\documentclass[11pt,oneside]{book}
\usepackage[T1]{fontenc}
\usepackage[utf8]{inputenc}
\usepackage{lmodern}
\usepackage[a4paper,margin=25mm]{geometry}
\usepackage{amsmath,amssymb,amsthm,enumitem,microtype}
\usepackage[hidelinks]{hyperref}
\hypersetup{pdftitle={Mathematical Foundations of Statistics and Learning},bookmarksnumbered=true}
\pagestyle{plain}
\setlength{\parindent}{0pt}
\setlength{\parskip}{4pt}
\setcounter{secnumdepth}{2}
\setcounter{tocdepth}{2}
\setlist[enumerate]{leftmargin=1.8em,itemsep=2pt,topsep=3pt}
\newtheorem{theorem}{Theorem}[chapter]
\newtheorem{proposition}[theorem]{Proposition}
\theoremstyle{definition}
\newtheorem{definition}[theorem]{Definition}
\newtheorem{example}[theorem]{Example}
\newtheorem{exercise}{Exercise}[section]
% Existing exercise blocks share the same heading and counter.
\let\practiceexercise\exercise
\let\endpracticeexercise\endexercise
\theoremstyle{remark}
\newtheorem*{remark}{Remark}
\DeclareMathOperator{\argmax}{argmax}
\DeclareMathOperator{\Law}{Law}
\newcommand{\Prob}{\mathsf{Prob}}
\newcommand{\R}{\mathbb R}

\newcommand{\notationentry}[3]{%
  \item[\ensuremath{#1}] #2\enspace\hyperref[#3]{(p.~\pageref*{#3})}}


% Declare notation at its definition; print the collected rows at the end.
\ExplSyntaxOn
\seq_new:N \g_book_notation_chapters_seq
\int_new:N \g_book_notation_serial_int
\tl_new:N \l_book_notation_chapter_tl
\tl_new:N \l_book_notation_label_tl

% A custom title is optional; new chapters are indexed automatically.
\NewDocumentCommand \notationgroup { m }
  { \cs_gset:cpn { book_notation_title_\theHchapter : } {#1} }

\NewDocumentCommand \indexnotation { m m }
  {
    \group_begin:
    \tl_set:Nx \l_book_notation_chapter_tl { \theHchapter }
    \cs_if_exist:cF { book_notation_title_\l_book_notation_chapter_tl : }
      {
        \cs_gset:cpx { book_notation_title_\l_book_notation_chapter_tl : }
          { Chapter~\thechapter }
      }
    \tl_if_exist:cF { g_book_notation_rows_\l_book_notation_chapter_tl _tl }
      {
        \tl_new:c { g_book_notation_rows_\l_book_notation_chapter_tl _tl }
        \seq_gput_right:NV \g_book_notation_chapters_seq
          \l_book_notation_chapter_tl
      }
    \int_gincr:N \g_book_notation_serial_int
    \tl_set:Nx \l_book_notation_label_tl
      { notation:\int_use:N \g_book_notation_serial_int }
    % Keep the enclosing definition's hyperlink anchor and capture this page.
    \exp_args:NV \label \l_book_notation_label_tl
    \tl_gput_right:cx { g_book_notation_rows_\l_book_notation_chapter_tl _tl }
      {
        \exp_not:N \notationentry
          { \exp_not:n {#1} } { \exp_not:n {#2} }
          { \l_book_notation_label_tl }
      }
    \group_end:
  }

\cs_new_protected:Npn \book_notation_print_group:n #1
  {
    \section*{\use:c { book_notation_title_#1: }}
    \begin{description}[style=multiline,leftmargin=4cm,
      labelwidth=3.7cm,labelsep=.3cm,itemsep=2pt]
    \tl_use:c { g_book_notation_rows_#1_tl }
    \end{description}
  }

\NewDocumentCommand \printnotationentries { }
  {
    \seq_map_function:NN \g_book_notation_chapters_seq
      \book_notation_print_group:n
  }
\ExplSyntaxOff

\begin{document}
\frontmatter
\tableofcontents
\mainmatter
\part{Statistical models and decisions}
\label{part:statistical-models-decisions}
\chapter{Statistical models and observations}
\label{chap:statistical-models}
\notationgroup{Experiments and sampling}
\section{Observation spaces, models, and parametrized families}
\label{sec:models-parametrizations}

\subsection{Observation spaces and data values}
\label{subsec:observation-spaces}

\begin{definition}[Observation space and data]
\label{def:observation-space}
\indexnotation{(\mathsf X,\mathcal A),\ X,\ x}{Observation space, random observation, and realized record.}
An \emph{observation space} is a nonempty measurable space
$(\mathsf X,\mathcal A)$. A \emph{data value} is an element $x\in\mathsf X$.
The law of a random observation taking values in this space is a probability
measure on $(\mathsf X,\mathcal A)$.
\end{definition}

Write $\Prob(\mathsf X,\mathcal A)$ for the set of these probability measures.
Once $P\in\Prob(\mathsf X,\mathcal A)$ is fixed, the canonical observation is
the identity map $X(x)=x$ on $(\mathsf X,\mathcal A,P)$; thus
$P(X\in A)=P(A)$ for $A\in\mathcal A$.
The notation $X$ denotes the random observation; $x$ denotes its realized value.

\begin{example}[The form of a data record]
\label{ex:observation-spaces}
A binary outcome uses $(\{0,1\},2^{\{0,1\}})$; a vector of $n\geq1$ real
measurements uses $(\R^n,\mathcal B(\R^n))$.
For a finite nonempty alphabet $V$, finite strings use
\[
\mathsf X=V^*:=\bigsqcup_{k=0}^{\infty}V^k,\qquad
\mathcal A=2^{V^*},\qquad V^0=\{\varepsilon\},
\]
where $\varepsilon$ is the empty string. A product observation space specifies
the form of the record; independence is a property of its law.
\end{example}

\subsection{Statistical models as collections of laws}
\label{subsec:statistical-models}

\begin{definition}[Statistical model]
\label{def:statistical-model}
\indexnotation{\Prob(\mathsf X,\mathcal A),\ \mathcal P}{Probability laws on the observation space, and the model of candidate laws.}
A \emph{statistical model} on $(\mathsf X,\mathcal A)$ is a nonempty set
\[
\mathcal P\subseteq\Prob(\mathsf X,\mathcal A).
\]
When the observation space must be displayed, write
$\mathcal M=(\mathsf X,\mathcal A,\mathcal P)$.
\end{definition}

The assertion that the data-generating law belongs to $\mathcal P$ is a
separate modeling assumption. On a fixed observation space, two models are
equal precisely when they contain the same probability measures.

\begin{example}[Bernoulli laws]
\label{ex:bernoulli}
\indexnotation{\delta_a}{Unit point mass at \(a\).}
On $\mathsf X=\{0,1\}$ with all subsets measurable, set
\[
P_p=(1-p)\delta_0+p\delta_1,\qquad
\mathcal P_{\mathrm B}=\{P_p:p\in[0,1]\},
\]
where $\delta_a$ denotes unit mass at $a$. Thus
$P_p(\{1\})=p$, $P_p(\{0\})=1-p$, and the model includes the deterministic
laws $P_0=\delta_0$ and $P_1=\delta_1$.
\end{example}

\begin{example}[Categorical laws]
\label{ex:categorical}
\indexnotation{\Delta_m}{Probability vectors on \(m\) outcomes.}
For $m\geq2$, take $\mathsf X=\{1,\ldots,m\}$ with all subsets measurable and
\[
\Delta_m:=\left\{p\in[0,1]^m:\sum_{i=1}^m p_i=1\right\}.
\]
Each $p\in\Delta_m$ gives $P_p=\sum_{i=1}^m p_i\delta_i$. Conversely, every
probability law $P$ on $\mathsf X$ satisfies
\[
P(A)=\sum_{i\in A}P(\{i\}),\qquad A\subseteq\mathsf X,
\]
so $p_i=P(\{i\})$ recovers its unique vector in $\Delta_m$.
Hence $\{P_p:p\in\Delta_m\}=\Prob(\mathsf X,2^{\mathsf X})$.
\end{example}

\begin{remark}[Assumed probability fact]
For probability measures $P,Q$ and nonnegative product-measurable $f$,
Tonelli's theorem gives
\[
\int f\,d(P\otimes Q)
=\int\left(\int f(u,v)\,Q(dv)\right)P(du).
\]
We use this result without proof.
\end{remark}

\begin{example}[Equal marginals, different joint models]
\label{ex:dependence}
\indexnotation{\Law(Z)}{Distribution of a random variable or tuple.}
Fix $\mu\in\R$ and let $\gamma_\mu=N(\mu,1)$. We construct two laws on
$(\R^2,\mathcal B(\R^2))$:
\[
P_\mu=\gamma_\mu\otimes\gamma_\mu=\Law(Z_1,Z_2),
\qquad Q_\mu=\Law(Z,Z),
\]
where $Z_1,Z_2$ are independent with law $\gamma_\mu$, and $Z$ has law
$\gamma_\mu$. In the second construction we draw $Z$ once and copy it
into both coordinates. For each Borel set $C\subseteq\R^2$, this means
\[
Q_\mu(C)=\mathbb P\big((Z,Z)\in C\big)
=\gamma_\mu\big(\{z\in\R:(z,z)\in C\}\big).
\]
The set on the right is Borel because $z\mapsto(z,z)$ is continuous.

For Borel sets $A,B\subseteq\R$, the probabilities of the rectangle
$A\times B$ are
\begin{align*}
P_\mu(A\times B)
&=\mathbb P(Z_1\in A,\ Z_2\in B)=\gamma_\mu(A)\gamma_\mu(B),\\
Q_\mu(A\times B)
&=\mathbb P(Z\in A,\ Z\in B)=\gamma_\mu(A\cap B).
\end{align*}
Taking $B=\R$ gives $\gamma_\mu(A)$ under both laws, so the first
coordinate marginals agree. Taking $A=\R$ gives $\gamma_\mu(B)$ under
both laws, so the second coordinate marginals agree too.

For a concrete difference, take $A=B=(-\infty,\mu]$. The normal law
is symmetric about $\mu$ and has no atom there, so $\gamma_\mu(A)=1/2$.
Thus the probability that both coordinates are at most $\mu$ is
$1/4$ under $P_\mu$ and $1/2$ under $Q_\mu$.

The diagonal $D=\{(u,v):u=v\}$ separates the two models for every
parameter value. It is closed, and $Q_\mu(D)=\mathbb P(Z=Z)=1$.
For the independent pair, Tonelli's theorem gives
\[
P_\mu(D)
=\int\left(\int \mathbf 1_{\{u=v\}}\,\gamma_\mu(dv)\right)\gamma_\mu(du)
=\int\gamma_\mu(\{u\})\,\gamma_\mu(du)
=0.
\]
In the inner integral $u$ is fixed. The only value of $v$ that contributes
is $v=u$, which has probability zero because $\gamma_\mu$ has a Lebesgue
density and a singleton has Lebesgue measure zero. Since these diagonal
probabilities hold for every $\mu$, no $P_\mu$ equals any $Q_\nu$.
\end{example}

\begin{example}[A parameter-dependent range on a fixed space]
\label{ex:uniform-endpoint}
On $(\R,\mathcal B(\R))$, consider
\[
\mathcal U=\{U_\theta:\theta>0\},\qquad
U_\theta(A)=\frac1\theta\int_A\mathbf 1_{[0,\theta]}(u)\,du.
\]
Here $U_\theta=\operatorname{Unif}[0,\theta]$.
The observation space remains $\R$ for every $\theta$, while
$U_\theta((\theta,\infty))=0$.
\end{example}

\subsection{Indexed families and reparametrizations}
\label{subsec:parametrizations}
\indexnotation{\pi,\ h}{Parametrization map and bijective change of labels.}

\begin{definition}[Parametrization]
\label{def:parametrization}
\indexnotation{\Theta,\ \theta,\ P_\theta}{Parameter space, label, and its observation law.}
A \emph{parametrization} of $\mathcal P$ is a surjection
\[
\pi:\Theta\longrightarrow\mathcal P,\qquad \theta\longmapsto P_\theta,
\]
where $\Theta$ is a nonempty set, called the \emph{parameter space}.
Thus $\mathcal P=\{P_\theta:\theta\in\Theta\}$.
\end{definition}

Injectivity, finite dimension, and measurable or topological structure on
$\Theta$ are additional conditions. Every model has the canonical
parametrization $\Theta=\mathcal P$, $\pi(P)=P$.

\begin{definition}[Reparametrization]
\label{def:reparametrization}
Let $(P_\theta)_{\theta\in\Theta}$ parametrize $\mathcal P$.
A \emph{reparametrization} is a bijection $h:\Lambda\to\Theta$ together with
the indexed family
\[
\widetilde P_\lambda=P_{h(\lambda)},\qquad \lambda\in\Lambda.
\]
\end{definition}

For a nonempty set $S$ and $f:S\to\overline{\R}:=[-\infty,\infty]$, write
\[
\argmax_{s\in S}f(s)
:=\{s\in S:f(s)=\sup_{u\in S}f(u)\}.
\]
This set may be empty.

\begin{proposition}[Change of parameter]
\label{prop:reparametrization}
Every reparametrization represents the same model. Its inverse recovers the
original indexed family, and compositions of reparametrizations are
reparametrizations.
\end{proposition}
\begin{proof}
The parametrization map of the new family is $\widetilde\pi=\pi\circ h$, so
$\widetilde\pi(\Lambda)=\pi(\Theta)=\mathcal P$. Also
$\widetilde P_{h^{-1}(\theta)}=P_\theta$.
If $g:\Gamma\to\Lambda$ is bijective, then
$\widetilde P_{g(\gamma)}=P_{(h\circ g)(\gamma)}$ and $h\circ g$ is bijective.
\end{proof}

\begin{proposition}[Optimization over a parametrized model]
\label{prop:parametrized-optimization}
\indexnotation{\argmax_{s\in S}f(s)}{Set of points attaining the supremum.}
Let $\pi:\Theta\to\mathcal P$, $\pi(\theta)=P_\theta$, be a parametrization,
and let $C:\mathcal P\to\overline{\R}$. Then
\begin{align}
\sup_{\theta\in\Theta}C(P_\theta)
&=\sup_{P\in\mathcal P}C(P), \label{eq:parametrized-supremum}\\
\argmax_{\theta\in\Theta}C(P_\theta)
&=\pi^{-1}\!\left(\argmax_{P\in\mathcal P}C(P)\right).
\label{eq:parametrized-argmax}
\end{align}
In particular, for a reparametrization $h:\Lambda\to\Theta$,
\[
\argmax_{\lambda\in\Lambda}C(\widetilde P_\lambda)
=h^{-1}\!\left(\argmax_{\theta\in\Theta}C(P_\theta)\right).
\]
\end{proposition}
\begin{proof}
Surjectivity gives
\[
\{C(P_\theta):\theta\in\Theta\}=\{C(P):P\in\mathcal P\},
\]
hence \eqref{eq:parametrized-supremum}. For $\theta\in\Theta$,
\[
\theta\in\argmax_{\vartheta\in\Theta}C(P_\vartheta)
\iff C(P_\theta)=\sup_{P\in\mathcal P}C(P)
\iff P_\theta\in\argmax_{P\in\mathcal P}C(P),
\]
which proves \eqref{eq:parametrized-argmax}. For the last assertion, use
$\widetilde\pi=\pi\circ h$ and
$(\pi\circ h)^{-1}(A)=h^{-1}(\pi^{-1}(A))$ for $A\subseteq\mathcal P$.
\end{proof}

\subsection{Worked parametrizations and exercises}
\label{subsec:parametrization-examples}

\begin{example}[Redundant Bernoulli labels]
\label{ex:redundant-labels}
For the model $\mathcal P_{\mathrm B}$ of Example~\ref{ex:bernoulli}, let
\[
\widetilde P_q=P_{q^2},\qquad q\in[-1,1].
\]
Every $P_p$ occurs, since $q=\sqrt p$ is allowed. For $p>0$ the distinct
labels $\sqrt p$ and $-\sqrt p$ describe the same law.
Thus $q\mapsto P_{q^2}$ is a parametrization, while $q\mapsto q^2$ is not
a reparametrization in Definition~\ref{def:reparametrization}.
\end{example}

\begin{example}[Counterexample: dropping boundary laws]
\label{ex:logistic-boundary}
Set
\[
h(\eta)=\frac{1}{1+e^{-\eta}},\qquad \eta\in\R.
\]
This is a bijection $\R\to(0,1)$ with inverse
$h^{-1}(p)=\log(p/(1-p))$. Consequently,
\[
\{P_{h(\eta)}:\eta\in\R\}
=\{P_p:0<p<1\}
\subsetneq\mathcal P_{\mathrm B}.
\]
It reparametrizes the interior Bernoulli model. It omits $P_0$ and $P_1$
from the closed model.
\end{example}

\begin{example}[Maximizers under the two changes of labels]
\label{ex:maximizers}
Fix $a\in[0,1]$ and define the criterion on Bernoulli laws by
\[
C_a(P)=-\bigl(P(\{1\})-a\bigr)^2.
\]
Since $C_a(P_p)=-(p-a)^2\leq0$, with equality exactly when $p=a$,
\[
\argmax_{P\in\mathcal P_{\mathrm B}}C_a(P)=\{P_a\},\qquad
\argmax_{q\in[-1,1]}C_a(P_{q^2})=\{-\sqrt a,\sqrt a\}.
\]
The second set is $\{0\}$ when $a=0$.
For the logistic parametrization,
\[
\argmax_{\eta\in\R}C_a(P_{h(\eta)})
=
\begin{cases}
\{\log(a/(1-a))\},&0<a<1,\\
\varnothing,&a\in\{0,1\}.
\end{cases}
\]
At either endpoint the supremum is still $0$, approached as
$\eta\to-\infty$ or $\eta\to+\infty$, respectively, but never attained.
The endpoint maximizer was lost when the model was restricted.
\end{example}

\subsection{Exercises}
\label{subsec:worked-examples-1-1}

\begin{practiceexercise}[A die roll and an event]
\label{wex:die-record}
On $\{1,\ldots,6\}$ with its power set and uniform law, distinguish record
$4$ from the event of an even record. Compute the probability of that event.
\end{practiceexercise}

\begin{practiceexercise}[Events on finite strings]
\label{wex:prefix-event}
Give $\varepsilon,a,ab$ masses $1/4,1/4,1/2$ on $\{a,b\}^*$. Find the
probability of a string starting with $a$ and of a string of length at most
one.
\end{practiceexercise}

\begin{practiceexercise}[An unordered pair as a record]
\label{wex:unordered-pair-space}
Draw two independent fair bits and report their values in increasing order.
Specify the observation space and the mass of every record.
\end{practiceexercise}

\begin{practiceexercise}[Reporting a threshold crossing]
\label{wex:threshold-report-space}
Let $U$ be uniform on $[0,1]$. Specify observation spaces and laws for the
records $U$ and $\mathbf1_{\{U>3/4\}}$.
\end{practiceexercise}

\begin{practiceexercise}[Lengths as observable events]
\label{wex:length-classes}
On $\{a,b\}^*$ compare $\delta_{ab}$ and $\delta_{ba}$ using first the sigma
algebra of unions of length classes, then the power set. Decide in each case
whether the laws differ.
\end{practiceexercise}

\begin{practiceexercise}[A record with two types of outcome]
\label{wex:sensor-status}
A sensor reports missing with probability $1/5$ and otherwise gives a
uniform reading on $[0,1]$. Specify a measurable observation space and
compute the probability of missing or a reading at most $1/2$.
\end{practiceexercise}

\begin{practiceexercise}[A restricted Bernoulli model]
\label{wex:restricted-bernoulli-model}
Decide whether $\delta_0$ and $P_{1/2}$ belong to $\{P_p:0\le p\le1/3\}$.
Justify your decisions using one observable event.
\end{practiceexercise}

\begin{practiceexercise}[Comparing event probabilities]
\label{wex:event-categorical}
Compare the categorical laws with vectors
\[
(1/8,1/8,1/4,1/2),\qquad (1/4,0,1/4,1/2).
\]
Compute both probabilities of $\{3,4\}$ and find an event on which the laws
differ.
\end{practiceexercise}

\begin{practiceexercise}[A simplex face]
\label{wex:categorical-face}
In the three-category model impose $P(\{3\})=0$. Parametrize the resulting
family by one coordinate. Which point masses belong? Does the uniform law on
three categories belong?
\end{practiceexercise}

\begin{practiceexercise}[Mixing Bernoulli laws]
\label{wex:two-point-convex}
Find the mass at one under $wP_{1/4}+(1-w)P_{3/4}$. Describe the model as
$w$ varies over $[0,1]$.
\end{practiceexercise}

\begin{practiceexercise}[A union of parameter ranges]
\label{wex:model-union}
Take the union of the Bernoulli submodels with $p\in[0,1/4]$ and
$p\in[3/4,1]$. Does the union contain the equal mixture of $P_0$ and $P_1$?
Identify that mixture.
\end{practiceexercise}

\begin{practiceexercise}[Comparing two normal laws]
\label{wex:normal-mean-separation}
Compare the probabilities of $(-\infty,0]$ under $N(0,1)$ and $N(1,1)$.
Justify the strict comparison without numerical tables.
\end{practiceexercise}

\begin{practiceexercise}[Complementary and independent bits]
\label{wex:same-marginals-fixed-sum}
Compare the marginals and the probability of a sum of one for $(Z,1-Z)$ with
$Z$ a fair bit and for two independent fair bits.
\end{practiceexercise}

\begin{practiceexercise}[A saturated measurement]
\label{wex:saturated-uniform}
For $U$ uniform on $[0,2]$, report $Y=\min(U,1)$. Compute $P(Y=1)$ and
$P(Y\le t)$ for $0<t<1$. Identify the atom and the continuous part.
\end{practiceexercise}

\begin{practiceexercise}[Two endpoint laws]
\label{wex:uniform-scale-separation}
For $0<a<b$, compare $\operatorname{Unif}[0,a]$ and
$\operatorname{Unif}[0,b]$ on $(a,b]$. Specify a common observation space.
\end{practiceexercise}

\begin{practiceexercise}[Signed labels for point masses]
\label{wex:signed-point-label}
Describe the image and fibers of $t\mapsto\delta_{t^2}$ on $\R$. Does
restricting to $t\ge0$ change the model? Is it a reparametrization in
Definition~\ref{def:reparametrization}?
\end{practiceexercise}

\begin{practiceexercise}[Reordering a finite family]
\label{wex:integer-label-reorder}
Let $P_i=\delta_i$, $i=1,2,3$, and let $h$ send $1,2,3$ to $2,3,1$. Check
whether $\widetilde P_j=P_{h(j)}$ is a reparametrization. Find the new label
of a criterion's unique maximizing law $\delta_3$.
\end{practiceexercise}

\begin{practiceexercise}[Variance or standard deviation]
\label{wex:normal-scale-label}
Compare the labels $v>0$ in $N(0,v)$ and $s$ in $N(0,s^2)$. Give the inverse
change for $s>0$ and explain what happens if negative labels are also
admitted.
\end{practiceexercise}

\begin{practiceexercise}[Rate or mean of an exponential law]
\label{wex:exponential-rate-mean}
For survival function $e^{-\lambda t}$, $t\ge0$, $\lambda>0$, check whether
$m=1/\lambda$ gives a reparametrization. Compute the mass of $(2,\infty)$ in
each parametrization.
\end{practiceexercise}

\begin{practiceexercise}[Odds as a Bernoulli label]
\label{wex:odds-reparam}
Invert $r=p/(1-p)$ on $0<p<1$. Find the law at odds two and decide which
endpoint laws finite positive odds represent.
\end{practiceexercise}

\begin{practiceexercise}[Centering a parameter or the data]
\label{wex:shifted-location-label}
For $P_\theta=N(\theta,1)$, find the law at label zero after relabelling by
$\eta=\theta-3$. Compare with the law of the transformed record $X-3$ at
original label zero.
\end{practiceexercise}

\begin{practiceexercise}[Maximizing in two parametrizations]
\label{wex:quadratic-tied-criterion}
Set $C(P_p)=-((p-1/4)(p-3/4))^2$ on the full Bernoulli model. Find all
maximizing laws and their labels under $\lambda=2p-1$.
\end{practiceexercise}

\begin{practiceexercise}[A criterion on an open range]
\label{wex:restricted-criterion-supremum}
For $C(P_p)=p$ on $0<p<1/2$, find the supremum and determine whether it is
attained. Repeat in labels $u=2p$.
\end{practiceexercise}

\begin{practiceexercise}[Maximizing an indicator criterion]
\label{wex:criterion-with-plateau}
For $p\in[0,1]$, take $C(P_p)=\mathbf1_{\{p\ge2/3\}}$. Find the maximizing
set first in labels $p$, then in labels $\lambda=1-p$.
\end{practiceexercise}

\begin{exercise}[Recording a count]
\label{exer:recording-a-count}
On $\mathsf X=\{0,1\}^2$, let $T(x_1,x_2)=x_1+x_2$.
List the fibers of $T$. Prove that no map
$d:\{0,1,2\}\to\mathsf X$ satisfies $d\circ T=\operatorname{id}_{\mathsf X}$.
Compare recording an ordered pair with recording only its count.
\end{exercise}

\begin{exercise}[Removing a simplex coordinate]
\label{exer:simplex-coordinates}
For $m\geq2$, set
\[
\Lambda=\left\{q\in[0,1]^{m-1}:\sum_{i=1}^{m-1}q_i\leq1\right\}.
\]
Prove that $h(q)=(q_1,\ldots,q_{m-1},1-\sum_{i=1}^{m-1}q_i)$ is a bijection
$\Lambda\to\Delta_m$. Write the categorical laws in these coordinates.
\end{exercise}

\begin{exercise}[Which Gaussian model?]
\label{exer:gaussian-labels}
Compare $\{N(\mu,1):\mu\in\R\}$ with each of
\[
\{N(e^\eta,1):\eta\in\R\},\qquad
\{N(\eta^2,1):\eta\in\R\},\qquad
\{N(\eta^3,1):\eta\in\R\}.
\]
Determine the represented models and whether each displayed change of labels
is bijective onto the original parameter space.
\end{exercise}

\begin{exercise}[Surjectivity at the correct level]
\label{exer:surjectivity}
Let $\pi:\Theta\to\mathcal P$ be a parametrization, $\Lambda$ nonempty,
and $h:\Lambda\to\Theta$ any map.
Show that $\pi\circ h$ parametrizes $\mathcal P$ if and only if
$\pi(h(\Lambda))=\mathcal P$.
Construct a finite example in which this holds although $h$ is not surjective.
\end{exercise}

\begin{exercise}[Changing the observable events]
\label{exer:coarse-sigma-algebra}
For the Bernoulli model, replace $2^{\{0,1\}}$ by
$\mathcal A_0=\{\varnothing,\{0,1\}\}$.
Compute the restricted measures $P_p|_{\mathcal A_0}$.
How many laws does the resulting model contain?
\end{exercise}

\begin{exercise}[Maximizers and derivatives]
\label{exer:derivatives}
Let $F(\theta)=-(\theta-1)^2$ on $\R$ and $h(\lambda)=\lambda^3$.
Prove that $h$ is bijective. Compute every zero of $F'$ and of $(F\circ h)'$,
and determine which are maximizers.
Verify that $h$ carries the maximizing sets bijectively, while zeros of the
derivatives need not correspond.
\end{exercise}
\clearpage
\section{Sampling schemes and product experiments}
\label{sec:sampling-schemes}

\subsection{Statistical experiments and joint sampling laws}
\label{subsec:experiments-sampling}

\begin{definition}[Statistical experiment]
\label{def:statistical-experiment}
\indexnotation{\mathcal E}{Statistical experiment with its parameter labels.}
A \emph{statistical experiment} indexed by a nonempty set $\Theta$ is
\[
\mathcal E=(\mathsf X,\mathcal A,(P_\theta)_{\theta\in\Theta}),
\qquad P_\theta\in\Prob(\mathsf X,\mathcal A).
\]
Its associated model is $\{P_\theta:\theta\in\Theta\}$.
The experiment retains the parameter labels.
\end{definition}

\begin{definition}[Fixed-size sampling scheme]
\label{def:fixed-size-sampling}
\indexnotation{Q_\theta,\ X_i}{Joint law of the sample and its coordinate observations.}
Fix $n\geq1$ and observation spaces $(\mathsf X_i,\mathcal A_i)$.
A \emph{fixed-size sampling scheme} is an experiment
\[
\left(\prod_{i=1}^n\mathsf X_i,\,
\bigotimes_{i=1}^n\mathcal A_i,\,
(Q_\theta)_{\theta\in\Theta}\right).
\]
The coordinate observations are $X_i(x_1,\ldots,x_n)=x_i$.
Each $Q_\theta$ specifies their joint law.
\end{definition}

\begin{remark}[Assumed probability fact: finite products]
For probability measures $P_i$ on $(\mathsf X_i,\mathcal A_i)$, there is a
unique probability measure $\bigotimes_{i=1}^nP_i$ on the product space with
\[
\left(\bigotimes_{i=1}^nP_i\right)
\left(\prod_{i=1}^n A_i\right)=\prod_{i=1}^nP_i(A_i),
\qquad A_i\in\mathcal A_i.
\]
A joint law equals this product exactly when its coordinate observations
are independent and have marginal laws $P_i$.
\end{remark}

\subsection{Independent composition and iid repetition}
\label{subsec:product-experiments}

\begin{definition}[Product experiment and iid repetition]
\label{def:product-experiment}
\indexnotation{\otimes_i P_{\theta,i},\ P_\theta^{\otimes n}}{Independent product and iid repetition.}
Let $\mathcal E_i=(\mathsf X_i,\mathcal A_i,
(P_{\theta,i})_{\theta\in\Theta})$ share the same parameter set.
Their \emph{independent product experiment} is
\[
\bigotimes_{i=1}^n\mathcal E_i
:=
\left(\prod_{i=1}^n\mathsf X_i,\,
\bigotimes_{i=1}^n\mathcal A_i,\,
\left(\bigotimes_{i=1}^nP_{\theta,i}\right)_{\theta\in\Theta}\right).
\]
For one experiment $\mathcal E$, its \emph{iid repetition}
$\mathcal E^{\otimes n}$ has laws $P_\theta^{\otimes n}$.
\end{definition}

All independence assertions concern a fixed $\theta$. In a product experiment
the observations may have different marginal laws. In an iid repetition they
are independent and identically distributed, with one shared parameter value.

\begin{proposition}[Equality of product laws]
\label{prop:product-law-equality}
For product experiments with $n\geq1$ and $\theta,\vartheta\in\Theta$,
\[
\bigotimes_{i=1}^nP_{\theta,i}
=\bigotimes_{i=1}^nP_{\vartheta,i}
\quad\Longleftrightarrow\quad
P_{\theta,i}=P_{\vartheta,i}\quad(1\leq i\leq n).
\]
In particular,
$P_\theta^{\otimes n}=P_\vartheta^{\otimes n}$ if and only if
$P_\theta=P_\vartheta$.
\end{proposition}
\begin{proof}
The $i$th coordinate marginal of a product law is its $i$th factor.
Equality of the joint laws therefore implies equality of every factor.
Conversely, equal factors give equal products. Apply this with
$P_{\theta,i}=P_\theta$ for the iid assertion.
\end{proof}
\begin{example}[Shared and separate parameter values]
\label{ex:shared-sampling-parameter}
For the Bernoulli laws of Example~\ref{ex:bernoulli}, the models
\[
\mathcal P_{\mathrm{same}}
=\{P_p\otimes P_p:p\in[0,1]\},\qquad
\mathcal P_{\mathrm{free}}
=\{P_p\otimes P_q:(p,q)\in[0,1]^2\}
\]
satisfy $\mathcal P_{\mathrm{same}}\subsetneq\mathcal P_{\mathrm{free}}$.
Indeed, $\delta_{(0,1)}=P_0\otimes P_1$ belongs to the second model.
If it equaled $P_p\otimes P_p$, its two marginals would force both $p=0$
and $p=1$. The two parameter sets specify different experiments and models.
\end{example}

\begin{example}[Independent but nonidentical measurements]
\label{ex:fixed-design-gaussian}
Fix known $z_1,\ldots,z_n\in\R$ and let
\[
Y_i=\theta z_i+\varepsilon_i,\qquad
\theta\in\R,\qquad
\varepsilon_1,\ldots,\varepsilon_n\ \text{iid }N(0,1).
\]
The observation law is $Q_\theta=\bigotimes_{i=1}^nN(\theta z_i,1)$.
For $n\geq2$, $z_1=0$, $z_2=1$, and $\theta\ne0$, the first two marginals differ.
This is an independent composition, but not an iid repetition.
\end{example}

\subsection{Empirical laws and averages}
\label{subsec:empirical-laws}

\begin{definition}[Empirical law]
\label{def:empirical-law}
\indexnotation{\widehat P_n,\ Pf}{Empirical law and the integral of \(f\) under \(P\).}
For $x=(x_1,\ldots,x_n)\in\mathsf X^n$, its \emph{empirical law} is
\[
\widehat P_n(x)=\frac1n\sum_{i=1}^n\delta_{x_i}.
\]
For a probability law $P$, write $Pf=\int f\,dP$ whenever defined.
For real-valued measurable $f$,
\[
\widehat P_n(x)f=\frac1n\sum_{i=1}^n f(x_i),\qquad
\widehat P_n(x)(A)=\frac1n\sum_{i=1}^n\mathbf1_A(x_i),
\quad A\in\mathcal A.
\]
\end{definition}

For a random sample, these scalar evaluations are measurable, being finite
sums of measurable coordinate functions.

\begin{proposition}[Finite-sample empirical averages]
\label{prop:empirical-moments}
Fix $\theta$. Let $Q_\theta$ be a sampling law on $\mathsf X^n$ under which
every $X_i$ has marginal law $P_\theta$.
For measurable $f:\mathsf X\to\R$ with $P_\theta f^2<\infty$,
\begin{align}
\mathbb E_{Q_\theta}[\widehat P_n f]
&=P_\theta f, \label{eq:empirical-expectation}\\
\operatorname{Var}_{Q_\theta}(\widehat P_n f)
&=\frac1{n^2}\sum_{i=1}^n\sum_{j=1}^n
\operatorname{Cov}_{Q_\theta}(f(X_i),f(X_j)).
\label{eq:empirical-covariance}
\end{align}
For the iid law $Q_\theta=P_\theta^{\otimes n}$, this yields
\begin{equation}
\mathbb E_{P_\theta^{\otimes n}}
\left[(\widehat P_n f-P_\theta f)^2\right]
=\frac{P_\theta f^2-(P_\theta f)^2}{n}.
\label{eq:iid-mean-square}
\end{equation}
\end{proposition}
\begin{proof}
Put $\mu=P_\theta f$ and $U_i=f(X_i)-\mu$. The second-moment assumption
ensures that all the following expectations are finite. Since
$\mathbb E_{Q_\theta}U_i=0$, linearity gives
\eqref{eq:empirical-expectation}, and
\[
\operatorname{Var}_{Q_\theta}(\widehat P_n f)
=\mathbb E_{Q_\theta}\left[\left(\frac1n\sum_iU_i\right)^2\right]
=\frac1{n^2}\sum_{i,j}\mathbb E_{Q_\theta}[U_iU_j],
\]
which is \eqref{eq:empirical-covariance}.
Under the iid law, off-diagonal expectations vanish and each diagonal term
is $P_\theta f^2-\mu^2$. Unbiasedness then gives
\eqref{eq:iid-mean-square}.
\end{proof}
\begin{example}[Empirical Bernoulli frequencies]
\label{ex:empirical-bernoulli}
For $x\in\{0,1\}^n$, put $\bar x_n=n^{-1}\sum_i x_i$. Then
\[
\widehat P_n(x)=(1-\bar x_n)\delta_0+\bar x_n\delta_1.
\]
For $X_1,\ldots,X_n$ iid with law $P_p$, take $f(u)=u$.
Since $P_pf=P_pf^2=p$, Proposition~\ref{prop:empirical-moments} gives
\[
\mathbb E[\bar X_n]=p,\qquad
\operatorname{Var}(\bar X_n)=\frac{p(1-p)}n,
\qquad \bar X_n=\frac1n\sum_iX_i.
\]
\end{example}

\begin{example}[Counterexample: copied observations]
\label{ex:copied-observations}
Let $Z\sim P$ and set $X_1=\cdots=X_n=Z$.
For real measurable $f$ with $Pf^2<\infty$, put
$\sigma^2=Pf^2-(Pf)^2$.
Every coordinate has law $P$, but
\[
\widehat P_n f=f(Z),\qquad
\operatorname{Cov}(f(X_i),f(X_j))=\sigma^2\quad\text{for every }i,j.
\]
Hence $\mathbb E[\widehat P_n f]=Pf$ and
$\operatorname{Var}(\widehat P_n f)=\sigma^2$.
For $n>1$ and $\sigma^2>0$, the iid variance formula fails.
\end{example}

\subsection{Sampling finite populations}
\label{subsec:finite-population-sampling}

\begin{example}[Sampling without replacement]
\label{ex:finite-population-variance}
Fix $N\geq2$, $1\leq n\leq N$, and a population
$a=(a_1,\ldots,a_N)\in\R^N$.
Choose $(I_1,\ldots,I_n)$ uniformly among ordered tuples of distinct
indices in $\{1,\ldots,N\}$, and observe $X_i=a_{I_i}$.
Write $Q_a$ for the joint observation law and set
\[
P_a=\frac1N\sum_{r=1}^N\delta_{a_r},\qquad
\mu=\frac1N\sum_{r=1}^N a_r,\qquad
v=\frac1N\sum_{r=1}^N(a_r-\mu)^2.
\]
Each $X_i$ has marginal law $P_a$. For $i\ne j$, each ordered pair of
distinct indices has probability $1/[N(N-1)]$. Since
$\sum_r(a_r-\mu)=0$,
\[
\begin{aligned}
\operatorname{Cov}_{Q_a}(X_i,X_j)
&=\frac1{N(N-1)}\sum_{r\ne s}(a_r-\mu)(a_s-\mu)\\
&=-\frac1{N(N-1)}\sum_{r=1}^N(a_r-\mu)^2
=-\frac v{N-1}.
\end{aligned}
\]
For $\bar X_n=n^{-1}\sum_iX_i$,
Proposition~\ref{prop:empirical-moments} therefore gives
\begin{equation}
\mathbb E_{Q_a}[\bar X_n]=\mu,\qquad
\operatorname{Var}_{Q_a}(\bar X_n)
=\frac1{n^2}\left(nv-\frac{n(n-1)v}{N-1}\right)
=\frac vn\frac{N-n}{N-1}.
\label{eq:finite-population-variance}
\end{equation}
The variance $v$ uses denominator $N$.
Relative to iid sampling from $P_a$, the variance is multiplied by
$(N-n)/(N-1)$. At $n=N$, every allowed sample has average $\mu$.
\end{example}

\begin{example}[Hypergeometric sampling]
\label{ex:hypergeometric}
In Example~\ref{ex:finite-population-variance}, take a population of $M$ ones
and $N-M$ zeros, where $M\in\{0,\ldots,N\}$ is the parameter.
Let $K=\sum_iX_i$. Each $n$-element subset of indices has $n!$ equally
likely orderings, so it has probability $\binom Nn^{-1}$.
Counting subsets with $k$ ones gives
\[
Q_M(K=k)=\frac{\binom Mk\binom{N-M}{n-k}}{\binom Nn},
\qquad 0\leq k\leq n,
\]
with $\binom ab=0$ when $b<0$ or $b>a$.
This is the hypergeometric law.
For $p=M/N$, the marginal law is $P_p$ and $v=p(1-p)$; hence
\[
\mathbb E_{Q_M}[K/n]=p,\qquad
\operatorname{Var}_{Q_M}(K/n)
=\frac{p(1-p)}n\frac{N-n}{N-1}.
\]
The formulas include $M=0$, $M=N$, and $n=N$.
\end{example}

\subsection{Exercises}
\label{subsec:worked-examples-1-2}

\begin{practiceexercise}[Independent coins with unequal probabilities]
\label{wex:different-coins}
For independent coins of success probabilities $1/4$ and $3/4$, compute the
chances of both successes and at least one success, and the expected number
of successes.
\end{practiceexercise}

\begin{practiceexercise}[Which coin produced the success?]
\label{wex:unequal-coins-given-count}
For independent coins of probabilities $1/4,3/4$, calculate the probability
that the first coin succeeded given exactly one success.
\end{practiceexercise}

\begin{practiceexercise}[A mean under a fixed design]
\label{wex:fixed-design-average}
Let $Y_i=\theta z_i+\varepsilon_i$, with $z=(-1,0,2)$ and independent
standard normal errors. Write the joint law and find the mean and variance
of $\bar Y$.
\end{practiceexercise}

\begin{practiceexercise}[Zero or one arrival in two intervals]
\label{wex:poisson-zero-or-one}
For independent $X_1\sim\operatorname{Pois}(a)$ and
$X_2\sim\operatorname{Pois}(b)$, compute the probabilities that $X_1+X_2$ is
zero or one by listing the contributing records.
\end{practiceexercise}

\begin{practiceexercise}[Two specified categories]
\label{wex:categorical-two-hits}
Take two iid draws from a law with category probabilities
$(1/2,1/3,1/6)$. Find the probability of one category-two and one
category-three record, accounting for both orders.
\end{practiceexercise}

\begin{practiceexercise}[Independent clusters of copies]
\label{wex:clustered-copies}
Let $Z_1,\ldots,Z_m$ be iid Bernoulli with parameter $p$, and record
$(Z_j,Z_j)$ for each $j$. Compare the variance of the $2m$-coordinate mean
with that of $2m$ independent draws.
\end{practiceexercise}

\begin{practiceexercise}[Partial copying]
\label{wex:partial-copy-covariance}
Let $X,Z$ be independent Bernoulli variables with parameter $p$.
Independently of them choose $Y=X$ with probability $r$ and $Y=Z$ otherwise.
Find the law of $Y$, $\mathbb E[XY]$, and $\operatorname{Var}((X+Y)/2)$.
\end{practiceexercise}

\begin{practiceexercise}[Averaging complementary bits]
\label{wex:paired-complement-mean}
For a fair bit $Z$, compute the covariance of $Z$ and $1-Z$. Use it to find
the variance of their average and verify directly from the possible records.
\end{practiceexercise}

\begin{practiceexercise}[Evaluating an empirical law]
\label{wex:empirical-feature-square}
Write the empirical law of $(-2,0,1,1)$. Compute its mass at one and its
integral of $f(x)=x^2$.
\end{practiceexercise}

\begin{practiceexercise}[Empirical mass of a union of categories]
\label{wex:empirical-category-subset}
Find the empirical vector of $(a,b,b,c,b,a)$ in order $a,b,c$. Use it to
compute the empirical mass of $\{a,c\}$.
\end{practiceexercise}

\begin{practiceexercise}[A category absent from the record]
\label{wex:empirical-unseen-category}
For $n$ iid draws and a category of true probability $q$, find the
probability that its empirical mass is zero. Evaluate at $n=5$, $q=1/2$.
\end{practiceexercise}

\begin{practiceexercise}[Mean-square error for a uniform sample]
\label{wex:uniform-mean-square}
For $n$ iid uniform draws on $[0,1]$, compute the first two moments of one
draw and the mean-square error of the sample mean about $1/2$.
\end{practiceexercise}

\begin{practiceexercise}[From mean-square error to error probability]
\label{wex:coin-error-probability}
Use the Bernoulli sample mean's mean-square error to bound $P(|\bar
X-p|\ge\epsilon)$. At $n=100$, $\epsilon=1/5$, give a bound valid for every
$p$.
\end{practiceexercise}

\begin{practiceexercise}[Two draws from four numerical values]
\label{wex:four-values-without-replacement}
Sample two distinct items uniformly from $\{1,2,3,4\}$. List the unordered
pairs and their sample means. Find the mean and variance of the resulting
distribution.
\end{practiceexercise}

\begin{practiceexercise}[At least one marked item]
\label{wex:rare-in-population}
In a population of ten items, three marked, draw two distinct items
uniformly. Find the chance of at least one marked item by counting the
complementary pairs.
\end{practiceexercise}

\begin{practiceexercise}[The empirical law of a census]
\label{wex:census-empirical-law}
Sample all four items of $(2,2,5,9)$ without replacement. Determine whether
the order and the empirical law vary. Compute the empirical law.
\end{practiceexercise}

\begin{practiceexercise}[Collision with and without replacement]
\label{wex:distinct-items-collision}
Draw twice from four distinct items, once with independent replacement and
once without. Compare the coordinate marginals and the probability of a
repeat.
\end{practiceexercise}

\begin{practiceexercise}[Adding a noise coordinate]
\label{wex:third-coordinate-marginal}
For $Q_\theta=P_\theta\otimes U$ with fixed $U$, prove
$Q_\theta=Q_\vartheta$ exactly when $P_\theta=P_\vartheta$. Use rectangle
probabilities for the product implication.
\end{practiceexercise}

\begin{exercise}[An ordered record and its count]
\label{exer:bernoulli-record-count}
For $n$ iid Bernoulli observations with parameter $p\in[0,1]$, compute the
probability of a specified binary record containing $k$ ones and the
probability of observing exactly $k$ ones. Derive the latter by counting
records. Treat $p=0$ and $p=1$ separately.
\end{exercise}

\begin{exercise}[Different marginal means]
\label{exer:nonidentical-bernoulli}
Let $X_i$ be independent Bernoulli variables with parameters $p_i$.
Compute the expectation and variance of $\bar X_n$.
Which assumption of Proposition~\ref{prop:empirical-moments} fails when the
$p_i$ are not all equal?
\end{exercise}

\begin{exercise}[A variance identity does not imply independence]
\label{exer:pairwise-independent}
Let $U,V$ be independent with
$\mathbb P(U=1)=\mathbb P(U=-1)=
\mathbb P(V=1)=\mathbb P(V=-1)=1/2$.
Set $(X_1,X_2,X_3)=(U,V,UV)$.
Show that each coordinate has the same law, every pair is independent,
and the three coordinates are not jointly independent.
Compute $\operatorname{Var}((X_1+X_2+X_3)/3)$.
\end{exercise}

\begin{exercise}[A finite population calculation]
\label{exer:finite-population-calculation}
Draw two distinct indices uniformly from the population $(0,0,1,1)$.
Find the complete law of the sample average and compute its variance directly.
Compare with \eqref{eq:finite-population-variance} and with iid sampling
from the same marginal law.
\end{exercise}

\begin{exercise}[Empirical laws discard order]
\label{exer:empirical-order}
Prove that $\widehat P_n$ is unchanged by permuting a record.
For the two laws $Q_0=\delta_{(0,1)}$ and $Q_1=\delta_{(1,0)}$ on
$\{0,1\}^2$, show that both possible records have the same empirical law.
Exhibit a measurable event separating $Q_0$ and $Q_1$.
\end{exercise}

\begin{exercise}[Mixtures and products]
\label{exer:mixture-products}
For the Bernoulli laws $P_0=\delta_0$ and $P_1=\delta_1$, compare
\[
Q=\tfrac12P_0^{\otimes2}+\tfrac12P_1^{\otimes2},
\qquad
R=\left(\tfrac12P_0+\tfrac12P_1\right)^{\otimes2}.
\]
Compute both laws on the four records, their marginal laws, and
$\operatorname{Cov}(X_1,X_2)$ under each.
\end{exercise}

\clearpage
\section{Parameters, statistical functionals, and identifiability}
\label{sec:identifiability}

Throughout this section, $\pi:\Theta\to\mathcal P$,
$\pi(\theta)=P_\theta$, is a parametrization on $(\mathsf X,\mathcal A)$.

\subsection{Targets and functionals}
\label{subsec:targets-functionals}

\begin{definition}[Target and statistical functional]
\label{def:target-functional}
\indexnotation{g,\ T}{Target as a function of the label, and functional as a function of the law.}
Let $S$ be a nonempty set. A \emph{target} is a map $g:\Theta\to S$.
A \emph{statistical functional} on $\mathcal P$ is a map
$T:\mathcal P\to S$. Its associated target is
\[
g(\theta)=T(P_\theta),\qquad \theta\in\Theta.
\]
The word functional does not require linearity.
\end{definition}

\begin{example}[Probabilities and moments as functionals]
\label{ex:law-functionals}
\indexnotation{\mu(P),\ v(P)}{Mean and variance as functionals of a law.}
For $A\in\mathcal A$, the map $T_A(P)=P(A)$ is a functional
$\mathcal P\to[0,1]$.
On $(\R,\mathcal B(\R))$, define
\[
\mathcal P_j=\left\{P\in\Prob(\R,\mathcal B(\R)):
\int |x|^j\,P(dx)<\infty\right\},\qquad j=1,2.
\]
The mean and variance functionals are
\[
\mu:\mathcal P_1\to\R,\quad \mu(P)=\int x\,P(dx),\qquad
v:\mathcal P_2\to[0,\infty),\quad
v(P)=\int(x-\mu(P))^2\,P(dx).
\]
Since $|x|\leq1+x^2$, $\mathcal P_2\subseteq\mathcal P_1$.
Expanding the square gives $v(P)=\int x^2\,P(dx)-\mu(P)^2$.
The laws $\tfrac12\delta_{-1}+\tfrac12\delta_1$ and $N(0,1)$ both have
mean $0$ and variance $1$, but assign probabilities $1$ and $0$ to
$\{-1,1\}$. Thus the pair $(\mu,v)$ does not determine an arbitrary law
in $\mathcal P_2$.
\end{example}

\subsection{Observational equivalence and factorization}
\label{subsec:observational-equivalence}

\begin{definition}[Observational equivalence]
\label{def:observational-equivalence}
\indexnotation{\theta\sim\vartheta,\ [\theta]}{Observational equivalence and the equivalence class of a label.}
Two labels $\theta,\vartheta\in\Theta$ are \emph{observationally equivalent}
if
\[
\theta\sim\vartheta\quad\Longleftrightarrow\quad P_\theta=P_\vartheta.
\]
Write $[\theta]=\{\vartheta\in\Theta:\vartheta\sim\theta\}$ and
$\Theta/{\sim}=\{[\theta]:\theta\in\Theta\}$.
\end{definition}

Reflexivity, symmetry, and transitivity of $\sim$ follow from equality of laws.
Thus $\Theta/{\sim}$ is the set of equivalence classes.

\begin{definition}[Identifiability]
\label{def:identifiability}
A target $g:\Theta\to S$ is \emph{identifiable} in $(P_\theta)_{\theta\in\Theta}$
if
\[
P_\theta=P_\vartheta\quad\Longrightarrow\quad g(\theta)=g(\vartheta)
\qquad(\theta,\vartheta\in\Theta).
\]
The parametrization is \emph{identifiable} if $\pi$ is injective, equivalently,
if the target $g(\theta)=\theta$ is identifiable.
\end{definition}

No measurable structure on $\Theta$, $\mathcal P$, $S$, or $\Theta/{\sim}$
is imposed here.

\begin{theorem}[Factorization through the law]
\label{thm:identifiable-factorization}
A target $g:\Theta\to S$ is identifiable if and only if there is a unique
statistical functional $T:\mathcal P\to S$ such that
\begin{equation}
g=T\circ\pi,\qquad\text{that is, }g(\theta)=T(P_\theta)
\quad(\theta\in\Theta).
\label{eq:identifiable-factorization}
\end{equation}
Moreover, the map
\[
\bar\pi:\Theta/{\sim}\longrightarrow\mathcal P,\qquad
\bar\pi([\theta])=P_\theta,
\]
is a well-defined bijection. Thus $[\theta]\mapsto P_\theta$ is an
identifiable parametrization of the same model.
\end{theorem}
\begin{proof}
If $g=T\circ\pi$, then $P_\theta=P_\vartheta$ implies
$g(\theta)=T(P_\theta)=T(P_\vartheta)=g(\vartheta)$.
Conversely, suppose $g$ is identifiable. For $P\in\mathcal P$, define
$T(P)=g(\theta)$ using any $\theta$ with $P_\theta=P$.
Surjectivity supplies such a $\theta$; identifiability makes its choice
irrelevant. The resulting $T$ satisfies \eqref{eq:identifiable-factorization},
and this equation prescribes its value at every $P\in\mathcal P$, proving
uniqueness.

Finally, $[\theta]=[\vartheta]$ if and only if $P_\theta=P_\vartheta$.
Hence $\bar\pi$ is well-defined and injective.
Surjectivity of $\pi$ implies surjectivity of $\bar\pi$.
\end{proof}

\begin{example}[Identifiable targets under redundant labels]
\label{ex:identifiable-bernoulli-targets}
For $q\in[-1,1]$, write $Q_q=P_{q^2}$ for the Bernoulli family of
Example~\ref{ex:redundant-labels}. Then
\[
Q_q=Q_r\iff q^2=r^2\iff r\in\{q,-q\},
\qquad [q]=\{q,-q\}.
\]
The class $[0]$ is a singleton.
A target $g:[-1,1]\to S$ is identifiable exactly when
$g(q)=g(-q)$ for every $q$; its functional is
\[
T(P_p)=g(\sqrt p),\qquad p\in[0,1].
\]
In particular, $q^2$ and $|q|$ are identifiable, with respective functionals
$P\mapsto P(\{1\})$ and $P\mapsto\sqrt{P(\{1\})}$.
The target $g(q)=q$ is not identifiable.
\end{example}

\begin{example}[Identifying a uniform endpoint]
\label{ex:identifiable-uniform-endpoint}
For $U_\theta=\operatorname{Unif}[0,\theta]$, $\theta>0$, as in
Example~\ref{ex:uniform-endpoint},
\[
\mu(U_\theta)=\frac1\theta\int_0^\theta x\,dx=\frac\theta2.
\]
Consequently, $U_\theta=U_\vartheta$ implies $\theta=\vartheta$.
The parametrization is identifiable, and its inverse functional is
$T(U)=2\mu(U)$ on $\mathcal U$.
\end{example}

\begin{example}[Mixture labels and an identifiable mean]
\label{ex:mixture-labels}
Let $\gamma_u=N(u,1)$. The two-component Gaussian mixture family
\[
M_{\alpha,u,w}=\alpha\gamma_u+(1-\alpha)\gamma_w,\qquad
(\alpha,u,w)\in[0,1]\times\R^2,
\]
satisfies
\[
M_{\alpha,u,w}=M_{1-\alpha,w,u},\qquad
M_{\alpha,u,u}=\gamma_u.
\]
For instance, $(1/3,0,1)$ and $(2/3,1,0)$ are distinct equivalent labels.
Even on the restricted parameter set $u\leq w$, the labels $(\alpha,u,u)$
give the same law for every $\alpha\in[0,1]$.
The full parameter triple is therefore not identifiable.
Every mixture has finite first moment, and linearity of integration gives
\[
\mu(M_{\alpha,u,w})=\alpha u+(1-\alpha)w.
\]
Thus the target $(\alpha,u,w)\mapsto\alpha u+(1-\alpha)w$ is identifiable.
\end{example}

\subsection{Products and observation schemes}
\label{subsec:identifiability-products}

\begin{proposition}[Identifiability in product experiments]
\label{prop:product-identifiability}
For $1\leq i\leq n$, $n\geq1$, let
$\mathcal E_i=(\mathsf X_i,\mathcal A_i,(P_{\theta,i})_{\theta\in\Theta})$
be statistical experiments.
The product parametrization
$Q_\theta=\bigotimes_{i=1}^nP_{\theta,i}$ is identifiable if and only if
\[
\text{for every }\theta\ne\vartheta
\text{ there is an }i\text{ with }P_{\theta,i}\ne P_{\vartheta,i}.
\]
For iid repetition, a target $g:\Theta\to S$ is identifiable under
$(P_\theta)_{\theta\in\Theta}$ if and only if it is identifiable under
$(P_\theta^{\otimes n})_{\theta\in\Theta}$.
\end{proposition}
\begin{proof}
Proposition~\ref{prop:product-law-equality} gives
\[
Q_\theta=Q_\vartheta
\iff P_{\theta,i}=P_{\vartheta,i}\text{ for every }i.
\]
Thus no two distinct labels have the same product law exactly under the
stated separation condition.
For iid repetition, the same proposition gives
$P_\theta^{\otimes n}=P_\vartheta^{\otimes n}\iff P_\theta=P_\vartheta$,
so the condition in Definition~\ref{def:identifiability} is unchanged.
\end{proof}

\begin{example}[Complementary measurements]
\label{ex:complementary-identification}
Let the parameter be $(u,v)\in\R^2$. The first experiment observes a law
$\gamma_u=N(u,1)$; the second observes $\gamma_v=N(v,1)$.
As parametrizations indexed by $\R^2$, the first does not distinguish
different $v$, and the second does not distinguish different $u$.
For independent joint observation, however,
\[
Q_{u,v}=\gamma_u\otimes\gamma_v,\qquad
Q_{u,v}=Q_{u',v'}
\iff \gamma_u=\gamma_{u'}\text{ and }\gamma_v=\gamma_{v'}
\iff (u,v)=(u',v').
\]
The last equivalence follows by taking means, with the converse immediate
from the formulas for the laws. The product parametrization is identifiable.
\end{example}

\subsection{Limits of recovery from data}
\label{subsec:identifiability-recovery}

\begin{proposition}[An obstruction from indistinguishable laws]
\label{prop:indistinguishable-error}
Let $g:\Theta\to\R$ and let
$t:(\mathsf X,\mathcal A)\to(\R,\mathcal B(\R))$ be measurable.
If $P_\theta=P_\vartheta$, then
\begin{equation}
\max\left\{
\mathbb E_{P_\theta}\!\left[(t(X)-g(\theta))^2\right],
\mathbb E_{P_\vartheta}\!\left[(t(X)-g(\vartheta))^2\right]
\right\}
\geq \frac{(g(\theta)-g(\vartheta))^2}{4}.
\label{eq:indistinguishable-square-error}
\end{equation}
The expectations are allowed to be infinite.
Consequently, if there is a measurable $t$ satisfying
\[
P_\theta\bigl(t(X)=g(\theta)\bigr)=1\qquad\text{for every }\theta\in\Theta,
\]
then $g$ is identifiable.
\end{proposition}
\begin{proof}
Put $a=g(\theta)$, $b=g(\vartheta)$, and $Q=P_\theta=P_\vartheta$.
For every $x\in\mathsf X$,
\[
(t(x)-a)^2+(t(x)-b)^2
=2\left(t(x)-\frac{a+b}{2}\right)^2+\frac{(a-b)^2}{2}.
\]
Integration under $Q$ gives a sum of squared errors at least $(a-b)^2/2$;
their maximum is at least half this sum.
This proves \eqref{eq:indistinguishable-square-error}, including infinite
expectations. If $t$ recovers both targets almost surely, both errors are
zero, so $a=b$. Applying this to every equivalent pair proves the last claim.
\end{proof}

\begin{example}[A sharp obstruction at every iid sample size]
\label{ex:signed-bernoulli-error}
In Example~\ref{ex:identifiable-bernoulli-targets}, observe
$n\geq1$ iid variables with joint law $Q_q^{\otimes n}$ and target $g(q)=q$.
For $r\in(0,1]$, the labels $r$ and $-r$ give the same joint law.
For every measurable $t_n:\{0,1\}^n\to\R$,
Proposition~\ref{prop:indistinguishable-error} yields
\[
\max\left\{
\mathbb E_{Q_r^{\otimes n}}\!\left[(t_n-r)^2\right],
\mathbb E_{Q_{-r}^{\otimes n}}\!\left[(t_n+r)^2\right]
\right\}\geq r^2.
\]
The constant map $t_n=0$ has squared error $r^2$ under both labels.
The bound is sharp for this pair, independently of $n$.
\end{example}

\begin{example}[Identifiability without exact finite-sample recovery]
\label{ex:bernoulli-no-exact-recovery}
The ordinary Bernoulli parameter $p\in[0,1]$ is identifiable because
$P_p(\{1\})=p$. It remains identifiable under $n\geq1$ iid observations.
For $0<p<1$, every record $x\in\{0,1\}^n$ has strictly positive probability:
\[
P_p^{\otimes n}(\{x\})=p^k(1-p)^{n-k}>0,
\qquad k=\sum_{i=1}^n x_i.
\]
Hence $P_p^{\otimes n}(t_n=p)=1$ forces $t_n(x)=p$ at every record.
A single $t_n$ cannot satisfy this for two distinct $p$ in $(0,1)$.
Thus no map recovers $p$ almost surely for every parameter value.
Nevertheless, \eqref{eq:iid-mean-square} gives
\[
\mathbb E_{P_p^{\otimes n}}\!\left[(\bar X_n-p)^2\right]
=\frac{p(1-p)}n.
\]
Identifiability is a property of the map from labels to laws.
\end{example}

\subsection{Exercises}
\label{subsec:worked-examples-1-3}

\begin{practiceexercise}[An additive normal model]
\label{wex:additive-mean-target}
Under $N(a+b,1)$, characterize when labels $(a,b)$ give the same law. Decide
whether $a$, $b$, and $a+b$ are identifiable and express any identifiable
target as a functional of the law.
\end{practiceexercise}

\begin{practiceexercise}[Recording a total count]
\label{wex:poisson-total-identification}
Independent Poisson counts have rates $a,b>0$ and only their sum is
recorded. Compute its law by convolution. Which of the individual rates and
their sum are identifiable?
\end{practiceexercise}

\begin{practiceexercise}[The sign of a scale]
\label{wex:signed-normal-scale}
For $N(m,s^2)$ with $s\ne0$, characterize observational equivalence and
decide whether $(m,s)$ and $(m,|s|)$ are identifiable.
\end{practiceexercise}

\begin{practiceexercise}[An atom in a mixture]
\label{wex:atom-mixture-weight}
For $P_\alpha=\alpha\delta_0+(1-\alpha)N(0,1)$, find a functional of
$P_\alpha$ recovering $\alpha$ for all $\alpha\in[0,1]$.
\end{practiceexercise}

\begin{practiceexercise}[A binary probability label]
\label{wex:arctangent-link}
A Bernoulli observation has success probability $1/2+\arctan(\theta)/\pi$.
Determine whether the real label is identifiable and, if so, express it in
terms of the law.
\end{practiceexercise}

\begin{practiceexercise}[A single bit from a location family]
\label{wex:threshold-location}
For $Z\sim N(\theta,1)$, record $\mathbf1_{\{Z>0\}}$. Express its success
probability via the standard normal distribution function and decide whether
it identifies $\theta$.
\end{practiceexercise}

\begin{practiceexercise}[Observing an absolute value]
\label{wex:folded-normal-sign}
For $Z\sim N(\theta,1)$, characterize which labels give the same law of
$|Z|$. Consider reflection and the second moment.
\end{practiceexercise}

\begin{practiceexercise}[A mixture of binary laws]
\label{wex:binary-mixture-confounding}
Compute the equal mixture of $P_u$ and $P_v$. Are the unordered component
probabilities identifiable? Find an identifiable scalar function of $(u,v)$.
\end{practiceexercise}

\begin{practiceexercise}[Sorted normal measurements]
\label{wex:unordered-normal-pair}
Sort two independent measurements $U\sim N(a,1)$ and $V\sim N(b,1)$. Decide whether the ordered and unordered pairs of means are
identifiable. Consider the sum and product of the records.
\end{practiceexercise}

\begin{practiceexercise}[A bit with known reporting noise]
\label{wex:known-flip-identification}
Flip a Bernoulli bit of probability $p$ independently with known probability
$\epsilon$. Compute the reported success probability and characterize the
$\epsilon$ for which $p$ is identifiable.
\end{practiceexercise}

\begin{practiceexercise}[A bit with unknown reporting noise]
\label{wex:unknown-flip-confounding}
In the bit-flip setup of Practice
exercise~\ref{wex:known-flip-identification}, take $p\in[0,1]$ and
$\epsilon\in[0,1/4]$ both unknown. Find two labels with the same reported
law but different $p$, and give an identifiable function of the pair.
\end{practiceexercise}

\begin{practiceexercise}[A uniform endpoint after capping]
\label{wex:censored-uniform-identification}
For $U$ uniform on $[0,\theta]$, report $Y=\min(U,c)$ with known $c>0$.
Compute $\mathbb E Y$ for $\theta\le c$ and $\theta>c$. Use it to decide
identifiability.
\end{practiceexercise}

\begin{practiceexercise}[A sign measurement of a normal scale]
\label{wex:sign-hides-variance}
For $Z\sim N(0,\sigma^2)$ with $\sigma>0$, only $\mathbf1_{\{Z>0\}}$ is
recorded. Find the equivalence classes and justify them.
\end{practiceexercise}

\begin{practiceexercise}[Identification with erasures]
\label{wex:erasure-identification}
At binary label $\theta$, report $\theta$ with probability $3/4$ and a
common erasure otherwise. Decide whether the label is identifiable and
whether one rule can recover it almost surely at both labels.
\end{practiceexercise}

\begin{practiceexercise}[Two measurements of two labels]
\label{wex:two-rate-contrasts}
For $a,b\in[0,1/4]$, record independent bits with probabilities $a+b$ and
$1/2+a-b$. Give recovery formulas from the joint law. Does the first bit
alone identify both labels?
\end{practiceexercise}

\begin{practiceexercise}[A median as a law functional]
\label{wex:lower-binary-median}
The lower median is the least outcome with cumulative mass at least $1/2$.
Compute it for $P_p$ on $\{0,1\}$ and give distinct laws with the same
value.
\end{practiceexercise}

\begin{practiceexercise}[Targets that distinguish laws]
\label{wex:translation-invariant-spread}
For $P_{a,b}=N(a+b,4)$, compare the targets $a$, $a+b$, and the variance.
Which are identifiable? Which separate every pair of distinct observed laws?
\end{practiceexercise}

\begin{practiceexercise}[Exact recovery on overlapping intervals]
\label{wex:nested-uniform-exact}
For $\operatorname{Unif}[0,1]$ and $\operatorname{Unif}[0,2]$, prove
identifiability of the endpoint and rule out exact recovery at both labels
from one record.
\end{practiceexercise}

\begin{practiceexercise}[Squared error at equivalent labels]
\label{wex:additive-mean-bound}
Under $N(a+b,1)$ compare labels $(0,1)$ and $(2,-1)$ for target $a$. Apply
Proposition~\ref{prop:indistinguishable-error} and find an action attaining
its lower bound.
\end{practiceexercise}

\begin{practiceexercise}[Identifying a sampling scheme]
\label{wex:diagonal-identifies-scheme}
In Example~\ref{ex:dependence}, take the target to be independent versus
copied draws. Compare identification and exact recovery from the pair and
from its first coordinate alone.
\end{practiceexercise}

\begin{practiceexercise}[Recovering a function of a label]
\label{wex:squared-record-exact}
The deterministic record is $x=\theta^2$. Compare identifiability and exact
recovery of $\theta$ and $|\theta|$ on $\R$.
\end{practiceexercise}

\begin{exercise}[Reparametrization and quotient labels]
\label{exer:identifiability-reparametrization}
Let $h:\Lambda\to\Theta$ be a reparametrization, with
$\widetilde P_\lambda=P_{h(\lambda)}$.
Prove that $g:\Theta\to S$ is identifiable in $(P_\theta)$ if and only if
$g\circ h$ is identifiable in $(\widetilde P_\lambda)$.
Show that $k:\Theta\to\Theta/{\sim}$, $k(\theta)=[\theta]$, is bijective
exactly when $\pi$ is injective. In that case, identify the bijection
required by Definition~\ref{def:reparametrization} to use quotient labels.
\end{exercise}

\begin{exercise}[Identifiability depends on the design]
\label{exer:fixed-design-identification}
For fixed $z_1,\ldots,z_n\in\R$, consider the family
$Q_\theta=\bigotimes_{i=1}^nN(\theta z_i,1)$, $\theta\in\R$, from
Example~\ref{ex:fixed-design-gaussian}.
Prove that it is identifiable if and only if $\sum_i z_i^2>0$.
In that case, construct a functional recovering $\theta$ from the coordinate
means of $Q_\theta$. Determine the equivalence classes when every $z_i=0$.
\end{exercise}

\begin{exercise}[Population labels and permutations]
\label{exer:population-identification}
Fix $N\geq2$ and $1\leq n\leq N$. Let $Q_a$ be the law of sampling without
replacement from $a\in\R^N$ in
Example~\ref{ex:finite-population-variance}.
For $a,b\in\R^N$, prove
\[
Q_a=Q_b
\iff \frac1N\sum_{r=1}^N\delta_{a_r}
     =\frac1N\sum_{r=1}^N\delta_{b_r}
\iff b\text{ is a permutation of }a.
\]
Which of the targets $a_1$, $N^{-1}\sum_r a_r$, and
$N^{-1}\sum_r(a_r-N^{-1}\sum_s a_s)^2$ are identifiable?
Include the case $n=N$.
\end{exercise}

\begin{exercise}[Mixture moments]
\label{exer:mixture-moments}
For the mixture in Example~\ref{ex:mixture-labels}, derive
\[
v(M_{\alpha,u,w})=1+\alpha(1-\alpha)(u-w)^2.
\]
Writing $\gamma_u=N(u,1)$, consider
\[
P=\tfrac12\gamma_{-1}+\tfrac12\gamma_1,\qquad
Q=\tfrac15\gamma_{-2}+\tfrac45\gamma_{1/2}.
\]
Show that $P$ and $Q$ have the same mean and variance, but different third
moments. Conclude that the first two moments do not identify every law
even within the two-component Gaussian mixture model.
\end{exercise}

\begin{exercise}[Exact recovery and singularity]
\label{exer:exact-recovery-singularity}
Let $P,Q$ be laws on $(\mathsf X,\mathcal A)$ and let $a\ne b$ be real.
Show that a measurable $t:\mathsf X\to\R$ with
$P(t=a)=Q(t=b)=1$ exists if and only if some $A\in\mathcal A$ satisfies
$P(A)=1$ and $Q(A)=0$.
Apply this criterion to $N(\mu,1)$ and $N(\nu,1)$ for $\mu\ne\nu$.
Compare the result with identifiability of the Gaussian location parameter.
\end{exercise}

\begin{exercise}[Vanishing error probabilities require identifiability]
\label{exer:consistency-requires-identifiability}
Let $g:\Theta\to\R$ and let $t_n:\mathsf X^n\to\R$ be measurable.
Suppose that, for every $\theta\in\Theta$ and every $\varepsilon>0$,
\[
P_\theta^{\otimes n}\bigl(|t_n-g(\theta)|>\varepsilon\bigr)
\longrightarrow0\qquad(n\to\infty).
\]
Prove that $g$ is identifiable. No moment assumption is imposed.
\end{exercise}

\clearpage
\section{Domination, densities, and likelihoods}
\label{sec:domination-likelihood}

\subsection{Common reference measures}
\label{subsec:domination}

\begin{definition}[Domination and density]
\label{def:dominated-experiment}
\indexnotation{\mu,\ f_\theta}{Dominating reference measure and a density of \(P_\theta\) with respect to it.}
An experiment on $(\mathsf X,\mathcal A)$ is \emph{dominated} if some
$\sigma$-finite measure $\mu$ on $(\mathsf X,\mathcal A)$ satisfies
$P_\theta\ll\mu$ for every $\theta\in\Theta$. Thus $\mu(A)=0$ implies
$P_\theta(A)=0$ for every $\theta$. A \emph{density} of $P_\theta$
with respect to $\mu$ is a nonnegative measurable function $f_\theta$ with
\[
P_\theta(A)=\int_A f_\theta\,d\mu
\qquad(A\in\mathcal A).
\]
We use versions that are finite everywhere.
\end{definition}

\begin{remark}[Assumed probability facts]
For a finite measure $P$ with $P\ll\mu$ and $\mu$ $\sigma$-finite,
the Radon--Nikodym theorem gives a density $dP/d\mu$, unique $\mu$-almost
everywhere. Its integral is $P(\mathsf X)<\infty$, so it can be set to zero
where it is infinite. If $f$ is a density of $P$ with respect to $\mu$,
then $\int u\,dP=\int uf\,d\mu$ for every nonnegative measurable $u$,
by approximation with simple functions. We also use the existence of
finite products of $\sigma$-finite measures and Tonelli's theorem for them.
\end{remark}

\begin{example}[A mixture of an atom and a continuous law]
\label{ex:mixed-reference}
On $[0,1]$ with its Borel sets, let $\lambda$ be Lebesgue measure restricted
to $[0,1]$. The laws $P_a=a\delta_0+(1-a)\lambda$, $0\leq a\leq1$, are
dominated by $\mu=\delta_0+\lambda$. A density is
\[
f_a(0)=a,\qquad f_a(x)=1-a\quad(0<x\leq1).
\]
Indeed, its integral on a Borel set $A$ is
$a\mathbf 1_A(0)+(1-a)\lambda(A)$. The formula includes $a=0$ and $a=1$.
Lebesgue measure alone does not dominate the family, since $\delta_0$ gives
mass one to the Lebesgue-null set $\{0\}$.
\end{example}

\begin{example}[A nondominated family]
\label{ex:nondominated-dirac}
The family $(\delta_\theta)_{\theta\in\R}$ on $(\R,\mathcal B(\R))$ has
no $\sigma$-finite dominating measure. If $\mu$ dominated it, each singleton
would have positive $\mu$-mass. Write $\R=\bigcup_{j\geq1}E_j$ with
$\mu(E_j)<\infty$. For any $j,m\geq1$, only finitely many $x\in E_j$
can satisfy $\mu(\{x\})\geq1/m$: the mass of any $k$ distinct such points
is at least $k/m$. The singletons of positive mass would therefore be
countable, a contradiction.
\end{example}

\subsection{Likelihood and density versions}
\label{subsec:likelihood-versions}

\begin{definition}[Likelihood and log likelihood]
\label{def:likelihood}
\indexnotation{L_x(\theta),\ \ell_x(\theta)}{Likelihood and log likelihood at record \(x\).}
Fix a common dominating measure and choose one density $f_\theta$ for each
$P_\theta$. For a fixed record $x\in\mathsf X$, the \emph{likelihood} and
\emph{log likelihood} are functions of $\theta$:
\[
L_x(\theta)=f_\theta(x),\qquad
\ell_x(\theta)=\log L_x(\theta),\qquad \log 0=-\infty.
\]
\end{definition}

No measure on $\Theta$ is needed. The normalization
$\int f_\theta\,d\mu=1$ integrates over records for fixed $\theta$.
No joint measurability in $(\theta,x)$ is imposed here.

\begin{proposition}[Change of reference measure]
\label{prop:likelihood-reference}
Let $\mu$ and $\nu$ be two $\sigma$-finite measures dominating the same
experiment. There exist density choices for both measures, a measurable set
$B$ with $P_\theta(B)=1$ for every $\theta$, and a positive finite function
$c$ on $B$ such that
\[
L_x^\nu(\theta)=c(x)L_x^\mu(\theta)
\qquad(x\in B,\ \theta\in\Theta).
\]
On $B$ the two likelihoods have the same maximizing set. Their log
likelihoods differ by $\log c(x)$.
\end{proposition}
\begin{proof}
Set $\lambda=\mu+\nu$, which is $\sigma$-finite by intersecting the
finite-measure covers for $\mu$ and $\nu$. Write
\[
h=\frac{d\mu}{d\lambda},\qquad k=\frac{d\nu}{d\lambda},\qquad
r_\theta=\frac{dP_\theta}{d\lambda},
\]
using finite nonnegative versions. For $h$ and $k$, which need not be
integrable on the whole space, the Radon--Nikodym theorem applies on a
disjoint countable partition into sets of finite $\lambda$-measure;
the resulting densities can be joined on the partition.

Since $\mu(\{h=0\})=0$ and $\nu(\{k=0\})=0$, domination gives
$P_\theta(B)=1$ for $B=\{h>0,\ k>0\}$. Define $f_\theta^\mu=r_\theta/h$
on $\{h>0\}$ and zero elsewhere. For measurable $A$,
\[
\int_A f_\theta^\mu\,d\mu
=\int_{A\cap\{h>0\}}r_\theta\,d\lambda
=P_\theta(A).
\]
The integration formula for $h$ holds also for these $\sigma$-finite
measures by the same simple-function argument. Define $f_\theta^\nu$
analogously with $k$. On $B$, their relation has factor $c=h/k$.
A positive factor preserves the order of likelihood values; taking logs
adds $\log c$, including when the value is $-\infty$.
\end{proof}

\begin{proposition}[Simultaneous agreement for countably many labels]
\label{prop:countable-density-agreement}
Suppose $\Theta$ is countable. If $(f_\theta)$ and $(g_\theta)$ are two
density choices with respect to the same $\mu$, there is a single
$\mu$-null set $N$ such that
\[
f_\theta(x)=g_\theta(x)\qquad(x\notin N,\ \theta\in\Theta).
\]
In particular, the two likelihood functions agree outside $N$, and
$P_\theta(N)=0$ for every $\theta$.
\end{proposition}
\begin{proof}
By uniqueness of the density, each set $\{f_\theta\ne g_\theta\}$ is
$\mu$-null. Their countable union is a single $\mu$-null set.
\end{proof}

With arbitrary density choices, observationally equivalent labels have
densities equal almost everywhere. Choosing a density once per law makes
their likelihood values equal at every record.

\subsection{Likelihoods for independent observations}
\label{subsec:product-likelihood}

\begin{proposition}[Product densities]
\label{prop:product-densities}
For $1\leq i\leq n$, let $f_{\theta,i}$ be a chosen density of
$P_{\theta,i}$ with respect to a $\sigma$-finite measure $\mu_i$.
For the independent product experiment of
Definition~\ref{def:product-experiment}, the product law $Q_\theta$
has density $\prod_{i=1}^n f_{\theta,i}(x_i)$ with respect to
$\bigotimes_{i=1}^n\mu_i$. With this choice,
\[
L_x(\theta)=\prod_{i=1}^n f_{\theta,i}(x_i),\qquad
\ell_x(\theta)=\sum_{i=1}^n\log f_{\theta,i}(x_i).
\]
For iid repetition one may take $f_{\theta,i}=f_\theta$ and $\mu_i=\mu$
for all $i$.
\end{proposition}
\begin{proof}
For a measurable rectangle $\prod_i A_i$, Tonelli gives
\[
\int_{\prod_i A_i}\prod_{i=1}^n f_{\theta,i}(x_i)\,
 d\left(\bigotimes_{i=1}^n\mu_i\right)(x)
=\prod_{i=1}^n\int_{A_i}f_{\theta,i}\,d\mu_i
=\prod_{i=1}^n P_{\theta,i}(A_i).
\]
Taking $A_i=\mathsf X_i$ shows that the density defines a probability
measure. Uniqueness of the probability product then gives $Q_\theta$.
The log identity holds also if any factor is zero.
\end{proof}

\begin{example}[Bernoulli records and counts]
\label{ex:bernoulli-likelihood}
Use counting measure on $\{0,1\}^n$. For $n\geq1$ iid observations from
$P_p$, write $k=\sum_i x_i$. Then
\[
L_x(p)=p^k(1-p)^{n-k},\qquad 0\leq p\leq1,
\]
where a power with exponent zero is an empty product, equal to one.
If the record is the count $k$, its likelihood with respect to counting
measure on $\{0,\ldots,n\}$ is $\binom nk L_x(p)$: all $\binom nk$ binary
records with that count have the same mass. The factor depends on the
record alone.

If $0<k<n$, the endpoints give likelihood zero, and for $0<p<1$,
\[
\frac{d}{dp}\ell_x(p)=\frac{k-np}{p(1-p)}.
\]
The derivative changes sign from positive to negative at $p=k/n$.
For $k=0$ the likelihood decreases on $[0,1]$; for $k=n$ it increases.
Thus in all cases the unique maximizer is $k/n$, for both records.
\end{example}

\begin{example}[Gaussian location with two references]
\label{ex:gaussian-likelihood}
For $n\geq1$ iid observations from $N(\theta,1)$, $\theta\in\R$,
the Lebesgue density gives
\[
\ell_x(\theta)=-\frac n2\log(2\pi)-\frac12\sum_{i=1}^n(x_i-\theta)^2.
\]
Writing $\bar x=n^{-1}\sum_i x_i$ and using $\sum_i(x_i-\bar x)=0$ gives
\[
\sum_i(x_i-\theta)^2=\sum_i(x_i-\bar x)^2+n(\theta-\bar x)^2.
\]
The unique maximizer is $\bar x$. Alternatively, use $N(0,1)^{\otimes n}$
as the reference. Its Lebesgue density is everywhere positive, so division
by that density gives a valid likelihood. The log likelihood is now
$\theta\sum_i x_i-n\theta^2/2$; the difference from the former log likelihood
depends only on $x$.
\end{example}

\begin{example}[Copied Bernoulli observations]
\label{ex:copied-likelihood}
If $X_i=Z$ for all $i$ with $Z\sim P_p$, the likelihood of the all-one
record is $p$, the likelihood of the all-zero record is $1-p$, and all
other records have likelihood zero, with counting measure as reference.
The iid experiment has the same coordinate marginals, yet it assigns
likelihood $p^n$ to the all-one record. For $n>1$ and $0<p<1$, these values
differ. The likelihood uses the joint law of the record.
\end{example}

\subsection{A parameter-dependent range}
\label{subsec:endpoint-likelihood}

\begin{example}[Uniform endpoint and attainment]
\label{ex:uniform-likelihood-versions}
For the uniform laws of Example~\ref{ex:uniform-endpoint}, the functions
\[
f_\theta^{\mathrm{cl}}(x)=\theta^{-1}\mathbf 1_{[0,\theta]}(x),
\qquad
f_\theta^{\mathrm{op}}(x)=\theta^{-1}\mathbf 1_{[0,\theta)}(x)
\qquad(\theta>0)
\]
are versions of the same Lebesgue density for each $\theta$. Consider an
iid record with $x_i\geq0$ for all $i$ and $m=\max_i x_i>0$.
The product formulas give
\[
L_x^{\mathrm{cl}}(\theta)=\theta^{-n}\mathbf 1_{\{\theta\geq m\}},
\qquad
L_x^{\mathrm{op}}(\theta)=\theta^{-n}\mathbf 1_{\{\theta>m\}}.
\]
Their common supremum is $m^{-n}$. It is attained only at $\theta=m$ in
the first choice and is unattained in the second, since $\theta^{-n}$ is
strictly decreasing on $(0,\infty)$. The records considered have
probability one under every member of the iid experiment. For each fixed
$\theta$, the differing values occur on a null set, but its location
depends on $\theta$.
\end{example}

\subsection{Exercises}
\label{subsec:worked-examples-1-4}

\begin{practiceexercise}[Countably many point masses]
\label{wex:countable-dirac-domination}
Find a common sigma-finite reference for the point masses on the positive
integers and give their densities. Why does the same argument not apply to
every point mass on $\R$?
\end{practiceexercise}

\begin{practiceexercise}[A reference for a countable model]
\label{wex:weighted-reference-mixture}
For laws $P_1,P_2,\ldots$ on one space, verify that
$\mu=\sum_{j\ge1}2^{-j}P_j$ is a probability law that dominates all the
$P_j$.
\end{practiceexercise}

\begin{practiceexercise}[Two atoms and a continuous tail]
\label{wex:atoms-and-exponential-reference}
For $a,b\ge0$, $a+b\le1$, take
$P=a\delta_0+b\delta_1+(1-a-b)\operatorname{Exp}(1)$. Find a density
relative to Lebesgue measure on $[0,\infty)$ plus $\delta_0+\delta_1$, and
integrate it on an arbitrary Borel set.
\end{practiceexercise}

\begin{practiceexercise}[Capping a uniform measurement]
\label{wex:capped-uniform-likelihood}
Observe $\min(U,c)$ for $U\sim\operatorname{Unif}[0,\theta]$ and known
$c>0$. Find the likelihood at and below $c$ relative to Lebesgue measure on
$[0,c]$ plus $\delta_c$. Verify normalization at $\theta=2c$.
\end{practiceexercise}

\begin{practiceexercise}[Rounding an exponential measurement]
\label{wex:rounded-exponential-masses}
For $T\sim\operatorname{Exp}(\lambda)$, find the counting density of
$\lfloor T\rfloor$. Sum its masses from $0$ to $k$ and verify the result
using the law of $T$.
\end{practiceexercise}

\begin{practiceexercise}[An exponential reference measure]
\label{wex:exponential-reference-change}
Find the density of a rate-$\lambda$ exponential law relative to the
rate-one law. Compare the two likelihoods relative to this reference and to
Lebesgue measure at a fixed record.
\end{practiceexercise}

\begin{practiceexercise}[Reference measures outside the support]
\label{wex:non-equivalent-references}
Put Bernoulli laws on $\{0,1,2,3\}$ with no mass at $2,3$. Check domination
for the references $\delta_0+\delta_1+\delta_2$ and
$2\delta_0+2\delta_1+\delta_3$. Relate the two likelihoods on a common
full-mass set.
\end{practiceexercise}

\begin{practiceexercise}[A record outside the model's support]
\label{wex:zero-likelihood-entire-model}
Use counting measure on $\{0,1,2\}$ for a family supported on $\{0,1\}$. At
record $2$, give the likelihood and log likelihood as functions of the label
and determine their maximizing sets.
\end{practiceexercise}

\begin{practiceexercise}[Order under unequal coin laws]
\label{wex:unequal-coin-likelihood-order}
Independent bits have probabilities $p$ and $p/2$, $0\le p\le1$. Find and
maximize the likelihood for $(1,0)$ and $(0,1)$ separately. Can the count
alone recover both answers?
\end{practiceexercise}

\begin{practiceexercise}[A Poisson likelihood]
\label{wex:poisson-rate-likelihood}
Derive and maximize the iid Poisson likelihood in its rate. Treat positive
and zero count totals separately, and distinguish whether rate zero is
admitted.
\end{practiceexercise}

\begin{practiceexercise}[An exponential waiting-time record]
\label{wex:exponential-rate-likelihood}
For iid exponential observations with rate $\lambda>0$, derive the log
likelihood and find its maximizer when the total time is positive.
\end{practiceexercise}

\begin{practiceexercise}[A normal scale with known mean]
\label{wex:normal-variance-likelihood}
For iid $N(0,v)$ observations, $v>0$, derive and maximize the log likelihood
in $v$. Treat the all-zero record separately.
\end{practiceexercise}

\begin{practiceexercise}[A moving uniform window]
\label{wex:moving-uniform-window}
For iid draws from $\operatorname{Unif}[\theta,\theta+1]$, use
closed-interval densities. Find the maximizing set from the minimum and
maximum, including a sample range greater than one.
\end{practiceexercise}

\begin{practiceexercise}[Likelihood with unequal variances]
\label{wex:precision-weighted-location}
For independent $X_1\sim N(\theta,1)$ and $X_2\sim N(\theta,4)$, find the
likelihood maximizer. Explain the relative weights of the two observations.
\end{practiceexercise}

\begin{practiceexercise}[A shifted exponential endpoint]
\label{wex:shifted-exponential-likelihood}
Let $X_i=\theta+E_i$ with iid rate-one exponentials. Use densities including
the endpoint, derive the joint likelihood and its positive set, and find the
maximizing label.
\end{practiceexercise}

\begin{practiceexercise}[Dependence and likelihood]
\label{wex:dependent-binary-joint-likelihood}
Start with a fair bit and repeat it with probability $p$, otherwise flip it.
Write the likelihood for each record. Compare it with the product of the
marginal densities as a function of $p$.
\end{practiceexercise}

\begin{practiceexercise}[Changing the units of a record]
\label{wex:exponential-units}
For $X\sim\operatorname{Exp}(\lambda)$ and known $c>0$, find the density of
$Y=cX$. Compare the iid likelihoods of transformed and original records.
\end{practiceexercise}

\begin{practiceexercise}[A Laplace location likelihood]
\label{wex:two-sided-location-likelihood}
For density $\tfrac12e^{-|x-\theta|}$ and iid record $(0,2,5)$, find the
likelihood maximizer by computing the slopes of the negative log likelihood
between observations.
\end{practiceexercise}

\begin{exercise}[Identifiability from densities]
\label{exer:identifiability-densities}
Show that $P_\theta=P_\vartheta$ if and only if
$f_\theta=f_\vartheta$ $\mu$-almost everywhere. Express identifiability of
the parametrization in terms of the sets $\{f_\theta\ne f_\vartheta\}$.
\end{exercise}

\begin{exercise}[Weighted counting measures]
\label{exer:weighted-counting-likelihood}
On a finite observation space with all subsets measurable, let $\mu$ be
counting measure and let $\nu(\{x\})=w_x>0$ for each $x$. Compute the
$\nu$-density in terms of the
$\mu$-density. For $n$ iid observations, find the factor relating the
likelihoods relative to $\mu^{\otimes n}$ and $\nu^{\otimes n}$.
\end{exercise}

\begin{exercise}[Gaussian measurements with fixed design]
\label{exer:gaussian-design-likelihood}
In Example~\ref{ex:fixed-design-gaussian}, derive the joint log likelihood
with respect to Lebesgue measure. If $\sum_i z_i^2>0$, compute its unique
maximizer by completing the square. Determine the maximizing set if
$z_i=0$ for every $i$.
\end{exercise}

\begin{exercise}[A reference for two different sampling schemes]
\label{exer:gaussian-diagonal-reference}
On $\R^2$ let $P_\theta=N(\theta,1)^{\otimes2}$ and let $Q_\theta$ be
the law of $(Z,Z)$ for $Z\sim N(\theta,1)$. Put $\mu=\lambda_2+Q_0$,
where $\lambda_2$ is planar Lebesgue measure. Show that $\mu$ dominates
both families. Give densities of $P_\theta$ and $Q_\theta$ by treating
the diagonal and its complement separately, and verify the formulas by
integration.
\end{exercise}

\begin{exercise}[The uniform sample maximum]
\label{exer:uniform-max-likelihood}
For $n\geq1$ iid uniform observations on $[0,\theta]$, compute the
distribution function of $M=\max_i X_i$. Find a Lebesgue density for
$M$ with the endpoint $\theta$ included. At a record with maximum $m>0$
and nonnegative coordinates, compare the resulting likelihood of $m$
with $L_x^{\mathrm{cl}}$ in Example~\ref{ex:uniform-likelihood-versions}.
\end{exercise}

\begin{exercise}[Joint measurability does not give simultaneous agreement]
\label{exer:diagonal-density-versions}
Let $\Theta=\mathsf X=[0,1]$, with Borel sets, and let every $P_\theta$
be the uniform law. Relative to Lebesgue measure, consider
\[
f_\theta(x)=1,\qquad g_\theta(x)=1+\mathbf 1_{\{x=\theta\}}.
\]
Verify that both choices are jointly measurable densities of this family.
Find their maximizing sets at each $x$ and show that no single null set
can make the likelihood functions agree outside it.
\end{exercise}

\chapter{Loss, risk, and decision rules}
\label{chap:statistical-decisions}
\notationgroup{Decision theory}
\section{Actions and decision rules}
\label{sec:actions-decision-rules}

Fix a statistical experiment
$\mathcal E=(\mathsf X,\mathcal A,(P_\theta)_{\theta\in\Theta})$.

\subsection{Actions and rules}
\label{subsec:actions-rules}

\begin{definition}[Action space and deterministic rule]
\label{def:deterministic-rule}
\indexnotation{(\mathsf D,\mathcal D),\ d}{Action space and deterministic rule.}
An \emph{action space} is a nonempty measurable space
$(\mathsf D,\mathcal D)$. A \emph{deterministic decision rule} is a
measurable map $d:(\mathsf X,\mathcal A)\to(\mathsf D,\mathcal D)$.
At record $x$ it takes action $d(x)$.
\end{definition}

\begin{example}[Two choices of actions for Bernoulli data]
\label{ex:bernoulli-decision-actions}
For a binary record $(x_1,\ldots,x_n)$, $n\geq1$, put $\bar x=n^{-1}\sum_i x_i$.
An estimate of the Bernoulli probability can take values in $[0,1]$
with its Borel sets; one rule is $d(x)=\bar x$. A choice between two
decisions can instead take values in $\{0,1\}$ with all subsets measurable;
one rule is $d(x)=\mathbf 1_{\{\bar x>1/2\}}$. The observation law can be
the same in both problems.
\end{example}

\begin{definition}[Randomized rule]
\label{def:randomized-rule}
\indexnotation{K(x,\cdot),\ K_d}{Action law at \(x\) and the kernel of a deterministic rule.}
A \emph{probability kernel} $K$ from $(\mathsf X,\mathcal A)$ to
$(\mathsf D,\mathcal D)$ assigns a probability measure $K(x,\cdot)$ to
each $x\in\mathsf X$, with $x\mapsto K(x,B)$ measurable for each
$B\in\mathcal D$. A \emph{randomized decision rule} is such a kernel:
at record $x$, it uses the action law $K(x,\cdot)$.
\end{definition}

The rule is fixed across parameter values. For a deterministic rule $d$,
\[
K_d(x,B)=\delta_{d(x)}(B)=\mathbf 1_B(d(x))
\]
is a kernel, by measurability of $d$. In particular, the definition of a
randomized rule includes deterministic rules.

\subsection{The joint observation and action law}
\label{subsec:joint-action-law}

\begin{proposition}[The law induced by a kernel]
\label{prop:joint-action-law}
\indexnotation{P\ltimes K,\ PK}{Joint observation and action law, and its action marginal.}
Let $P$ be a probability law on $(\mathsf X,\mathcal A)$ and $K$ a
probability kernel to $(\mathsf D,\mathcal D)$. There is a unique probability
law $P\ltimes K$ on $(\mathsf X\times\mathsf D,\mathcal A\otimes\mathcal D)$
such that
\begin{equation}
(P\ltimes K)(A\times B)=\int_A K(x,B)\,P(dx)
\qquad(A\in\mathcal A,\ B\in\mathcal D).
\label{eq:joint-action-rectangles}
\end{equation}
For each nonnegative product-measurable $g$, the inner integral below is
measurable and
\begin{equation}
\int g\,d(P\ltimes K)
=\int_{\mathsf X}\left(\int_{\mathsf D} g(x,a)\,K(x,da)\right)P(dx).
\label{eq:kernel-iterated-integral}
\end{equation}
The observation marginal is $P$, and the action marginal, denoted by
$PK$, satisfies $(PK)(B)=\int K(x,B)\,P(dx)$.
For $K_d$ this action marginal is $P\circ d^{-1}$.
\end{proposition}
\begin{proof}
For $C\in\mathcal A\otimes\mathcal D$, put
$C_x=\{a:(x,a)\in C\}$. For fixed $x$, the class of sets $C$ with
$C_x\in\mathcal D$ is a sigma algebra containing the measurable rectangles.
The class of $C$ for which $x\mapsto K(x,C_x)$ is measurable contains the
rectangles and is closed under complements and countable disjoint unions.
For complements use $K(x,C_x^{\mathrm c})=1-K(x,C_x)$; for disjoint
unions use countable additivity of $K(x,\cdot)$.
The $\pi$--$\lambda$ theorem therefore gives measurability for every $C$.

Define $(P\ltimes K)(C)=\int K(x,C_x)\,P(dx)$. This has total mass one.
For disjoint sets $C_j$, their sections are disjoint for each $x$,
so countable additivity follows from monotone convergence. The rectangle
formula is immediate, and the rectangles determine a probability law
by the $\pi$--$\lambda$ theorem.

Equation~\eqref{eq:kernel-iterated-integral} now holds for indicators
by construction. Linearity extends it to nonnegative simple functions;
monotone convergence extends it to all nonnegative measurable $g$ and
gives measurability of the inner integral. The marginal formulas follow
by taking $B=\mathsf D$ or $A=\mathsf X$ in
\eqref{eq:joint-action-rectangles}. In particular,
$(PK_d)(B)=\int\mathbf 1_B(d(x))\,P(dx)=P(d^{-1}(B))$.
\end{proof}

\begin{example}[Randomization preserves indistinguishability]
\label{ex:randomized-indistinguishability}
If $P_\theta=P_\vartheta$, then $P_\theta\ltimes K=P_\vartheta\ltimes K$
for every fixed $K$, by \eqref{eq:joint-action-rectangles}.
In Example~\ref{ex:identifiable-bernoulli-targets}, the labels $r$ and
$-r$, $0<r\leq1$, have the same law $P_{r^2}$. A real-action kernel must
give the same action distribution at both labels, so it cannot take
action $r$ almost surely at the first and action $-r$ almost surely at
the second. The same argument holds with any finite iid repetition.
\end{example}

\subsection{Mixtures of rules}
\label{subsec:rule-mixtures}

\begin{proposition}[Mixing kernels]
\label{prop:mixture-rules}
If $K_0,K_1$ are kernels and $t:\mathsf X\to[0,1]$ is measurable, then
\[
K_t(x,B)=(1-t(x))K_0(x,B)+t(x)K_1(x,B)
\]
is a kernel. If $t$ is constant, equal to $s$, then for every $P$
\[
P\ltimes K_t=(1-s)(P\ltimes K_0)+s(P\ltimes K_1),\qquad
PK_t=(1-s)PK_0+sPK_1.
\]
\end{proposition}
\begin{proof}
For fixed $x$, $K_t(x,\cdot)$ is a convex combination of probability
measures. For fixed $B$, the formula is measurable in $x$.
When $t\equiv s$, linearity in \eqref{eq:joint-action-rectangles} gives
the equality of joint laws on rectangles, hence everywhere. Taking
action marginals gives the second equality.
\end{proof}

\begin{example}[Randomly flipping a binary observation]
\label{ex:binary-flip-rule}
For $X\sim P_p$, mix the identity rule $d_0(x)=x$ with the flip
$d_1(x)=1-x$, using $t\equiv1/4$. Then
\[
K_t(0,\{1\})=\tfrac14,\qquad K_t(1,\{1\})=\tfrac34.
\]
Under the joint law, the probability of observing zero and taking action
one is $(1-p)/4$. The probability of taking action one is
$(1-p)/4+3p/4=1/4+p/2$. No deterministic rule on the binary record can
give this action marginal for every $p$: at $p=0$ its action is constant,
whereas $P_0K_t$ gives both actions positive probability.
\end{example}

\begin{example}[A mixture and an average of actions]
\label{ex:mixing-versus-averaging-actions}
For the Borel action space $[0,1]$, an equal mixture of the rules $0$ and $1$
has action law $\tfrac12\delta_0+\tfrac12\delta_1$, for every observation law.
Averaging the two actions pointwise instead gives the constant rule $1/2$.
Both action laws have mean $1/2$; their second moments are respectively
$1/2$ and $1/4$.
\end{example}

\subsection{Uniform simulation for finite action spaces}
\label{subsec:finite-rule-simulation}

\begin{proposition}[Simulation by one uniform draw]
\label{prop:finite-kernel-simulation}
Let $\mathsf D=\{a_1,\ldots,a_m\}$, $m\geq1$, with all subsets measurable.
For each kernel $K$ to $\mathsf D$ there is a measurable map
$d:\mathsf X\times[0,1)\to\mathsf D$ with
\begin{equation}
K(x,B)=\lambda\{u:d(x,u)\in B\}
\qquad(x\in\mathsf X,\ B\subseteq\mathsf D),
\label{eq:uniform-kernel-simulation}
\end{equation}
where $\lambda$ is uniform probability on $[0,1)$ with its Borel sets.
Conversely, this formula defines a kernel for each such map $d$.
For every $P$, the law of $(X,d(X,U))$ under $P\otimes\lambda$ is
$P\ltimes K$.
\end{proposition}
\begin{proof}
Write $c_0(x)=0$ and $c_j(x)=\sum_{i=1}^jK(x,\{a_i\})$. These functions
are measurable and satisfy $0=c_0\leq\cdots\leq c_m=1$. Set $d(x,u)=a_j$
on the set
\[
\{(x,u):c_{j-1}(x)\leq u<c_j(x)\}.
\]
These measurable sets partition $\mathsf X\times[0,1)$, with empty
intervals when an action has probability zero. For each $x$, the interval
for $a_j$ has length $c_j(x)-c_{j-1}(x)=K(x,\{a_j\})$. Summing over actions
in $B$ gives \eqref{eq:uniform-kernel-simulation}.

Conversely, a measurable $d$ gives, for fixed $x$, the distribution
of $d(x,U)$ under $\lambda$. Its probability of $B$ is measurable in $x$
by the inner-integral assertion of Proposition~\ref{prop:joint-action-law}
applied to the constant kernel $x\mapsto\lambda$. Thus it is a kernel.
Under $P\otimes\lambda$, Tonelli gives
\[
(P\otimes\lambda)\{X\in A,\ d(X,U)\in B\}
=\int_A K(x,B)\,P(dx).
\]
The rectangle formula identifies the joint law.
\end{proof}

\begin{example}[Randomizing at a tie]
\label{ex:count-tie-rule}
Observe the count $S=X_1+X_2$ of two iid Bernoulli observations with
parameter $p$. Take action one at $S=2$, zero at $S=0$, and $\mathbf 1_{\{U<1/2\}}$ at
$S=1$, where $U$ is independent and uniform on $[0,1)$.
The probability of taking action one is
\[
p^2+\tfrac12\bigl(2p(1-p)\bigr)=p.
\]
A deterministic rule on the count with this property for all $p$ must take
action zero at $S=0$ and one at $S=2$, as $p=0$ and $p=1$ show. At $p=1/2$,
its probability of taking action one is then either $1/4$ or $3/4$, according
to the action at $S=1$. If the ordered record is available, the deterministic
rule $d(x_1,x_2)=x_1$ has action law $P_p$ for every $p$.
\end{example}

\subsection{Exercises}
\label{subsec:worked-examples-2-1}

\begin{practiceexercise}[A deterministic three-level rule]
\label{wex:three-level-rule}
Let $d(x)$ be $-1,0,1$ on $x<-1$, $-1\le x\le1$, and $x>1$, respectively.
Check measurability and give its action law at record $x=1$.
\end{practiceexercise}

\begin{practiceexercise}[A joint law from two kernel rows]
\label{wex:finite-kernel-joint}
Suppose $P(\{1\})=1/4$, $K(0,\{1\})=1/3$, and $K(1,\{1\})=2/3$. Compute the
joint masses at $(0,1)$ and $(1,1)$, the action-one probability, and the
probability the record was one given action one.
\end{practiceexercise}

\begin{practiceexercise}[Clipping a real action]
\label{wex:clipping-action-marginal}
Clip a uniform draw from $[-1,2]$ to the interval $[0,1]$.
Determine the atoms and continuous part of the resulting action law.
\end{practiceexercise}

\begin{practiceexercise}[A kernel constant in the record]
\label{wex:constant-kernel-independence}
Suppose $K(x,\cdot)=Q$ for every record. Use the rectangle formula to
identify the joint law and decide whether the record and action are
independent.
\end{practiceexercise}

\begin{practiceexercise}[Two rules on a fair bit]
\label{wex:opposite-rules-fair-bit}
For a fair observation $X$, compare the action marginals of $d_0(X)=X$ and
$d_1(X)=1-X$. Compare their joint probabilities of $\{X=A\}$, with $A$ the
action.
\end{practiceexercise}

\begin{practiceexercise}[Rules at an impossible record]
\label{wex:kernel-null-record}
All model laws on $\{0,1,2\}$ are supported on $\{0,1\}$. Two kernels agree
at $0,1$ but may differ at $2$. Compare their joint observation and action
laws at each label.
\end{practiceexercise}

\begin{practiceexercise}[Simulating a nonlinear kernel]
\label{wex:quadratic-action-probability}
For $x\in[0,1]$ let $K(x,\{1\})=x^2$. Implement the rule with an independent
uniform draw and find the marginal action-one probability for uniform $X$.
\end{practiceexercise}

\begin{practiceexercise}[An action from two uniform variables]
\label{wex:scaled-uniform-action}
Let $A=XU$ for independent uniform $X,U$ on $[0,1]$. Describe the kernel at
each $x$ including zero, justify measurability, and compute $P(A\le t)$ for
$0<t<1$.
\end{practiceexercise}

\begin{practiceexercise}[Adding noise to an action]
\label{wex:symmetric-action-noise}
At record $x$ take action $x+U$ with independent
$U\sim\operatorname{Unif}[-1,1]$. Find the probability of a nonnegative
action as a function of $x$.
\end{practiceexercise}

\begin{practiceexercise}[Selecting one of two measurements]
\label{wex:select-a-coordinate-kernel}
Independent bits $X_1,X_2$ have probabilities $1/4,3/4$. Select either
coordinate with probability $1/2$. Write the kernel and find the action
marginal.
\end{practiceexercise}

\begin{practiceexercise}[A record-dependent mixing weight]
\label{wex:gated-rule-mixture}
Let $X$ have masses $1/4,1/2,1/4$ at $-1,0,1$. Mix $d_0(x)=0$ and
$d_1(x)=x$, choosing $d_1$ with respective probabilities $1/4,1/2,3/4$. Find
the action law. Can a constant weight reproduce it?
\end{practiceexercise}

\begin{practiceexercise}[A tie between categories]
\label{wex:category-tie-kernel}
Choose uniformly among categories with maximal count. Give the action law
for $(a,b,a,b)$, on actions $a,b,c$. Compare with always breaking ties in
favor of $a$ when the labels $a,b$ are interchanged.
\end{practiceexercise}

\begin{practiceexercise}[The intervals of a simulator]
\label{wex:three-interval-simulation}
Implement probabilities $(1/6,1/3,1/2)$ on $(a,b,c)$ by successive half-open
intervals in $[0,1)$. What action does the value $1/2$ give? Does changing
one endpoint assignment change the action law?
\end{practiceexercise}

\begin{practiceexercise}[Two successive flips]
\label{wex:two-random-flips}
Flip a bit independently twice, each time with probability $1/4$. Find the
chance the final bit differs from the original and give the resulting
single-step kernel.
\end{practiceexercise}

\begin{practiceexercise}[Uniform draws and dependence]
\label{wex:dependent-auxiliary-draw}
For uniform $X$, compare $\mathbf1_{\{U<X\}}$ when $U$ is independent
uniform and when $U=X$. Why does a uniform marginal for $U$ not determine
the action law?
\end{practiceexercise}

\begin{exercise}[Same action law, different joint laws]
\label{exer:action-versus-joint-law}
Let $X\sim P_{1/2}$ and use binary actions. Compare $d(x)=x$ with the
kernel $K(x,\{1\})=1/2$ at both records. Compute their action marginals and
their joint probabilities of the event $\{X=D\}$, where $D$ is the action
coordinate.
\end{exercise}

\begin{exercise}[Equality of kernels under a fixed law]
\label{exer:joint-law-kernel-equality}
For a finite action space with all subsets measurable, prove that
$P\ltimes K=P\ltimes H$ if and only if a single $P$-null set $N$ satisfies
$K(x,\cdot)=H(x,\cdot)$ for every $x\notin N$. For the forward direction,
use \eqref{eq:joint-action-rectangles} with $B$ a singleton, and integrate
the difference over the sets where it is positive or negative.
\end{exercise}

\begin{exercise}[A weight that depends on the record]
\label{exer:record-dependent-rule-mixture}
For the identity and flip rules in Example~\ref{ex:binary-flip-rule},
use $t(x)=x$ instead of $t\equiv1/4$. Compute the action distribution
when $X\sim P_p$. Compare it with the distribution from the constant
weight $t\equiv\mathbb E_{P_p}[t(X)]$ for fixed $p\in(0,1)$.
\end{exercise}

\begin{exercise}[A squared-error bound for randomized rules]
\label{exer:randomized-recovery-bound}
Suppose $P_\theta=P_\vartheta$, let $g:\Theta\to\R$, and let $K$ be a
kernel to $\R$ with its Borel sets. Write
\[
e_\eta=\int_{\R}(a-g(\eta))^2\,(P_\eta K)(da).
\]
Prove that $\max\{e_\theta,e_\vartheta\}\geq
(g(\theta)-g(\vartheta))^2/4$, allowing infinite integrals. Find a constant
rule attaining the bound for these two labels.
\end{exercise}

\begin{exercise}[Three actions from one uniform draw]
\label{exer:three-action-simulation}
For $x\in[0,1]$, the respective probabilities of actions $0,1,2$ are
$(1-x)^2$, $2x(1-x)$, and $x^2$. Give an explicit uniform simulation as
in \eqref{eq:uniform-kernel-simulation}, including the endpoints $x=0,1$.
Compute the action marginal for a uniform observation on $[0,1]$.
\end{exercise}

\begin{exercise}[Countably many actions]
\label{exer:countable-action-simulation}
Take a countable action space $\{a_1,a_2,\ldots\}$ with all subsets
measurable. Extend Proposition~\ref{prop:finite-kernel-simulation} to this space. Show that every $u\in[0,1)$
belongs to one of the intervals in the proof, even if some probabilities
are zero, and verify the kernel identity for every subset of actions.
\end{exercise}

\clearpage
\section{Loss, risk, and comparison of procedures}
\label{sec:loss-risk}

\subsection{Loss and risk}
\label{subsec:loss-risk-definitions}

Fix an experiment $\mathcal E=(\mathsf X,\mathcal A,(P_\theta)_{\theta\in\Theta})$
and an action space $(\mathsf D,\mathcal D)$.

\begin{definition}[Loss]
\label{def:decision-loss}
\indexnotation{L(\theta,a)}{Loss of action \(a\) at parameter \(\theta\).}
A \emph{loss function} is a map $L:\Theta\times\mathsf D\to[0,\infty]$
such that $a\mapsto L(\theta,a)$ is $\mathcal D$-measurable for each
$\theta$. The value $L(\theta,a)$ is the cost of taking action $a$
when the parameter is $\theta$.
\end{definition}

The parameter set has no specified sigma algebra here. For a real target
$g:\Theta\to\R$, the functions $(a-g(\theta))^2$ and $|a-g(\theta)|$ on
real actions are called squared and absolute loss, respectively.

\begin{definition}[Risk]
\label{def:decision-risk}
\indexnotation{R_L(\theta,K)}{Risk of rule \(K\) at \(\theta\).}
The \emph{risk} of a rule $K$ at $\theta$ under loss $L$ is
\begin{equation}
R_L(\theta,K)=\int_{\mathsf D} L(\theta,a)\,(P_\theta K)(da).
\label{eq:risk-action-law}
\end{equation}
The map $\theta\mapsto R_L(\theta,K)$ is its \emph{risk function}.
For a deterministic rule $d$, write $R_L(\theta,d)=R_L(\theta,K_d)$.
\end{definition}

By Proposition~\ref{prop:joint-action-law}, integrating the loss against
the joint law gives the equivalent expressions
\begin{align}
R_L(\theta,K)
&=\int L(\theta,a)\,(P_\theta\ltimes K)(d(x,a)) \notag\\
&=\int_{\mathsf X}\left(\int_{\mathsf D}L(\theta,a)\,K(x,da)\right)P_\theta(dx).
\label{eq:risk-iterated}
\end{align}
In particular, $R_L(\theta,d)=\int L(\theta,d(x))\,P_\theta(dx)$.
Nonnegativity makes these integrals well-defined, with $+\infty$ allowed.

\begin{example}[Asymmetric binary costs]
\label{ex:asymmetric-binary-risk}
Let the parameter and action sets be $\{0,1\}$, with all subsets of the
action set measurable. Take $L(0,0)=L(1,1)=0$,
$L(0,1)=c_{01}>0$, and $L(1,0)=c_{10}>0$, with finite costs.
If $D$ is the action coordinate, then
\[
R_L(0,K)=c_{01}(P_0\ltimes K)\{D=1\},\qquad
R_L(1,K)=c_{10}(P_1\ltimes K)\{D=0\}.
\]
For instance, take a binary observation whose success probabilities
at labels $0$ and $1$ are $1/4$ and $3/4$. The rule $d(x)=x$ has risk
pair $(c_{01}/4,c_{10}/4)$. The constant rules $0$ and $1$ have pairs
$(0,c_{10})$ and $(c_{01},0)$, respectively.
\end{example}

\subsection{Pointwise comparison}
\label{subsec:risk-comparison}

\begin{definition}[Risk comparison]
\label{def:risk-comparison}
\indexnotation{K\preceq_L H}{Pointwise comparison of the risks of two rules.}
For a fixed experiment and loss, write $K\preceq_L H$ when
$R_L(\theta,K)\leq R_L(\theta,H)$ for every $\theta$.
Say that $K$ \emph{dominates} $H$ if $K\preceq_L H$ and the inequality
is strict for at least one $\theta$. The rules are \emph{risk equivalent}
if their risk functions agree everywhere.
\end{definition}

\begin{example}[Crossing risk functions]
\label{ex:bernoulli-crossing-risks}
For $n\geq1$ iid Bernoulli observations with parameter $p\in[0,1]$,
take actions in the Borel space $[0,1]$ and loss $L(p,a)=(a-p)^2$.
By \eqref{eq:iid-mean-square}, the sample mean and the constant $1/2$
have risks
\[
R_L(p,\bar X_n)=\frac{p(1-p)}n,\qquad
R_L(p,1/2)=(p-1/2)^2.
\]
Writing $u=p-1/2$, their difference is
\[
R_L(p,\bar X_n)-R_L(p,1/2)
=\frac{1/4-(n+1)u^2}{n}.
\]
The constant has lower risk when $|u|<1/(2\sqrt{n+1})$; the sample mean
has lower risk when $|u|>1/(2\sqrt{n+1})$. They tie at equality.
Neither rule dominates the other.
\end{example}

\subsection{Risk of mixtures}
\label{subsec:risk-mixtures}

\begin{proposition}[Affinity under constant mixing]
\label{prop:affine-risk}
Let $K_s=(1-s)K_0+sK_1$, with $0<s<1$ constant. Then
\[
R_L(\theta,K_s)=(1-s)R_L(\theta,K_0)+sR_L(\theta,K_1)
\qquad(\theta\in\Theta).
\]
The identity holds in $[0,\infty]$. At $s=0$ and $s=1$ the rule is
$K_0$ and $K_1$, respectively.
\end{proposition}
\begin{proof}
By Proposition~\ref{prop:mixture-rules}, the action law is
$P_\theta K_s=(1-s)P_\theta K_0+sP_\theta K_1$. Integrate the nonnegative
function $L(\theta,\cdot)$ against this measure. Linearity for simple
functions and monotone convergence give the formula, including infinite
risk.
\end{proof}

\begin{example}[Domination by a randomized rule]
\label{ex:randomized-domination}
Let the observation space consist of one point, with the same law at
both labels $0,1$. Use three actions $a,b,c$, with all subsets measurable,
and losses
\[
(L(0,a),L(1,a))=(0,4),\quad
(L(0,b),L(1,b))=(4,0),\quad
(L(0,c),L(1,c))=(3,3).
\]
A deterministic rule is constant, so its risk pair is one of the three
displayed pairs. Neither $a$ nor $b$ dominates $c$. The rule choosing
$a$ and $b$ with equal probabilities has risk pair $(2,2)$ and dominates
$c$. The equal observation laws do not force equal risks at the two
labels: the loss can depend on the label.
\end{example}

\subsection{Squared loss and convex losses}
\label{subsec:risk-convex}

\begin{proposition}[Squared risk, bias, and variance]
\label{prop:bias-variance-risk}
Take real actions with their Borel sets and $L(\theta,a)=(a-g(\theta))^2$
for a real target $g$. Fix $\theta$ with
$\int a^2\,(P_\theta K)(da)<\infty$. If $m_\theta$ and $v_\theta$ are the
mean and variance of the action law $P_\theta K$, then
\[
R_L(\theta,K)=v_\theta+\bigl(m_\theta-g(\theta)\bigr)^2.
\]
The difference $m_\theta-g(\theta)$ is the \emph{bias} at $\theta$.
\end{proposition}
\begin{proof}
Put $Q=P_\theta K$. The finite second moment implies
$\int |a|\,Q(da)<\infty$ by Cauchy--Schwarz. Expand
\[
(a-g(\theta))^2
=(a-m_\theta)^2+2(a-m_\theta)(m_\theta-g(\theta))
 +(m_\theta-g(\theta))^2.
\]
The middle term integrates to zero under $Q$, since $m_\theta$ is its
mean. The other two terms give the stated variance and squared bias.
\end{proof}

\begin{example}[A shifted normal sample mean]
\label{ex:shifted-normal-risk}
For $n\geq1$ iid observations from $N(\theta,1)$ and $g(\theta)=\theta$,
take the real-action rule $d_b(x)=\bar x+b$ for fixed $b\in\R$.
Its action law has mean $\theta+b$ and variance $1/n$, by independence.
Under squared loss,
\[
R_L(\theta,d_b)=\frac1n+b^2.
\]
If $b\ne0$, the rule $d_0$ has strictly smaller risk at every $\theta$.
\end{example}

\begin{proposition}[Taking the mean action under convex loss]
\label{prop:convex-mean-action}
Let $\mathsf D=[b,c]$, $b<c$, with its Borel sets, and suppose
$L(\theta,\cdot)$ is finite and convex on $[b,c]$ for every $\theta$.
For any kernel $K$ define
\[
m_K(x)=\int_{[b,c]}a\,K(x,da).
\]
Then $m_K$ is a deterministic rule into $[b,c]$, and $m_K\preceq_L K$.
\end{proposition}
\begin{proof}
The function $m_K-b$ is measurable by
\eqref{eq:kernel-iterated-integral}, applied to the nonnegative function
$(x,a)\mapsto a-b$. Since $b\leq a\leq c$, we have $b\leq m_K(x)\leq c$.
For each fixed $\theta,x$, scalar Jensen's inequality for a convex
function on an interval gives
\[
L(\theta,m_K(x))\leq\int_{[b,c]}L(\theta,a)\,K(x,da).
\]
Integrating with respect to $P_\theta$ and using
\eqref{eq:risk-iterated} proves the comparison for every $\theta$.
\end{proof}

\begin{example}[Selecting one observation or taking the mean]
\label{ex:random-selection-risk}
For $n\geq2$ iid Bernoulli observations with parameter $p$, take actions
in the Borel space $[0,1]$. Consider $K(x,\cdot)=n^{-1}\sum_i\delta_{x_i}$,
which selects a coordinate uniformly. Its action law is Bernoulli with
probability $p$ of one, since
\[
(P_p^{\otimes n}K)(\{1\})=\mathbb E_{P_p^{\otimes n}}[\bar X_n]=p.
\]
With $g(p)=p$, its squared risk is $p(1-p)$. The mean action is
$m_K(x)=\bar x$, whose risk is $p(1-p)/n$. Thus the mean rule dominates
$K$, with strict improvement for $0<p<1$ and equality at $p=0,1$.
\end{example}

\subsection{Exercises}
\label{subsec:worked-examples-2-2}

\begin{practiceexercise}[A directional numerical loss]
\label{wex:asymmetric-numeric-loss}
Let $L(\theta,a)=2(\theta-a)_++(a-\theta)_+$ with $t_+=\max(t,0)$. At
$\theta=1$, compute the risks of actions zero and two and of their equal
mixture.
\end{practiceexercise}

\begin{practiceexercise}[Comparing two losses]
\label{wex:loss-reverses-ranking}
At target zero, compare the rule taking one with probability $1/10$ and zero
otherwise with the constant $1/5$. Compute both risks under absolute and
squared loss and compare the rankings.
\end{practiceexercise}

\begin{practiceexercise}[Error probabilities of a threshold rule]
\label{wex:symmetric-normal-classification-risk}
At labels $-1,1$ observe $X\sim N(\theta,1)$ and choose one when $X>0$. Find
both unit-error risks and compare them.
\end{practiceexercise}

\begin{practiceexercise}[Moving a threshold]
\label{wex:shifted-threshold-tradeoff}
At labels $-1,1$ observe $X\sim N(\theta,1)$ and choose one when $X>t$.
Express the two unit-error risks in $t$ and determine how each changes as
$t$ increases.
\end{practiceexercise}

\begin{practiceexercise}[Averaging or selecting a measurement]
\label{wex:weighted-normal-risk}
Independent $X_1,X_2$ have means $\theta$ and variances one and four.
Compute squared risk for $wX_1+(1-w)X_2$ and for selecting $X_1$ with
probability $w$ and $X_2$ otherwise. Compare at $w=4/5$.
\end{practiceexercise}

\begin{practiceexercise}[Shrinkage and the parameter range]
\label{wex:bounded-normal-shrinkage}
For $X\sim N(\theta,1)$, compare the squared risks of $X$ and $X/2$. Decide
whether $X/2$ dominates $X$ on $[-1,1]$ and on $\R$.
\end{practiceexercise}

\begin{practiceexercise}[Clipping to a known interval]
\label{wex:clip-reduces-square}
For $X\sim N(\theta,1)$, $\theta\in[0,1]$, compare pointwise squared losses
before and after clipping $X$ to $[0,1]$. Determine whether the risk
comparison is strict at every label.
\end{practiceexercise}

\begin{practiceexercise}[Independent action noise]
\label{wex:noise-adds-square-risk}
A rule $d$ has finite squared risk for a real target. Independent noise $U$
has mean zero and variance $\sigma^2$. Express the risk of $d(X)+U$ in terms
of that of $d(X)$.
\end{practiceexercise}

\begin{practiceexercise}[An infinite error penalty]
\label{wex:infinite-error-penalty}
Give loss zero to a correct binary label and infinite loss to a mistake.
Characterize finite risk at a label. With a constant observation at both
labels, can any rule have finite risk at both?
\end{practiceexercise}

\begin{practiceexercise}[Moments of a countable action law]
\label{wex:countable-infinite-variance}
An observation has mass proportional to $j^{-3}$ on the positive integers.
For action $d(j)=j$ at target zero, determine whether its mean and squared
risk are finite.
\end{practiceexercise}

\begin{practiceexercise}[Comparing action laws]
\label{wex:equal-square-risk-different-laws}
On a constant record compare the fair law on $\{-1,1\}$ with masses
$1/4,1/2,1/4$ at $-\sqrt2,0,\sqrt2$. Find their squared risks at every real
target and absolute risks at zero.
\end{practiceexercise}

\begin{practiceexercise}[A mixture using the record]
\label{wex:data-dependent-risk-mixture}
Observe $X=\theta$ at binary label $\theta$. Under unit-error loss find the
risk pair of the mixture that takes constant action one with weight $X$ and
zero otherwise. Can a constant weight give the same pair?
\end{practiceexercise}

\begin{practiceexercise}[A mean action under nonconvex loss]
\label{wex:nonconvex-mean-action}
On $[-1,1]$ use $L(a)=(a^2-1)^2$. Compare an equal mixture of $-1$ and $1$
with its mean action. Which hypothesis of
Proposition~\ref{prop:convex-mean-action} fails?
\end{practiceexercise}

\begin{practiceexercise}[A mean action under absolute loss]
\label{wex:absolute-loss-jensen-tie}
At target zero a rule takes zero or two with equal probabilities. Compare
its absolute risk with the loss of its mean action. What property of the
loss explains the comparison?
\end{practiceexercise}

\begin{practiceexercise}[Missing a tolerance band]
\label{wex:tolerance-loss-normal-mean}
For iid $N(\theta,1)$ measurements and
$L(\theta,a)=\mathbf1_{\{|a-\theta|>\epsilon\}}$, $\epsilon>0$, compute the
sample mean's risk and decide whether it depends on $\theta$.
\end{practiceexercise}

\begin{exercise}[Asymmetric risk after a random flip]
\label{exer:binary-flip-risk}
In Example~\ref{ex:asymmetric-binary-risk}, compare the identity rule
with the binary flip mixture of Example~\ref{ex:binary-flip-rule}.
Use the two specified success probabilities $1/4$ and $3/4$ to compute
both risk pairs. Which rule dominates?
\end{exercise}

\begin{exercise}[Changes of loss units]
\label{exer:risk-comparison-rescaling}
Let $s(\theta)>0$ and $b(\theta)\geq0$ be finite for every $\theta$.
For $L'(\theta,a)=s(\theta)L(\theta,a)+b(\theta)$, express
$R_{L'}(\theta,K)$ in terms of $R_L(\theta,K)$, including when the latter
is infinite. Show that risk equivalence and domination are unchanged.
\end{exercise}

\begin{exercise}[Bernoulli shrinkage]
\label{exer:bernoulli-shrinkage-risk}
In Example~\ref{ex:bernoulli-crossing-risks}, fix $0<\alpha<1$ and take
$d_\alpha(x)=\alpha\bar x+(1-\alpha)/2$. Compute its bias, variance,
and squared risk. Determine for which $p$ its risk is strictly smaller
than that of the sample mean and for which $p$ it is strictly larger.
\end{exercise}

\begin{exercise}[Finite mean, infinite squared risk]
\label{exer:infinite-squared-risk}
Take one parameter label, with target zero, and observe a uniform variable
on $(0,1]$ with its Borel sets. Use the real-action rule $d(x)=x^{-1/2}$.
Compute the mean of its action law and its squared risk. Check the moment
hypothesis in Proposition~\ref{prop:bias-variance-risk}.
\end{exercise}

\begin{exercise}[The exact gap under squared loss]
\label{exer:mean-action-risk-gap}
In Proposition~\ref{prop:convex-mean-action}, take
$L(\theta,a)=(a-g(\theta))^2$ for a real target $g$. Prove
\[
R_L(\theta,K)-R_L(\theta,m_K)
=\int_{\mathsf X}\int_{[b,c]}(a-m_K(x))^2\,K(x,da)\,P_\theta(dx).
\]
Use the zero-integral condition for nonnegative functions to characterize
equality at a given $\theta$.
\end{exercise}

\begin{exercise}[Risk functions and action laws]
\label{exer:risks-identify-action-marginals}
For rules $K,H$ in a fixed experiment, prove that they have the same
risk function for every loss if and only if $P_\theta K=P_\theta H$ for
every $\theta$. For necessity, use losses $L(\theta,a)=\mathbf 1_B(a)$.
Does equality of risks for one squared loss imply equality of action laws?
\end{exercise}

\clearpage
\section{Bayes risk and optimal actions}
\label{sec:bayes-risk}

\subsection{Averaging risk over the parameter}
\label{subsec:prior-risk}
\indexnotation{\mathcal T,\ \pi}{Parameter sigma algebra and prior law.}

Give $\Theta$ a sigma algebra $\mathcal T$. In this section assume that
$\theta\mapsto P_\theta$ is a probability kernel to $(\mathsf X,\mathcal A)$,
and that $L$ is $\mathcal T\otimes\mathcal D$-measurable. These are
additional hypotheses on the experiment and loss.

\begin{definition}[Prior and predictive law]
\label{def:prior-predictive}
\indexnotation{J_\pi,\ M_\pi}{Joint parameter and observation law, and prior predictive law.}
A \emph{prior} $\pi$ is a probability law on $(\Theta,\mathcal T)$.
The joint parameter and observation law is $J_\pi=\pi\ltimes P$, where
the kernel $P$ sends $\theta$ to $P_\theta$. Its observation marginal
$M_\pi$, the \emph{prior predictive law}, is
\[
M_\pi(A)=\int_\Theta P_\theta(A)\,\pi(d\theta)
\qquad(A\in\mathcal A).
\]
\end{definition}

For any rule $K$, the function
$h(\theta,x)=\int L(\theta,a)\,K(x,da)$ is product-measurable:
apply the inner-integral assertion of Proposition~\ref{prop:joint-action-law}
to the kernel $(\theta,x)\mapsto K(x,\cdot)$. Applying it again to $P$
shows that $\theta\mapsto\int h(\theta,x)\,P_\theta(dx)=R_L(\theta,K)$
is measurable.

\begin{definition}[Bayes risk and Bayes rule]
\label{def:bayes-risk}
\indexnotation{r_\pi(K),\ r_\pi^\star}{Prior average risk and Bayes value.}
The \emph{Bayes risk of $K$ under $\pi$} is
\[
r_\pi(K)=\int_\Theta R_L(\theta,K)\,\pi(d\theta).
\]
The \emph{Bayes value} is $r_\pi^\star=\inf_K r_\pi(K)$, over all kernels
to the specified action space. A \emph{Bayes rule} under $\pi$ is a rule
attaining this infimum. As before, a deterministic rule is identified
with its Dirac kernel.
\end{definition}

\subsection{Posterior laws and expected loss}
\label{subsec:posterior-loss}

\begin{definition}[Posterior kernel]
\label{def:posterior-kernel}
\indexnotation{\Pi(x,\cdot)}{Posterior law at record \(x\).}
A \emph{posterior kernel} $\Pi$ from the observation space to
$(\Theta,\mathcal T)$ satisfies
\begin{equation}
\int_A\Pi(x,B)\,M_\pi(dx)
=\int_B P_\theta(A)\,\pi(d\theta)
\qquad(A\in\mathcal A,\ B\in\mathcal T).
\label{eq:posterior-rectangles}
\end{equation}
Thus $M_\pi\ltimes\Pi$ is the law obtained by reversing the coordinates
of $J_\pi$.
\end{definition}

\begin{proposition}[A posterior from densities]
\label{prop:posterior-from-densities}
\indexnotation{m(x)}{Density of the prior predictive law.}
Suppose the experiment is dominated by a $\sigma$-finite $\mu$ and
admits a jointly measurable choice $f(\theta,x)$ of densities. Then
$M_\pi$ has density
\[
m(x)=\int_\Theta f(\theta,x)\,\pi(d\theta)
\]
with respect to $\mu$. A posterior kernel is
\begin{equation}
\Pi(x,B)=\frac{\int_B f(\theta,x)\,\pi(d\theta)}{m(x)}
\qquad\text{when }0<m(x)<\infty.
\label{eq:posterior-density-formula}
\end{equation}
On the complementary $M_\pi$-null set, one may set $\Pi(x,\cdot)=\pi$.
\end{proposition}
\begin{proof}
Tonelli's theorem gives measurability of $m$ and, for $A\in\mathcal A$,
\[
M_\pi(A)=\int_\Theta\int_A f(\theta,x)\,\mu(dx)\,\pi(d\theta)
=\int_A m(x)\,\mu(dx).
\]
Since $\int m\,d\mu=1$, the set where $m=\infty$ is $\mu$-null, and the
set where $m=0$ is $M_\pi$-null. On $\{0<m<\infty\}$ the numerator in
\eqref{eq:posterior-density-formula} defines a measure in $B$ with total
mass $m(x)$. For fixed $B$ it is measurable in $x$, by Tonelli. Hence the
formula and the fallback prior define a probability kernel. Integrating
it over $A$ against $M_\pi$ cancels the factor $m$ and gives
$\int_A\int_B f(\theta,x)\,\pi(d\theta)\,\mu(dx)$: the inner integral
vanishes where $m=0$. Another use of Tonelli gives the right side of
\eqref{eq:posterior-rectangles}.
\end{proof}

For finitely many labels with all subsets measurable one can always use
$\mu=\sum_{\theta\in\Theta}P_\theta$. Joint measurability then follows
from measurability of the finitely many densities.

\begin{proposition}[Risk as averaged posterior loss]
\label{prop:posterior-risk-averaging}
\indexnotation{C(x,a)}{Posterior expected loss of action \(a\).}
Given a posterior kernel $\Pi$, the \emph{posterior expected loss}
\[
C(x,a)=\int_\Theta L(\theta,a)\,\Pi(x,d\theta)
\]
is product-measurable. For each rule $K$,
\begin{equation}
r_\pi(K)=\int_{\mathsf X}\int_{\mathsf D} C(x,a)\,K(x,da)\,M_\pi(dx).
\label{eq:posterior-risk-averaging}
\end{equation}
If a measurable rule $d$ minimizes $C(x,\cdot)$ outside a single
$M_\pi$-null set, then $d$ is a Bayes rule.
\end{proposition}
\begin{proof}
The kernel $(x,a)\mapsto\Pi(x,\cdot)$ and
Proposition~\ref{prop:joint-action-law} give measurability of $C$.
Using the function $h$ defined above, we have
\[
r_\pi(K)=\int h(\theta,x)\,J_\pi(d\theta,dx)
=\int_{\mathsf X}\int_\Theta h(\theta,x)\,\Pi(x,d\theta)\,M_\pi(dx).
\]
The second equality uses \eqref{eq:posterior-rectangles} and the
kernel integral formula. For each $x$, apply Tonelli to
$\Pi(x,\cdot)\otimes K(x,\cdot)$ to interchange the two integrals in
\[
\int_\Theta h(\theta,x)\,\Pi(x,d\theta)
=\int_\Theta\int_{\mathsf D}L(\theta,a)\,K(x,da)\,\Pi(x,d\theta).
\]
This gives \eqref{eq:posterior-risk-averaging}. If $d$ minimizes
pointwise outside a null set, then $C(x,d(x))\leq\int C(x,a)\,K(x,da)$
there, for every $K$. Integrating proves optimality.
\end{proof}

\subsection{Optimal decisions among finitely many actions}
\label{subsec:finite-bayes-actions}

\begin{proposition}[Finite action spaces]
\label{prop:finite-bayes-action}
Suppose a posterior kernel exists and $\mathsf D=\{a_1,\ldots,a_m\}$,
$m\geq1$, with all subsets measurable. Put $c(x)=\min_j C(x,a_j)$.
Then
\[
r_\pi^\star=\int_{\mathsf X}c(x)\,M_\pi(dx),
\]
and the rule choosing the first minimizing action is measurable and
attains this value. If the value is finite, a rule is Bayes exactly
when it assigns probability one to the minimizing actions for
$M_\pi$-almost every $x$.
\end{proposition}
\begin{proof}
The minimum of the finitely many measurable functions $C_j(x)=C(x,a_j)$
is measurable. The set on which $a_j$ is the first minimizer is
$\{C_j=c\}\cap\bigcap_{i<j}\{C_i>c\}$, so the selected rule is measurable.
Its Bayes risk is $\int c\,dM_\pi$ by
\eqref{eq:posterior-risk-averaging}, and every rule has risk at least
this value.

If the value is finite, $c(x)$ is finite almost everywhere. On that
set the excess over $c(x)$ is the nonnegative integral of
$C(x,a)-c(x)$ against $K(x,\cdot)$. Its integral over $x$ vanishes
exactly when the excess vanishes almost everywhere. For each such
$x$, this holds exactly when $K(x,\cdot)$ assigns probability one
to the set on which $C(x,a)=c(x)$.
\end{proof}

\begin{example}[Asymmetric binary decisions and a tie]
\label{ex:binary-bayes-actions}
Use the two binary observation laws from
Example~\ref{ex:asymmetric-binary-risk}, with prior masses $1/2,1/2$.
Both records have predictive probability $1/2$. Write $q_x=\Pi(x,\{1\})$.
Bayes' formula gives $q_0=1/4$ and $q_1=3/4$.
The posterior losses of actions zero and one are
$c_{10}q_x$ and $c_{01}(1-q_x)$. Thus action one is optimal when
$q_x\geq c_{01}/(c_{01}+c_{10})$, with a tie at equality.

For $c_{01}=3$ and $c_{10}=1$, action zero is uniquely optimal at
$x=0$, and either action is optimal at $x=1$. The Bayes value is
$\tfrac12(1/4+3/4)=1/2$. Any randomization at $x=1$ attains the same value
if the rule takes action zero at $x=0$.
\end{example}

\begin{example}[Ties without information]
\label{ex:bayes-ties-no-information}
In Example~\ref{ex:randomized-domination}, give both labels prior mass
$1/2$. The observation is constant, so the posterior is the prior.
The posterior losses of $a,b,c$ are $2,2,3$. Every mixture supported on
$a,b$ is Bayes. For example, the constant rule $a$ and the equal mixture
of $a,b$ both have Bayes risk $2$, but have risk pairs $(0,4)$ and $(2,2)$.
\end{example}

\begin{example}[A record with predictive probability zero]
\label{ex:bayes-prior-null-record}
Let $\Theta=\mathsf X=\mathsf D=\{0,1\}$, with all subsets measurable,
$P_\theta=\delta_\theta$, and $L(\theta,a)=\mathbf 1_{\{a\ne\theta\}}$.
Take $\pi=\delta_0$. The predictive law is $\delta_0$, so a rule is
Bayes precisely when $K(0,\{0\})=1$. In particular, the constant rule zero is
Bayes, with risk pair $(0,1)$. The identity rule is also Bayes and
dominates it, with pair $(0,0)$. The prior assigns zero mass to the
label at which the two risks differ.
\end{example}

\subsection{The posterior mean under squared loss}
\label{subsec:posterior-mean}

\begin{proposition}[The posterior mean is optimal for squared loss]
\label{prop:bayes-squared-mean}
\indexnotation{d_\pi(x),\ v(x)}{Posterior mean and variance of the target under squared loss.}
Suppose a posterior kernel exists, take real Borel actions, and let
$L(\theta,a)=(a-g(\theta))^2$ for a measurable real target $g$ with
$\int g^2\,d\pi<\infty$. For $M_\pi$-almost every $x$, let $d_\pi(x)$
and $v(x)$ be the mean and variance of $g$ under $\Pi(x,\cdot)$.
They admit measurable choices, and
\begin{align}
r_\pi^\star&=\int_{\mathsf X}v(x)\,M_\pi(dx), \label{eq:bayes-value-square}\\
r_\pi(K)&=r_\pi^\star+
\int_{\mathsf X}\int_{\R}(a-d_\pi(x))^2\,K(x,da)\,M_\pi(dx).
\label{eq:bayes-square-gap}
\end{align}
Thus $d_\pi$ is a Bayes rule. Every Bayes rule takes the action
$d_\pi(x)$ with probability one for $M_\pi$-almost every $x$.
\end{proposition}
\begin{proof}
The posterior rectangle identity implies
\[
\int_{\mathsf X}\int_\Theta g(\theta)^2\,\Pi(x,d\theta)\,M_\pi(dx)
=\int_\Theta g(\theta)^2\,\pi(d\theta)<\infty.
\]
The posterior second moment is measurable and finite on a set of full
$M_\pi$-measure. On this set the mean exists by Cauchy--Schwarz and is
measurable by integrating the positive and negative parts of $g$. Extend
it by zero elsewhere. The variance is measurable on the same set by the
kernel integral formula for $(x,\theta)\mapsto(g(\theta)-d_\pi(x))^2$;
extend it by zero as well.

At every record in the full-measure set, the square expansion of
Proposition~\ref{prop:bias-variance-risk}, now for the posterior law of
$g(\theta)$, gives $C(x,a)=v(x)+(a-d_\pi(x))^2$ for each $a$.
The mean rule attains $\int v\,dM_\pi$, which is finite because the
variance is bounded by the posterior second moment. Equation
\eqref{eq:posterior-risk-averaging} now yields
\eqref{eq:bayes-square-gap}. Its nonnegative excess term vanishes exactly
when $K(x,\{d_\pi(x)\})=1$ almost everywhere.
\end{proof}

\begin{example}[A two-point prior for a Bernoulli probability]
\label{ex:bayes-two-point-bernoulli}
Observe one Bernoulli variable with unknown success probability $p$.
Give $p=1/4$ and $p=3/4$ prior mass $1/2$ each, and use squared loss for
$g(p)=p$. The probability of $p=3/4$ given $x$ is $q_x$ as in
Example~\ref{ex:binary-bayes-actions}. Thus the mean action is
$1/4+q_x/2$, equal to $3/8$ at $x=0$ and $5/8$ at $x=1$.
At either record the posterior variance is
\[
(3/4-1/4)^2q_x(1-q_x)=3/64.
\]
The Bayes value is $3/64$. For comparison, the constant $1/2$ has
squared risk $1/16$ at both labels and hence Bayes risk $1/16$.
\end{example}

\begin{example}[A uniform prior for a Bernoulli probability]
\label{ex:bayes-uniform-bernoulli}
Keep one Bernoulli observation, and put the uniform prior on $p\in[0,1]$.
Relative to counting measure on the records, the density is $p$ at
$x=1$ and $1-p$ at $x=0$. Each record has predictive probability $1/2$.
The posterior densities on $[0,1]$ are $2p$ at $x=1$ and $2(1-p)$ at
$x=0$. Their means are respectively $2/3$ and $1/3$.
At $x=1$ the variance is $\int_0^1 2p^3\,dp-(2/3)^2=1/18$;
reflection $p\mapsto1-p$ gives the same variance at $x=0$.
The Bayes value under squared loss is therefore $1/18$.
\end{example}

\subsection{Exercises}
\label{subsec:worked-examples-2-3}

\begin{practiceexercise}[Updating an unequal prior]
\label{wex:unequal-coin-prior}
Give $p=1/4,3/4$ prior masses $2/3,1/3$. Compute the predictive success
probability and the posterior mass of the larger rate after success and
after failure.
\end{practiceexercise}

\begin{practiceexercise}[Updating posterior odds]
\label{wex:coin-posterior-odds}
For labels $p=1/4,3/4$ with prior odds $1/2$ for the larger rate, compute
the likelihood ratios for a success and a failure. Find the posterior odds
after two successes and one failure.
\end{practiceexercise}

\begin{practiceexercise}[Independence given the label]
\label{wex:prior-induced-dependence}
Give $p=1/4,3/4$ equal prior mass and observe two conditionally iid bits.
Compute their predictive marginals and joint success probability. Are the
observations independent under the predictive law?
\end{practiceexercise}

\begin{practiceexercise}[Predicting the next bit]
\label{wex:next-bit-prediction}
After a success in Exercise~\ref{wex:unequal-coin-prior}, find the
predictive success probability of a second draw independent of the first
conditional on $p$.
\end{practiceexercise}

\begin{practiceexercise}[A posterior on two endpoints]
\label{wex:endpoint-posterior-update}
Give endpoints one and two equal prior mass and observe one uniform draw on
$[0,\theta]$. Compute the posterior weights at $x=1/2$ and $x=3/2$.
\end{practiceexercise}

\begin{practiceexercise}[A nonuniform continuous prior]
\label{wex:increasing-prior-binary}
Use prior density $2p$ on $[0,1]$ and observe a Bernoulli bit. Compute the
predictive success probability and the two posterior densities and means.
\end{practiceexercise}

\begin{practiceexercise}[A prior with an atom]
\label{wex:prior-atom-survives}
Mix $\delta_0$ and the uniform prior on $p$ with equal weights. Compute the
posterior atom at zero after failure and after success. Find the posterior
mean after failure.
\end{practiceexercise}

\begin{practiceexercise}[A count and an ordered record]
\label{wex:count-posterior-same}
For two iid Bernoulli observations and a uniform prior, compare the
posterior densities from record $(0,1)$ and from only the information that
there was one success. Explain how the likelihood factors enter
normalization.
\end{practiceexercise}

\begin{practiceexercise}[Prior odds and asymmetric costs]
\label{wex:asymmetric-posterior-threshold}
In Exercise~\ref{wex:unequal-coin-prior}, predicting the larger
rate incorrectly costs two and missing it costs one. Find the decision
threshold, the optimal action at each record, and the Bayes risk.
\end{practiceexercise}

\begin{practiceexercise}[Unit error loss on three labels]
\label{wex:three-class-posterior-mode}
With posterior probabilities $(1/2,1/3,1/6)$, compute the posterior losses
of the three label predictions under unit error loss and find the optimum.
\end{practiceexercise}

\begin{practiceexercise}[A decision with an abstention option]
\label{wex:reject-option}
Predict either of two labels under unit error loss, or abstain for loss
$1/3$. Find the optimal action at posterior probability vectors $(3/5,2/5)$
and $(4/5,1/5)$.
\end{practiceexercise}

\begin{practiceexercise}[Squared loss on a grid]
\label{wex:grid-constrained-posterior-mean}
The target is binary with posterior mass $13/20$ at one. Find the
squared-loss minimizer over all real actions and over $\{0,1/2,1\}$.
\end{practiceexercise}

\begin{practiceexercise}[Asymmetric absolute loss]
\label{wex:asymmetric-absolute-posterior}
A posterior puts mass $3/5$ at zero and the remainder at two. Find the
optimal action under absolute loss and when underestimation costs twice as
much per unit as overestimation. Justify why actions outside $[0,2]$ cannot
improve the optimum.
\end{practiceexercise}

\begin{practiceexercise}[Bayes risk with erasures]
\label{wex:erasure-bayes-value}
Give the two binary labels equal prior mass. Report the label with
probability $r\in(0,1)$ and erase otherwise. Find the posterior variance on
a reveal and on an erasure, then the squared-loss Bayes value.
\end{practiceexercise}

\begin{practiceexercise}[Prior average versus worst case]
\label{wex:unequal-prior-risk-tradeoff}
With no data and squared loss, the target is $1$ with prior probability
$3/4$ and $-1$ otherwise. Find the Bayes action. Compare its average and
worst-case risks with those of action zero.
\end{practiceexercise}

\begin{exercise}[Binary costs with an arbitrary prior]
\label{exer:binary-bayes-threshold}
Use the observation laws and positive costs from
Example~\ref{ex:asymmetric-binary-risk}. Give label one prior mass
$s\in[0,1]$. For each record, derive the condition on $s$ under which
action one minimizes posterior loss. Include ties and the endpoints
$s=0,1$. Express the Bayes value as a sum of two minima, one for each
record.
\end{exercise}

\begin{exercise}[Changing density versions]
\label{exer:posterior-density-versions}
Suppose $f$ and $\widetilde f$ are jointly measurable density choices
for the same experiment and the same dominating measure $\mu$.
Show that the kernels constructed by
\eqref{eq:posterior-density-formula} agree as probability measures for
$M_\pi$-almost every $x$. Apply Tonelli to the indicator of
$\{(\theta,x):f(\theta,x)\ne\widetilde f(\theta,x)\}$ to find a single
exceptional set of records.
\end{exercise}

\begin{exercise}[Positive prior mass and domination]
\label{exer:finite-positive-prior-domination}
Suppose $\Theta$ is finite, every label has positive prior mass, and
a Bayes rule has finite Bayes risk. Prove that no rule dominates it.
Explain why Example~\ref{ex:bayes-prior-null-record} does not
contradict this conclusion.
\end{exercise}

\begin{exercise}[Absolute loss and posterior medians]
\label{exer:bayes-absolute-median}
Let $g$ be measurable with finitely many real values, and use absolute
loss on real actions. At a record $x$, show that an action $a$ minimizes posterior
expected loss exactly when
\[
\Pi(x,\{\theta:g(\theta)<a\})\leq\tfrac12,\qquad
\Pi(x,\{\theta:g(\theta)>a\})\leq\tfrac12.
\]
Use the left and right slopes of the finite weighted sum of absolute
values. Show that the smallest target value at which the posterior
cumulative mass reaches $1/2$ gives a measurable Bayes rule.
\end{exercise}

\begin{exercise}[A uniform prior and an iid record]
\label{exer:bayes-uniform-bernoulli-iid}
Extend Example~\ref{ex:bayes-uniform-bernoulli} to $n\geq1$ iid
observations. For a record with $k$ ones, derive the posterior density
up to its normalizing constant and prove that the mean is
$(k+1)/(n+2)$. To compute the ratio of the integrals, integrate the
derivative of $p^{k+1}(1-p)^{n-k+1}$ on $[0,1]$.
\end{exercise}

\begin{exercise}[Rescaling loss and tilting the prior]
\label{exer:bayes-rescaling-prior}
Let $s(\theta)>0$ be measurable and finite, with $Z=\int s\,d\pi<\infty$.
Set $L'(\theta,a)=s(\theta)L(\theta,a)$ and
$\pi_s(B)=Z^{-1}\int_B s\,d\pi$. Show that the Bayes risk under $L'$
and $\pi$ is $Z$ times that under $L$ and $\pi_s$.
Give an example in which rescaling changes the set of Bayes rules
while the prior is held fixed.
\end{exercise}

\clearpage
\section{Minimax risk, admissibility, and existence}
\label{sec:minimax-admissibility}

\subsection{Worst-case risk and admissibility}
\label{subsec:worst-case-risk}

Fix an experiment, action space, and loss $L$. The class of rules consists
of all kernels to the action space, including deterministic rules.

\begin{definition}[Worst-case risk and minimax rule]
\label{def:minimax-rule}
\indexnotation{W(K),\ V}{Worst-case risk and minimax value.}
The \emph{worst-case risk} of $K$ and the \emph{minimax value} are
\[
W(K)=\sup_{\theta\in\Theta}R_L(\theta,K),\qquad
V=\inf_K W(K).
\]
A rule is \emph{minimax} if $W(K)=V$.
\end{definition}

\begin{definition}[Admissibility]
\label{def:admissibility}
A rule is \emph{admissible} if no rule dominates it in the sense of
Definition~\ref{def:risk-comparison}.
\end{definition}

\begin{example}[A dominated minimax rule]
\label{ex:dominated-minimax}
Let the observation be constant at both labels $0,1$. Use two actions
$a,b$ with risk pairs $(1,0)$ and $(1,1)$, respectively. These pairs are
their loss pairs because the observation is constant. Every rule
has risk one at label zero and at most one at label one, so every rule
is minimax. The rule taking action $a$ dominates the rule taking $b$.
\end{example}

\subsection{Bayes lower bounds}
\label{subsec:bayes-minimax-bound}

When averaging over a prior, use the measurability assumptions of
Section~\ref{sec:bayes-risk}.

\begin{definition}[Least favorable prior]
\label{def:least-favorable-prior}
A prior is \emph{least favorable} if it attains $\sup_\pi r_\pi^\star$,
where the supremum is over all probability priors on the given
measurable parameter space.
\end{definition}

\begin{proposition}[A Bayes bound and a minimax certificate]
\label{prop:bayes-minimax-certificate}
For every prior $\pi$ and rule $K$,
\[
r_\pi^\star\leq V\leq W(K).
\]
If $W(K)\leq r_\pi^\star$, then all three values agree, $K$ is minimax,
and $\pi$ is least favorable. In particular, a Bayes rule with finite
risk constant in $\theta$ is minimax, and its prior is least favorable.
\end{proposition}
\begin{proof}
For each rule $H$, the average of its risk under $\pi$ is at most
$W(H)$, so $r_\pi^\star\leq r_\pi(H)\leq W(H)$. Taking the infimum
over $H$ gives $r_\pi^\star\leq V$. The other inequality follows from
the definition of $V$. If the outer values are equal, $K$ attains $V$
and no other prior can have Bayes value above $r_\pi^\star=V$.
For a Bayes rule with constant finite risk, the prior average and
supremum of its risk both equal that constant.
\end{proof}

\begin{proposition}[Strict improvement and prior mass]
\label{prop:bayes-admissibility-mass}
Let $K$ be Bayes under $\pi$ with $r_\pi(K)<\infty$.
If $H\preceq_L K$, then
\[
\pi\{\theta:R_L(\theta,H)<R_L(\theta,K)\}=0.
\]
In particular, if $\Theta$ is finite and every label has positive prior
mass, then $K$ is admissible.
\end{proposition}
\begin{proof}
The inequality of risks implies $r_\pi(H)\leq r_\pi(K)$; Bayes
optimality gives equality. The nonnegative difference
$R_L(\theta,K)-R_L(\theta,H)$ is defined wherever $R_L(\theta,K)$
is finite, a set of full $\pi$-measure. It integrates to zero,
so the strict improvement set is $\pi$-null. If every label has
positive mass, a nonempty strict improvement set cannot be null.
\end{proof}

\subsection{Finite problems}
\label{subsec:finite-minimax}

\begin{proposition}[Existence in a finite problem]
\label{prop:finite-minimax-existence}
Suppose $\Theta=\{\theta_1,\ldots,\theta_m\}$, $\mathsf X$, and
$\mathsf D$ are finite with all subsets measurable, and all losses are
finite. The set
\[
S=\{(R_L(\theta_1,K),\ldots,R_L(\theta_m,K)):K\text{ a rule}\}
\subseteq\R^m
\]
is compact and convex. There exists an admissible minimax rule.
\end{proposition}
\begin{proof}
A kernel is a point of a product of simplexes, one probability vector
$(K(x,\{a\}))_{a\in\mathsf D}$ for each $x\in\mathsf X$. This product is
closed and bounded in a finite-dimensional Euclidean space, hence
compact, and it is convex. Each risk coordinate is linear in the
kernel probabilities:
\[
R_L(\theta_i,K)
=\sum_{x\in\mathsf X}\sum_{a\in\mathsf D}
P_{\theta_i}(\{x\})K(x,\{a\})L(\theta_i,a).
\]
Thus the image $S$ is compact and convex. The continuous function
$r\mapsto\max_i r_i$ attains its minimum on $S$. Among the minimizing
risk vectors, a nonempty compact set, choose one minimizing
$\sum_i r_i$. A rule dominating this vector would still be minimax
and would have a strictly smaller sum. Hence the chosen vector
comes from an admissible minimax rule.
\end{proof}

\begin{theorem}[Finite minimax equality]
\label{thm:finite-minimax-equality}
Under the hypotheses of Proposition~\ref{prop:finite-minimax-existence},
\[
V=\max_{\pi}r_\pi^\star.
\]
For any minimax rule $K$ and least favorable prior $\pi$, the rule is
Bayes under $\pi$, and every label of positive $\pi$-mass has risk $V$
under $K$.
\end{theorem}
\begin{proof}
The prior can be identified with its probability vector
$(\pi_1,\ldots,\pi_m)$, and $r_\pi^\star=\inf_{r\in S}\pi\cdot r$.
We construct priors that give lower bounds approaching $V$.

Put $C=S+[0,\infty)^m$. This set is convex and closed: if
$r_j+h_j$ converges with $r_j\in S$, $h_j\geq0$ coordinatewise,
a subsequence has $r_j\to r\in S$, and the corresponding $h_j$
converges to a vector with nonnegative coordinates. Fix $t<V$ and
set $z=(t,\ldots,t)$. Then $z\notin C$, since $r\leq z$ for some
$r\in S$ would give $\max_i r_i\leq t$.

Choose $c\in C$ nearest to $z$ in Euclidean distance. Such a point
exists by restricting the search to the compact intersection of
$C$ with a closed ball about $z$ containing at least one point of $C$.
Write $u=c-z\ne0$. For any $r\in C$, the points
$c+s(r-c)$, $0\leq s\leq1$, belong to $C$. Expanding their squared
distances to $z$ and letting $s\downarrow0$ gives
$u\cdot(r-c)\geq0$. Applying this inequality to $r=c+e_i$, where
$e_i$ is the $i$th coordinate unit vector, gives $u_i\geq0$.
Hence $\pi=u/\sum_i u_i$ is a prior. For every $r\in S$,
\[
\pi\cdot r\ \geq\ \pi\cdot c
=t+\frac{\|u\|^2}{\sum_i u_i}>t.
\]
Choose $t_j=V-1/j$ and take a convergent subsequence of the resulting
priors, using compactness of the simplex. Its limit $\pi^\star$
satisfies $\pi^\star\cdot r\geq V$ for every $r\in S$. Therefore
$r_{\pi^\star}^\star\geq V$, and the Bayes bound gives equality.

If $K$ is minimax and $\pi$ is least favorable, then
$V=r_\pi^\star\leq r_\pi(K)\leq W(K)=V$. Thus $K$ is Bayes under
$\pi$. Its risks are at most $V$, so their average can equal $V$
only if each coordinate of positive prior mass equals $V$.
\end{proof}

\begin{example}[Randomization reduces worst-case risk]
\label{ex:finite-randomized-minimax}
In Example~\ref{ex:randomized-domination}, the deterministic risk
pairs for $a,b,c$ are $(0,4),(4,0),(3,3)$. Their worst-case risks
are $4,4,3$. Giving both labels prior mass $1/2$ gives Bayes value
$2$: actions $a,b$ both have average risk $2$, and any mixture
with $c$ raises that average. The equal mixture of $a,b$ has risk
pair $(2,2)$, so the Bayes lower bound certifies minimax value $2$.
This mixture is the only minimax rule: optimality under the equal
prior excludes $c$, and both risks can be at most $2$ only with
equal weights on $a,b$.

The pure action $a$ is admissible, since a rule with risk zero at
label zero must take action $a$ with probability one. Its worst-case
risk is $4$, so admissibility alone does not imply minimaxity.
\end{example}

\subsection{An equalizer under squared loss}
\label{subsec:bernoulli-minimax}

\begin{example}[Minimax estimation of a Bernoulli probability]
\label{ex:bernoulli-minimax}
Observe $n\geq1$ iid Bernoulli variables with parameter $p\in[0,1]$,
and use squared loss $L(p,a)=(a-p)^2$ with Borel actions in $[0,1]$.
Put $S=\sum_i X_i$, $\alpha=\sqrt n/2$, and define
\[
d(X)=\frac{S+\alpha}{n+2\alpha}.
\]
Since $S$ has mean $np$ and variance $np(1-p)$, the bias of $d$ is
$\alpha(1-2p)/(n+2\alpha)$. Its variance is $np(1-p)/(n+2\alpha)^2$.
The squared risk is therefore
\begin{align*}
R_L(p,d)
&=\frac{np(1-p)+\alpha^2(1-2p)^2}{(n+2\alpha)^2}\\
&=\frac{1}{4(\sqrt n+1)^2},
\end{align*}
using $\alpha^2=n/4$ and $p(1-p)+(1-2p)^2/4=1/4$.

Take the prior with density proportional to
$p^{\alpha-1}(1-p)^{\alpha-1}$ on $(0,1)$, with no mass at the endpoints.
The integral is finite because $\alpha>0$. For a record with $S=k$,
the likelihood of Example~\ref{ex:bernoulli-likelihood} gives posterior
density proportional to $p^{b-1}(1-p)^{c-1}$, with
$b=k+\alpha$, $c=n-k+\alpha$. For $b,c>0$, integration of the derivative
of $p^b(1-p)^c$ gives
\[
b\int_0^1 p^{b-1}(1-p)^c\,dp
=c\int_0^1 p^b(1-p)^{c-1}\,dp.
\]
Adding the two integrals without their coefficients gives the
normalizing integral. The posterior mean is thus $b/(b+c)$,
which is $d(X)$. The mean is in $[0,1]$, so
Proposition~\ref{prop:bayes-squared-mean} gives optimality also for
this restricted action space. Its constant risk makes it minimax
by Proposition~\ref{prop:bayes-minimax-certificate}.

The rule is also admissible. Any dominating rule would be Bayes
under the same prior. Each binary record has positive predictive
probability because its likelihood is positive on $(0,1)$.
The equality criterion for squared loss therefore forces any Bayes
rule to take action $d(x)$ with probability one at every record.
\end{example}

\subsection{Failures of attainment and equality}
\label{subsec:minimax-failures}

\begin{example}[A finite value without an optimal rule]
\label{ex:minimax-nonattainment}
Use one label and one observation point, Borel actions in $(0,1)$,
and loss $L(a)=a$. The constant actions $1/j$, $j\geq2$, have risks
tending to zero, so $V=0$. Every rule has strictly positive risk:
its action is positive with probability one, so its mean cannot
be zero. Hence no rule is minimax. Indeed, a rule of risk $r>0$
is dominated by the constant action $r/2$.
\end{example}

\begin{example}[A gap between Bayes and minimax values]
\label{ex:infinite-minimax-gap}
Use one observation point and take both labels and actions in
$\{1,2,\ldots\}$ with all subsets measurable. Set
$L(j,a)=\mathbf 1_{\{a\leq j\}}$. A rule is an action law $Q$, with
risk $Q(\{1,\ldots,j\})$ at label $j$. As $j\to\infty$ this risk
increases to one, so $W(K)=1$ for every rule and $V=1$.
For every prior $\pi$, the constant action $a$ has Bayes risk
$\pi(\{j:j\geq a\})$, which tends to zero as $a\to\infty$.
Nonnegativity then gives $r_\pi^\star=0$ for every $\pi$. Hence
$\sup_\pi r_\pi^\star=0<V$.
\end{example}

\subsection{Exercises}
\label{subsec:worked-examples-2-4}

\begin{practiceexercise}[Guessing one of three labels]
\label{wex:three-label-minimax}
With no data and unit error loss on three labels, find a minimax action law
and its value. Compare with the best deterministic value.
\end{practiceexercise}

\begin{practiceexercise}[Minimaxity with an erasure symbol]
\label{wex:erasure-minimax}
At either binary label, report the label correctly with probability
$r\in(0,1)$ and erase otherwise. Find a minimax rule for unit error loss
using the equal prior as a certificate.
\end{practiceexercise}

\begin{practiceexercise}[A bounded target without observations]
\label{wex:interval-no-data-minimax}
For target $\theta\in[1,5]$ and squared loss, use an equal prior on
endpoints to find a Bayes lower bound. Find a constant action whose
worst-case risk meets it.
\end{practiceexercise}

\begin{practiceexercise}[A minimax problem with equivalent labels]
\label{wex:equivalent-signs-minimax}
Let $\theta\in\{-2,0,2\}$ and $X\sim N(\theta^2,1)$, with squared loss for
$\theta$. Find a prior supported on the two indistinguishable labels and use
it to certify a minimax rule.
\end{practiceexercise}

\begin{practiceexercise}[Worst-case risk for four Bernoulli draws]
\label{wex:mean-versus-bernoulli-equalizer}
Compute the sample mean's worst squared risk and compare it with the
equalizer of Example~\ref{ex:bernoulli-minimax} at $n=4$. Explain why good
endpoint performance alone does not settle the comparison.
\end{practiceexercise}

\begin{practiceexercise}[Balancing endpoint and interior risk]
\label{wex:affine-one-bit-minimax}
With one Bernoulli observation and squared loss, derive the risk of
$d_w=wX+(1-w)/2$ for $0\le w\le1$. Find its worst case as a function of $w$
and minimize within this class.
\end{practiceexercise}

\begin{practiceexercise}[A symmetric signal]
\label{wex:symmetric-bit-minimax}
Under either binary label report it correctly with probability $r>1/2$ and
flip it otherwise. Under unit error loss and equal priors, find a Bayes
rule, compute its two risks, and determine the minimax value.
\end{practiceexercise}

\begin{practiceexercise}[Randomization at an asymmetric record]
\label{wex:asymmetric-bit-minimax}
At label zero let $X=a$ always; at label one let $X=a,b$ be equally likely.
Choose label one at $b$ and with probability $s$ at $a$. Under unit error
loss choose $s$ to equalize risks and find a prior that certifies the
resulting rule.
\end{practiceexercise}

\begin{practiceexercise}[Duplicating a label]
\label{wex:duplicated-label-least-favorable}
With no data, use labels $0_a,0_b,1$ and actions zero and one with unit
error loss for the underlying bit. Express the Bayes value in terms of the
prior and find all least favorable priors.
\end{practiceexercise}

\begin{practiceexercise}[Admissibility versus minimaxity]
\label{wex:admissible-nonminimax-posterior-mean}
With no data and squared loss for $\theta\in\{0,1\}$, give one prior mass
$1/4$. Find the Bayes action and show that it is admissible. Compare its
worst-case risk with that of action $1/2$.
\end{practiceexercise}

\begin{practiceexercise}[Bayes comparisons with infinite risk]
\label{wex:infinite-bayes-risk-warning}
With one record, two actions have risk vectors $(+\infty,0)$ and
$(+\infty,1)$. Use a prior giving positive mass to the first label. Find the
Bayes value and determine whether a rule attaining it can be dominated.
\end{practiceexercise}

\begin{practiceexercise}[A null set of strict improvement]
\label{wex:full-support-null-improvement}
With no data and uniform prior on $[0,1]$, let one action have zero loss
everywhere and the other have loss $\mathbf1_{\{\theta=0\}}$. Compare Bayes
risks, check domination, and identify the strict improvement set.
\end{practiceexercise}

\begin{practiceexercise}[A noncompact action space]
\label{wex:coercive-loss-attains}
With one label and record, use $L(a)=1+(a-2)^2$ on $\R$. Find the minimax
value and characterize every randomized rule attaining it.
\end{practiceexercise}

\begin{practiceexercise}[A continuous loss on the real line]
\label{wex:continuous-loss-at-infinity}
With one label and record, use $L(a)=e^{-a}$ on $\R$. Find the minimax value
and determine whether any probability law on actions attains it.
\end{practiceexercise}

\begin{practiceexercise}[A sequence of prior bounds]
\label{wex:countable-squared-least-favorable}
With no data, estimate $1/j$, $j\ge1$, by squared loss on $[0,1]$. Use
priors with equal mass on $1,m$ to find the minimax value. Determine whether
a least favorable prior attains it.
\end{practiceexercise}

\begin{exercise}[Restricting to deterministic rules]
\label{exer:deterministic-minimax-gap}
In Example~\ref{ex:finite-randomized-minimax}, restrict both the
minimax and Bayes optimizations to deterministic rules. Compute
the minimax value and the supremum of Bayes values over all priors.
Identify the property of $S$ in the proof of
Theorem~\ref{thm:finite-minimax-equality} that fails.
\end{exercise}

\begin{exercise}[Unequal least favorable weights]
\label{exer:unequal-least-favorable}
There is one observation point, two labels, and two actions with
loss pairs $(0,2)$ and $(3,0)$. Find the minimax mixing probabilities
and a least favorable prior. Verify both optimalities by computing
the risk pair of the mixture and the Bayes value under the prior.
\end{exercise}

\begin{exercise}[Minimaxity and prior support]
\label{exer:minimax-prior-support}
In Example~\ref{ex:dominated-minimax}, compute the Bayes value as a
function of the prior mass of label zero. Find all least favorable
priors and characterize their Bayes rules. Which rules are admissible?
\end{exercise}

\begin{exercise}[Gaussian location and a sequence of priors]
\label{exer:normal-minimax-priors}
For $n\geq1$ iid $N(\theta,1)$ observations, $\theta\in\R$, use
squared loss for estimating $\theta$. Under prior $N(0,\tau^2)$,
$\tau>0$, complete the square in the product of prior density and
likelihood. Show that the posterior variance is
$(n+\tau^{-2})^{-1}$ and hence equals the Bayes value.
Let $\tau\to\infty$ and use the risk of the sample mean to find $V$.
\end{exercise}

\begin{exercise}[Compact actions with a discontinuous loss]
\label{exer:compact-actions-nonattainment}
Change the action space in Example~\ref{ex:minimax-nonattainment}
to $[0,1]$, and set $L(0)=1$ and $L(a)=a$ for $a>0$.
Show that there is still no minimax rule, even with randomization.
\end{exercise}

\begin{exercise}[A supremum without a least favorable prior]
\label{exer:least-favorable-nonattainment}
Use one observation point, one available action, and labels
$j\in\{1,2,\ldots\}$ with all subsets measurable. At label $j$ the
loss of the single action is $1-1/j$. Compute $V$ and
$\sup_\pi r_\pi^\star$. Show that no prior attains the latter value.
\end{exercise}

\chapter{Information, divergence, and scoring rules}
\label{chap:information-scoring}
\notationgroup{Information and prediction}
\section{Entropy, cross-entropy, and logarithmic loss}
\label{sec:entropy-log-loss}

\subsection{Logarithmic loss on a countable space}
\label{subsec:discrete-log-loss}
\indexnotation{p(y),\ q(y),\ S_P}{Masses under \(P,Q\) and the discrete support of \(P\).}

Let $\mathsf Y$ be a nonempty finite or countable set with all subsets
measurable. For laws $P,Q$ on $\mathsf Y$, write $p(y)=P(\{y\})$ and
$q(y)=Q(\{y\})$, and let $S_P=\{y:p(y)>0\}$ be the support in this
discrete setting. All logarithms are natural; the unit is the nat.
Dividing by $\log 2$ gives the corresponding quantities in bits.

\begin{definition}[Logarithmic loss]
\label{def:log-loss}
\indexnotation{L_{\log}(y,Q)}{Logarithmic loss of forecast \(Q\) at outcome \(y\).}
A probability forecast $Q$ incurs \emph{logarithmic loss}
\[
L_{\log}(y,Q)=-\log q(y)
\]
when the outcome is $y$. We set $-\log 0=+\infty$. With the sigma
algebra generated by the coordinate maps $Q\mapsto q(y)$ on the
forecast space, this is a measurable action section for each $y$.
\end{definition}

\begin{definition}[Entropy and cross-entropy]
\label{def:discrete-entropy-cross-entropy}
\indexnotation{H(P),\ H(P,Q)}{Discrete entropy and cross-entropy.}
The \emph{cross-entropy} of $P$ relative to $Q$, and the \emph{entropy}
of $P$, are
\[
H(P,Q)=\mathbb E_P[L_{\log}(Y,Q)]
=\sum_{y\in S_P}p(y)(-\log q(y)),
\qquad H(P)=H(P,P).
\]
Here $Y$ has law $P$. The sums are nonnegative and may equal $+\infty$.
Outcomes with $p(y)=0$ contribute zero, including when $q(y)=0$.
If $p(y)>0$ and $q(y)=0$ for some $y$, then $H(P,Q)=+\infty$.
\end{definition}

\begin{example}[Forecasting a binary outcome]
\label{ex:bernoulli-cross-entropy}
For the Bernoulli laws of Example~\ref{ex:bernoulli},
\[
H(P_p,P_q)=-p\log q-(1-p)\log(1-q),\qquad
H(P_p)=-p\log p-(1-p)\log(1-p),
\]
with zero-weight terms omitted. For a fair outcome, $p=1/2$, the
forecast $q=1/2$ has expected loss $\log 2$. The forecast $q=1/4$
has expected loss $\tfrac12\log(16/3)>\log 2$. At $q=0$ the loss
is infinite on outcome one, which has probability $1/2$, so the
expected loss is infinite.
\end{example}

\begin{example}[Infinite expected loss with full support]
\label{ex:infinite-cross-entropy-full-support}
On $\{1,2,\ldots\}$ let $p(j)=1/(Zj^2)$, with
$Z=\sum_{j\geq1}j^{-2}$, and $q(j)=2^{-j}$. Both laws are positive
at every outcome. The entropy of $P$ is finite, since
\[
H(P)=\log Z+\frac{2}{Z}\sum_{j\geq1}\frac{\log j}{j^2}<\infty
\]
by the integral test for the tail. On the other hand,
\[
H(P,Q)=\frac{\log 2}{Z}\sum_{j\geq1}\frac1j=+\infty,
\]
since the loss at outcome $j$ grows as $j\log 2$.
\end{example}

\subsection{Entropy bounds and product forecasts}
\label{subsec:entropy-bounds-products}

\begin{proposition}[Elementary bounds for entropy]
\label{prop:discrete-entropy-bounds}
For any discrete law $P$, $H(P)\geq0$, with equality exactly for point
masses. If $|S_P|=s<\infty$, then $H(P)\leq\log s$, with equality
exactly when $P$ is uniform on $S_P$. On an alphabet of size $m$ this
gives $H(P)\leq\log m$, with equality exactly for the uniform law
on the whole alphabet.
\end{proposition}
\begin{proof}
Each positive-mass term $p(y)(-\log p(y))$ is nonnegative and vanishes
exactly when $p(y)=1$. This proves the assertion about zero entropy.

For $t>0$, the function $t-1-\log t$ decreases up to $1$ and increases
after $1$, so
\begin{equation}
\log t\leq t-1,\qquad\text{with equality exactly at }t=1.
\label{eq:log-tangent-bound}
\end{equation}
If $|S_P|=s<\infty$, apply this bound at $t=1/(sp(y))$ to get
\[
H(P)-\log s
=\sum_{y\in S_P}p(y)\log\frac1{sp(y)}
\leq\sum_{y\in S_P}\left(\frac1s-p(y)\right)=0.
\]
Equality requires $p(y)=1/s$ at each $y\in S_P$. For an alphabet of
size $m$, equality at $\log m$ also requires $s=m$.
\end{proof}

\begin{proposition}[Expected log loss for a product forecast]
\label{prop:product-forecast-cross-entropy}
Let $R$ be a law on two countable outcome spaces with marginals
$P_1,P_2$. For forecasts $Q_1,Q_2$ on the respective spaces,
\[
H(R,Q_1\otimes Q_2)=H(P_1,Q_1)+H(P_2,Q_2).
\]
This holds also when a term is infinite. In particular,
$H(P_1\otimes P_2)=H(P_1)+H(P_2)$.
\end{proposition}
\begin{proof}
The mass assigned by the product forecast at $(y_1,y_2)$ is
$q_1(y_1)q_2(y_2)$. Its negative log is
$-\log q_1(y_1)-\log q_2(y_2)$, including the case of a zero factor.
Taking expectations under $R$, each summand depends on only one
coordinate. Nonnegative integration gives the stated sum even for
infinite values. Taking $R=Q_1\otimes Q_2=P_1\otimes P_2$ gives
the entropy identity.
\end{proof}

\begin{example}[Two independent bits or one copied bit]
\label{ex:entropy-copied-bit}
Let $U$ be the fair law on $\{0,1\}$. For two independent bits the joint
law is $Q=U\otimes U$ and $H(Q)=2\log 2$. If a single fair bit is
copied into both coordinates, the joint law $R$ assigns mass $1/2$
to each of $(0,0)$ and $(1,1)$, so $H(R)=\log 2$. The two laws have
the same marginals, as in Example~\ref{ex:dependence}.

Forecasting the copied pair with $Q$ incurs $2\log 2$ on each record
with positive $R$-mass, so $H(R,Q)=2\log 2$. Reversing the roles gives
$H(Q,R)=+\infty$, because $Q$ gives positive probability to unequal
bits and $R$ assigns those records zero mass.
\end{example}

\subsection{Empirical cross-entropy and likelihood}
\label{subsec:empirical-cross-entropy}

\begin{proposition}[Average log loss on a record]
\label{prop:empirical-cross-entropy}
For a record $y_{1:n}=(y_1,\ldots,y_n)$, $n\geq1$, let
$\widehat P_n$ be its empirical law. Then
\begin{equation}
H(\widehat P_n,Q)
=\frac1n\sum_{i=1}^n L_{\log}(y_i,Q)
=-\frac1n\log Q^{\otimes n}(\{y_{1:n}\}).
\label{eq:empirical-cross-entropy-likelihood}
\end{equation}
Over a parametrized family $(Q_\theta)$, the minimizers of the left
side are precisely the likelihood maximizers for the iid experiment.
\end{proposition}
\begin{proof}
By Definition~\ref{def:empirical-law}, the mass of $y$ under
$\widehat P_n$ is the number of occurrences of $y$ divided by $n$.
Grouping the summands of average log loss by outcome gives the
cross-entropy. The iid probability of the ordered record is
$\prod_i q(y_i)$, so taking its negative log gives the other equality.
If a forecast assigns zero probability to an observed outcome, both
sides are infinite. The map $t\mapsto -n^{-1}\log t$ is strictly
decreasing on $[0,1]$ with the extended value at zero, so it reverses
the order of all likelihood values, including at the boundary.
\end{proof}

\begin{example}[Four binary observations]
\label{ex:empirical-binary-cross-entropy}
For the record $(1,0,1,1)$, the empirical probability of one is $3/4$.
Under an iid Bernoulli forecast with parameter $q$, the empirical
cross-entropy is
\[
-\tfrac34\log q-\tfrac14\log(1-q).
\]
The likelihood of the record is $q^3(1-q)$, so the unique minimizing
forecast is $q=3/4$ by Example~\ref{ex:bernoulli-likelihood}.
At this forecast the cross-entropy is
$\tfrac34\log(4/3)+\tfrac14\log 4$. The fair forecast gives $\log 2$,
and each of the forecasts $q=0,1$ gives infinite loss.
\end{example}

\subsection{Density-based logarithmic loss}
\label{subsec:density-log-loss}

\begin{definition}[Entropy relative to a reference measure]
\label{def:density-entropy}
\indexnotation{h_\mu(P),\ H_\mu(P,Q)}{Entropy and cross-entropy relative to a reference measure.}
Let $P,Q$ have finite density versions $f,g$ with respect to a common
$\sigma$-finite measure $\mu$. Define
\[
H_\mu(P,Q)=\int(-\log g)\,dP,\qquad
h_\mu(P)=H_\mu(P,P),
\]
when the indicated integral is defined. Here it may be finite,
$+\infty$, or $-\infty$; it is undefined if its positive and negative
parts both have infinite integral. For Lebesgue measure,
$h_\mu(P)$ is called \emph{differential entropy}.
\end{definition}

Changing a density on a $\mu$-null set does not change these integrals
because $P\ll\mu$. A density can exceed one, so $-\log g$ can be negative.
This differs from logarithmic loss for probability masses, which
always takes values in $[0,\infty]$.

\begin{proposition}[Changing the scale of the reference measure]
\label{prop:rescale-entropy-reference}
Let $\nu=c\mu$, where $0<c<\infty$. Then $H_\nu(P,Q)$ is defined
exactly when $H_\mu(P,Q)$ is defined, and in that case
\[
H_\nu(P,Q)=H_\mu(P,Q)+\log c.
\]
The same statements hold for $h_\nu(P)$ and $h_\mu(P)$.
\end{proposition}
\begin{proof}
A density $g$ with respect to $\mu$ becomes $g/c$ with respect to $\nu$.
Hence $-\log(g/c)=-\log g+\log c$, also on $\{g=0\}$. Adding a
finite constant bounds each part of the new integrand by the
corresponding part of the old one plus the absolute value of the
constant, and conversely. Since $P$ is a
probability measure, the integral is defined in exactly the same
cases, and its value shifts by $\log c$ whenever it is defined.
Taking $Q=P$ gives the assertion for entropy.
\end{proof}

\begin{example}[Differential entropy of a uniform law]
\label{ex:uniform-differential-entropy}
Let $P$ be uniform on $[0,a]$, $a>0$, as in
Example~\ref{ex:uniform-endpoint}. Relative to Lebesgue measure
$\lambda$, its density is $1/a$ on the interval. Thus
\[
h_\lambda(P)=-\int_{[0,a]}\log(1/a)\,P(dx)=\log a.
\]
For $a<1$ this is negative. With reference measure
$\nu=a^{-1}\lambda$ the density is one on the interval, so
$h_\nu(P)=0$. The probability law $P$ is the same in both calculations.
\end{example}

\subsection{Exercises}
\label{subsec:worked-examples-3-1}

\begin{practiceexercise}[Comparing losses for one outcome]
\label{wex:rare-outcome-bits}
Two forecasts give the observed outcome probabilities $1/8$ and $1/2$.
Compute both logarithmic losses in nats and bits. Express their difference
as the log of a probability ratio.
\end{practiceexercise}

\begin{practiceexercise}[Adding unused symbols]
\label{wex:uniform-entropy-unused-alphabet}
Find the entropy of the uniform law on eight symbols. Embed it in an
alphabet of sixteen symbols with zero mass on the new ones. Compare its
entropy and the alphabet-size bound before and after.
\end{practiceexercise}

\begin{practiceexercise}[Forecasting a deterministic outcome]
\label{wex:deterministic-truth-forecast}
For $P=\delta_a$, compute $H(P)$ and $H(P,Q)$ when $q(a)=1/100$ and when
$q(a)=0$.
\end{practiceexercise}

\begin{practiceexercise}[Entropy of three unequal masses]
\label{wex:entropy-in-bits}
Compute the entropy in bits of $(1/2,1/4,1/4)$. Compare it with the entropy
of a uniform law on three symbols.
\end{practiceexercise}

\begin{practiceexercise}[Cross-entropy in both directions]
\label{wex:finite-cross-entropy-asymmetry}
For $P=(1/2,1/4,1/4)$ and $Q$ uniform on three symbols, compute and compare
$H(P,Q)$ and $H(Q,P)$.
\end{practiceexercise}

\begin{practiceexercise}[Appending a known coordinate]
\label{wex:deterministic-coordinate-entropy}
For $Y\sim P$ and fixed $c$, compute the entropy of $(Y,c)$ from its joint
masses and from the product entropy identity.
\end{practiceexercise}

\begin{practiceexercise}[Parity and joint entropy]
\label{wex:parity-entropy}
Draw two independent fair bits and let the third be their sum modulo two.
Find the joint entropy in bits, check whether pairs are independent, and
compute the cross-entropy against the product of all three marginals.
\end{practiceexercise}

\begin{practiceexercise}[Forecasting the second bit from the first]
\label{wex:conditional-bit-forecast}
The true law copies a fair bit. Compare the independent fair forecast with a
forecast assigning mass $9/20$ to each equal pair and $1/20$ to each unequal
pair. Find both cross-entropies and the second forecast's conditional match
probability.
\end{practiceexercise}

\begin{practiceexercise}[A forecast built from the correct marginals]
\label{wex:biased-marginal-product-entropy}
Give $(0,0),(1,0),(1,1)$ masses $1/2,1/4,1/4$, forming law $R$. Compute
$H(R)$ and its cross-entropy against the product of its marginals. Find
their difference.
\end{practiceexercise}

\begin{practiceexercise}[Smoothing an empirical forecast]
\label{wex:smoothing-zero-masses}
On $m$ symbols use $(1-\epsilon)\widehat P_n+\epsilon U$ for uniform $U$ and
$0<\epsilon<1$. Give a uniform upper bound on log loss and compute the loss
on any unobserved symbol.
\end{practiceexercise}

\begin{practiceexercise}[Six observations in three categories]
\label{wex:categorical-record-cross-entropy}
Use the record $(a,b,a,c,a,b)$ and forecast $(1/2,1/4,1/4)$ in
order $a,b,c$. Compute average log loss directly and from the iid likelihood
of the ordered record.
\end{practiceexercise}

\begin{practiceexercise}[Average loss for a Poisson forecast]
\label{wex:poisson-record-cross-entropy}
For count record $(0,2)$, express average log loss as a function of the
positive Poisson forecast rate. Find its minimizer and compare it with the
sample mean.
\end{practiceexercise}

\begin{practiceexercise}[Forecasting a waiting time]
\label{wex:exponential-cross-entropy}
The true and forecast laws are exponential with rates $\lambda,\rho>0$.
Compute their cross-entropy relative to Lebesgue measure and minimize it
over $\rho$. Derive the differential entropy and find when it is negative.
\end{practiceexercise}

\begin{practiceexercise}[Uniform laws with different supports]
\label{wex:uniform-support-cross-entropy}
For $P=\operatorname{Unif}[0,2]$ and $Q=\operatorname{Unif}[0,1]$, compute
both directions of cross-entropy relative to Lebesgue measure. Explain the
role of the support.
\end{practiceexercise}

\begin{practiceexercise}[Scaling a continuous random variable]
\label{wex:change-coordinate-differential-entropy}
For $X$ with density $f$ and $Y=cX$, $c>0$, derive the density of $Y$ and
the differential entropy relation when defined. Evaluate for uniform $X$ on
$[0,1/2]$ and $c=10$.
\end{practiceexercise}

\begin{exercise}[The binary entropy curve]
\label{exer:binary-entropy}
For $p\in[0,1]$, compute the derivative of $H(P_p)$ on $(0,1)$.
Show that the curve is symmetric about $1/2$, continuous at $0,1$,
and uniquely maximized at $1/2$.
\end{exercise}

\begin{exercise}[A law with infinite entropy]
\label{exer:infinite-discrete-entropy}
Take masses $p(j)$ proportional to $1/(j(\log j)^2)$
for $j\geq2$. Show that they can be normalized and that the resulting
entropy is infinite. Compare the entropy
summands for large $j$ with $1/(j\log j)$.
\end{exercise}

\begin{exercise}[Geometric forecasts]
\label{exer:geometric-cross-entropy}
On $\{1,2,\ldots\}$ let $P$ assign mass $2^{-j}$ to outcome $j$.
For $0<r<1$ let $Q_r$ assign mass $(1-r)r^{j-1}$. Compute
$H(P,Q_r)$ using $\mathbb E_P[Y-1]$, find its minimizer in $r$,
and give $H(P)$ in nats and in bits.
\end{exercise}

\begin{exercise}[Empirical categorical forecasts]
\label{exer:empirical-categorical-log-loss}
On an alphabet of size $m\geq2$, show that the empirical law
$\widehat P_n$ uniquely minimizes $H(\widehat P_n,Q)$ over all
forecasts. For $Q$ positive at every observed outcome, apply
\eqref{eq:log-tangent-bound} to $q(y)/\widehat p_n(y)$ there.
If some outcomes are unobserved, compute the infimum over forecasts
that are strictly positive on the whole alphabet and determine
whether it is attained.
\end{exercise}

\begin{exercise}[Gaussian cross-entropy]
\label{exer:gaussian-cross-entropy}
For $P=N(u,s^2)$ and $Q=N(v,t^2)$ with $s,t>0$, compute
$H_\lambda(P,Q)$ relative to Lebesgue measure. Use the normal density
and $\mathbb E_P[(Y-v)^2]=s^2+(u-v)^2$. By taking $Q=P$, derive
$h_\lambda(P)$ and determine when it is negative.
\end{exercise}

\begin{exercise}[A density with undefined differential entropy]
\label{exer:undefined-differential-entropy}
Set $f(x)=1/(x|\log x|^2)$ for
$x\in(0,e^{-2})\cup(e^2,\infty)$, and set $f(x)=0$ elsewhere.
Show that this is a probability density and that both the positive
and negative parts of $-\log f$ have infinite integral under this
law. Use the substitution $t=|\log x|$ separately on each interval.
\end{exercise}


\clearpage
\section{Relative entropy and divergence inequalities}
\label{sec:relative-entropy}

\subsection{Relative entropy and excess logarithmic loss}
\label{subsec:relative-entropy-definition}

\begin{definition}[Relative entropy]
\label{def:relative-entropy}
\indexnotation{D(P\|Q)}{Relative entropy of \(P\) with respect to \(Q\).}
Let $P,Q$ be probability laws on the same measurable space. Put
$\varphi(t)=t\log t-t+1$ for $t\geq0$, with $0\log 0=0$.
The \emph{relative entropy}, or \emph{Kullback--Leibler divergence},
of $P$ with respect to $Q$ is
\[
D(P\|Q)=
\begin{cases}
\displaystyle\int\varphi\!\left(\frac{dP}{dQ}\right)dQ,
 & P\ll Q,\\[4pt]
+\infty, & P\not\ll Q.
\end{cases}
\]
Here $P\ll Q$ denotes absolute continuity. On $(0,\infty)$ the
derivative of $\varphi$ is $\log t$, so its minimum is $\varphi(1)=0$.
At zero $\varphi(0)=1$. The integral is therefore nonnegative and
well defined in $[0,\infty]$.
\end{definition}

\begin{proposition}[Likelihood-ratio and density formulas]
\label{prop:relative-entropy-density}
One has $D(P\|Q)=0$ if and only if $P=Q$. If $P\ll Q$ and
$r=dP/dQ$, then $\log r$ is defined $P$-almost surely and
\[
D(P\|Q)=\int\log r\,dP,\qquad \int(\log r)_-\,dP\leq 1/e,
\]
where $u_-=\max\{-u,0\}$. More generally, if $f,g$ are densities of
$P,Q$ with respect to a common $\sigma$-finite measure $\mu$, then
\[
D(P\|Q)=\int f\log(f/g)\,d\mu .
\]
In this integrand a zero numerator contributes zero, and a positive
numerator with zero denominator contributes $+\infty$. Its negative
part has integral at most $1/e$.
\end{proposition}
\begin{proof}
If $P\ll Q$, then $\int r\,dQ=1$ and $P(r=0)=0$.
On $0<t<1$, the maximum of $-t\log t$ is $1/e$. Hence
\[
\int(\log r)_-\,dP
=\int_{\{0<r<1\}}-r\log r\,dQ\leq 1/e.
\]
Changing measure and integrating $\varphi(r)=r\log r-r+1$
gives the likelihood-ratio formula. Since $\varphi$ vanishes only at
one, the integral of $\varphi(r)$ is zero exactly when $r=1$
$Q$-almost surely, or equivalently $P=Q$.

For the density formula, the negative part of $f\log(f/g)$ is
bounded by $g/e$, by applying the same bound to $f/g$ on
$\{0<f<g\}$. If $P(g=0)=0$, then $f/g$ on $\{g>0\}$ is a
version of $dP/dQ$, as integration over any measurable set shows.
If $P(g=0)>0$, then $\{f>0=g\}$ has positive $\mu$-measure and
zero $Q$-measure. Both the integral and $D(P\|Q)$ are then infinite.
\end{proof}

\begin{proposition}[Excess logarithmic loss]
\label{prop:relative-entropy-cross-entropy}
For laws on a finite or countable set,
\[
H(P,Q)=H(P)+D(P\|Q).
\]
For laws with densities relative to $\mu$ as in
Definition~\ref{def:density-entropy}, if $h_\mu(P)$ is finite, then
$H_\mu(P,Q)$ is defined and
\[
H_\mu(P,Q)=h_\mu(P)+D(P\|Q).
\]
Thus $Q=P$ uniquely minimizes expected logarithmic loss when the
entropy of $P$ is finite. In the discrete case, if $H(P)=+\infty$,
then every forecast has infinite expected loss.
\end{proposition}
\begin{proof}
For the discrete claim take counting measure and write
\[
-\log q(y)=-\log p(y)+\log(p(y)/q(y))
\]
on $\{p>0\}$. The first term on the right is nonnegative.
The second has $P$-integrable negative part by
Proposition~\ref{prop:relative-entropy-density}, so the integrals add,
including at $+\infty$. For general densities, the same decomposition
holds $P$-almost surely with $p,q$ replaced by $f,g$. The term
$-\log f$ is $P$-integrable by hypothesis, so the negative part of
the sum is again integrable and the same argument applies.
The minimizer assertions follow from the identity and the equality
condition for $D$.
\end{proof}

\begin{example}[Bernoulli forecasts]
\label{ex:bernoulli-relative-entropy}
For the laws of Example~\ref{ex:bernoulli-cross-entropy}, the excess
loss from using $P_{1/4}$ when the true law is $P_{1/2}$ is
\[
D(P_{1/2}\|P_{1/4})
=\tfrac12\log 2+\tfrac12\log(2/3)=\tfrac12\log(4/3).
\]
Reversing the laws gives
\[
D(P_{1/4}\|P_{1/2})
=\tfrac14\log(1/2)+\tfrac34\log(3/2)=\tfrac34\log 3-\log 2.
\]
The first divergence exceeds the second by $\tfrac14\log(256/243)$.
\end{example}

\begin{example}[Nested uniform laws]
\label{ex:nested-uniform-relative-entropy}
Let $P$ be uniform on $[0,1]$ and $Q$ uniform on $[0,2]$.
The density ratio $dP/dQ$ equals two $P$-almost surely, so
$D(P\|Q)=\log 2$. Since $P((1,2])=0$ and $Q((1,2])=1/2$,
one has $D(Q\|P)=+\infty$.
\end{example}

\begin{example}[Normal laws]
\label{ex:normal-relative-entropy}
Let $P=N(u,s^2)$ and $Q=N(v,t^2)$, where $s,t>0$. Their log density
ratio at $y$ is
\[
\log(t/s)+\frac{(y-v)^2}{2t^2}-\frac{(y-u)^2}{2s^2}.
\]
Under $P$ the two squared terms have expectations
$s^2+(u-v)^2$ and $s^2$, respectively. Consequently
\[
D(P\|Q)=\log(t/s)+\frac{s^2+(u-v)^2}{2t^2}-\tfrac12.
\]
In particular, at common variance $s^2=t^2$ this is
$(u-v)^2/(2s^2)$.
\end{example}

\begin{example}[Infinite divergence with full support]
\label{ex:infinite-relative-entropy-full-support}
Use the laws $p(j)=1/(Zj^2)$ and $q(j)=2^{-j}$ from
Example~\ref{ex:infinite-cross-entropy-full-support}. They are
mutually absolutely continuous. Since $H(P)<\infty$ and
$H(P,Q)=+\infty$, Proposition~\ref{prop:relative-entropy-cross-entropy}
gives $D(P\|Q)=+\infty$.
\end{example}

\subsection{The log-sum inequality and joint convexity}
\label{subsec:log-sum-joint-convexity}

\begin{proposition}[Log-sum inequality]
\label{prop:log-sum-inequality}
For nonnegative numbers $a_i,b_i$, $1\leq i\leq m$, put
$A=\sum_i a_i$ and $B=\sum_i b_i$. With the conventions of
Proposition~\ref{prop:relative-entropy-density},
\[
\sum_{i=1}^m a_i\log(a_i/b_i)\ \geq\ A\log(A/B).
\]
If $B>0$ and $a_i=0$ whenever $b_i=0$, equality holds precisely when
$a_i/b_i$ is constant over the indices with $b_i>0$.
\end{proposition}
\begin{proof}
If some $a_i>0=b_i$, the left side is infinite. In the remaining
case $B=0$ implies that all terms are zero. Suppose $B>0$.
The function $t\mapsto t\log t$ is strictly convex on $[0,\infty)$:
its second derivative is positive on $(0,\infty)$, and for
$0<\alpha<1$, $t>0$,
$(\alpha t)\log(\alpha t)<\alpha t\log t$.
Apply Jensen's inequality to the ratios $a_i/b_i$ with weights
$b_i/B$ on the indices with $b_i>0$. Their weighted mean is $A/B$.
Multiplying by $B$ yields the inequality, with equality exactly
when the ratios coincide by strict convexity.
\end{proof}

\begin{proposition}[Joint convexity]
\label{prop:relative-entropy-joint-convexity}
Let $P_i,Q_i$ be probability laws on one measurable space and let
$w_i>0$, $1\leq i\leq m$, with $\sum_i w_i=1$. Then
\[
D\!\left(\sum_i w_i P_i\,\middle\|\,\sum_i w_i Q_i\right)
\leq\sum_i w_iD(P_i\|Q_i).
\]
\end{proposition}
\begin{proof}
Use the common dominating measure $\mu=\sum_i(P_i+Q_i)$ and
write $f_i,g_i$ for the densities. Apply the log-sum inequality
pointwise to $a_i=w_if_i$ and $b_i=w_ig_i$. Integrate and use
Proposition~\ref{prop:relative-entropy-density}. All the signed
terms have integrable negative parts.
\end{proof}

\begin{example}[Strict inequality under mixing]
\label{ex:strict-relative-entropy-convexity}
For Bernoulli laws, $\tfrac12P_{1/4}+\tfrac12P_{3/4}=P_{1/2}$.
With $Q=P_{1/2}$, both $D(P_{1/4}\|Q)$ and $D(P_{3/4}\|Q)$
are positive. Their average strictly exceeds the divergence
$D(\tfrac12P_{1/4}+\tfrac12P_{3/4}\|Q)=0$.
\end{example}

\subsection{Total variation and Pinsker's inequality}
\label{subsec:pinsker}

\begin{definition}[Total variation distance]
\label{def:total-variation-distance}
\indexnotation{d_{\mathrm{TV}}(P,Q)}{Total variation distance between two laws.}
For laws $P,Q$ on a measurable space $(\mathsf Y,\mathcal Y)$, define
\[
d_{\mathrm{TV}}(P,Q)=\sup_{A\in\mathcal Y}|P(A)-Q(A)|.
\]
\end{definition}

\begin{proposition}[Density formula for total variation]
\label{prop:total-variation-density}
If $P,Q$ have densities $f,g$ relative to a common $\sigma$-finite
measure $\mu$, then
\[
d_{\mathrm{TV}}(P,Q)=\tfrac12\int |f-g|\,d\mu.
\]
The set $\{f\geq g\}$ attains the supremum in
Definition~\ref{def:total-variation-distance}.
\end{proposition}
\begin{proof}
The integrals of $(f-g)_+$ and $(f-g)_-$ agree since
$\int(f-g)\,d\mu=0$. Each equals $\tfrac12\int |f-g|\,d\mu$.
For any $A$, the absolute integral of $f-g$ over $A$ is bounded
by that common value. On $\{f\geq g\}$ the integral attains it.
\end{proof}

\begin{theorem}[Pinsker's inequality]
\label{thm:pinsker}
For probability laws $P,Q$ on the same measurable space,
\[
D(P\|Q)\geq 2d_{\mathrm{TV}}(P,Q)^2.
\]
\end{theorem}
\begin{proof}
The claim is immediate for $D(P\|Q)=+\infty$. Otherwise $P\ll Q$.
Fix a measurable $A$ and put $p=P(A)$ and $q=Q(A)$. If $q=0$
or $q=1$, then $p=q$. For $0<q<1$, define
\[
b_q(p)=p\log(p/q)+(1-p)\log((1-p)/(1-q)),\qquad 0\leq p\leq1 .
\]
Apply Jensen's inequality to the convex function $\varphi$ under
$Q$ conditioned on $A$ and on $A^c$, with argument $r=dP/dQ$.
Its two conditional means are $p/q$ and $(1-p)/(1-q)$. Thus
\[
D(P\|Q)\geq q\varphi(p/q)+(1-q)\varphi((1-p)/(1-q))=b_q(p).
\]
On $0<p<1$ we have $b_q''(p)=1/[p(1-p)]\geq4$, while
$b_q(q)=b_q'(q)=0$. Hence $b_q(p)-2(p-q)^2$ is convex and
has a minimum of zero at $p=q$. By continuity the bound also
holds at $p=0,1$. Taking the supremum over $A$ gives the result.
\end{proof}

\begin{example}[Sharpness of the constant]
\label{ex:pinsker-sharpness}
For $0<|\varepsilon|<1/2$, take $P=P_{1/2+\varepsilon}$ and
$Q=P_{1/2}$. The two singleton mass differences are $\varepsilon$
and $-\varepsilon$, so $d_{\mathrm{TV}}(P,Q)=|\varepsilon|$.
The function $b_{1/2}$ in the proof of Theorem~\ref{thm:pinsker}
has value and derivative zero at $1/2$ and second derivative four
there. Taylor's formula gives
\[
D(P\|Q)=b_{1/2}(1/2+\varepsilon)
=2\varepsilon^2+o(\varepsilon^2).
\]
Taking the ratio to $d_{\mathrm{TV}}(P,Q)^2$ and letting
$\varepsilon\to0$ shows that the constant two is optimal.
\end{example}

\subsection{Exercises}
\label{subsec:practice-3-2}

\begin{practiceexercise}[Point mass]
\label{pex:relative-entropy-point-mass}
For a countable alphabet and $P=\delta_y$, compute $D(P\|Q)$ and
$D(Q\|P)$ in terms of $q(y)$, including $q(y)=0$ and $q(y)=1$.
\end{practiceexercise}

\begin{practiceexercise}[Three outcomes]
\label{pex:relative-entropy-three-outcomes}
Let $P=(1/2,1/3,1/6)$ and $Q=(1/3,1/3,1/3)$. Compute both
directions of the divergence.
\end{practiceexercise}

\begin{practiceexercise}[Zero coordinates]
\label{pex:relative-entropy-zero-coordinates}
For $P=(2/3,0,1/3,0)$ and $Q=(1/2,1/4,1/4,0)$, determine which
direction of absolute continuity holds and compute both divergences.
\end{practiceexercise}

\begin{practiceexercise}[Geometric laws]
\label{pex:relative-entropy-geometric}
On the positive integers let $P_r(j)=(1-r)r^{j-1}$ for $0<r<1$.
Use the mean of $j-1$ under $P_r$ to compute $D(P_r\|P_s)$.
\end{practiceexercise}

\begin{practiceexercise}[Exponential rates]
\label{pex:relative-entropy-exponential}
For $\lambda,\eta>0$, compute the divergence from the law with
density $\lambda e^{-\lambda x}$ on $[0,\infty)$ to the law with
density $\eta e^{-\eta x}$ on the same half-line.
\end{practiceexercise}

\begin{practiceexercise}[Relabeling]
\label{pex:relative-entropy-relabeling}
Prove that applying the same bijection of a countable alphabet
to both laws preserves their divergence.
\end{practiceexercise}

\begin{practiceexercise}[Rescaling the reference]
\label{pex:relative-entropy-reference-scale}
If $\nu=c\mu$ with $c>0$, show directly in the density formula
that replacing $\mu$ by $\nu$ leaves $D(P\|Q)$ unchanged.
\end{practiceexercise}

\begin{practiceexercise}[Uniform forecast]
\label{pex:relative-entropy-uniform-forecast}
For a law $P$ on an alphabet of size $m$ and the uniform law $U$,
express $D(P\|U)$ in terms of $H(P)$ and $m$.
\end{practiceexercise}

\begin{practiceexercise}[Applying log-sum]
\label{pex:log-sum-two-terms}
Evaluate the two sides of Proposition~\ref{prop:log-sum-inequality}
for $a=(2,1)$ and $b=(1,3)$. Determine whether equality holds.
\end{practiceexercise}

\begin{practiceexercise}[Log-sum with zero masses]
\label{pex:log-sum-zero-masses}
With $b=(1,0,2)$, find all $a\in[0,\infty)^3$ for which the two
sides of the log-sum inequality are finite and equal.
\end{practiceexercise}

\begin{practiceexercise}[Convexity in the forecast]
\label{pex:relative-entropy-forecast-mixture}
Let $P=P_{1/2}$, $Q_1=P_{1/4}$, and $Q_2=P_{1/3}$. Compute the
divergence to the equal mixture of $Q_1,Q_2$ and compare it to
the average of the two divergences.
\end{practiceexercise}

\begin{practiceexercise}[An event as a binary outcome]
\label{pex:relative-entropy-binary-event}
For $P=(1/2,1/3,1/6)$ and $Q=(1/3,1/3,1/3)$, let
$A=\{1,2\}$. Compute the divergence of the laws of $\mathbf 1_A$
under $P$ and $Q$. Compare it to $D(P\|Q)$ using log-sum.
\end{practiceexercise}

\begin{practiceexercise}[A maximizing event]
\label{pex:tv-three-outcomes}
Compute $d_{\mathrm{TV}}$ between $(1/2,1/3,1/6)$ and
$(1/4,1/4,1/2)$. List every event attaining the supremum.
\end{practiceexercise}

\begin{practiceexercise}[Uniform intervals]
\label{pex:tv-nested-intervals}
For the uniform laws on $[0,a]$ and $[0,b]$ with $0<a<b$,
compute their total variation distance and find an interval
attaining it.
\end{practiceexercise}

\begin{practiceexercise}[A Poisson event]
\label{pex:tv-poisson-zero}
Use the event $\{0\}$ to give a lower bound on the total
variation distance between Poisson laws with rates one and two.
\end{practiceexercise}

\begin{practiceexercise}[From divergence to event probabilities]
\label{pex:pinsker-event-bound}
Suppose $D(P\|Q)\leq 1/50$ and $Q(A)=1/4$. Give the interval
for $P(A)$ implied by Theorem~\ref{thm:pinsker}.
\end{practiceexercise}

\begin{practiceexercise}[Normal mean shift]
\label{pex:pinsker-normal-mean-shift}
Apply Theorem~\ref{thm:pinsker} to $N(u,s^2)$ and $N(v,s^2)$.
For which values of $|u-v|/s$ does the resulting upper bound
improve on the bound one?
\end{practiceexercise}

\begin{practiceexercise}[Disjoint supports]
\label{pex:disjoint-relative-entropy-tv}
Suppose $P(A)=1$ and $Q(A)=0$ for a measurable $A$.
Compute $d_{\mathrm{TV}}(P,Q)$ and both directions of relative entropy.
\end{practiceexercise}

\begin{exercise}[Infinite entropy and finite relative entropy]
\label{exer:relative-entropy-infinite-entropy}
Let $P$ have masses $p(j)=c/[j(\log j)^2]$ for $j\geq2$, with
$c$ the normalizing constant. Prove $H(P)=+\infty$. For
$0<\alpha<1$ set $Q=(1-\alpha)P+\alpha\delta_2$. Show that
$0<D(P\|Q)\leq-\log(1-\alpha)$ and compute it in terms of
$p(2)$ and $\alpha$. Determine $H(P,Q)$.
\end{exercise}

\begin{exercise}[Undefined entropy and defined divergence]
\label{exer:relative-entropy-undefined-entropy}
Use the law $P$ of Exercise~\ref{exer:undefined-differential-entropy}.
Let $A=(0,e^{-2})$ and define $Q$ by $dQ/dP=1/2$ on $A$ and
$dQ/dP=3/2$ on $A^c$. Show that $Q$ is a probability law, compute
$D(P\|Q)$, and determine whether $H_\lambda(P,Q)$ is defined
for Lebesgue measure $\lambda$.
\end{exercise}

\begin{exercise}[Best variance at a prescribed mean]
\label{exer:relative-entropy-normal-variance}
For fixed $P=N(u,s^2)$ and fixed $v$, minimize
$D(P\|N(v,t^2))$ over $t>0$. Compute the minimum and show
that the minimizing variance is unique.
\end{exercise}

\begin{exercise}[Likelihood regret]
\label{exer:relative-entropy-likelihood-regret}
Let $y_1,\ldots,y_n$ be a nonempty sample on a finite alphabet,
and let $\widehat P_n$ be its empirical law. Express the difference
between the log-likelihoods at $\widehat P_n$ and at $Q$ as a
multiple of a divergence. Treat forecasts assigning zero
probability to an observed outcome.
\end{exercise}

\begin{exercise}[Finite partitions]
\label{exer:relative-entropy-finite-partitions}
Let $P\ll Q$, $D(P\|Q)<\infty$, and let $(A_i)_{i=1}^m$ be a
measurable partition. Prove that
\[
D(P\|Q)\geq\sum_i P(A_i)\log\frac{P(A_i)}{Q(A_i)}.
\]
Show that equality holds precisely when $dP/dQ$ is
$Q$-almost surely constant on each cell of positive $Q$-mass.
\end{exercise}

\begin{exercise}[Equality in Pinsker's inequality]
\label{exer:pinsker-equality}
Prove that equality in Theorem~\ref{thm:pinsker} holds only
when $P=Q$. Use the maximizing event from
Proposition~\ref{prop:total-variation-density} with $\mu=P+Q$,
and establish the strict binary inequality for $p\ne q$.
\end{exercise}



\clearpage
\section{Conditional information, chain rules, and data processing}
\label{sec:conditional-information}

\subsection{Conditional laws and entropy}
\label{subsec:conditional-entropy}

For random elements $X,Y,Z$ on one probability space, write $P_X$ for
the law of $X$ and $P_{XY}$ for the law of the ordered pair $(X,Y)$.
Use the same convention for longer tuples.

\begin{definition}[Conditional law]
\label{def:conditional-law}
\indexnotation{P_X,\ P_{XY},\ P_{Y\mid X}}{Marginal law, joint law, and conditional law kernel.}
A \emph{conditional law} of $Y$ given $X$ is a probability kernel
$P_{Y\mid X}$ from the value space of $X$ to that of $Y$ such that
\[
P_{XY}=P_X\ltimes P_{Y\mid X}.
\]
The joint law induced by a kernel is defined in
Proposition~\ref{prop:joint-action-law}.
\end{definition}

A sigma algebra is countably generated if it is generated by a countable
family of sets.

\begin{definition}[Conditional entropy] of a countable outcome]
\label{def:conditional-entropy}
\indexnotation{H(Y\mid X)}{Mean entropy of the conditional law of a countable outcome.}
If $Y$ takes values in a countable set with all subsets measurable, define
\[
H(Y\mid X)=\int H(P_{Y\mid X}(x,\cdot))\,P_X(dx).
\]
Use $H(Y)=H(P_Y)$ and, when both variables are countable-valued,
$H(X,Y)=H(P_{XY})$. The integrand in the conditional entropy is a
countable sum of nonnegative measurable terms. Existence and uniqueness
up to $P_X$-null sets follow from Proposition~\ref{prop:conditional-law-versions}.
\end{definition}

\begin{proposition}[Versions of a conditional law]
\label{prop:conditional-law-versions}
If the sigma algebra of the value space of $Y$ is countably
generated, any two conditional laws of $Y$ given $X$ agree as
measures for $P_X$-almost every $x$. If $Y$ takes values in a
countable set with all subsets measurable, a conditional law
exists for every $X$.
\end{proposition}
\begin{proof}
For each set $B$ in a countable generating $\pi$-system containing
the whole value space of $Y$, the rectangle identity implies that
the two functions $x\mapsto P_{Y\mid X}(x,B)$ agree $P_X$-almost
everywhere. Off the union of these null sets the two measures agree
on the generating system, and hence agree on the sigma algebra.

For existence in the countable case, let $k_y$ be the nonnegative
Radon--Nikodym density of the measure $A\mapsto P_{XY}(A\times\{y\})$
with respect to $P_X$. Then $\sum_y k_y=1$ $P_X$-almost everywhere,
by summing the measures and using uniqueness of the density.
Use these masses on the set where they sum to one and a fixed point
mass elsewhere. Summing the rectangle identities over $y\in B$
gives the required kernel.
\end{proof}

\begin{proposition}[Entropy chain rule]
\label{prop:discrete-entropy-chain-rule}
For countable-valued $X,Y$,
\[
H(X,Y)=H(X)+H(Y\mid X)
\]
in $[0,\infty]$.
\end{proposition}
\begin{proof}
At every pair $(x,y)$ of positive joint mass, that mass is
$p_X(x)P_{Y\mid X}(x,\{y\})$. Its negative log is the sum of the
two negative logs. Take expectations and integrate the second
term first over $y$. Both terms are nonnegative.
\end{proof}

\begin{example}[A selector and a fair bit]
\label{ex:conditional-entropy-selector}
Let $X,B$ be independent fair bits and set $Y=XB$.
At $X=0$ the conditional law of $Y$ is a point mass, and at $X=1$
it is fair. Thus
\[
H(Y\mid X)=\tfrac12\log 2,\qquad
H(X,Y)=\tfrac32\log 2.
\]
Indeed, the three positive joint masses are $1/2,1/4,1/4$.
\end{example}

\subsection{The relative-entropy chain rule}
\label{subsec:relative-entropy-chain-rule}

\begin{theorem}[Relative-entropy chain rule]
\label{thm:relative-entropy-chain-rule}
Let $P,Q$ be laws on a common measurable space and $K,L$ kernels from
that space to $(\mathsf Y,\mathcal Y)$, with $\mathcal Y$ countably
generated. Write $K_x=K(x,\cdot)$ and $L_x=L(x,\cdot)$. The function
$x\mapsto D(K_x\|L_x)$ has a measurable version up to $(P+Q)$-null sets,
and
\[
D(P\ltimes K\|Q\ltimes L)=D(P\|Q)+\int D(K_x\|L_x)\,P(dx).
\]
All terms are in $[0,\infty]$.
\end{theorem}
\begin{proof}
Put $S=(P+Q)/2$, $N=(K+L)/2$, and $M=S\ltimes N$. Choose densities
$f=d(S\ltimes K)/dM$ and $g=d(S\ltimes L)/dM$ in $[0,2]$.
For each $B$ in a countable generating $\pi$-system containing
$\mathsf Y$, the rectangle identity gives
\[
\int_B f(x,y)\,N(x,dy)=K(x,B)
\qquad\text{for }S\text{-almost every }x.
\]
The same assertion holds for $g$ and $L$. Taking a common
full-measure set and using uniqueness on the generating system,
$f(x,\cdot),g(x,\cdot)$ are densities of $K_x,L_x$ relative to $N_x$
there. By Proposition~\ref{prop:relative-entropy-density},
the inner integral of $f\log(f/g)$ under $N(x,\cdot)$ represents
$D(K_x\|L_x)$ on this set. Its negative part has integral at most
$2/e$ for every $x$, since it is bounded pointwise by $g/e$.
The two parts have measurable integrals by
\eqref{eq:kernel-iterated-integral}. Setting the inner integral to zero
off the full-measure set gives the measurable version.

If $P\not\ll Q$, a rectangle $A\times\mathsf Y$ with $P(A)>0=Q(A)$
makes the joint divergence infinite, and the marginal term is
also infinite. Suppose $P\ll Q$. Write $a=dP/dS$ and $b=dQ/dS$.
The joint densities of $P\ltimes K,Q\ltimes L$ relative to $M$
are $a(x)f(x,y),b(x)g(x,y)$, by the iterated integral formula.
The log density ratio under $P\ltimes K$ is therefore
\[
\log(a(x)/b(x))+\log(f(x,y)/g(x,y))
\qquad (P\ltimes K)\text{-almost surely}.
\]
Here $a(x),b(x),f(x,y)>0$ almost surely under $P\ltimes K$,
and the second term is $+\infty$ when $g(x,y)=0$.
The first term has integrable negative part under $P$.
The second has integrable negative part under $P\ltimes K$
by the fiberwise bound $1/e$. Integrating the two terms gives
the stated sum, by Proposition~\ref{prop:relative-entropy-density}
and \eqref{eq:kernel-iterated-integral}.
\end{proof}

\begin{example}[A marginal discrepancy and a conditional discrepancy]
\label{ex:relative-entropy-chain-bits}
Let $X$ be binary with $P(X=1)=1/2$ and $Q(X=1)=1/4$.
Under $Q$, the conditional law of $Y$ is fair at each $x$.
Under $P$, it is fair at $x=0$ and concentrated at one at $x=1$.
The chain rule and Example~\ref{ex:bernoulli-relative-entropy} give
\[
D(P_{XY}\|Q_{XY})
=\tfrac12\log(4/3)+\tfrac12\log 2.
\]
Only the row $x=1$ contributes to the conditional term. Its
divergence from the fair law is $\log 2$.
\end{example}

\subsection{Mutual information and conditional mutual information}
\label{subsec:mutual-information}

\begin{definition}[Mutual information]
\label{def:mutual-information}
\indexnotation{I(X;Y)}{Mutual information between two random elements.}
For arbitrary random elements $X,Y$ define
\[
I(X;Y)=D(P_{XY}\|P_X\otimes P_Y).
\]
\end{definition}

\begin{definition}[Conditional mutual information]
\label{def:conditional-mutual-information}
\indexnotation{I(X;Y\mid Z)}{Mean mutual information under the conditional law of the pair.}
Suppose the value spaces of $X,Y$ have countably generated sigma algebras
and a conditional law $J=P_{(X,Y)\mid Z}$ exists. Let $A_z,B_z$ be the
marginals of $J(z,\cdot)$. Define
\[
I(X;Y\mid Z)=\int D(J(z,\cdot)\|A_z\otimes B_z)\,P_Z(dz).
\]
The kernel $z\mapsto A_z\otimes B_z$ is measurable on rectangles and
hence on the product sigma algebra by the $\pi$--$\lambda$ theorem.
Use the measurable version of the divergence supplied by
Theorem~\ref{thm:relative-entropy-chain-rule}. The integral is
independent of the version of $J$ by
Proposition~\ref{prop:conditional-law-versions}.
The variables are \emph{conditionally independent given $Z$} if
$J(z,\cdot)=A_z\otimes B_z$ for $P_Z$-almost every $z$.
\end{definition}

\begin{proposition}[Information, independence, and entropy]
\label{prop:mutual-information-entropy}
Mutual information is symmetric and vanishes exactly when the
variables are independent. Conditional mutual information vanishes
exactly when the variables are conditionally independent.
If the value space of $Y$ has a countably generated sigma algebra
and $Y$ admits a conditional law given $X$, then
\[
I(X;Y)=\int D(P_{Y\mid X}(x,\cdot)\|P_Y)\,P_X(dx).
\]
If $Y$ is countable-valued and $X$ is arbitrary, then
\[
H(Y)=H(Y\mid X)+I(X;Y)
\]
in $[0,\infty]$.
\end{proposition}
\begin{proof}
A measurable bijection with measurable inverse preserves relative
entropy. It preserves absolute continuity, and in the absolutely
continuous case the density after the bijection is the original
density composed with its inverse. This follows by integrating over
inverse images. Applying this fact to coordinate swap gives symmetry.
The zero criteria follow from Proposition~\ref{prop:relative-entropy-density}
and, in the conditional case, nonnegativity of the integrand.

The integral formula follows from the relative-entropy chain rule
against $P_X\otimes P_Y$. For countable $Y$, apply Proposition~\ref{prop:relative-entropy-cross-entropy} to each
conditional law against $P_Y$ and integrate. The integrated
cross-entropy is $H(Y)$, since the conditional masses average to
the marginal masses. The entropy term is $H(Y\mid X)$.
\end{proof}

\begin{theorem}[Mutual-information chain rule]
\label{thm:mutual-information-chain-rule}
Under the hypotheses of Definition~\ref{def:conditional-mutual-information},
assume also that $Y$ admits a conditional law given $(Z,X)$. Then
\[
I(X;(Z,Y))=I(X;Z)+I(X;Y\mid Z).
\]
\end{theorem}
\begin{proof}
Work in coordinate order $(Z,X,Y)$. Let $C=P_{Y\mid(Z,X)}$
and $\widetilde B((z,x),\cdot)=B_z$, using the marginal kernel
in Definition~\ref{def:conditional-mutual-information}. Put
\[
R=P_{ZX}\ltimes C,\qquad
Q=(P_Z\otimes P_X)\ltimes\widetilde B,\qquad
S=P_{ZX}\ltimes\widetilde B .
\]
The law $Q$ is the reordering of $P_X\otimes P_{ZY}$.
Also $S=P_Z\ltimes (z\mapsto A_z\otimes B_z)$.
Both identities follow on rectangles, using $A_z,B_z$ as conditional
marginals. The relative-entropy chain rule applied to $R,Q$ and
to $R,S$ with first coordinate $(Z,X)$ yields
\[
D(R\|Q)=D(P_{ZX}\|P_Z\otimes P_X)+D(R\|S).
\]
Applied to $R=P_Z\ltimes J$ and $S$ with first coordinate $Z$,
it gives $D(R\|S)=I(X;Y\mid Z)$. Reordering $R,Q$ gives
$D(R\|Q)=I(X;(Z,Y))$.
\end{proof}

\begin{example}[Conditioning on parity]
\label{ex:mutual-information-parity}
Let $X,Y$ be independent fair bits and $Z=(X+Y)\bmod 2$.
Then $I(X;Y)=0$. Given either parity, the pair is uniform on two
points, with both marginals fair. The density ratio to the product
of those marginals equals two on the two points, so
$I(X;Y\mid Z)=\log 2$. Also $X$ is independent of $Z$, hence
Theorem~\ref{thm:mutual-information-chain-rule} gives
$I(X;(Z,Y))=\log 2$.
\end{example}

\begin{example}[A common cause]
\label{ex:mutual-information-common-cause}
Let $Z$ be a fair bit. Given $Z$, let $X,Y$ be independent bits,
each equal to $Z$ with probability $3/4$. Then $I(X;Y\mid Z)=0$.
Unconditionally both bits are fair and
\[
\Pr(X=Y)=(3/4)^2+(1/4)^2=5/8 .
\]
By the symmetry under simultaneous bit flips, each diagonal
pair has mass $5/16$ and each off-diagonal pair has mass $3/16$.
Thus $I(X;Y)=\tfrac58\log(5/4)+\tfrac38\log(3/4)>0$.
\end{example}

\begin{example}[Copying a continuous variable]
\label{ex:mutual-information-continuous-copy}
Let $X$ be uniform on $[0,1]$ and $Y=X$. The diagonal has joint
mass one. Under $P_X\otimes P_Y$ its mass is
$\int P_Y(\{x\})\,P_X(dx)=0$ by Tonelli's theorem. The joint law
is therefore singular with respect to the product law and
$I(X;Y)=+\infty$.
\end{example}

\subsection{Data processing}
\label{subsec:information-data-processing}

\begin{definition}[Markov chain of three random elements]
\label{def:markov-chain-three}
\indexnotation{X\to Y\to Z}{Markov chain in which the last conditional law depends on the middle variable.}
Write $X\to Y\to Z$ if there is a probability kernel $K$ from the
value space of $Y$ to that of $Z$ with
$P_{XYZ}=P_{XY}\ltimes\widetilde K$, where
$\widetilde K((x,y),\cdot)=K(y,\cdot)$.
\end{definition}

\begin{theorem}[Data processing for relative entropy]
\label{thm:relative-entropy-data-processing}
For laws $P,Q$ on a common measurable space and any probability
kernel $K$ from that space, one has
\[
D(PK\|QK)\leq D(P\|Q).
\]
In particular, for a measurable map $T$,
$D(P\circ T^{-1}\|Q\circ T^{-1})\leq D(P\|Q)$.
\end{theorem}
\begin{proof}
It suffices to treat $D(P\|Q)<\infty$. First consider a measurable
map $T$ and put $P_T=P\circ T^{-1}$ and $Q_T=Q\circ T^{-1}$.
Then $P_T\ll Q_T$. For $r=dP/dQ$ and $s=dP_T/dQ_T$,
$s\circ T$ is a version of $\mathbb E_Q[r\mid\sigma(T)]$, as is
seen by integrating over inverse images under $T$. We use
conditional Jensen's inequality for the integrable $r$ and
convex $\varphi$ of Definition~\ref{def:relative-entropy},
with $\varphi(r)$ integrable under $Q$. It gives
\[
D(P_T\|Q_T)=\int\varphi(s\circ T)\,dQ
\leq\int\varphi(r)\,dQ=D(P\|Q).
\]
For a kernel $K$, $P\ltimes K$ has density $(x,y)\mapsto r(x)$
relative to $Q\ltimes K$, by \eqref{eq:kernel-iterated-integral}.
Their divergence is thus $D(P\|Q)$. Apply the map case to projection
onto the second coordinate.
\end{proof}

\begin{proposition}[Data processing for mutual information]
\label{prop:mutual-information-data-processing}
If $X\to Y\to Z$, then $I(X;Z)\leq I(X;Y)$.
\end{proposition}
\begin{proof}
Use the kernel from $(X,Y)$ to $(X,Z)$ that keeps $x$ and draws
$z$ from $K(y,\cdot)$ in Definition~\ref{def:markov-chain-three}.
Its kernel property follows on rectangles and then on the product
sigma algebra by the $\pi$--$\lambda$ theorem. It sends $P_{XY}$
to $P_{XZ}$ and $P_X\otimes P_Y$ to $P_X\otimes P_Z$, as
seen on rectangles. Apply Theorem~\ref{thm:relative-entropy-data-processing}
to those two input laws.
\end{proof}

\begin{example}[Two binary channels]
\label{ex:mutual-information-binary-channels}
Let $X$ be fair, and let independent noise bits $E,F$ also be
independent of $X$, with $\Pr(E=1)=\varepsilon$ and
$\Pr(F=1)=\delta$. Set $Y=(X+E)\bmod2$ and $Z=(Y+F)\bmod2$.
Then $X\to Y\to Z$, and $Y,Z$ are both fair. Given $X$, the
probabilities that $Y$ and $Z$ differ from $X$ are
$\varepsilon$ and $\rho=\varepsilon+\delta-2\varepsilon\delta$.
For $h(p)=H(P_p)$, Proposition~\ref{prop:mutual-information-entropy}
gives
\[
I(X;Y)=\log 2-h(\varepsilon),\qquad I(X;Z)=\log 2-h(\rho).
\]
If $0<\varepsilon,\delta<1/2$, then
$\varepsilon<\rho<1/2$. Since $h'(p)=\log((1-p)/p)>0$ on
$(0,1/2)$, the second channel strictly decreases the information.
\end{example}

\subsection{Exercises}
\label{subsec:practice-3-3}

\begin{practiceexercise}[Reading a joint table]
\label{pex:conditional-law-joint-table}
The rows and columns of a joint mass table are indexed by $0,1$.
The row for $X=0$ is $(1/2,1/4)$ and that for $X=1$ is $(1/8,1/8)$.
Find $P_X$ and the two rows of $P_{Y\mid X}$.
\end{practiceexercise}

\begin{practiceexercise}[A null row]
\label{pex:conditional-law-null-row}
Let $X=0$ almost surely and let $Y$ be a fair bit. Treat the
value space of $X$ as $\{0,1\}$. Describe all kernels
$P_{Y\mid X}$ and compute $H(Y\mid X)$.
\end{practiceexercise}

\begin{practiceexercise}[Entropy in both orders]
\label{pex:conditional-entropy-both-orders}
Use the table of Exercise~\ref{pex:conditional-law-joint-table}
to compute $H(Y\mid X)$ and $H(X\mid Y)$. Verify the chain rule in
both orders.
\end{practiceexercise}

\begin{practiceexercise}[Independent outcomes]
\label{pex:conditional-entropy-independent}
For independent $X$ and countable-valued $Y$, give an explicit
kernel for $Y$ given $X$. Show that $H(Y\mid X)=H(Y)$.
\end{practiceexercise}

\begin{practiceexercise}[A deterministic countable outcome]
\label{pex:conditional-entropy-deterministic}
If $Y=T(X)$ for a measurable map $T$ into a countable set, prove
that $H(Y\mid X)=0$ and compute $I(X;Y)$ in terms of $H(Y)$.
\end{practiceexercise}

\begin{practiceexercise}[Entropy of an unbounded count]
\label{pex:conditional-entropy-count}
Let $X$ be a fair bit. Given $X=0$, put $Y=1$, and given $X=1$
let $Y$ have masses $2^{-j}$ on the positive integers.
Compute $H(Y\mid X)$ and $H(X,Y)$.
\end{practiceexercise}

\begin{practiceexercise}[Two independent coordinates]
\label{pex:relative-entropy-products}
Use the relative-entropy chain rule to compute
$D(P_1\otimes P_2\|Q_1\otimes Q_2)$ for countable second
coordinates. Include the case in which a marginal divergence
is infinite.
\end{practiceexercise}

\begin{practiceexercise}[A singular row]
\label{pex:relative-entropy-singular-row}
Let $P_X=Q_X$ be fair. Under $P$, $Y=X$, and under $Q$, $Y=0$
regardless of $X$. Use the chain rule to compute the joint
divergence in both directions.
\end{practiceexercise}

\begin{practiceexercise}[An absent row under the true law]
\label{pex:relative-entropy-absent-row}
Let $P_X=\delta_0$ and let $Q_X$ be fair on $\{0,1\}$. For a binary
second coordinate, set $K_0=L_0$ and take $K_1,L_1$ to be
different point masses. Compute $D(P_X\ltimes K\|Q_X\ltimes L)$.
\end{practiceexercise}

\begin{practiceexercise}[A common second-coordinate kernel]
\label{pex:relative-entropy-common-kernel}
For $P=(3/4,1/4)$ and $Q=(1/4,3/4)$, let $K$ output the same
fair law at either input. Compute $D(P\ltimes K\|Q\ltimes K)$
and $D(PK\|QK)$.
\end{practiceexercise}

\begin{practiceexercise}[An information bound]
\label{pex:mutual-information-alphabet-bound}
If $Y$ takes at most $m$ values, prove $I(X;Y)\leq\log m$.
Characterize equality using $P_Y$ and $P_{Y\mid X}$.
\end{practiceexercise}

\begin{practiceexercise}[A constant conditioning variable]
\label{pex:conditional-mutual-constant}
For finite $X,Y$ and constant $Z$, compute $I(X;Y\mid Z)$ in
terms of the unconditional joint and marginal laws.
\end{practiceexercise}

\begin{practiceexercise}[Conditioning on the observed variable]
\label{pex:conditional-mutual-observed}
For finite $X,Y$, prove $I(X;Y\mid X)=0$ by factoring the
conditional joint law.
\end{practiceexercise}

\begin{practiceexercise}[A biased copied bit]
\label{pex:mutual-information-biased-copy}
Let $X$ be a Bernoulli bit with parameter $p$ and $Y=X$.
Compute $I(X;Y)$ by summing the two diagonal contributions.
Include $p=0,1$.
\end{practiceexercise}

\begin{practiceexercise}[A quantized copy]
\label{pex:mutual-information-quantized-copy}
Let $X$ be uniform on $[0,1)$ and $Y=\lfloor mX\rfloor$ for an
integer $m\geq2$. Compute $I(X;Y)$. If $Z=Y\bmod2$ and $m$ is
even, compute $I(X;Z)$.
\end{practiceexercise}

\begin{practiceexercise}[Applying log-sum to a map]
\label{pex:information-processing-log-sum}
On a finite space, prove the map case of
Theorem~\ref{thm:relative-entropy-data-processing} by applying
log-sum separately to each fiber of the map.
\end{practiceexercise}

\begin{practiceexercise}[A noisy bit near complete noise]
\label{pex:mutual-information-complete-noise}
For a fair input bit and independent flip probability
$\varepsilon$, determine where $I(X;Y)$ attains its maximum
and minimum on $[0,1]$. Include the meaning of $\varepsilon=1$.
\end{practiceexercise}

\begin{practiceexercise}[Extra independent randomness]
\label{pex:mutual-information-independent-noise}
Let $X,Y,U$ be finite-valued with $U$ independent of $(X,Y)$.
Prove $I(X;(Y,U))=I(X;Y)$ using the mutual-information chain rule.
\end{practiceexercise}

\begin{exercise}[The best product forecast]
\label{exer:mutual-information-product-forecast}
For a joint law $R=P_{XY}$ on a finite alphabet and forecasts
$Q_X,Q_Y$, derive
\[
D(R\|Q_X\otimes Q_Y)
=I(X;Y)+D(P_X\|Q_X)+D(P_Y\|Q_Y)
\]
from the cross-entropy identities in
Sections~\ref{sec:entropy-log-loss} and \ref{sec:relative-entropy}.
Find the unique minimizing pair $(Q_X,Q_Y)$.
\end{exercise}

\begin{exercise}[Chain identities with infinite entropy]
\label{exer:mutual-information-infinite-copy}
Let $X$ have the law of
Exercise~\ref{exer:relative-entropy-infinite-entropy}, and put $Y=X$.
Compute $H(X,Y)$, $H(Y\mid X)$, and $I(X;Y)$, and determine
which differences of joint and marginal entropies are undefined.
\end{exercise}

\begin{exercise}[Additive Gaussian noise]
\label{exer:mutual-information-normal-channel}
Let $X\sim N(0,s^2)$ and $E\sim N(0,t^2)$ be independent with
$s,t>0$, and put $Y=X+E$. Use the conditional law $N(x,t^2)$
and the marginal $N(0,s^2+t^2)$ of $Y$ to compute $I(X;Y)$.
Determine its limits as $t\to0$ and $t\to\infty$ with $s$ fixed.
\end{exercise}

\begin{exercise}[Equality under finite coarsening]
\label{exer:relative-entropy-processing-equality}
For strictly positive laws $P,Q$ on a finite space and a map $T$,
show that $D(P\circ T^{-1}\|Q\circ T^{-1})=D(P\|Q)$ exactly
when $p/q$ is constant on each nonempty fiber of $T$. Give a
nonconstant, noninjective $T$ and distinct laws attaining equality.
\end{exercise}

\begin{exercise}[Information lost along a Markov chain]
\label{exer:mutual-information-processing-gap}
For a Markov chain $X\to Y\to Z$ of finite-valued variables, show that
\[
I(X;Y)=I(X;Z)+I(X;Y\mid Z).
\]
Derive an equality criterion for $I(X;Z)\leq I(X;Y)$ by applying
the mutual-information chain rule to the pair $(Y,Z)$ in both orders.
\end{exercise}

\begin{exercise}[A sequence of observations]
\label{exer:mutual-information-sequence}
For finite $X,Y_1,\ldots,Y_n$, derive the sum of successive conditional
information terms for $I(X;(Y_1,\ldots,Y_n))$. Apply it when the
$Y_i$ are independent fair bits and $X$ is their parity. Determine
which summands vanish and compute the remaining one.
\end{exercise}



\clearpage
\section{Proper scoring rules and probability estimation}
\label{sec:proper-scoring-rules}

\subsection{Propriety and regret}
\label{subsec:score-propriety}

Let $\mathsf Y$ be a nonempty finite or countable alphabet with all
subsets measurable. Equip its laws with the sigma algebra generated
by the coordinate maps $Q\mapsto q(y)$ and each forecast class with
the induced sigma algebra.

\begin{definition}[Scoring rule and propriety]
\label{def:proper-scoring-rule}
\indexnotation{S(y,Q),\ s_S(P,Q)}{Score as a loss at outcome \(y\) and expected score under \(P\).}
A \emph{scoring rule} on a class $\mathcal P$ of laws on $\mathsf Y$
is a loss $S(y,Q)\in[0,\infty]$ for $y\in\mathsf Y$ and
$Q\in\mathcal P$, measurable in $Q$ for each $y$. For any law $P$ write
\[
s_S(P,Q)=\sum_{y:p(y)>0}p(y)S(y,Q).
\]
The rule is \emph{proper} on $\mathcal P$ if
$s_S(P,P)\leq s_S(P,Q)$ for all $P,Q\in\mathcal P$. It is
\emph{strictly proper} if the inequality is strict for $Q\ne P$.
\end{definition}

\begin{definition}[Truthful score and regret]
\label{def:scoring-regret}
\indexnotation{B_S(P),\ \Delta_S(P,Q)}{Expected score at the true law and excess expected score at the forecast.}
For a proper rule set $B_S(P)=s_S(P,P)$. When $B_S(P)<\infty$,
the \emph{regret} of $Q$ at $P$ is
\[
\Delta_S(P,Q)=s_S(P,Q)-B_S(P)\in[0,\infty].
\]
\end{definition}

\begin{proposition}[Concavity of the truthful score]
\label{prop:proper-score-concavity}
Suppose $\mathcal P$ is convex, $S$ is proper on $\mathcal P$,
and $B_S$ is finite on $\mathcal P$. Then $B_S$ is concave.
If $S$ is strictly proper, $B_S$ is strictly concave.
\end{proposition}
\begin{proof}
For $0<t<1$ put $P_t=(1-t)P_0+tP_1$. By linearity in the true law
and propriety,
\[
B_S(P_t)
=(1-t)s_S(P_0,P_t)+t\,s_S(P_1,P_t)
\geq (1-t)B_S(P_0)+t B_S(P_1).
\]
If $P_0\ne P_1$, then $P_t$ differs from both endpoints, so both
comparisons are strict under strict propriety.
\end{proof}

\begin{proposition}[Affine supports on a finite alphabet]
\label{prop:proper-score-affine-support}
Let $\mathsf Y$ be finite, and let $S$ be a finite score on the whole
probability simplex. Put $A_Q(P)=s_S(P,Q)$ and $B(P)=s_S(P,P)$.
Then $S$ is proper exactly when, for every $Q$, the affine function
$A_Q$ satisfies $A_Q\geq B$ on the simplex and $A_Q(Q)=B(Q)$.
It is strictly proper exactly when each such support touches $B$
only at $Q$.

Conversely, let $B$ be a finite nonnegative concave function on the
simplex and let $A_Q$ be a finite affine support from above at each
$Q$. If $Q\mapsto A_Q(\delta_y)$ is measurable for every $y$, then
$S(y,Q)=A_Q(\delta_y)$ is a proper score with truthful score $B$.
\end{proposition}
\begin{proof}
For each $Q$, the expectation $s_S(P,Q)$ is affine in $P$.
Propriety is the inequality $B(P)\leq A_Q(P)$, and equality
holds at $P=Q$ by definition. Strict propriety gives exactly the
stated equality condition.

Conversely, $A_Q(\delta_y)\geq B(\delta_y)\geq0$. By affinity,
$\sum_y p(y)A_Q(\delta_y)=A_Q(P)$, so the score has expectation
$A_Q(P)$ and truthful score $A_P(P)=B(P)$. The support inequality
proves propriety.
\end{proof}

\begin{example}[Reporting only the mode]
\label{ex:proper-score-mode}
For binary outcomes define $m(Q)=\mathbf 1_{\{q(1)>1/2\}}$ and
$S(y,Q)=\mathbf 1_{\{y\ne m(Q)\}}$. Predicting zero gives expected
loss $p(1)$ and predicting one gives $1-p(1)$. The forecast $Q=P$
selects a minimizing label, so the score is proper. At $p(1)=3/4$,
every forecast with $q(1)>1/2$ gives the same minimum $1/4$.
\end{example}

\subsection{Logarithmic and quadratic scores}
\label{subsec:logarithmic-quadratic-scores}

The logarithmic score is the loss $L_{\log}$ from
Definition~\ref{def:log-loss}.

\begin{definition}[Quadratic score]
\label{def:quadratic-score}
\indexnotation{S_{\mathrm{Br}}(y,Q)}{Quadratic probability score, also called the Brier score.}
The \emph{quadratic score}, or \emph{Brier score}, is
\[
S_{\mathrm{Br}}(y,Q)=
\sum_{a\in\mathsf Y}\bigl(q(a)-\mathbf 1_{\{a=y\}}\bigr)^2.
\]
In the binary case this convention gives $2(q(1)-y)^2$ for
$y\in\{0,1\}$.
\end{definition}

\begin{proposition}[Propriety of the logarithmic and quadratic scores]
\label{prop:log-quadratic-strictly-proper}
The logarithmic score is proper on all laws on $\mathsf Y$ and
strictly proper on the class of laws of finite entropy. At any
$P$ with $H(P)<\infty$ its regret against any forecast $Q$ is
$D(P\|Q)$.

The quadratic score lies in $[0,2]$ and is strictly proper on all
laws on $\mathsf Y$. Its truthful score and regret are
\[
B_{\mathrm{Br}}(P)=1-\sum_y p(y)^2,\qquad
\Delta_{\mathrm{Br}}(P,Q)=\sum_y(p(y)-q(y))^2.
\]
Here the subscript $\mathrm{Br}$ denotes $S_{\mathrm{Br}}$.
\end{proposition}
\begin{proof}
For the logarithmic score, $s_S(P,Q)=H(P,Q)$ and
$B_S(P)=H(P)$. Apply Proposition~\ref{prop:relative-entropy-cross-entropy}
and the zero criterion for relative entropy.

For the quadratic score, expansion gives
$S_{\mathrm{Br}}(y,Q)=1+\sum_a q(a)^2-2q(y)$. The sum is at most
one because $q(a)^2\leq q(a)$, which gives the bound and
measurability. Taking expectations and completing the square gives
\[
s_{S_{\mathrm{Br}}}(P,Q)
=1-\sum_y p(y)^2+\sum_y(p(y)-q(y))^2.
\]
All sums in the expansion are finite. The squared difference
vanishes exactly when the masses agree at every outcome.
\end{proof}

\begin{example}[A uniform forecast for three outcomes]
\label{ex:brier-score-three-outcomes}
Let $P=(1/2,1/3,1/6)$ and $Q=(1/3,1/3,1/3)$. The truthful
quadratic score is
\[
B_{\mathrm{Br}}(P)=1-\left(\tfrac14+\tfrac19+\tfrac1{36}\right)
=\tfrac{11}{18}.
\]
The regret is $(1/6)^2+0^2+(-1/6)^2=1/18$, giving an
expected score of $2/3$ at the uniform forecast.
\end{example}

\subsection{Conditional forecasts}
\label{subsec:proper-conditional-forecasts}

\begin{definition}[Risk of a probability forecast]
\label{def:probability-forecast-risk}
\indexnotation{\mathcal R_S(Q)}{Population score risk of a conditional forecast kernel.}
For a random pair $(X,Y)$ with $Y$ taking values in $\mathsf Y$,
a \emph{probability forecast based on $X$} is a probability kernel
$Q$ from the value space of $X$ to $\mathsf Y$. Write
$Q_x=Q(x,\cdot)$ and $P_x=P_{Y\mid X}(x,\cdot)$.
For a score whose forecast class contains the values of $Q_x$, set
\[
\mathcal R_S(Q)=\mathbb E[S(Y,Q(X,\cdot))].
\]
The integrand is measurable because $Y$ is countable and
$x\mapsto Q_x$ is measurable in each coordinate.
\end{definition}

\begin{theorem}[Optimal conditional forecasts]
\label{thm:proper-conditional-forecast}
Let $S$ be proper on $\mathcal P$. Suppose the kernels
$x\mapsto P_x,Q_x$ take values in $\mathcal P$ and that
$B_S(P_x)<\infty$ for $P_X$-almost every $x$. Then
\[
\mathcal R_S(Q)=\mathcal R_S(P_{Y\mid X})
 +\int\Delta_S(P_x,Q_x)\,P_X(dx).
\]
Thus $P_{Y\mid X}$ minimizes population risk. If $S$ is strictly
proper and $\mathcal R_S(P_{Y\mid X})<\infty$, the minimizers are
exactly the forecasts equal to $P_{Y\mid X}$ for $P_X$-almost every $x$.
\end{theorem}
\begin{proof}
By \eqref{eq:kernel-iterated-integral}, the risk of $Q$ is the integral
of $s_S(P_x,Q_x)$ under $P_X$. For almost every $x$ this is
$B_S(P_x)+\Delta_S(P_x,Q_x)$. Each expected score is measurable as
a countable sum, with zero-weight terms contributing zero. The
regret is measurable on the set where $B_S(P_x)<\infty$ and can be
set to zero elsewhere. Integration gives the identity.
If the optimal risk is finite, equality holds precisely when the
nonnegative regret vanishes almost surely. For a strictly proper
score, this is exactly the asserted agreement of forecasts.
\end{proof}

For the logarithmic score, Proposition~\ref{prop:relative-entropy-cross-entropy}
gives the decomposition on the whole forecast space,
\[
\mathcal R_{L_{\log}}(Q)
=H(Y\mid X)+\int D(P_x\|Q_x)\,P_X(dx).
\]
In particular, if $H(Y\mid X)$ is finite, equality with the optimal
risk requires $Q_x=P_x$ almost surely. If $H(Y\mid X)=+\infty$,
every forecast has infinite risk.

\begin{example}[A feature that improves a probability forecast]
\label{ex:brier-score-conditional-selector}
In Example~\ref{ex:conditional-entropy-selector}, $Y=XB$ for
independent fair bits $X,B$. The conditional probabilities of one are
$0$ at $X=0$ and $1/2$ at $X=1$. The optimal conditional quadratic
risk is therefore $\tfrac12\cdot0+\tfrac12\cdot\tfrac12=1/4$.
Among constant forecasts, the optimal probability of one is $1/4$,
the marginal probability. Its quadratic risk is
$2(1/4)(3/4)=3/8$, an excess of $1/8$.
\end{example}

\subsection{Empirical fitting and prediction}
\label{subsec:proper-score-estimation}

\begin{proposition}[Empirical fitting by a proper score]
\label{prop:proper-score-empirical}
For a record $y_{1:n}$ of length $n\geq1$ with empirical law
$\widehat P_n$, the mean score of a forecast $Q$ on the record is
$s_S(\widehat P_n,Q)$. If $\widehat P_n$ lies in the forecast class
and $S$ is strictly proper on that class, $\widehat P_n$ is the
unique minimizer of the mean score.

Over any forecast family, the mean logarithmic and quadratic scores
have the same minimizers as, respectively,
\[
D(\widehat P_n\|Q)
\qquad\text{and}\qquad
\sum_y\bigl(\widehat p_n(y)-q(y)\bigr)^2.
\]
\end{proposition}
\begin{proof}
Group the terms $S(y_i,Q)/n$ by outcome and apply
Definition~\ref{def:empirical-law}. Strict propriety gives the
first minimizer claim. The empirical law has finite support, so
the logarithmic and quadratic regret identities apply with a
finite truthful term that is constant in $Q$.
\end{proof}

\begin{example}[Two scores on a restricted forecast family]
\label{ex:proper-scores-restricted-fit}
On three outcomes, suppose a record has counts $(4,1,3)$.
Take forecasts $Q=(q,3/4-q,1/4)$ with $0\leq q\leq3/4$.
The part of the mean log loss depending on $q$ is
\[
-\tfrac12\log q-\tfrac18\log(3/4-q).
\]
This is strictly convex with derivative zero at $q=3/5$ and
infinite at the endpoints. By the quadratic regret formula, the
quadratic fit minimizes
\[
(q-\tfrac12)^2+(\tfrac58-q)^2+(\tfrac18)^2.
\]
Its unique minimum is at $q=9/16$.
\end{example}

\begin{theorem}[Empirical forecasts on a fresh observation]
\label{thm:empirical-score-fresh-observation}
Let $Y_1,\ldots,Y_{n+1}$ be iid with law $P$ on $\mathsf Y$,
and let $\widehat P_n$ be the empirical law of the first $n$ draws,
$n\geq1$. Then
\[
\mathbb E\bigl[S_{\mathrm{Br}}(Y_{n+1},\widehat P_n)\bigr]
=\left(1+\frac1n\right)B_{\mathrm{Br}}(P).
\]
If $P$ is nondegenerate, then
$\mathbb E[L_{\log}(Y_{n+1},\widehat P_n)]=+\infty$.
If $P$ is a point mass, both scores are zero almost surely.
\end{theorem}
\begin{proof}
Given the training record, the fresh outcome has law $P$.
By Proposition~\ref{prop:log-quadratic-strictly-proper} its
conditional expected quadratic score is $B_{\mathrm{Br}}(P)$
plus $\sum_y(\widehat p_n(y)-p(y))^2$. For each $y$,
$\widehat p_n(y)$ is the mean of $n$ independent indicators
of mean $p(y)$, hence has mean $p(y)$ and variance $p(y)(1-p(y))/n$.
By nonnegative integration,
\[
\mathbb E\sum_y(\widehat p_n(y)-p(y))^2
=\frac1n\sum_y p(y)(1-p(y))
=\frac1n B_{\mathrm{Br}}(P).
\]

For nondegenerate $P$ there is a $y$ with $0<p(y)<1$.
With probability $(1-p(y))^n p(y)>0$, the training record omits
$y$ and the fresh outcome equals $y$. The logarithmic loss is
infinite on this event. For a point mass the training and fresh
outcomes coincide almost surely and both scores vanish.
\end{proof}

\begin{example}[One observation of a fair bit]
\label{ex:proper-scores-one-fair-observation}
Let $Y_1,Y_2$ be independent fair bits and forecast with
$\widehat P_1=\delta_{Y_1}$. Both training scores are zero.
On the fresh outcome $Y_2$, both scores are zero when $Y_2=Y_1$.
When $Y_2\ne Y_1$, the quadratic score is two and the logarithmic
score is infinite. These two events each have probability $1/2$,
so the expected fresh quadratic score is one and the expected fresh
logarithmic score is infinite.
\end{example}

\begin{example}[Smoothing one observation]
\label{ex:proper-scores-smoothed-fair-observation}
In the same fair-bit experiment let $U$ be the fair law and
$Q_\alpha=(1-\alpha)\delta_{Y_1}+\alpha U$ for $0<\alpha\leq1$.
The matching outcome gets probability $1-\alpha/2$ and the other
gets $\alpha/2$. The training log loss is $-\log(1-\alpha/2)$
and the expected fresh log loss is
\[
-\tfrac12\log\bigl((1-\alpha/2)(\alpha/2)\bigr).
\]
The fresh loss is minimized at $\alpha=1$, where the product
of the two probabilities is $1/4$. The expected fresh quadratic
score is $1/2+(1-\alpha)^2/2$, by the quadratic regret identity
at the fair law.
\end{example}

\subsection{Exercises}
\label{subsec:practice-3-4}

\begin{practiceexercise}[Binary normalization]
\label{pex:quadratic-score-binary-normalization}
For $y\in\{0,1\}$ verify the binary formula
$S_{\mathrm{Br}}(y,Q)=2(q(1)-y)^2$. Compute the score for
$q(1)=1/4$ at each outcome.
\end{practiceexercise}

\begin{practiceexercise}[Quadratic score of a point forecast]
\label{pex:quadratic-score-point-forecast}
For any law $P$ and outcome $a$, compute $s_{S_{\mathrm{Br}}}(P,\delta_a)$.
Which point forecasts minimize this expectation on a finite alphabet?
\end{practiceexercise}

\begin{practiceexercise}[Two regrets for the same forecast]
\label{pex:proper-score-log-and-quadratic-regrets}
Let $P$ be fair on $\{0,1\}$ and let $q(1)=1/4$. Compute the
logarithmic and quadratic regrets in their respective normalizations.
\end{practiceexercise}

\begin{practiceexercise}[Infinite truthful logarithmic loss]
\label{pex:proper-score-infinite-log-entropy}
Let $P$ be the infinite-entropy law of
Exercise~\ref{exer:relative-entropy-infinite-entropy}. Determine the
set of minimizers of $s_{L_{\log}}(P,Q)$ over all laws $Q$.
Explain why the logarithmic regret from
Definition~\ref{def:scoring-regret} is undefined at this $P$.
\end{practiceexercise}

\begin{practiceexercise}[Changing units and shifting outcomes]
\label{pex:proper-score-shift-and-scale}
For a finite alphabet, let $S'(y,Q)=cS(y,Q)+b(y)$ with
$c>0$ and $b(y)$ finite and nonnegative. Show that propriety
and strict propriety are preserved, and relate the two regrets
where they are defined.
\end{practiceexercise}

\begin{practiceexercise}[Absolute loss for a reported probability]
\label{pex:proper-score-absolute-probability}
For binary $y$ let $S(y,Q)=|y-q(1)|$. Determine all minimizers
of the expected score as a function of $p(1)$ and decide whether
the score is proper on all binary laws.
\end{practiceexercise}

\begin{practiceexercise}[Concavity gap for the quadratic score]
\label{pex:quadratic-score-concavity-gap}
For $P_t=(1-t)P_0+tP_1$ with $0<t<1$, express the difference
between $B_{\mathrm{Br}}(P_t)$ and
$(1-t)B_{\mathrm{Br}}(P_0)+tB_{\mathrm{Br}}(P_1)$ in terms of
the two mass vectors.
\end{practiceexercise}

\begin{practiceexercise}[The affine support of a quadratic score]
\label{pex:quadratic-score-affine-support}
For binary outcomes and fixed $q=q(1)$, write the expected
quadratic score as an affine function of $p=p(1)$. Show that
its graph is tangent to $B_{\mathrm{Br}}(P)$ as a function of $p$
at $p=q$.
\end{practiceexercise}

\begin{practiceexercise}[Strict propriety on a smaller class]
\label{pex:proper-score-restrict-domain}
Show that restricting a strictly proper score to a nonempty
subclass preserves strict propriety. For the score in
Example~\ref{ex:proper-score-mode}, give a two-element subclass
on which it is strictly proper.
\end{practiceexercise}

\begin{practiceexercise}[A forecast that ignores the feature]
\label{pex:proper-score-constant-log-penalty}
Suppose $Y$ takes finitely many values. Show that the optimal
constant forecast under log loss has risk $H(Y)$ and exceeds
the optimal conditional risk by $I(X;Y)$.
\end{practiceexercise}

\begin{practiceexercise}[Conditional quadratic regret]
\label{pex:proper-score-conditional-regret}
Let $X$ be a fair bit, with conditional probabilities of $Y=1$
equal to $1/4$ at $X=0$ and $3/4$ at $X=1$. Compute the
quadratic regret in population risk of the constant fair forecast.
\end{practiceexercise}

\begin{practiceexercise}[An unobserved feature value]
\label{pex:proper-score-null-feature}
Suppose $X=0$ almost surely but its value space is $\{0,1\}$, and
$Y$ is a fair bit. Describe all conditional forecast kernels
minimizing the population quadratic score.
\end{practiceexercise}

\begin{practiceexercise}[Fitting a binary record]
\label{pex:proper-score-fit-binary-record}
For the record $(1,0,1,1)$, compute the mean quadratic score
as a function of $q(1)$ and find its unique minimizer.
\end{practiceexercise}

\begin{practiceexercise}[Clipping a binary forecast]
\label{pex:proper-score-clipped-binary-fit}
A binary record has empirical frequency $3/4$ of ones.
Minimize each of the mean logarithmic and quadratic scores
over $q(1)\in[1/5,3/5]$.
\end{practiceexercise}

\begin{practiceexercise}[An unseen outcome in a fresh draw]
\label{pex:proper-score-unseen-fair-bit}
For $n$ training observations of a fair bit and one independent
fresh observation, compute the probability that the empirical
forecast incurs infinite fresh log loss.
\end{practiceexercise}

\begin{practiceexercise}[Uniform law on a finite alphabet]
\label{pex:quadratic-score-fresh-uniform}
For $P$ uniform on $m$ outcomes, compute the expected fresh
quadratic score of $\widehat P_n$ and its excess over the
truthful score.
\end{practiceexercise}

\begin{practiceexercise}[A uniform smoothing component]
\label{pex:proper-score-smoothed-log-bound}
Let $U$ be uniform on $m$ outcomes and
$Q_\alpha=(1-\alpha)\widehat P_n+\alpha U$ with $0<\alpha\leq1$.
Give a finite upper bound on $L_{\log}(y,Q_\alpha)$ that holds
for every record and outcome.
\end{practiceexercise}

\begin{practiceexercise}[Pooling two categories]
\label{pex:proper-score-pool-two-categories}
For $P=(1/2,1/3,1/6)$, minimize expected quadratic score
over forecasts with $q(2)=q(3)$. Give the minimizing law
and its regret.
\end{practiceexercise}

\begin{exercise}[Strict concavity and strict propriety]
\label{exer:proper-score-strict-concavity-converse}
Let $S$ be proper on a convex class $\mathcal P$ and $B_S$
finite and strictly concave. Show that $S$ is strictly proper.
If $P\ne Q$ and $s_S(P,Q)=B_S(P)$, apply affinity in the true
law along the segment joining $P$ and $Q$.
\end{exercise}

\begin{exercise}[A score built from the Euclidean norm]
\label{exer:spherical-score-propriety}
On a finite alphabet, set
$S(y,Q)=1-q(y)/(\sum_a q(a)^2)^{1/2}$.
Prove that the score lies in $[0,1]$ and is strictly proper.
Compute its truthful score and identify the equality condition
in the inequality used to minimize its expectation.
\end{exercise}

\begin{exercise}[Value of a feature under the quadratic score]
\label{exer:quadratic-score-feature-value}
Show that the minimum risk over constant forecasts exceeds the
minimum conditional risk by
\[
\sum_y \operatorname{Var}\!\bigl(P_{Y\mid X}(X,\{y\})\bigr).
\]
Allow a countable outcome space. Characterize the case of zero
improvement in terms of independence.
\end{exercise}

\begin{exercise}[Training score and fresh score]
\label{exer:quadratic-score-training-fresh-gap}
For iid observations with law $P$, show that the expected mean
quadratic score of $\widehat P_n$ on its own training record is
$(1-1/n)B_{\mathrm{Br}}(P)$. Compute the difference between the
expected fresh and training scores. Include $n=1$.
\end{exercise}

\begin{exercise}[Shrinking the empirical probability]
\label{exer:quadratic-score-optimal-shrinkage}
For iid Bernoulli data with parameter $p$, forecast with
$q_\alpha=(1-\alpha)\widehat p_n+\alpha/2$,
$0\leq\alpha\leq1$. Use the variance and bias of $q_\alpha$
to compute the expected fresh quadratic regret and its
minimizing $\alpha$ as a function of $p,n$. Include $p=0,1/2,1$.
\end{exercise}

\begin{exercise}[No forecast has finite log loss for every law]
\label{exer:proper-score-no-universal-countable-forecast}
On a countably infinite alphabet, fix any forecast $Q$.
Construct a true law $P$ of finite entropy with $H(P,Q)=+\infty$.
If all masses of $Q$ are positive, choose distinct outcomes $y_j$
with $q(y_j)\leq e^{-j^2}$ and use masses proportional to $j^{-2}$
on that subsequence.
\end{exercise}



\chapter{Statistics, sufficiency, and information reduction}
\label{chap:statistics-sufficiency}
\notationgroup{Statistics and reductions}
\section{Statistics and induced experiments}
\label{sec:statistics-induced-experiments}

Fix an experiment
$\mathcal E=(\mathsf X,\mathcal A,(P_\theta)_{\theta\in\Theta})$.

\subsection{Statistics and induced laws}
\label{subsec:statistics-induced-laws}

\begin{definition}[Statistic and induced experiment]
\label{def:statistic-induced-experiment}
\indexnotation{T,\ P_\theta^T,\ \mathcal E^T}{Statistic, induced law, and induced experiment.}
A \emph{statistic} is a measurable map
$T:(\mathsf X,\mathcal A)\to(\mathsf S,\mathcal S)$.
The induced law $P_\theta^T$ on $(\mathsf S,\mathcal S)$ is
\[
P_\theta^T(B)=P_\theta(T^{-1}(B)),\qquad B\in\mathcal S.
\]
The experiment $\mathcal E^T=(\mathsf S,\mathcal S,
(P_\theta^T)_{\theta\in\Theta})$ is the \emph{induced experiment}.
For any measure $\nu$ on $\mathsf X$, use the same pushforward
notation $\nu^T=\nu\circ T^{-1}$.
\end{definition}

\begin{proposition}[Integration and composition]
\label{prop:statistic-integration-composition}
For each nonnegative measurable $v$ on $\mathsf S$,
\[
\int v\,dP_\theta^T=\int v(T(x))\,P_\theta(dx).
\]
The formula also holds for $P_\theta^T$-integrable $v$.
For a measurable map $h$ on $\mathsf S$ into another measurable
space, $(P_\theta^T)^h=P_\theta^{h\circ T}$.
\end{proposition}
\begin{proof}
Preimages preserve disjoint unions and the whole space, so
$P_\theta^T$ is a probability law. The integration formula holds
for indicators by definition. Extend it by linearity to simple
functions and by monotone convergence to nonnegative functions.
Use positive and negative parts in the integrable case. The
composition formula follows from
$(h\circ T)^{-1}(C)=T^{-1}(h^{-1}(C))$ for every measurable $C$.
\end{proof}

\begin{example}[The Bernoulli count]
\label{ex:statistic-bernoulli-count}
Under $P_p^{\otimes n}$ for the Bernoulli family, take
$T(x)=\sum_{i=1}^n x_i$. There are $\binom nk$ records with count
$k$, each of probability $p^k(1-p)^{n-k}$ for $0<p<1$. Thus the
induced mass function is
\[
P_p^{\otimes n}(T=k)=\binom nk p^k(1-p)^{n-k},
\qquad k=0,\ldots,n.
\]
At $p=0$ and $p=1$ the induced laws are point masses at $0$ and
$n$, respectively. Since $\mathbb E_p T=np$, the count
experiment identifies $p$ for $n\geq1$.
\end{example}

\begin{example}[The mean of a normal sample]
\label{ex:statistic-normal-mean}
For independent $X_i\sim N(\theta,\sigma^2)$ with fixed $\sigma>0$,
take $T(x)=n^{-1}\sum_i x_i$. A linear combination of independent
normal variables is normal. The mean and variance of $T(X)$
are $\theta$ and $\sigma^2/n$, so the induced experiment has
laws $N(\theta,\sigma^2/n)$. The sum and the mean determine
each other by multiplication and division by $n$.
\end{example}

\begin{example}[The maximum of a uniform sample]
\label{ex:statistic-uniform-maximum}
Let $X_1,\ldots,X_n$ be independent and uniform on $[0,\theta]$
with $\theta>0$. For $T(x)=\max_i x_i$, independence gives
\[
P_\theta(T\leq t)=\prod_{i=1}^n P_\theta(X_i\leq t)
=(t/\theta)^n,\qquad 0\leq t\leq\theta,
\]
where $P_\theta$ denotes the joint law in this example.
The distribution function is zero for $t<0$ and one for
$t>\theta$. Its density on $(0,\theta)$ is
$nt^{n-1}/\theta^n$. The right endpoint of the interval of
positive density identifies $\theta$.
\end{example}

\subsection{Information in a statistic}
\label{subsec:information-in-a-statistic}

\begin{definition}[Sigma algebra generated by a statistic]
\label{def:statistic-sigma-algebra}
\indexnotation{\sigma(T)}{Sub-sigma algebra of observation events determined by the statistic \(T\).}
The sigma algebra generated by $T$ is
\[
\sigma(T)=\{T^{-1}(B):B\in\mathcal S\}\subseteq\mathcal A .
\]
Thus $\sigma(T)$ is the smallest sigma algebra on $\mathsf X$
for which $T$ is measurable.
\end{definition}

\begin{proposition}[Real-valued functions of a statistic]
\label{prop:real-statistic-factorization}
A function $U:\mathsf X\to\R$ is measurable with respect to
$\sigma(T)$ if and only if $U=h\circ T$ for some measurable
$h:\mathsf S\to\R$.
\end{proposition}
\begin{proof}
If $U=h\circ T$, then the inverse image under $U$ of any Borel
set belongs to $\sigma(T)$. Conversely, take finite-valued
$\sigma(T)$-measurable simple functions $U_n$ converging pointwise
to $U$. Each level set of $U_n$ is $T^{-1}(B)$ for some
$B\in\mathcal S$. Replacing those level-set indicators by
$\mathbf 1_B$ in the simple function gives a measurable $h_n$
with $U_n=h_n\circ T$. Set $h$ to the limit of $h_n$ wherever
the limit exists and is finite, and to zero elsewhere. This is
measurable, and $h(T(x))=U(x)$ for every $x$.
\end{proof}

\begin{proposition}[Identifiability and distances after reduction]
\label{prop:statistic-information-contraction}
Every target identifiable in $\mathcal E^T$ is identifiable in
$\mathcal E$. For any two laws $P,Q$ on the observation space,
\[
D(P^T\|Q^T)\leq D(P\|Q),\qquad
d_{\mathrm{TV}}(P^T,Q^T)\leq d_{\mathrm{TV}}(P,Q).
\]
If $U=h\circ T$ for a measurable $h$, the same conclusions
hold with the roles of $\mathcal E^T$ and $\mathcal E^U$ as the
original and reduced experiments, respectively.
\end{proposition}
\begin{proof}
Equality $P_\theta=P_\vartheta$ implies $P_\theta^T=P_\vartheta^T$
by the preimage formula. This proves the identifiability claim
by Definition~\ref{def:identifiability}. The relative-entropy
bound is Theorem~\ref{thm:relative-entropy-data-processing}.
For total variation, each
$|P^T(B)-Q^T(B)|=|P(T^{-1}(B))-Q(T^{-1}(B))|$ is bounded by
$d_{\mathrm{TV}}(P,Q)$, so the supremum satisfies the same bound.
For $U=h\circ T$, apply these arguments to $h$ on the induced
experiment and use Proposition~\ref{prop:statistic-integration-composition}.
\end{proof}

\begin{example}[The absolute value of a normal observation]
\label{ex:statistic-folded-normal}
Let $P_\theta=N(\theta,1)$ with $\theta\in\R$, and take $T(x)=|x|$.
Reflection sends $P_\theta$ to $P_{-\theta}$ and leaves $|x|$
fixed, so $P_\theta^T=P_{-\theta}^T$. Moreover
$\mathbb E_\theta[T(X)^2]=1+\theta^2$, so equality of the induced
laws implies equality of $|\theta|$. The induced experiment
identifies $|\theta|$ and loses the sign. For $a>0$,
Example~\ref{ex:normal-relative-entropy} gives
\[
D(P_a^T\|P_{-a}^T)=0,\qquad D(P_a\|P_{-a})=2a^2.
\]
\end{example}

\subsection{Domination and induced densities}
\label{subsec:statistic-induced-densities}

\begin{proposition}[Domination and densities after reduction]
\label{prop:statistic-induced-density}
If $\mathcal E$ is dominated in the sense of
Definition~\ref{def:dominated-experiment}, then $\mathcal E^T$
is dominated. More precisely, if $\mu$ is a dominating
probability measure on $\mathsf X$, then $\mu^T$ dominates the
induced laws. For $f_\theta=dP_\theta/d\mu$ and
$g_\theta=dP_\theta^T/d\mu^T$,
\[
g_\theta\circ T=\mathbb E_\mu[f_\theta\mid\sigma(T)]
\qquad \mu\text{-almost surely}.
\]
\end{proposition}
\begin{proof}
Let $\nu$ be a $\sigma$-finite dominating measure. Take a countable
measurable partition $(E_j)_{j\geq1}$ of $\mathsf X$ into sets of
finite $\nu$-measure. The function
$w=\sum_j 2^{-j}(1+\nu(E_j))^{-1}\mathbf 1_{E_j}$ is strictly
positive and has finite, positive integral. Normalizing $w\,d\nu$
gives a probability measure $\mu$ with the same null sets as $\nu$.
If $\mu^T(B)=0$, then $\mu(T^{-1}(B))=0$ and hence
$P_\theta^T(B)=0$. Thus $\mu^T$ is a dominating probability measure.

For every $B\in\mathcal S$, the definition of the densities gives
\[
\int_{T^{-1}(B)}g_\theta(T(x))\,\mu(dx)
=\int_B g_\theta\,d\mu^T
=P_\theta(T^{-1}(B))
=\int_{T^{-1}(B)}f_\theta\,d\mu.
\]
The first equality is the integration formula. The sets $T^{-1}(B)$
are precisely those in $\sigma(T)$, which proves the conditional
expectation identity.
\end{proof}

\begin{example}[Pooling two intervals]
\label{ex:statistic-domination-pooling}
Let $P_p=(1-p)\operatorname{Unif}[-1,0]+p\operatorname{Unif}[0,1]$
on $\R$ with $0\leq p\leq1$, and put $T(x)=\mathbf 1_{\{x>0\}}$.
Lebesgue measure dominates the original laws, but its pushforward
assigns infinite mass to both outputs. The probability measure
$\mu=\operatorname{Unif}[-1,1]$ also dominates the original laws.
Its induced law is fair, while $P_p^T$ has mass $p$ at one.
The induced density at $0,1$ is therefore $2(1-p),2p$,
respectively. These are the averages of the original density
relative to $\mu$ on the two fibers.
\end{example}

\subsection{Decision rules based on a statistic}
\label{subsec:statistic-decision-rules}

\begin{theorem}[Lifting rules from an induced experiment]
\label{thm:statistic-risk-lifting}
Fix an action space and a loss $L$ on the parameter and action sets.
If $K$ is a decision kernel for $\mathcal E^T$, then the kernel
$\widetilde K(x,\cdot)=K(T(x),\cdot)$ for $\mathcal E$ satisfies
\[
P_\theta\widetilde K=P_\theta^T K
\qquad\text{for every }\theta.
\]
Consequently, its risk in $\mathcal E$ equals the risk of $K$
in $\mathcal E^T$ at each parameter. The minimax value in
$\mathcal E^T$ is at least the value in $\mathcal E$. The same
ordering holds for the infimum of average risk under any prior
for which the average risks are defined.

If there is a measurable $h:\mathsf S\to\mathsf X$ with
$h\circ T=\mathrm{id}_{\mathsf X}$, then the two experiments
have the same set of attainable risk functions.
\end{theorem}
\begin{proof}
The kernel property of $\widetilde K$ follows from measurability
of $T$ and of the action probabilities under $K$. For each
measurable action set $B$,
\[
(P_\theta\widetilde K)(B)
=\int K(T(x),B)\,P_\theta(dx)
=\int K(t,B)\,P_\theta^T(dt)
=(P_\theta^T K)(B).
\]
By \eqref{eq:risk-action-law}, the risks agree pointwise in
$\theta$. The inequalities for minimax and prior average risk
follow by taking infima over rules.

If $h\circ T=\mathrm{id}_{\mathsf X}$ and $J$ is a rule on the
original space, the kernel $K(t,\cdot)=J(h(t),\cdot)$ lifts to
$J$. Thus every original risk function is attained in the
induced experiment as well.
\end{proof}

\begin{example}[Identification with increased testing risk]
\label{ex:statistic-increased-testing-risk}
Consider two labels, with laws on $\{1,2,3,4\}$ given by
\[
P=(1/2,1/8,1/4,1/8),\qquad Q=(1/8,1/4,1/8,1/2).
\]
Let $T$ record whether the outcome lies in $\{1,2\}$ or
$\{3,4\}$. The induced laws $(5/8,3/8)$ and $(3/8,5/8)$
are distinct, so the label is still identifiable.

Under equal prior weights and zero-one loss for choosing a label,
the contribution of outcome $i$ to average risk is minimized
by choosing the label of larger mass. Its minimum is
$\tfrac12\min\{P(\{i\}),Q(\{i\})\}$, since any randomization
averages the two possible contributions. The minimum average
risk is $1/4$ with the full observation and $3/8$ after applying
the same calculation to the two induced laws.
\end{example}

\subsection{Exercises}
\label{subsec:practice-4-1}

\begin{practiceexercise}[A two-cell statistic]
\label{pex:statistic-two-cell-law}
Let $P$ have masses $(1/2,1/3,1/6)$ on $\{1,2,3\}$.
For $T(x)=\mathbf 1_{\{x\geq2\}}$, compute $P^T$ and list all
events in $\sigma(T)$.
\end{practiceexercise}

\begin{practiceexercise}[Indicators and event probabilities]
\label{pex:statistic-indicator-law}
For a measurable event $A$, compute $P^T$ for $T=\mathbf 1_A$.
For a second law $Q$, express $d_{\mathrm{TV}}(P^T,Q^T)$ in terms
of $P(A),Q(A)$.
\end{practiceexercise}

\begin{practiceexercise}[Grouping consecutive values]
\label{pex:statistic-grouped-law}
On $\{0,1,2,3,4,5\}$ put $T(x)=\lfloor x/2\rfloor$ and
$U(x)=\mathbf 1_{\{x\geq2\}}$. Find $h$ with $U=h\circ T$,
then list the events in $\sigma(U)$.
\end{practiceexercise}

\begin{practiceexercise}[A statistic with unused values]
\label{pex:statistic-unused-values}
Take $\mathsf X=\{0,1\}$ and $T(x)=2x$ with value space
$\mathsf S=\{0,1,2\}$. For $P(\{1\})=p$ compute $P^T$.
Describe all functions $h:\mathsf S\to\R$ with $h(T(x))=x$.
\end{practiceexercise}

\begin{practiceexercise}[A single retained coordinate]
\label{pex:statistic-first-coordinate}
Let $X_1,\ldots,X_n$ be iid Bernoulli with parameter $p$.
Find the induced law of $T(x)=x_1$ and show that it identifies $p$.
\end{practiceexercise}

\begin{practiceexercise}[Detecting at least one success]
\label{pex:statistic-any-success}
For a Bernoulli sample of size $n\geq1$, let
$T(x)=\mathbf 1_{\{\sum_i x_i>0\}}$.
Find the induced law and determine whether it identifies $p$ on $[0,1]$.
\end{practiceexercise}

\begin{practiceexercise}[The parity of a count]
\label{pex:statistic-count-parity}
For two iid Bernoulli observations, take
$T(x)=(x_1+x_2)\bmod2$. Compute its mass at one and describe
all pairs of parameter values with the same induced law.
\end{practiceexercise}

\begin{practiceexercise}[Thresholding a normal observation]
\label{pex:statistic-normal-threshold}
For $X\sim N(\theta,1)$ and $T(x)=\mathbf 1_{\{x>0\}}$,
express $P_\theta^T(\{1\})$ using the standard normal distribution
function. Show that the induced family identifies $\theta$.
\end{practiceexercise}

\begin{practiceexercise}[Censoring at a fixed level]
\label{pex:statistic-censored-exponential}
Let $X$ have density $\lambda e^{-\lambda x}$ on $[0,\infty)$
and take $T(x)=\min\{x,c\}$ for $c>0$. Give the induced density
on $(0,c)$ and the mass at $c$.
\end{practiceexercise}

\begin{practiceexercise}[The minimum of a uniform sample]
\label{pex:statistic-uniform-minimum}
For $n$ iid draws from $\operatorname{Unif}[0,\theta]$, derive
the distribution function of the minimum using the probability
that all draws exceed $t$.
\end{practiceexercise}

\begin{practiceexercise}[An expectation from the induced law]
\label{pex:statistic-count-moments}
Use the induced law of the Bernoulli count with $n=2$ to compute
$\mathbb E_p[T^2]$. Verify the result by expanding $(X_1+X_2)^2$.
\end{practiceexercise}

\begin{practiceexercise}[Equal information by reversible scaling]
\label{pex:statistic-mean-sum-distances}
For laws $P,Q$ on $\R^n$, let $T$ be the sample sum and $U$ the
sample mean. Prove equality of the relative entropies of their
induced law pairs and equality of their total variation distances.
\end{practiceexercise}

\begin{practiceexercise}[A finite induced density]
\label{pex:statistic-finite-induced-density}
Let $\mu$ be uniform on $\{1,2,3\}$ and $P=(1/2,1/3,1/6)$.
Take $T(x)=\mathbf 1_{\{x=1\}}$. Compute $dP/d\mu$ and
$dP^T/d\mu^T$. Check the conditional expectation identity
on both fibers.
\end{practiceexercise}

\begin{practiceexercise}[A sigma-finite pushforward]
\label{pex:statistic-floor-lebesgue}
For Lebesgue measure $\lambda$ on $\R$ and $T(x)=\lfloor x\rfloor$
with value space $\mathbb Z$, compute $\lambda^T$ on singletons.
Is $\lambda^T$ sigma-finite?
\end{practiceexercise}

\begin{practiceexercise}[Lifting a randomized rule]
\label{pex:statistic-lift-randomized-rule}
In the Bernoulli experiment with $n=2$, let a rule output action
one with probability $k/2$ when the count is $k$. Write the
lifted action-one probability at each full record. Compute the
unconditional probability of action one.
\end{practiceexercise}

\begin{practiceexercise}[Risk of a constant statistic]
\label{pex:statistic-constant-testing-risk}
Consider two labels with prior masses $\pi$ and $1-\pi$ and
zero-one loss. Compute the smallest average risk using a
constant statistic, allowing randomized actions.
\end{practiceexercise}

\begin{practiceexercise}[Distances in the four-outcome example]
\label{pex:statistic-testing-example-distances}
Compute $d_{\mathrm{TV}}(P,Q)$ and $d_{\mathrm{TV}}(P^T,Q^T)$ in
Example~\ref{ex:statistic-increased-testing-risk}. Show directly
from log-sum on the first fiber that relative entropy decreases
strictly.
\end{practiceexercise}

\begin{practiceexercise}[Reconstructing a record]
\label{pex:statistic-record-left-inverse}
For $T(x)=(x,x^2)$ on $\R$, give a measurable left inverse on
$\R^2$. Express any randomized rule based on $x$ as a rule
on the value space of $T$ with the same risk.
\end{practiceexercise}

\begin{exercise}[Counts on a finite alphabet]
\label{exer:statistic-multinomial-counts}
For $n$ iid draws from a law with masses $(p_1,\ldots,p_m)$,
take $T(x)$ to be the vector of category counts. Derive its
joint mass function by counting ordered records. Show that
the induced experiment identifies the mass vector for $n\geq1$.
\end{exercise}

\begin{exercise}[The law of an order statistic]
\label{exer:statistic-order-distribution}
Let $X_1,\ldots,X_n$ be iid real variables with distribution
function $F$. For $1\leq k\leq n$, let $X_{(k)}$ be the $k$th
value in the sample written in nondecreasing order, including
ties. Show that the event $X_{(k)}\leq t$ is measurable and
derive its probability in terms of $F(t)$ by counting
observations at most $t$.
\end{exercise}

\begin{exercise}[A normal statistic that discards location]
\label{exer:statistic-normal-difference}
For independent $X_1,X_2\sim N(\theta,\sigma^2)$ with
$\sigma>0$, find the law of $T(X)=X_1-X_2$. With $\sigma$
known, determine which targets of $\theta$ are identifiable
from $T$. With both $\theta$ and $\sigma$ unknown, determine
which pairs of labels are observationally equivalent under $T$.
\end{exercise}

\begin{exercise}[Relative entropy of a Bernoulli count]
\label{exer:statistic-bernoulli-count-preserves-kl}
For $0<p,q<1$, show that the relative entropy between the
induced count laws for $n$ iid draws equals that between the
full sample laws. Compute the ratio of the two probability
masses at a count $k$ and compare it with the ratio at any
record with that count.
\end{exercise}

\begin{exercise}[Equal-prior testing and total variation]
\label{exer:statistic-tv-testing-equality}
On a finite space, show that the optimal equal-prior average
zero-one risk for testing $P$ against $Q$ is
$\tfrac12(1-d_{\mathrm{TV}}(P,Q))$, allowing randomization.
For a statistic with finitely many values, prove that this
risk is preserved exactly when the differences $p(x)-q(x)$
have a common weak sign on each fiber.
\end{exercise}

\begin{exercise}[A likelihood ratio determined by a statistic]
\label{exer:statistic-ratio-preserves-kl}
Suppose $P\ll Q$ and $dP/dQ$ has a $\sigma(T)$-measurable
real-valued version. Use
Propositions~\ref{prop:real-statistic-factorization} and
\ref{prop:statistic-induced-density} to identify the induced
likelihood ratio. Deduce $D(P^T\|Q^T)=D(P\|Q)$, including
when the common value is infinite.
\end{exercise}



\clearpage
\section{Sufficiency and conditional formulations}
\label{sec:sufficiency-conditional}

Fix an experiment $\mathcal E=(\mathsf X,\mathcal A,
(P_\theta)_{\theta\in\Theta})$ and a statistic
$T:(\mathsf X,\mathcal A)\to(\mathsf S,\mathcal S)$.

\subsection{Common conditional probabilities}
\label{subsec:sufficient-statistic}

\begin{definition}[Sufficient statistic]
\label{def:sufficient-statistic}
\indexnotation{c_A}{Common conditional probability of an observation event \(A\) given the statistic.}
The statistic $T$ is \emph{sufficient} for $\mathcal E$ if for every
$A\in\mathcal A$ there is a measurable $c_A:\mathsf S\to[0,1]$
such that
\begin{equation}
P_\theta(A\cap T^{-1}(B))=\int_B c_A(t)\,P_\theta^T(dt)
\qquad(\theta\in\Theta,\ B\in\mathcal S).
\label{eq:sufficiency-common-conditional}
\end{equation}
Equivalently, $c_A\circ T$ is a version of the conditional expectation
of $\mathbf 1_A$ given $\sigma(T)$ under every $P_\theta$.
\end{definition}

\begin{definition}[Common conditional kernel]
\label{def:common-conditional-kernel}
\indexnotation{K(t,\cdot)}{Conditional law of the full observation at statistic value \(t\), common to the experiment.}
A \emph{common conditional kernel} for $T$ is a probability kernel
$K$ from $\mathsf S$ to $\mathsf X$ such that $K(\cdot,A)$ satisfies
\eqref{eq:sufficiency-common-conditional} for every $A$.
Equivalently, under each $P_\theta$, the joint law of $(T(X),X)$
is $P_\theta^T\ltimes K$.
\end{definition}

\begin{proposition}[When common conditional probabilities form a kernel]
\label{prop:sufficiency-common-kernel-existence}
A common conditional kernel implies sufficiency. Conversely,
if $T$ is sufficient, a common conditional kernel exists in either
of the following cases.
\begin{enumerate}
\item The observation space $\mathsf X$ is countable with all
      subsets measurable.
\item A countable mixture $M$ of model laws dominates every
      $P_\theta$, and under $M$ a conditional law of $X$ given
      $T(X)$ exists.
\end{enumerate}
In the second case, any such conditional kernel under $M$ is common.
\end{proposition}
\begin{proof}
The forward implication follows by taking $c_A(t)=K(t,A)$.

For countable $\mathsf X$, choose a common version $c_{\{x\}}$
for each singleton. Under every $P_\theta^T$ the sum of these
functions is one almost surely, since its integral on each $B$
is $P_\theta(T^{-1}(B))$ by monotone convergence. On the measurable
set where they sum to one, use them as the masses of $K(t,\cdot)$.
Use a fixed point mass elsewhere. For any event $A$, summing the
singleton identities over $x\in A$ verifies
\eqref{eq:sufficiency-common-conditional} for this kernel.

For the mixture case, the common version $c_A$ satisfies the
conditional identity under $M$ by averaging the identities
under the component laws. If $K$ is a conditional kernel under
$M$, uniqueness of the conditional expectation gives
$K(t,A)=c_A(t)$ for $M^T$-almost every $t$. Because
$P_\theta^T\ll M^T$ for each $\theta$, this equality also holds
under every induced law. This argument applies to each $A$.
\end{proof}

\begin{proposition}[Reconstruction and the fibers]
\label{prop:sufficiency-reconstruction-fibers}
For a kernel $K$ from $\mathsf S$ to $\mathsf X$, the following
are equivalent.
\begin{enumerate}
\item $K$ is a common conditional kernel for $T$.
\item For every $\theta$, $P_\theta^T K=P_\theta$, and for every
      $B\in\mathcal S$,
      \[
      K(t,T^{-1}(B))=\mathbf 1_B(t)
      \qquad P_\theta^T\text{-almost surely in }t.
      \]
\end{enumerate}
The null set in the second condition may depend on $\theta$ and $B$.
\end{proposition}
\begin{proof}
For a common conditional kernel, put $B=\mathsf S$ in
\eqref{eq:sufficiency-common-conditional} to get reconstruction.
For a fixed $B$, apply the same identity to $A=T^{-1}(B)$
and integrate over an arbitrary $C\in\mathcal S$. The integral is
$P_\theta^T(B\cap C)$, so uniqueness gives
$K(t,T^{-1}(B))=\mathbf 1_B(t)$ almost surely.

Conversely, fix $\theta,A,B$. The fiber condition implies
$K(t,A\cap T^{-1}(B))=\mathbf 1_B(t)K(t,A)$ for $P_\theta^T$-almost
every $t$. Integrate and use $P_\theta^T K=P_\theta$ to obtain
\eqref{eq:sufficiency-common-conditional}.
\end{proof}

\subsection{Conditional laws in sampling models}
\label{subsec:sufficiency-conditional-examples}

\begin{example}[A Bernoulli sample given its count]
\label{ex:sufficiency-bernoulli-count}
In Example~\ref{ex:statistic-bernoulli-count}, fix a count $k$
having positive probability under $p$. Every record of count $k$
has the same probability under $p$. The conditional law is
therefore uniform on the $\binom nk$ records of count $k$.
Use that uniform law at each $k=0,\ldots,n$, including counts
of zero probability under $p=0$ or $p=1$. This defines a
single conditional kernel for all $p\in[0,1]$.
\end{example}

\begin{example}[A normal sample given its mean]
\label{ex:sufficiency-normal-mean}
For independent $X_i\sim N(\theta,\sigma^2)$ with known $\sigma>0$,
write $R_i=X_i-\overline X$ and $R=(R_1,\ldots,R_n)$.
The law of $R$ is independent of $\theta$, since subtracting
$\overline X$ cancels the common location. The vector
$(\overline X,R)$ is jointly Gaussian, and
\[
\operatorname{Cov}(\overline X,R_i)
=\sigma^2/n-\sigma^2/n=0 .
\]
The joint Gaussian characteristic function then factors, so
$\overline X$ and $R$ are independent. With
$\mathbf 1=(1,\ldots,1)$, the law of $t\mathbf 1+R$ gives a
kernel $K(t,\cdot)$. Its measurability follows by integrating
the Borel function $\mathbf 1_A(t\mathbf 1+r)$ against the
fixed law of $R$. Independence verifies that it is a common
conditional kernel given $\overline X$.
\end{example}

\begin{example}[Two uniform draws given their maximum]
\label{ex:sufficiency-uniform-maximum-pair}
Let $X_1,X_2$ be iid uniform on $[0,\theta]$, $\theta>0$,
and $T=\max(X_1,X_2)$. At $t\geq0$, draw $U$ uniformly on
$[0,1]$ and place $t,tU$ in either order with equal probability.
This defines a measurable kernel $K(t,\cdot)$. To verify the
conditional law, use the induced density $2t/\theta^2$ from
Example~\ref{ex:statistic-uniform-maximum}. For a Borel record
event $A$ and statistic event $B$,
\[
\int_B K(t,A)\,P_\theta^T(dt)
=\frac{1}{\theta^2}\int_{B\cap(0,\theta)}
  \int_0^t \bigl(\mathbf 1_A(t,u)+\mathbf 1_A(u,t)\bigr)\,du\,dt.
\]
The two terms integrate over the two triangles of the sample
square. The diagonal has probability zero, so the integral
equals $P_\theta(A\cap T^{-1}(B))$.
\end{example}

\begin{example}[Identification from one coordinate]
\label{ex:sufficiency-first-bernoulli-coordinate}
Let $X_1,X_2$ be iid Bernoulli with $0<p<1$ and retain
$T=X_1$. Its law determines $p$. However, the conditional
probability of $A=\{X_2=1\}$ given $T=0$ is $p$ by independence.
The value $T=0$ has positive probability at every parameter,
so no single $c_A(0)$ can satisfy
\eqref{eq:sufficiency-common-conditional} for the family.
\end{example}

\subsection{Posterior characterization for finite parameter sets}
\label{subsec:sufficiency-finite-posterior}

\begin{theorem}[Sufficiency and posterior probabilities]
\label{thm:sufficiency-finite-posterior}
Suppose $\Theta$ is finite with all subsets measurable. Fix prior
masses $\pi_\theta>0$ with $\sum_\theta\pi_\theta=1$, and write
$M=\sum_\theta\pi_\theta P_\theta$.
Let $\eta_\theta(x)$ and $\eta_\theta^T(t)$ be the posterior
probabilities of label $\theta$ from the original and induced
experiments, respectively. Then $T$ is sufficient if and only if
\[
\eta_\theta(x)=\eta_\theta^T(T(x))
\qquad M\text{-almost surely, for every }\theta.
\]
\end{theorem}
\begin{proof}
By the posterior rectangle identity
\eqref{eq:posterior-rectangles}, the posterior masses are
$\pi_\theta f_\theta$ and $\pi_\theta g_\theta$, where
$f_\theta=dP_\theta/dM$ and $g_\theta=dP_\theta^T/dM^T$.
Thus the claimed posterior equality is equivalent to
$f_\theta=g_\theta\circ T$ $M$-almost surely for every $\theta$.

If $T$ is sufficient, $c_A\circ T$ is a conditional expectation
of $\mathbf 1_A$ under $M$, by averaging
\eqref{eq:sufficiency-common-conditional} over the prior. Hence
\[
P_\theta(A)=\int c_A\,dP_\theta^T
=\int (c_A\circ T)(g_\theta\circ T)\,dM
=\int_A(g_\theta\circ T)\,dM.
\]
In the last equality the weight $g_\theta\circ T$ is nonnegative
and $\sigma(T)$-measurable. Uniqueness of the density gives the
posterior equality.

Conversely, suppose $f_\theta=g_\theta\circ T$ almost surely
for every $\theta$. Take a $[0,1]$-valued conditional expectation
of $\mathbf 1_A$ given $\sigma(T)$ under $M$. By
Proposition~\ref{prop:real-statistic-factorization}, write it as
$c_A\circ T$ with $c_A$ measurable, clipping $c_A$ to $[0,1]$
if necessary. Integrate the conditional expectation identity
against $\mathbf 1_{T^{-1}(B)}(g_\theta\circ T)$. This gives
\[
P_\theta(A\cap T^{-1}(B))
=\int_{T^{-1}(B)}(c_A\circ T) f_\theta\,dM
=\int_B c_A\,dP_\theta^T .
\]
Thus $c_A$ is a common version for every $A$.
\end{proof}

\begin{example}[A prior on two Bernoulli probabilities]
\label{ex:sufficiency-bernoulli-posterior}
Give $p=1/4$ and $p=3/4$ equal prior mass and observe two iid
bits. At any record of count $k$, the ratio of the likelihood
at $3/4$ to that at $1/4$ is $3^{2k-2}$. Hence the posterior
probability of $p=3/4$ is
\[
\frac{3^{2k-2}}{1+3^{2k-2}}.
\]
The likelihoods of the count have the same ratio, since each
record likelihood is multiplied by the same $\binom{2}{k}$.
Both observations therefore give posterior probabilities
$1/10,1/2,9/10$ at counts $0,1,2$, respectively.
\end{example}

\subsection{Decision risk and information}
\label{subsec:sufficiency-decision-information}

\begin{theorem}[Decision and information preservation]
\label{thm:sufficiency-risk-information}
Suppose a single kernel $K$ from $\mathsf S$ to $\mathsf X$
satisfies $P_\theta^T K=P_\theta$ for all $\theta$, as a common
conditional kernel does. Then the two experiments have the same
attainable risk functions for every action space and loss.
For all $\theta,\vartheta\in\Theta$,
\[
D(P_\theta^T\|P_\vartheta^T)=D(P_\theta\|P_\vartheta),\qquad
d_{\mathrm{TV}}(P_\theta^T,P_\vartheta^T)
=d_{\mathrm{TV}}(P_\theta,P_\vartheta).
\]
\end{theorem}
\begin{proof}
Given a full-observation decision kernel $J$, use the reduced
rule $t\mapsto\int J(x,\cdot)\,K(t,dx)$. This is a kernel by
measurability of kernel integrals and monotone convergence.
Its action law at $\theta$ is $(P_\theta^T K)J=P_\theta J$ by
iterated integration, so the risks agree. Conversely,
Theorem~\ref{thm:statistic-risk-lifting} lifts every reduced rule
with the same risk.

Relative entropy contracts along $T$ and along $K$ by
Theorem~\ref{thm:relative-entropy-data-processing}. As $K$
reconstructs both laws, these two inequalities force equality.
For total variation, a kernel also contracts distances.
Indeed, for laws with densities $f,g$ relative to their sum $\mu$,
the difference of their output probabilities on a set $A$ is
$\int K(t,A)(f(t)-g(t))\,\mu(dt)$. Since $0\leq K(t,A)\leq1$,
its absolute value is at most
$\tfrac12\int |f-g|\,d\mu$ by equality of the positive and negative
masses of $f-g$. Take the supremum over $A$ and use
Propositions~\ref{prop:total-variation-density}
and~\ref{prop:statistic-information-contraction} for the two directions.
\end{proof}

\begin{example}[A reconstruction kernel that changes the statistic]
\label{ex:sufficiency-reconstruction-without-fibers}
Let $X=(U,V)$, with $U$ Bernoulli of parameter $p$ and $V$ an
independent fair bit. Use the identity statistic $T(X)=X$.
The kernel sending $(u,v)$ to $(u,1-v)$ reconstructs $P_p$
from $P_p^T$ because flipping the fair bit preserves the joint law.
For $B=\{(u,v):v=0\}$, it gives
$K(t,T^{-1}(B))=1-\mathbf 1_B(t)$, so it fails the fiber condition
of Proposition~\ref{prop:sufficiency-reconstruction-fibers}.
The point-mass kernel $t\mapsto\delta_t$ is common conditional.
\end{example}

\subsection{Exercises}
\label{subsec:practice-4-2}

\begin{practiceexercise}[A particular count fiber]
\label{pex:sufficiency-bernoulli-three-fiber}
For three iid Bernoulli draws with $0<p<1$, list the records
with count two and their conditional masses given that count.
Compute the conditional probability that the first draw is one.
\end{practiceexercise}

\begin{practiceexercise}[Counts at the endpoints]
\label{pex:sufficiency-bernoulli-null-fibers}
For two Bernoulli draws and parameter set $\{0,1\}$, describe
all common conditional kernels for the count. Specify the
freedom at count one.
\end{practiceexercise}

\begin{practiceexercise}[One coordinate given the count]
\label{pex:sufficiency-bernoulli-one-coordinate}
For $n$ Bernoulli draws and any count $k$ of positive probability,
derive $\Pr_p(X_1=1\mid T=k)=k/n$ by counting records.
Include $k=0$ and $k=n$.
\end{practiceexercise}

\begin{practiceexercise}[A discrete uniform maximum]
\label{pex:sufficiency-discrete-uniform-maximum}
Let two independent draws be uniform on $\{1,\ldots,m\}$
with unknown positive integer $m$. Show that, conditional
on their maximum $t$, their joint law is uniform on
$2t-1$ ordered records. Prove sufficiency.
\end{practiceexercise}

\begin{practiceexercise}[A constant statistic]
\label{pex:sufficiency-constant-statistic}
Show that a constant statistic is sufficient if and only
if all model laws are equal.
\end{practiceexercise}

\begin{practiceexercise}[An invertible encoding]
\label{pex:sufficiency-invertible-statistic}
Suppose $T$ has a measurable left inverse $h$. Show that
$K(t,\cdot)=\delta_{h(t)}$ is a common conditional kernel.
\end{practiceexercise}

\begin{practiceexercise}[The same generated sigma algebra]
\label{pex:sufficiency-same-sigma-algebra}
Suppose $\sigma(U)=\sigma(T)$ and $T$ is sufficient.
Use Proposition~\ref{prop:real-statistic-factorization}
to show that $U$ is sufficient.
\end{practiceexercise}

\begin{practiceexercise}[A parameter-free extra measurement]
\label{pex:sufficiency-independent-extra-measurement}
Let $P_\theta=Q_\theta\otimes R$ on a product space, with $R$
fixed. For $T$ the first projection, give a common conditional
kernel and verify the defining identity on rectangles.
\end{practiceexercise}

\begin{practiceexercise}[Two finite fibers]
\label{pex:sufficiency-test-two-fibers}
On four ordered outcomes, take laws
\[
P=(1/4,1/4,1/8,3/8),\qquad Q=(1/8,1/8,3/8,3/8).
\]
Let $T$ group the first two and last two outcomes into two cells.
Compute both conditional laws on each cell and decide whether
$T$ is sufficient for $\{P,Q\}$.
\end{practiceexercise}

\begin{practiceexercise}[Unequal prior weights]
\label{pex:sufficiency-bernoulli-unequal-prior}
In Example~\ref{ex:sufficiency-bernoulli-posterior}, give
$p=3/4$ prior mass $1/3$. Compute its posterior mass at
counts zero, one, and two.
\end{practiceexercise}

\begin{practiceexercise}[A prior supported on one label]
\label{pex:sufficiency-posterior-zero-prior}
For a two-label experiment with distinct laws, take a prior
concentrated on one label and take $T$ constant. Check the
posterior equality of Theorem~\ref{thm:sufficiency-finite-posterior}
and determine whether $T$ is sufficient.
\end{practiceexercise}

\begin{practiceexercise}[A mixture conditional law]
\label{pex:sufficiency-mixture-conditional-law}
For the two-draw Bernoulli family, take a prior mixture of the
laws at $p=1/4$ and $p=3/4$ with equal weights. Compute the
conditional probabilities of $(1,0)$ and $(0,1)$ given
count one under the mixture.
\end{practiceexercise}

\begin{practiceexercise}[Checking a fiber directly]
\label{pex:sufficiency-flip-fiber-failure}
In Example~\ref{ex:sufficiency-reconstruction-without-fibers},
choose a record event $A$ and statistic event $B$ for which
the flipping kernel gives a different value from
$P_p(A\cap T^{-1}(B))$. Compute both values.
\end{practiceexercise}

\begin{practiceexercise}[Simulating a decision from the count]
\label{pex:sufficiency-reduce-first-bit-rule}
For $n$ iid Bernoulli draws, the full rule outputs its
first bit. Use the common count kernel to give a reduced
randomized rule with the same action law at every $p$.
\end{practiceexercise}

\begin{practiceexercise}[Testing with two Bernoulli draws]
\label{pex:sufficiency-bernoulli-tv}
For parameters $1/4$ and $3/4$, compute the total variation
distance between the two count laws when $n=2$.
Use sufficiency to infer the distance for the full records.
\end{practiceexercise}

\begin{practiceexercise}[A normal residual for two draws]
\label{pex:sufficiency-normal-two-draw-kernel}
For two iid $N(\theta,\sigma^2)$ draws with known $\sigma$,
write $(X_1,X_2)=(\overline X+D,\overline X-D)$. Find the
law of $D$ and write the common kernel given $\overline X=t$.
\end{practiceexercise}

\begin{practiceexercise}[Conditional moments under a uniform maximum]
\label{pex:sufficiency-uniform-maximum-conditional-moment}
In Example~\ref{ex:sufficiency-uniform-maximum-pair}, compute
$\mathbb E_\theta[X_1\mid T=t]$ using the displayed kernel.
\end{practiceexercise}

\begin{practiceexercise}[A kernel on unused statistic values]
\label{pex:sufficiency-unused-statistic-values}
Suppose $K$ is a common conditional kernel and
$N\in\mathcal S$ satisfies $P_\theta^T(N)=0$ for every $\theta$.
Show that replacing $K$ on $N$ by any fixed probability law
on $\mathsf X$ gives another common conditional kernel.
\end{practiceexercise}

\begin{exercise}[Poisson observations given their total]
\label{exer:sufficiency-poisson-total}
For $n$ iid Poisson observations with rate $\lambda>0$, derive
the conditional joint mass at $(x_1,\ldots,x_n)$ given
$\sum_i X_i=k$. Show that the total is sufficient. One may
obtain the law of the total by summing the record masses
with the multinomial formula.
\end{exercise}

\begin{exercise}[A uniform sample given its maximum]
\label{exer:sufficiency-uniform-maximum-general}
Extend Example~\ref{ex:sufficiency-uniform-maximum-pair} to
$n\geq2$ iid uniform draws. At maximum $t$, put $t$ in a
uniformly chosen coordinate and fill the remaining coordinates
with independent uniform draws on $[0,t]$. Verify the
conditional identity by partitioning the sample cube according
to the coordinate attaining the unique maximum.
\end{exercise}

\begin{exercise}[Conditional independence of label and observation]
\label{exer:sufficiency-conditional-independence}
Let $\Theta,\mathsf X,\mathsf S$ be finite with all subsets
measurable and give every label positive prior mass. Under the
joint prior and observation law, show that $T$ is sufficient if
and only if the label and the full observation are conditionally
independent given $T(X)$. Express the equivalence in terms of
conditional mutual information.
\end{exercise}

\begin{exercise}[Pairwise sufficiency on a countable space]
\label{exer:sufficiency-pairwise-countable}
Let $\mathsf X,\mathsf S$ be countable with all subsets
measurable. Suppose $T$ is sufficient in each two-label
subexperiment. Prove that it is sufficient for the whole
family, allowing an arbitrary parameter set. For each value
$t$, compare only the labels that assign it positive mass.
\end{exercise}

\begin{exercise}[Averaging a square-loss rule]
\label{exer:sufficiency-condition-square-loss}
Let $K$ be a common conditional kernel for $T$ and $d$ a bounded
real-valued rule on $\mathsf X$. Define
$d_T(t)=\int d(x)\,K(t,dx)$. Show that $d_T(T(X))$ has the same
mean as $d(X)$ at every $\theta$. For squared loss at a real
target $g(\theta)$, express the risk decrease as an expected
conditional variance and characterize equality at each label.
\end{exercise}

\begin{exercise}[Reconstruction implies sufficiency in a finite experiment]
\label{exer:sufficiency-reconstruction-finite-converse}
Suppose $\Theta,\mathsf X,\mathsf S$ are finite with all subsets
measurable on the observation and statistic spaces, and
$P_\theta^TK=P_\theta$ for a single kernel $K$. Take a mixture
$M$ with positive weight on every label. Use the same kernel
to reconstruct $M$, and use data processing in both directions
to show $D(P_\theta^T\|M^T)=D(P_\theta\|M)$. Apply the equality
condition in log-sum on each fiber to construct a common
conditional kernel from $M$.
\end{exercise}



\clearpage
\section{Factorization criteria under explicit assumptions}
\label{sec:sufficiency-factorization}

Fix an experiment $\mathcal E=(\mathsf X,\mathcal A,
(P_\theta)_{\theta\in\Theta})$ and a statistic
$T:(\mathsf X,\mathcal A)\to(\mathsf S,\mathcal S)$.

\subsection{A reference law from the model}
\label{subsec:sufficiency-reference-mixture}

\begin{proposition}[A dominating countable mixture]
\label{prop:dominating-countable-mixture}
If $\mathcal E$ is dominated, there is a probability law
$M=\sum_{j\geq1}w_jP_{\theta_j}$ with $w_j>0$ and $\sum_jw_j=1$
such that $P_\theta\ll M$ for every $\theta$.
Consequently, an event is $M$-null exactly when it is null
under every model law.
\end{proposition}
\begin{proof}
Choose a dominating probability measure $\nu$, as in the proof of
Proposition~\ref{prop:statistic-induced-density}, and put
$A_\theta=\{dP_\theta/d\nu>0\}$ using one density version per label.
Let $b$ be the supremum of $\nu(\bigcup_{\theta\in F}A_\theta)$ over
nonempty finite subsets $F\subseteq\Theta$. Choose finite $F_j$ for
which these measures tend to $b$, and put
$A=\bigcup_j\bigcup_{\theta\in F_j}A_\theta$.
By continuity from below applied to the increasing finite unions,
$\nu(A)=b$. For every $\theta$, the same argument applied after
adjoining $A_\theta$ gives $\nu(A\cup A_\theta)\leq b$, so
$\nu(A_\theta\setminus A)=0$.

List the selected labels in a sequence $(\theta_j)_{j\geq1}$,
repeating labels if necessary. For any strictly positive weights
summing to one, the density of the resulting mixture is positive
$\nu$-almost everywhere on $A$. This follows by summing the
component densities by monotone convergence. If $M(B)=0$, then
$\nu(B\cap A)=0$, so $P_\theta(B)=0$ for every $\theta$.
Conversely, an event null under all model laws is null under their
mixture.
\end{proof}

\subsection{The factorization criterion}
\label{subsec:fisher-neyman-factorization}

\begin{definition}[Factorization through a statistic]
\label{def:density-factorization-statistic}
\indexnotation{g_\theta(T(x))h(x)}{Factorization of a model density through a statistic.}
Let $\mu$ be a $\sigma$-finite measure dominating the model.
Densities $f_\theta=dP_\theta/d\mu$ \emph{factor through $T$}
if there are finite nonnegative measurable functions
$h$ on $\mathsf X$ and $g_\theta$ on $\mathsf S$ such that
\[
f_\theta(x)=g_\theta(T(x))h(x)
\qquad \mu\text{-almost surely for each }\theta .
\]
The factor $h$ is common to the whole family. The exceptional
set for this equality may depend on $\theta$.
\end{definition}

\begin{theorem}[Fisher--Neyman factorization criterion]
\label{thm:fisher-neyman-factorization}
Suppose $\mathcal E$ is dominated and let $\mu$ be any
$\sigma$-finite dominating measure. Then the following are equivalent.
\begin{enumerate}
\item $T$ is sufficient.
\item For any dominating countable mixture $M$ of model laws,
      every $dP_\theta/dM$ has a $\sigma(T)$-measurable version.
\item The densities relative to $\mu$ factor through $T$.
\end{enumerate}
\end{theorem}
\begin{proof}
First suppose $T$ is sufficient, and let $M$ be any dominating
countable mixture. For each $A\in\mathcal A$, the function
$c_A\circ T$ from Definition~\ref{def:sufficient-statistic}
is a conditional expectation of $\mathbf 1_A$ under $M$ by averaging
\eqref{eq:sufficiency-common-conditional} over the mixture.
Put $r_\theta=dP_\theta^T/dM^T$, which exists by domination.
Using the conditional expectation identity with the nonnegative
$\sigma(T)$-measurable weight $r_\theta\circ T$ gives
\[
P_\theta(A)=\int c_A\,dP_\theta^T
=\int (c_A\circ T)(r_\theta\circ T)\,dM
=\int_A r_\theta\circ T\,dM .
\]
Thus $r_\theta\circ T$ is a version of $dP_\theta/dM$.

Conversely, it suffices that such measurable versions exist for
one dominating countable mixture $M$. Write $v_\theta$ for the
$\sigma(T)$-measurable density. Choose a $[0,1]$-valued version of
$\mathbb E_M[\mathbf 1_A\mid\sigma(T)]$ and write it as $c_A\circ T$
by Proposition~\ref{prop:real-statistic-factorization}, clipping
$c_A$ to $[0,1]$ if needed. For $B\in\mathcal S$,
\[
\int_B c_A\,dP_\theta^T
=\int_{T^{-1}(B)}(c_A\circ T)v_\theta\,dM
=\int_{A\cap T^{-1}(B)}v_\theta\,dM .
\]
This is \eqref{eq:sufficiency-common-conditional}.

To deduce factorization from the measurable versions, let
$h=dM/d\mu$ and use Proposition~\ref{prop:real-statistic-factorization}
to write $v_\theta=g_\theta\circ T$. Choose finite nonnegative
versions, taking the positive part of $g_\theta$ if needed.
The change-of-measure formula gives $f_\theta=(g_\theta\circ T)h$
$\mu$-almost surely.

Finally, suppose $f_\theta=(g_\theta\circ T)h$ $\mu$-almost surely
for each $\theta$. Choose a dominating countable mixture with
labels $\theta_j$ and weights $w_j$. Put
$G(t)=\sum_jw_jg_{\theta_j}(t)$. On $\{h>0\}$ a density of $M$
is $h(G\circ T)$, and it is zero on $\{h=0\}$.
The countably many component factorizations hold together
$\mu$-almost surely. Since the density of $M$ is integrable,
$0<G(T(x))<\infty$ for $M$-almost every $x$.
For each $\theta$ the ratio of $f_\theta$ to the mixture density
on its positive finite set is $g_\theta(T(x))/G(T(x))$ there,
$M$-almost surely. By domination this ratio gives
$dP_\theta/dM$. Defining $g_\theta/G$ to be zero where
$G\notin(0,\infty)$ supplies the required measurable version.
\end{proof}

\begin{proposition}[Likelihood ratios to a model law]
\label{prop:sufficiency-likelihood-ratios-reference}
Suppose a model law $P_{\theta_0}$ dominates every $P_\theta$.
Then $T$ is sufficient exactly when for every $\theta$ the
likelihood ratio $dP_\theta/dP_{\theta_0}$ has a
$\sigma(T)$-measurable version.
\end{proposition}
\begin{proof}
The model is dominated by $P_{\theta_0}$. Use the mixture that
repeats $\theta_0$ in Theorem~\ref{thm:fisher-neyman-factorization}.
The proof of that theorem shows that one dominating mixture
suffices for the converse.
\end{proof}

\subsection{Applications to sampling models}
\label{subsec:factorization-sampling-models}

\begin{example}[Category counts]
\label{ex:factorization-category-counts}
Let $X_1,\ldots,X_n$ be iid on $\{1,\ldots,m\}$ with probability
vector $p=(p_1,\ldots,p_m)$. With respect to counting measure,
the record density is $\prod_{j=1}^m p_j^{N_j(x)}$, where $N_j$
is the count of category $j$. Use $0^0=1$ when $p_j=N_j(x)=0$.
This factors through $(N_1,\ldots,N_m)$ with $h=1$, so the counts
are sufficient on the whole probability simplex.
\end{example}

\begin{example}[Poisson observations with different exposures]
\label{ex:factorization-poisson-exposures}
Let $X_i$ be independent Poisson variables of means $a_i\lambda$,
with known $a_i>0$ and unknown $\lambda>0$. Write $A=\sum_i a_i$
and $T(x)=\sum_i x_i$. The record mass is
\[
f_\lambda(x)=e^{-A\lambda}\lambda^{T(x)}
\prod_i \frac{a_i^{x_i}}{x_i!}.
\]
The last product is the common factor, so $T$ is sufficient even
when the observations have different means.
\end{example}

\begin{example}[Normal samples with unknown mean and variance]
\label{ex:factorization-normal-two-summaries}
Let $X_i$ be iid $N(u,v)$, with $u\in\R$ and variance $v>0$
both unknown. With $s_1=\sum_i x_i$ and $s_2=\sum_i x_i^2$,
the density relative to Lebesgue measure is
\[
f_{u,v}(x)=(2\pi v)^{-n/2}
\exp\!\left(-\frac{s_2-2us_1+nu^2}{2v}\right).
\]
Thus $(s_1,s_2)$ is sufficient. For $n\geq2$ the sample mean
and sample variance are also sufficient, since the same density factors through them using
$s_1=n\overline x$ and $s_2=(n-1)s^2+n\overline x^{\,2}$, where
$s^2=(n-1)^{-1}\sum_i(x_i-\overline x)^2$.
\end{example}

\begin{example}[A uniform interval with unknown endpoints]
\label{ex:factorization-uniform-two-endpoints}
Let $X_i$ be iid uniform on $[a,b]$, with $a<b$ both unknown.
Their joint density is
\[
(b-a)^{-n}\mathbf 1_{\{a\leq\min_i x_i,\ \max_i x_i\leq b\}} .
\]
The minimum and maximum are jointly sufficient. The entire
parameter-dependent support condition appears in the factor
that depends on the statistic.
\end{example}

\begin{example}[A likelihood ratio that distinguishes signs]
\label{ex:factorization-folded-normal-failure}
For $P_\theta=N(\theta,1)$, the likelihood ratio relative to
$P_0$ is $\exp(\theta x-\theta^2/2)$. If $T(x)=|x|$, a
$\sigma(T)$-measurable version would have the form $g(|x|)$
by Proposition~\ref{prop:real-statistic-factorization}.
Since $P_0$ is invariant under reflection, agreement with
that version almost surely at $x$ also gives agreement at
$-x$ on a set of full $P_0$-measure. For $\theta\ne0$ the
ratios at $x$ and $-x$ differ whenever $x\ne0$. Thus $T$ is
insufficient by Proposition~\ref{prop:sufficiency-likelihood-ratios-reference}.
\end{example}

\begin{example}[Density versions on parameter-dependent null sets]
\label{ex:factorization-parameterwise-null-sets}
On $[0,1]$ let every $P_\theta$ be uniform, with $\theta\in[0,1]$.
Use the density version $f_\theta(x)=1+\mathbf 1_{\{x=\theta\}}$
relative to Lebesgue measure. The constant statistic gives
the factorization $g_\theta=1$, $h=1$ almost everywhere for
each $\theta$. For these versions and factors there is no
single point outside all exceptional sets, since at each
$x$ the equality fails for $\theta=x$.
\end{example}

\subsection{Exercises}
\label{subsec:practice-4-3}

\begin{practiceexercise}[Two model laws with different supports]
\label{pex:factorization-two-law-mixture}
Let $P=(1/2,1/2,0)$ and $Q=(0,1/2,1/2)$ on three outcomes.
Compute $M=(P+Q)/2$ and the densities of $P,Q$ relative to $M$.
Determine whether either law dominates the other.
\end{practiceexercise}

\begin{practiceexercise}[A countable family of point masses]
\label{pex:factorization-countable-point-mixture}
On the positive integers, take the family $(\delta_j)_{j\geq1}$.
Find a dominating mixture and compute $d\delta_j/dM$.
\end{practiceexercise}

\begin{practiceexercise}[Uniform laws with no largest endpoint]
\label{pex:factorization-uniform-mixture}
For $\{\operatorname{Unif}[0,\theta]:\theta>0\}$, give a
countable mixture of model laws that dominates the family.
Show that no single model law dominates it.
\end{practiceexercise}

\begin{practiceexercise}[A parameter-free change of reference]
\label{pex:factorization-reference-rescaling}
Suppose $f_\theta=(g_\theta\circ T)h$ relative to $\mu$ and
$\nu=c\mu$ for $c>0$. Write a factorization relative to $\nu$.
\end{practiceexercise}

\begin{practiceexercise}[Multiplying the common factor]
\label{pex:factorization-positive-rescaling}
For a positive finite measurable function $a$ on the statistic
space, replace $h$ by $(a\circ T)h$. Find new $g_\theta$ that
give the same densities.
\end{practiceexercise}

\begin{practiceexercise}[A finer statistic]
\label{pex:factorization-refinement}
If $T$ is sufficient in a dominated model and $T=b\circ U$
for a measurable $b$, show by factorization that $U$ is sufficient.
\end{practiceexercise}

\begin{practiceexercise}[Binary counts]
\label{pex:factorization-bernoulli-count}
Factor the Bernoulli sample mass through the count directly.
Specify the factors at parameters zero and one.
\end{practiceexercise}

\begin{practiceexercise}[Dropping one category count]
\label{pex:factorization-one-count-redundant}
For a fixed-size sample on $m\geq2$ categories, express the
last count using the first $m-1$. Deduce sufficiency of the
first $m-1$ counts.
\end{practiceexercise}

\begin{practiceexercise}[Unequal Poisson exposures]
\label{pex:factorization-poisson-two-exposures}
In Example~\ref{ex:factorization-poisson-exposures}, take
$a_1=1$ and $a_2=2$. At $\lambda=2$ and $\lambda=1$, compute
the likelihood ratio on the fiber with total three.
\end{practiceexercise}

\begin{practiceexercise}[An exponential rate]
\label{pex:factorization-exponential-total}
For iid exponential draws of rate $\lambda>0$, factor the
joint density through their sum, including the support indicator.
\end{practiceexercise}

\begin{practiceexercise}[A shifted exponential model]
\label{pex:factorization-shifted-exponential}
Let $X_i=\theta+E_i$, where the $E_i$ are iid rate-one
exponential variables. Factor the density through the sample
minimum, keeping the parameter-dependent support explicit.
\end{practiceexercise}

\begin{practiceexercise}[Normal mean with known variance]
\label{pex:factorization-normal-known-variance}
In Example~\ref{ex:factorization-normal-two-summaries}, fix
$v$. Write a factorization through $s_1$ alone.
\end{practiceexercise}

\begin{practiceexercise}[Normal variance with known mean]
\label{pex:factorization-normal-known-mean}
For iid $N(0,v)$ with unknown $v>0$, exhibit a sufficient
statistic by factoring the density.
\end{practiceexercise}

\begin{practiceexercise}[A uniform scale with known left endpoint]
\label{pex:factorization-uniform-one-endpoint}
For iid uniform draws on $[0,\theta]$, factor the density
through the maximum. Specify the common factor on $\R^n$.
\end{practiceexercise}

\begin{practiceexercise}[The range of two uniform draws]
\label{pex:factorization-uniform-range-failure}
For two draws from $\operatorname{Unif}[a,a+1]$ with
$a\in\{0,1/2\}$, compare the range laws and the full record laws.
Decide whether the range is sufficient using
Theorem~\ref{thm:sufficiency-finite-posterior}.
\end{practiceexercise}

\begin{practiceexercise}[Likelihood ratios on finite cells]
\label{pex:factorization-finite-ratio-cells}
On four outcomes, let $Q$ be uniform and $P=(1/8,1/8,3/8,3/8)$.
Check sufficiency for $\{P,Q\}$ when $T$ groups consecutive
pairs. Repeat for the grouping $\{1,3\},\{2,4\}$.
\end{practiceexercise}

\begin{practiceexercise}[Fixed conditional weights]
\label{pex:factorization-split-fixed-weights}
Take $P_p=(p/3,2p/3,1-p)$ for $0\leq p\leq1$. Factor the
record mass through the indicator that the observation lies
in $\{1,2\}$. Use the factorization to recover the two
conditional masses within that cell.
\end{practiceexercise}

\begin{practiceexercise}[Choosing pointwise density versions]
\label{pex:factorization-pointwise-versions}
In Example~\ref{ex:factorization-parameterwise-null-sets},
choose versions that give a pointwise factorization at every
$x,\theta$. Explain why this choice leaves the model unchanged.
\end{practiceexercise}

\begin{exercise}[Likelihood vectors for a finite experiment]
\label{exer:factorization-finite-likelihood-vectors}
Let the parameter set, observation space, and statistic space be
finite, with all subsets measurable in the observation and
statistic spaces. Under an equally weighted mixture of the model laws,
show that the statistic is sufficient exactly when the nonzero
likelihood vectors $(P_\theta(\{x\}))_{\theta\in\Theta}$ of any two
records of positive mixture mass in the same fiber are proportional. Explain how records null under all model
laws enter this criterion.
\end{exercise}

\begin{exercise}[A sufficient statistic with no finite dominating subfamily]
\label{exer:factorization-infinite-dominating-mixture}
For $\theta\in\{1,2,\ldots\}$, observe $(J,V)$ with $J=\theta$
almost surely and $V$ an independent fair bit. Show that no
finite mixture of model laws dominates the family. Give a
dominating countable mixture and use the resulting densities
to show that $J$ is sufficient.
\end{exercise}

\begin{exercise}[Order statistics for iid densities]
\label{exer:factorization-order-statistics}
Let $X_1,\ldots,X_n$ be iid real observations with densities
$f_\theta$ relative to Lebesgue measure. Prove that the
sorted vector is a sufficient statistic. Your factorization
should include models under which some factors can vanish.
\end{exercise}

\begin{exercise}[A two-parameter gamma family]
\label{exer:factorization-gamma-two-parameters}
For $\alpha,\beta>0$, take the density on $(0,\infty)$ given by
\[
f_{\alpha,\beta}(x)=\frac{\beta^\alpha x^{\alpha-1}e^{-\beta x}}
{\Gamma(\alpha)},\qquad
\Gamma(\alpha)=\int_0^\infty t^{\alpha-1}e^{-t}\,dt.
\]
For iid samples find a two-component sufficient statistic.
Then find a one-component sufficient statistic in each
subfamily obtained by fixing one parameter.
\end{exercise}

\begin{exercise}[The mean alone with unknown normal variance]
\label{exer:factorization-normal-mean-unknown-variance}
For two iid $N(0,v)$ draws with $v\in\{1,2\}$, use the
likelihood ratio to show that the sample mean is insufficient.
To justify the almost-everywhere claim, use the coordinates
$(\overline X,X_1-X_2)$ under $v=1$.
\end{exercise}

\begin{exercise}[A nondominated model with a sufficient reduction]
\label{exer:factorization-nondominated-sufficiency}
On $\R\times\{0,1\}$ take $P_\theta=\delta_\theta\otimes U$,
where $U$ is a fixed fair law. Construct a common conditional
kernel for the first coordinate. Show that no $\sigma$-finite
measure dominates the model by considering the disjoint sets
$\{\theta\}\times\{0,1\}$.
\end{exercise}



\clearpage
\section{Minimal sufficiency, completeness, and ancillarity}
\label{sec:minimal-complete-ancillary}

Fix an experiment $\mathcal E=(\mathsf X,\mathcal A,
(P_\theta)_{\theta\in\Theta})$ with $\Theta\ne\varnothing$.

\subsection{Minimal sufficiency}
\label{subsec:minimal-sufficiency}

\begin{definition}[Minimal sufficient statistic]
\label{def:minimal-sufficient-statistic}
\indexnotation{T\preceq U}{Every event determined by \(T\) is determined by \(U\) up to a common model null set.}
For statistics $T,U$, write $T\preceq U$ if for every
$A\in\sigma(T)$ there is a $B\in\sigma(U)$ with
$P_\theta(A\mathbin{\triangle} B)=0$ for every $\theta$, where
$\triangle$ denotes symmetric difference. A sufficient statistic
$T$ is \emph{minimal sufficient} if $T\preceq U$ for every
sufficient statistic $U$.
\end{definition}

\begin{proposition}[Comparison up to common null sets]
\label{prop:statistic-comparison-versions}
The relation $\preceq$ is reflexive and transitive. Moreover,
$T\preceq U$ if and only if every bounded real $\sigma(T)$-measurable
function agrees with a bounded $\sigma(U)$-measurable function
almost surely under every $P_\theta$.
\end{proposition}
\begin{proof}
Reflexivity is immediate. For transitivity, the symmetric difference
between the first and third events is contained in the union of
the two successive symmetric differences. For the function statement,
one direction follows by taking indicators and thresholding the
corresponding function at $1/2$.

For the other direction, approximate a bounded $\sigma(T)$-measurable
function $a$ by simple functions pointwise. Replace the finitely many
level sets at each step by their $\sigma(U)$ versions. The resulting
$\sigma(U)$-measurable simple functions agree with the original ones
off a countable union of common null sets. Take their limit where it
exists, and use zero elsewhere. Clipping to $[-\sup|a|,\sup|a|]$
gives the desired bounded version.
\end{proof}

\begin{theorem}[Minimal sufficiency in a dominated experiment]
\label{thm:minimal-sufficient-density-profile}
Let $M$ be a dominating countable mixture of model laws and choose
finite nonnegative densities $v_\theta=dP_\theta/dM$. The statistic
\[
V(x)=(v_\theta(x))_{\theta\in\Theta},
\]
with value space $[0,\infty)^\Theta$ and its product sigma algebra,
is minimal sufficient. Any two minimal sufficient statistics
determine the same events up to common model null sets.
\end{theorem}
\begin{proof}
Each coordinate density is $\sigma(V)$-measurable, so $V$ is sufficient
by Theorem~\ref{thm:fisher-neyman-factorization}. If $U$ is sufficient,
each coordinate has a $\sigma(U)$-measurable version by the same
theorem. Thus every set $v_\theta^{-1}(C)$ for Borel $C$ has a
$\sigma(U)$ version modulo an $M$-null set. The events in $\mathcal A$
with such versions form a sigma algebra: complements preserve the
version property, and the exceptional set for a countable union is
contained in the union of the exceptional sets. This sigma algebra
contains all the coordinate events, hence contains $\sigma(V)$.
The $M$-null sets are exactly the common model null sets by
Proposition~\ref{prop:dominating-countable-mixture}. Therefore
$V\preceq U$. Mutual comparison of two minimal sufficient statistics
follows from the definition.
\end{proof}

\begin{example}[Two likelihood-ratio levels]
\label{ex:minimal-sufficient-finite-cells}
On four ordered outcomes, let $Q$ be uniform and let $P$ have
masses $(1/8,1/8,3/8,3/8)$. The ratio $dP/dQ$ takes the value
$1/2$ on the first two outcomes and $3/2$ on the last two.
Thus $T(x)=\mathbf 1_{\{x\in\{3,4\}\}}$ is sufficient by
Proposition~\ref{prop:sufficiency-likelihood-ratios-reference}.
For any sufficient $U$, the ratio is $\sigma(U)$-measurable up
to a $Q$-null set. Here every outcome has positive $Q$ mass, so
$T$ itself is $\sigma(U)$-measurable by thresholding the ratio
at one. Hence $T$ is minimal sufficient.
\end{example}

\begin{example}[Minimal normal summaries]
\label{ex:minimal-sufficient-normal-summaries}
In Example~\ref{ex:factorization-normal-two-summaries},
the pair $(s_1,s_2)=(\sum_i x_i,\sum_i x_i^2)$ is sufficient.
Let $P_{u,v}$ be the law of the sample, and use the dominating
model law $P_{0,1}$. Two log likelihood ratios are
\[
\log\frac{dP_{1,1}}{dP_{0,1}}=s_1-n/2,\qquad
\log\frac{dP_{0,2}}{dP_{0,1}}=s_2/4-(n/2)\log 2 .
\]
Each ratio has a version measurable with respect to any
sufficient $U$. Taking logarithms on the positive values gives
versions of $s_1,s_2$ under $P_{0,1}$, hence under all model laws.
Both versions agree together off a common null set, so every
Borel event of the pair has a $\sigma(U)$ version. Thus the
pair is minimal sufficient.
\end{example}

\subsection{Completeness}
\label{subsec:statistic-completeness}

\begin{definition}[Complete statistic]
\label{def:complete-statistic}
A statistic $T$ is \emph{complete} if every measurable real
function $a$ on its value space that satisfies
\[
\mathbb E_\theta |a(T)|<\infty,\qquad \mathbb E_\theta a(T)=0
\qquad(\theta\in\Theta)
\]
also satisfies $P_\theta(a(T)=0)=1$ for every $\theta$.
It is \emph{boundedly complete} if this implication holds for
every bounded measurable real function $a$.
\end{definition}

\begin{proposition}[Uniqueness of an unbiased function of a complete statistic]
\label{prop:complete-unbiased-uniqueness}
Let $T$ be complete. If measurable functions $a(T),b(T)$ are integrable
at each label and have the same expectation at every label, then
$a(T)=b(T)$ almost surely under every $P_\theta$.
For boundedly complete $T$, the same holds for bounded $a,b$.
\end{proposition}
\begin{proof}
Apply completeness to $a-b$. For the bounded case the difference
is bounded.
\end{proof}

\begin{theorem}[Bounded completeness and minimality]
\label{thm:bounded-complete-sufficient-minimal}
In a dominated experiment, every boundedly complete sufficient
statistic is minimal sufficient.
\end{theorem}
\begin{proof}
Choose a minimal sufficient statistic $V$ by
Theorem~\ref{thm:minimal-sufficient-density-profile}. Then $V\preceq T$.
For $A\in\sigma(T)$, sufficiency of $V$ gives a bounded common
conditional probability $c_A(V)$ of $A$ given $V$.
Proposition~\ref{prop:statistic-comparison-versions} gives a bounded
$\sigma(T)$-measurable version of $c_A(V)$, and
Proposition~\ref{prop:real-statistic-factorization} writes this
version as $a(T)$ for bounded measurable $a$ (clipping if needed).
The function $\mathbf 1_A-a(T)$ is a bounded measurable function
of $T$, by the same factorization. Its expectation is zero under
each $P_\theta$. Bounded completeness implies
$\mathbf 1_A=c_A(V)$ almost surely under every label.
The event $\{c_A(V)>1/2\}$ is in $\sigma(V)$ and is a version of $A$.
Thus $T\preceq V$, and transitivity gives minimal sufficiency of $T$.
\end{proof}

\begin{example}[Completeness of the Bernoulli count]
\label{ex:complete-bernoulli-count}
Let $T$ be the count in an iid Bernoulli sample of size $n$ with
$0<p<1$. If $\mathbb E_p a(T)=0$ for every $p$, then
\[
0=(1-p)^{-n}\mathbb E_p a(T)
=\sum_{k=0}^n a(k)\binom nk \left(\frac{p}{1-p}\right)^k .
\]
This polynomial vanishes at every positive argument, so all
coefficients vanish. Thus $T$ is complete. By
Example~\ref{ex:sufficiency-bernoulli-count} it is sufficient,
so it is also minimal sufficient. The same conclusions hold
if $p=0$ and $p=1$ are included.
\end{example}

\begin{example}[Completeness of the uniform maximum]
\label{ex:complete-uniform-maximum}
For $n$ iid draws from $\operatorname{Unif}[0,\theta]$ with
$\theta>0$, take $T=\max_i X_i$. Its induced density is
$nt^{n-1}/\theta^n$ on $(0,\theta)$ by
Example~\ref{ex:statistic-uniform-maximum}. If $a(T)$ is integrable
with mean zero at every $\theta$, then
\[
\int_0^\theta a(t)t^{n-1}\,dt=0\qquad(\theta>0).
\]
For each $m\geq1$, the integrable function $a(t)t^{n-1}$ on $(0,m]$
defines a finite signed measure. Its integrals on $(0,\theta]$
for $0<\theta\leq m$ vanish. These intervals determine the measure
by the $\pi$--$\lambda$ theorem applied to its positive and negative
parts. Thus $a(t)t^{n-1}=0$ almost everywhere on $(0,m]$ for every
$m$. It follows that $a(T)=0$ almost surely at every label.
This proves completeness. The density factors through $T$ with
common factor $\mathbf 1_{\{\min_i x_i\geq0\}}$, so $T$ is also
sufficient and minimal sufficient.
\end{example}

\begin{example}[Minimal sufficiency with an unbiased function of zero]
\label{ex:minimal-sufficient-incomplete-three}
On three ordered outcomes let $P=(1/2,1/3,1/6)$ and
$Q=(1/6,1/3,1/2)$. The likelihood ratio takes the three distinct
values $3,1,1/3$. Its value determines the outcome, so the identity
statistic is minimal sufficient: any sufficient statistic
determines the likelihood ratio by
Proposition~\ref{prop:sufficiency-likelihood-ratios-reference}.
The function with values $(1,-2,1)$ has mean zero under both
laws and is nonzero at every outcome. Thus the statistic fails
bounded completeness.
\end{example}

\subsection{Ancillarity and independence}
\label{subsec:statistic-ancillarity}

\begin{definition}[Ancillary statistic]
\label{def:ancillary-statistic}
A statistic $V$ is \emph{ancillary} if its induced law $P_\theta^V$
is the same for all $\theta\in\Theta$.
\end{definition}

\begin{theorem}[Basu's theorem]
\label{thm:basu-independence}
A boundedly complete sufficient statistic $T$ and an ancillary
statistic $V$ are independent under every $P_\theta$.
\end{theorem}
\begin{proof}
Fix a measurable set $B$ in the value space of $V$. The probability
$p_B=P_\theta(V\in B)$ is independent of $\theta$. By sufficiency
of $T$ there is a common conditional probability $c(T)$ for the
event $\{V\in B\}$ given $\sigma(T)$. Its expectation is $p_B$
for every $\theta$. Bounded completeness applied to $c-p_B$
gives $c(T)=p_B$ almost surely under each $P_\theta$.
For any event $A\in\sigma(T)$, integration over $A$ gives
$P_\theta(A\cap\{V\in B\})=P_\theta(A)p_B$. This proves independence.
\end{proof}

\begin{example}[An ancillary ratio in the uniform scale model]
\label{ex:basu-uniform-ratio}
For two iid $\operatorname{Unif}[0,\theta]$ draws, let
$T=\max(X_1,X_2)$ and $V=\min(X_1,X_2)/T$, defining $V=0$ when
$T=0$. Writing $(X_1,X_2)$ in law as $\theta(U_1,U_2)$ for iid
uniform $U_i$ on $[0,1]$ shows that $V$ is ancillary.
By Example~\ref{ex:complete-uniform-maximum} and Basu's theorem,
$V$ and $T$ are independent under every $\theta$. The common
kernel in Example~\ref{ex:sufficiency-uniform-maximum-pair}
gives the law of $V$ directly: at each $t>0$ the smaller
coordinate is $tU$ with $U$ uniform on $[0,1]$, so $V$ is uniform
on $[0,1]$.
\end{example}

\subsection{Exercises}
\label{subsec:practice-4-4}

\begin{practiceexercise}[A three-point likelihood profile]
\label{pex:minimal-profile-three-points}
On three ordered outcomes take $P=(1/2,1/2,0)$ and $Q=(0,1/2,1/2)$.
Compute the profile $(dP/dM,dQ/dM)$ for $M=(P+Q)/2$.
Which outcomes can a minimal sufficient statistic group together?
\end{practiceexercise}

\begin{practiceexercise}[An outcome null under every label]
\label{pex:minimal-profile-null-outcome}
Append a fourth outcome of zero mass under both laws in
Exercise~\ref{pex:minimal-profile-three-points}.
Explain how the value of the profile at that outcome can be chosen.
\end{practiceexercise}

\begin{practiceexercise}[Transitivity with common null sets]
\label{pex:minimal-statistic-comparison}
On $\{1,2,3\}$ suppose all laws are supported on $\{1,2\}$.
Let $T(x)=\mathbf 1_{\{x=1\}}$ and $U(x)=\mathbf 1_{\{x=2\}}$.
Show $T\preceq U$ and $U\preceq T$. Compare the exact sigma algebras.
\end{practiceexercise}

\begin{practiceexercise}[A one-law experiment]
\label{pex:minimal-one-law}
Show that a constant statistic is minimal sufficient when all
the model laws are equal.
\end{practiceexercise}

\begin{practiceexercise}[Keeping an independent extra bit]
\label{pex:minimal-sufficient-extra-bit}
Let $X=(Y,V)$, where $Y$ is Bernoulli of parameter $p\in(0,1)$
and $V$ is an independent fair bit. Find a minimal sufficient
statistic and explain why the identity statistic is sufficient
but fails minimality.
\end{practiceexercise}

\begin{practiceexercise}[A polynomial on a smaller parameter set]
\label{pex:complete-bernoulli-subinterval}
Show that the count in $n$ Bernoulli draws is complete if the
parameter ranges over a nonempty open subinterval of $(0,1)$.
\end{practiceexercise}

\begin{practiceexercise}[A count at two parameter values]
\label{pex:complete-two-parameter-counterexample}
For two Bernoulli draws with $p\in\{1/4,3/4\}$, find a nonzero
function of the count with expectation zero at both labels.
\end{practiceexercise}

\begin{practiceexercise}[Completeness of a constant]
\label{pex:complete-constant}
Show that a constant statistic is complete for every experiment.
Give a model in which it fails sufficiency.
\end{practiceexercise}

\begin{practiceexercise}[Finite range]
\label{pex:complete-finite-range}
Show that completeness and bounded completeness coincide when
a statistic takes finitely many values.
\end{practiceexercise}

\begin{practiceexercise}[Functions of complete statistics]
\label{pex:complete-measurable-function}
If $T$ is complete and $U=h(T)$ for a measurable map $h$,
show that $U$ is complete. Give a case where $T$ is sufficient
and $U$ is insufficient.
\end{practiceexercise}

\begin{practiceexercise}[Completeness of a uniform maximum on a bounded family]
\label{pex:complete-uniform-bounded-parameter}
In Example~\ref{ex:complete-uniform-maximum}, restrict $\theta$ to
$(0,1)$. Determine whether the maximum is still complete.
\end{practiceexercise}

\begin{practiceexercise}[An ancillary sign]
\label{pex:ancillary-normal-sign}
If $X\sim N(0,\sigma^2)$ with $\sigma>0$ unknown, show that
$\mathbf 1_{\{X>0\}}$ is ancillary. Show directly that it is
independent of $|X|$ under every label.
\end{practiceexercise}

\begin{practiceexercise}[Ancillarity under translation]
\label{pex:ancillary-location-range}
Let $X_i=\theta+Z_i$ with a fixed joint law for $(Z_1,\ldots,Z_n)$.
Show that the sample range is ancillary.
\end{practiceexercise}

\begin{practiceexercise}[Independent fair randomization]
\label{pex:ancillary-fair-randomization}
In Exercise~\ref{pex:minimal-sufficient-extra-bit},
show that $V$ is ancillary. Find a bounded nonzero function of
the full observation with expectation zero for all $p$.
\end{practiceexercise}

\begin{practiceexercise}[Dependent ancillary statistics]
\label{pex:ancillary-dependent-statistics}
Let $Y$ be a Bernoulli bit of parameter $p$ and $V$ an independent fair bit.
Show that $V$ and $1-V$ are ancillary and determine their joint law.
Are they independent?
\end{practiceexercise}

\begin{practiceexercise}[Ancillary coordinates with an informative pair]
\label{pex:ancillary-pair-informative}
With $Y,V$ as in the previous exercise, let $W=(Y+V)\bmod2$.
Show that $W$ is ancillary and compute $P_p(V\ne W)$.
For $p\in(0,1)$, decide whether $(V,W)$ is ancillary.
\end{practiceexercise}

\begin{practiceexercise}[A nonancillary event given a complete statistic]
\label{pex:basu-nonancillary-event}
For an iid Bernoulli sample of size $n\geq2$, let $T$ be its count
and $V=X_1$. Compute the conditional probability of $V=1$ given
$T=k$. Show that $T,V$ are dependent when $0<p<1$.
\end{practiceexercise}

\begin{practiceexercise}[An ancillary function of a complete statistic]
\label{pex:ancillary-function-complete}
Suppose $T$ is boundedly complete and $V=h(T)$ is ancillary.
For a measurable $B$ in the value space of $V$, show that the
common probability $P_\theta(V\in B)$ is either zero or one.
\end{practiceexercise}

\begin{exercise}[Proportional likelihood vectors]
\label{exer:minimal-finite-likelihood-vectors}
Consider a finite experiment with all observation events measurable.
On outcomes with positive mass under the equally weighted mixture
of the model laws, group the outcomes whose likelihood vectors are
proportional. Show that the class label gives a minimal sufficient
statistic. Assign any value to outcomes of zero mixture mass and
justify that freedom.
\end{exercise}

\begin{exercise}[A complete count at finitely many parameter values]
\label{exer:complete-binomial-finite-parameter-set}
For $n\geq1$, show that the count in $n$ Bernoulli draws with
parameter set $\Theta\subset(0,1)$ is complete exactly when
$\Theta$ contains at least $n+1$ distinct values. Use polynomial
uniqueness in one direction and construct a polynomial vanishing
at all the transformed parameter values in the other.
\end{exercise}

\begin{exercise}[Bounded completeness without completeness]
\label{exer:bounded-complete-not-complete}
On the positive integers let $P_j$ give masses $2/3$ and $1/3$ to
$j$ and $j+1$, respectively, for every $j\geq1$. Show that the
identity statistic is boundedly complete. Construct a real function
with mean zero and finite absolute expectation under each $P_j$
that demonstrates failure of completeness.
\end{exercise}

\begin{exercise}[Minimal sufficiency with two unknown uniform endpoints]
\label{exer:minimal-uniform-endpoints}
Let $n\geq2$ and let the model consist of iid samples from
$\operatorname{Unif}[a,b]$ for all real $a<b$. Prove minimal
sufficiency of the sample minimum and maximum. Use a countable
mixture of the model laws with rational endpoints, giving each
positive weight. The positivity sets of the corresponding
densities relative to the mixture give versions of the events
that the entire sample lies in each rational interval. Use these
events to recover both extrema up to a common null set.
\end{exercise}

\begin{exercise}[The normalized range in a uniform scale family]
\label{exer:basu-uniform-normalized-range}
For $n\geq2$ iid draws from $\operatorname{Unif}[0,\theta]$,
let $T$ be the maximum and $V=(T-\min_i X_i)/T$ when $T>0$,
with $V=0$ otherwise. Show that $V$ and $T$ are independent
at every label. Derive the distribution function of $V$ on
$[0,1]$ from the common conditional kernel in
Exercise~\ref{exer:sufficiency-uniform-maximum-general}.
\end{exercise}

\begin{exercise}[Why bounded completeness matters in Basu's theorem]
\label{exer:basu-incomplete-sufficiency}
Take $X=(Y,V)$ with independent $Y\sim\operatorname{Bernoulli}(p)$
and a fair bit $V$, and let $T=X$. Verify sufficiency of $T$
and ancillarity of $V$. Show that $T$ and $V$ are dependent and
identify a bounded function that violates completeness of $T$.
\end{exercise}



\part{Statistical inference}
\label{part:statistical-inference}
\chapter{Point estimation and finite-sample optimality}
\label{chap:point-estimation}
\notationgroup{Point estimation}
\section{Estimators, bias, variance, and mean squared error}
\label{sec:estimators-bias-variance}

Fix an experiment $\mathcal E=(\mathsf X,\mathcal A,
(P_\theta)_{\theta\in\Theta})$ and a real target $g:\Theta\to\R$.
Write $\mathbb E_\theta$ and $\operatorname{Var}_\theta$ for
expectation and variance under $P_\theta$.

\subsection{Estimators and squared risk}
\label{subsec:point-estimators}

\begin{definition}[Point estimator and estimate]
\label{def:point-estimator}
\indexnotation{\widehat g,\ \widehat g(x)}{Estimator of a target \(g\) and its value at the observed record \(x\).}
A \emph{point estimator} of $g$ is a measurable function
$\widehat g:\mathsf X\to\R$. The value $\widehat g(x)$ at an
observed record $x$ is the \emph{estimate}. The estimator is a
deterministic decision rule with real actions.
\end{definition}

\begin{definition}[Bias, unbiasedness, and mean squared error]
\label{def:bias-mse}
\indexnotation{b_\theta(\widehat g),\ \operatorname{MSE}_\theta(\widehat g)}{Bias and mean squared error at parameter \(\theta\) for the target under study.}
At a label where $\widehat g$ is integrable, its bias is
$b_\theta(\widehat g)=\mathbb E_\theta\widehat g-g(\theta)$.
The estimator is \emph{unbiased} for $g$ if it is integrable
and its bias is zero at every label. Its \emph{mean squared error}
is the squared-loss risk
\[
\operatorname{MSE}_\theta(\widehat g)
=\mathbb E_\theta[(\widehat g-g(\theta))^2]\in[0,\infty].
\]
Its variance at a label with finite second moment is
$\operatorname{Var}_\theta(\widehat g)$.
\end{definition}

\begin{proposition}[Finite mean squared error and its decomposition]
\label{prop:estimator-mse-decomposition}
At each label, $\operatorname{MSE}_\theta(\widehat g)$ is finite
exactly when $\mathbb E_\theta[\widehat g^2]<\infty$. In this case
\[
\operatorname{MSE}_\theta(\widehat g)
=\operatorname{Var}_\theta(\widehat g)+b_\theta(\widehat g)^2.
\]
For an unbiased estimator of finite second moment the mean squared
error equals the variance.
\end{proposition}
\begin{proof}
At a fixed label, each of $\widehat g^2$ and
$(\widehat g-g(\theta))^2$ is bounded by twice the other plus
$2g(\theta)^2$, by $(u+v)^2\leq2u^2+2v^2$. This proves the
equivalence. Apply Proposition~\ref{prop:bias-variance-risk} to the
deterministic rule $\widehat g$.
\end{proof}

\begin{proposition}[Identifiability is necessary for unbiased estimation]
\label{prop:unbiased-implies-identifiable}
If $g$ admits an unbiased estimator in $\mathcal E$, then $g$ is
identifiable. If $\widehat g=h\circ T$ is unbiased for $g$ for a
statistic $T$ and a measurable real $h$, then $g$ is identifiable
in the induced experiment $\mathcal E^T$.
\end{proposition}
\begin{proof}
If $P_\theta=P_\vartheta$ then the expectations of $\widehat g$ agree.
By unbiasedness they equal $g(\theta)$ and $g(\vartheta)$, respectively.
For the reduced experiment use
$\mathbb E_\theta[h(T)]=\int h\,dP_\theta^T$ from
Proposition~\ref{prop:statistic-integration-composition} and apply
the same argument to $h$.
\end{proof}

\begin{example}[An identifiable target with no unbiased estimator]
\label{ex:identifiable-no-unbiased-single-bit}
Observe one Bernoulli bit with parameter $p\in[0,1]$ and let
$g(p)=p^2$. The target is identifiable since the law determines $p$.
An estimator with values $a$ at zero and $b$ at one has mean
$(1-p)a+pb$, which is affine in $p$. Matching $p^2$ at the
endpoints forces $a=0$ and $b=1$, giving mean $1/2$ at $p=1/2$
where the target is $1/4$. Thus no estimator is unbiased for this
target with one observation.
\end{example}

\subsection{Estimating population moments}
\label{subsec:estimating-population-moments}

\begin{proposition}[Estimating a mean and a variance]
\label{prop:sample-mean-variance-moments}
Let $X_1,\ldots,X_n$ be iid real random variables with mean $\mu$
and finite variance $\sigma^2$. The sample mean $\overline X$ has
mean $\mu$ and mean squared error $\sigma^2/n$ for estimating $\mu$.
For $n\geq2$,
\[
S_n^2=\frac1{n-1}\sum_{i=1}^n(X_i-\overline X)^2
\]
is unbiased for $\sigma^2$. With denominator $n$ instead, the
estimator has bias $-\sigma^2/n$.
\end{proposition}
\begin{proof}
The mean and its squared risk follow from
Proposition~\ref{prop:empirical-moments} with $f(x)=x$. Expanding
about $\overline X$ gives
\[
\sum_i(X_i-\mu)^2
=\sum_i(X_i-\overline X)^2+n(\overline X-\mu)^2,
\]
since $\sum_i(X_i-\overline X)=0$. The first sum on the left has
expectation $n\sigma^2$ and the last term has expectation $\sigma^2$.
Thus the residual sum of squares has expectation $(n-1)\sigma^2$,
which yields both claims about denominators.
\end{proof}

\begin{example}[Estimating a squared probability with two bits]
\label{ex:unbiased-squared-probability-two-bits}
For two iid Bernoulli bits, $\widehat g=X_1X_2$ is a Bernoulli
variable of mean $p^2$ by independence. Its squared risk for
$p^2$ is $p^2(1-p^2)$. The estimator $\overline X^{\,2}$ instead
has mean
\[
\mathbb E_p[\overline X^{\,2}]=p^2+\tfrac12p(1-p)
\]
by the mean and variance of $\overline X$. It has positive bias
at every $0<p<1$.
\end{example}

\begin{example}[Two observations and the sample variance]
\label{ex:sample-variance-two-bits}
For two iid Bernoulli bits, $S_2^2=(X_1-X_2)^2/2$.
It equals $1/2$ with probability $2p(1-p)$ and is zero otherwise.
Its mean is $p(1-p)$ and its squared risk for that target is
\[
\operatorname{Var}_p(S_2^2)
=\tfrac12p(1-p)-p^2(1-p)^2.
\]
The estimator with denominator two is $S_2^2/2$, so it has
bias $-p(1-p)/2$.
\end{example}

\subsection{Comparing squared risks}
\label{subsec:comparing-estimator-squared-risks}

\begin{proposition}[Projection to a known interval]
\label{prop:estimator-interval-projection}
Suppose $g(\theta)\in[a,b]$ at every label, with $a\leq b$ finite.
Define $\widehat g_{[a,b]}=\min\{b,\max\{a,\widehat g\}\}$.
Its squared error is at most that of $\widehat g$ at every
observation and label. If $\widehat g$ has finite mean squared
error at a label and falls outside $[a,b]$ with positive probability
there, then projection strictly reduces the mean squared error
at that label.
\end{proposition}
\begin{proof}
For $z\in[a,b]$ the projection fixes $z$. For $z<a$, replacing $z$
by $a$ strictly decreases its distance to any $g(\theta)\in[a,b]$.
For $z>b$ the same follows by replacing $z$ by $b$.
Squaring and integrating proves the claims. In the strict case
the nonnegative difference of the squared errors is integrable
and positive on an event of positive probability.
\end{proof}

\begin{proposition}[Affine shrinkage toward a constant]
\label{prop:estimator-affine-shrinkage}
Suppose $U$ is unbiased for $g$ and square integrable at every
label. For $c\in\R$ and $0\leq\lambda\leq1$, put
$U_\lambda=(1-\lambda)U+\lambda c$. Then
\[
\operatorname{MSE}_\theta(U_\lambda)
=(1-\lambda)^2\operatorname{Var}_\theta(U)
+\lambda^2(c-g(\theta))^2.
\]
Writing $v=\operatorname{Var}_\theta(U)$ and $d=c-g(\theta)$,
the risk at a fixed label is minimized at
$\lambda=v/(v+d^2)$ when $v+d^2>0$. When $v=d=0$ all weights
give zero risk.
\end{proposition}
\begin{proof}
The bias is $\lambda d$ and the variance is $(1-\lambda)^2v$.
Apply Proposition~\ref{prop:estimator-mse-decomposition}.
If $v+d^2>0$, completing the square gives a quadratic with
positive coefficient $v+d^2$ and vertex at $v/(v+d^2)$, which
lies in $[0,1]$.
\end{proof}

The minimizing weight is computed at a fixed label. A rule using
one fixed weight is evaluated by substituting that same weight
at each label.

\begin{example}[Estimating a uniform endpoint]
\label{ex:mse-uniform-endpoint-comparison}
For $n$ iid draws from $\operatorname{Unif}[0,\theta]$ with
$\theta>0$, let $T$ be the sample maximum. Integrating its
density from Example~\ref{ex:statistic-uniform-maximum} gives
$\mathbb E_\theta[T^k]=n\theta^k/(n+k)$ for $k=1,2$.
Thus $T$ has bias $-\theta/(n+1)$ and variance
$n\theta^2/((n+1)^2(n+2))$. The rescaled estimator
$U=(n+1)T/n$ is unbiased. The two risks are
\[
\operatorname{MSE}_\theta(T)=\frac{2\theta^2}{(n+1)(n+2)},\qquad
\operatorname{MSE}_\theta(U)=\frac{\theta^2}{n(n+2)}.
\]
Their ratio $\operatorname{MSE}_\theta(U)/\operatorname{MSE}_\theta(T)$
is $(n+1)/(2n)$, which equals one at $n=1$ and is smaller for $n>1$.
\end{example}

\begin{example}[Clipping a normal mean]
\label{ex:mse-clipped-normal-mean}
For $n$ iid draws from $N(\theta,1)$ with $0\leq\theta\leq1$,
$\overline X$ is unbiased with squared risk $1/n$.
The clipped mean $\min\{1,\max\{0,\overline X\}\}$ has strictly
smaller squared risk at each label by
Proposition~\ref{prop:estimator-interval-projection}, since the
normal law of $\overline X$ gives positive mass outside $[0,1]$.
At $\theta=0$ the clipped estimator is nonnegative and is positive
with positive probability, so its bias is strictly positive there.
\end{example}

\begin{example}[Shrinkage with crossing risks]
\label{ex:mse-normal-shrinkage-crossing}
For $X\sim N(\theta,1)$ with $\theta\in\R$, the estimator $X$
has squared risk one. Taking $\lambda=1/2$ and $c=0$ in
Proposition~\ref{prop:estimator-affine-shrinkage}, the risk of
$X/2$ is $(1+\theta^2)/4$. It is smaller for $|\theta|<\sqrt 3$,
equal at $|\theta|=\sqrt 3$, and larger for $|\theta|>\sqrt 3$.
\end{example}

\subsection{Exercises}
\label{subsec:practice-5-1}

\begin{practiceexercise}[A constant estimate]
\label{pex:mse-constant-estimate}
For a real target $g$ and constant estimate $c$, compute the bias,
variance, and mean squared error at each label.
\end{practiceexercise}

\begin{practiceexercise}[Changing the units of a target]
\label{pex:mse-affine-units}
Suppose $\widehat g$ has finite second moment and estimates $g$.
Express the bias, variance, and mean squared error of
$a\widehat g+b$ for estimating $ag+b$ in terms of those of $\widehat g$.
\end{practiceexercise}

\begin{practiceexercise}[A shifted sample mean]
\label{pex:mse-shifted-sample-mean}
For iid observations of mean $\mu$ and variance $\sigma^2$, take
$\overline X+c$ for a fixed $c\in\R$. Compute its squared risk for
estimating $\mu$ and find the minimizing $c$.
\end{practiceexercise}

\begin{practiceexercise}[Finite mean and infinite mean squared error]
\label{pex:mse-uniform-power}
Let $X$ be uniform on $(0,1]$ and use $X^{-a}$ to estimate zero,
with $a>0$. Determine the values of $a$ for which the bias is
finite and the values for which the mean squared error is finite.
\end{practiceexercise}

\begin{practiceexercise}[Indistinguishable labels]
\label{pex:mse-same-law-two-targets}
Suppose $P_0=P_1$ but $g(0)=0$ and $g(1)=1$. For an integrable
estimator, express both biases in terms of its common mean.
What is the smallest possible maximum of their absolute values?
\end{practiceexercise}

\begin{practiceexercise}[Estimating a cubic probability]
\label{pex:mse-cubic-bernoulli-three-bits}
For three iid Bernoulli observations, find an unbiased estimator
of $p^3$ taking values in $\{0,1\}$. Compute its squared risk.
\end{practiceexercise}

\begin{practiceexercise}[A biased plug-in estimator]
\label{pex:mse-squared-frequency-bias}
For $n$ iid Bernoulli observations, compute the bias of
$\overline X^{\,2}$ for estimating $p^2$. At which parameters
does the bias vanish?
\end{practiceexercise}

\begin{practiceexercise}[Comparing means of nested samples]
\label{pex:mse-two-sample-sizes}
Let $X_1,\ldots,X_n$ be iid with variance $\sigma^2$. For
$1\leq m<n$, compare the mean squared error of the mean of
the first $m$ draws to that of the full sample mean.
\end{practiceexercise}

\begin{practiceexercise}[Repeated observations]
\label{pex:mse-repeated-single-draw}
Suppose $X_i=Z$ for all $i$, with $\mathbb E_\theta Z=g(\theta)$
and $\operatorname{Var}_\theta(Z)=v_\theta<\infty$.
Compute the mean squared error of $\overline X$ and explain
which covariance terms contribute to it.
\end{practiceexercise}

\begin{practiceexercise}[Two forms of the sample variance]
\label{pex:mse-sample-variance-raw-moments}
Show that for $n\geq2$,
$(n-1)S_n^2=\sum_i X_i^2-n\overline X^{\,2}$.
Use this expression to compute the expectation of $S_n^2$ for
iid real observations with finite variance.
\end{practiceexercise}

\begin{practiceexercise}[Variance with a known mean]
\label{pex:mse-known-mean-variance}
For iid real observations with known mean $\mu$, show that
$n^{-1}\sum_i(X_i-\mu)^2$ is unbiased for $\sigma^2$. If the
fourth central moment is finite, express its mean squared
error in terms of that moment and $\sigma^2$.
\end{practiceexercise}

\begin{practiceexercise}[A nonnegative target]
\label{pex:mse-project-half-line}
Suppose $g(\theta)\geq0$ for every $\theta$. Compare the pointwise
squared errors of $\widehat g$ and $\max\{0,\widehat g\}$.
Give a condition for strict improvement in mean squared error.
\end{practiceexercise}

\begin{practiceexercise}[Clipping at a boundary]
\label{pex:mse-clipping-boundary-bias}
Let $X\sim N(\theta,1)$ with $\theta\in[0,1]$. Determine the sign
of the bias of the clipped estimator at $\theta=1$ and show
that its bias vanishes at $\theta=1/2$ by reflection.
\end{practiceexercise}

\begin{practiceexercise}[A known bound on a normal mean]
\label{pex:mse-bounded-normal-shrinkage-weight}
For $X\sim N(\theta,1)$ with $|\theta|\leq M$, $M>0$, express
the worst-case risk of $(1-\lambda)X$ over this interval and
minimize it over $0\leq\lambda\leq1$.
\end{practiceexercise}

\begin{practiceexercise}[When shrinkage improves risk]
\label{pex:mse-shrinkage-improvement}
In Proposition~\ref{prop:estimator-affine-shrinkage}, assume
$v>0$ at a fixed label. Find the weights $\lambda\in(0,1]$
for which $U_\lambda$ has strictly smaller risk than $U$.
\end{practiceexercise}

\begin{practiceexercise}[A sample mean for a uniform endpoint]
\label{pex:mse-uniform-double-mean}
For $n$ iid draws from $\operatorname{Unif}[0,\theta]$, show
that $2\overline X$ is unbiased for $\theta$. Compute its
squared risk and compare it with that of the unbiased rescaled
maximum in Example~\ref{ex:mse-uniform-endpoint-comparison}.
\end{practiceexercise}

\begin{practiceexercise}[Averaging two independent estimates]
\label{pex:mse-average-independent-estimates}
Let $U,V$ be independent at a fixed label, both unbiased for
the same real target there, with variances $v_U,v_V>0$.
Compute the squared risk of $wU+(1-w)V$ and find the minimizing
$w\in[0,1]$ at that label.
\end{practiceexercise}

\begin{practiceexercise}[Zero squared risk]
\label{pex:mse-zero-risk}
Show that zero mean squared error at a label is equivalent to
the estimate equaling the target almost surely there. If two
labels give mutually absolutely continuous laws and distinct
target values, can the same estimator have zero risk at both?
\end{practiceexercise}

\begin{exercise}[Which Bernoulli targets admit unbiased estimation?]
\label{exer:unbiased-bernoulli-polynomial-targets}
For $n\geq1$ iid Bernoulli draws with $p\in[0,1]$, show that
a target admits an unbiased estimator exactly when it is the
restriction of a polynomial of degree at most $n$. For the
converse, construct an unbiased estimator of each monomial
of degree $0$ through $n$.
\end{exercise}

\begin{exercise}[Best rescaling of the uniform maximum]
\label{exer:mse-optimal-uniform-maximum-rescaling}
In Example~\ref{ex:mse-uniform-endpoint-comparison}, minimize
the squared risk of $cT$ over $c\in\R$ at every $\theta>0$.
Show that one common $c$ is optimal throughout the model.
Compute its bias and risk, and compare its risk with both
estimators in the example for every $n\geq1$.
\end{exercise}

\begin{exercise}[Correlated observations and an error floor]
\label{exer:mse-common-noise-floor}
Let $X_i=\theta+Z+\varepsilon_i$, where $Z,\varepsilon_1,\ldots,
\varepsilon_n$ are independent, centered, and have variances
$\tau^2,\sigma^2,\ldots,\sigma^2$. Compute the squared risk of
$\overline X$ for $\theta$ and the expectation of $S_n^2$
for $n\geq2$. Determine which variance $S_n^2$ estimates.
\end{exercise}

\begin{exercise}[Gaussian variance with a variable denominator]
\label{exer:mse-normal-variance-denominator}
Let $X_i$ be iid $N(\mu,\sigma^2)$, $n\geq2$, and put
$Q=\sum_i(X_i-\overline X)^2$. Use the probability fact
$Q/\sigma^2\sim\chi^2_{n-1}$, with mean $n-1$ and variance
$2(n-1)$, to compute the squared risk of $Q/d$ for
$d>0$. Find the minimizing $d$ and compare it with $n-1$ and $n$.
\end{exercise}

\begin{exercise}[One shrinkage weight on a restricted Bernoulli model]
\label{exer:mse-bernoulli-fixed-shrinkage-interval}
For $n$ iid Bernoulli observations, take
$U_\lambda=(1-\lambda)\overline X+\lambda/2$.
Fix $0<a<1/2$ and find all weights $\lambda\in(0,1]$
that strictly improve on $\overline X$ in squared risk at
every $p\in[a,1-a]$. Determine whether any such weight
improves on the whole closed unit interval.
\end{exercise}

\begin{exercise}[Minimum squared risk among linear unbiased averages]
\label{exer:mse-heteroscedastic-weighted-average}
Suppose independent $X_1,\ldots,X_n$ have common mean $\theta$
and known positive finite variances $v_1,\ldots,v_n$. Among
estimators $\sum_i w_iX_i$ with $\sum_i w_i=1$, prove that the
unique minimum-risk weights are proportional to $1/v_i$.
Compute the minimum and characterize when equal weights attain it.
\end{exercise}



\clearpage
\section{Maximum likelihood and other estimating constructions}
\label{sec:estimator-constructions}

Fix a dominated experiment $\mathcal E=(\mathsf X,\mathcal A,
(P_\theta)_{\theta\in\Theta})$. Fix a common reference measure and
density versions as in Definition~\ref{def:likelihood}.

\subsection{Maximum likelihood as a selection}
\label{subsec:maximum-likelihood-selection}

\begin{definition}[Maximizing set and maximum likelihood estimator]
\label{def:maximum-likelihood-estimator}
\indexnotation{\mathcal M(x),\ \widehat\theta_{\mathrm{ML}}}{Likelihood maximizing set and a measurable choice of maximizer.}
At a record $x$ put
\[
\mathcal M(x)=\argmax_{\theta\in\Theta}L_x(\theta).
\]
If $\Theta$ is equipped with a sigma algebra, a \emph{maximum
likelihood estimator} is a measurable map
$\widehat\theta_{\mathrm{ML}}:\mathsf X\to\Theta$ such that
$\widehat\theta_{\mathrm{ML}}(x)\in\mathcal M(x)$ on a measurable
set of full probability under every model law. A measurable
real-valued target $g$ then gives the estimate
$g(\widehat\theta_{\mathrm{ML}}(x))$.
\end{definition}

By the convention for $\argmax$, $\mathcal M(x)$ is empty when the
supremum is unattained, and equals $\Theta$ when all likelihood
values are zero. Maximizing $\ell_x$ gives the same set.

\begin{proposition}[Measurable maximizers in two cases]
\label{prop:ml-measurable-maximizers}
If $\Theta$ is finite with all subsets measurable, the first
maximizer in any fixed ordering is a maximum likelihood estimator.
If $\Theta$ is a nonempty compact subset of $\R^d$ with its Borel
sets, suppose that for each record $x$ the map $\theta\mapsto L_x(\theta)$
is continuous and has a unique maximizer. The unique maximizing
choice is measurable.
\end{proposition}
\begin{proof}
For finite $\Theta$, the event that the first maximizing index is $i$
is given by strict likelihood comparisons with all earlier indices and
weak comparisons with all later ones. Each comparison is measurable
because the likelihood at a fixed label is a measurable function of $x$.

For compact $\Theta$, choose a countable dense sequence $(q_j)$ in
$\Theta$, allowing repetitions. Continuity and compactness give
$s(x):=\sup_j L_x(q_j)=\max_\theta L_x(\theta)$, so $s$ is measurable.
For each $k\geq1$, choose the first $q_j$ with
$L_x(q_j)>s(x)-1/k$. This choice is measurable by countably many
measurable comparisons. At fixed $x$, every convergent subsequence
of the chosen points has limit maximizing $L_x$, by continuity.
Compactness and uniqueness therefore give convergence to the
maximizer. The pointwise limit is measurable.
\end{proof}

\subsection{Reduction and attainment}
\label{subsec:ml-reduction-attainment}

\begin{proposition}[Likelihood maximization after a sufficient reduction]
\label{prop:ml-sufficient-reduction}
Suppose $T$ is sufficient and $M$ is a dominating countable mixture
of model laws. There are density choices relative to $M$ and $M^T$
for which the full likelihood at $x$ equals the induced likelihood
at $T(x)$ for every $x,\theta$. With these choices, a maximum
likelihood estimator for the induced experiment lifts by composition
with $T$. For countable $\Theta$, arbitrary density choices relative
to these same references give the same conclusion up to common
model null sets.
\end{proposition}
\begin{proof}
Choose finite nonnegative induced densities
$r_\theta=dP_\theta^T/dM^T$. By the proof of
Theorem~\ref{thm:fisher-neyman-factorization}, $r_\theta\circ T$
is a full density of $P_\theta$ relative to $M$. With these choices
the maximizing sets at $x$ and $T(x)$ agree. The preimage of a set
of full probability under every $P_\theta^T$ has full probability
under every $P_\theta$, and composition preserves measurability.
For countable $\Theta$, apply
Proposition~\ref{prop:countable-density-agreement} in each
experiment and take the union of the full-experiment exceptional
set and the preimage of the induced-experiment exceptional set.
\end{proof}

\begin{example}[An open parameter interval for Bernoulli trials]
\label{ex:ml-open-bernoulli}
For $n$ iid Bernoulli trials with $p\in(0,1)$, the sample
proportion maximizes the likelihood on records with $0<k<n$
successes by Example~\ref{ex:bernoulli-likelihood}. On the all-zero
record the likelihood is $(1-p)^n$, with supremum one approached
as $p\downarrow0$. On the all-one record it is $p^n$, with supremum
one approached as $p\uparrow1$. Each record has positive probability
at every allowed $p$, so there is no maximum likelihood estimator
on a common full-probability set. With parameter space $[0,1]$
the same supremum is attained at the corresponding endpoint.
\end{example}

\begin{example}[Gaussian location and variance]
\label{ex:ml-normal-location-variance}
For iid observations from $N(u,v)$ with $u\in\R$ and $v>0$, write
$Q(x)=\sum_i(x_i-\overline x)^2$. For fixed $v$, the decomposition in
Example~\ref{ex:gaussian-likelihood} gives the unique optimum
$u=\overline x$. The log likelihood there is
\[
-\frac n2\log(2\pi v)-\frac{Q(x)}{2v},
\]
whose derivative is $(Q(x)-nv)/(2v^2)$. If $Q(x)>0$ it increases
up to $v=Q(x)/n$ and decreases afterwards, giving the unique optimum.
If $Q(x)=0$ it tends to $+\infty$ as $v\downarrow0$.
For $n\geq2$ the latter event has probability zero under every model
law since $X_1=X_2$ has probability zero. Thus the formula defines
an estimator after assigning, for example, $(\overline x,1)$ on that
event. For $n=1$ the log likelihood is unbounded at every record.
For $n\geq2$, the estimator of $v$ has bias $-v/n$ by
Proposition~\ref{prop:sample-mean-variance-moments}.
\end{example}

\begin{example}[Uniform endpoint and density versions]
\label{ex:ml-sample-uniform-versions}
Let $X_1,\ldots,X_n$ be iid uniform on $[0,\theta]$, $\theta>0$.
On the common full-probability set of nonnegative records with
maximum $T>0$, the closed-support density versions in
Example~\ref{ex:uniform-likelihood-versions} give likelihood
$\theta^{-n}\mathbf 1_{\{\theta\geq T\}}$. It is maximized at
$\theta=T$. Using the half-open versions for each observation
gives $\theta^{-n}\mathbf 1_{\{\theta>T\}}$, whose supremum is
$T^{-n}$ and is unattained. For each fixed $\theta$ the two
product densities agree almost everywhere. The maximizing
set changes because the excluded endpoint in the half-open
version depends on $\theta$.
\end{example}

\subsection{Moment matching and minimum distance}
\label{subsec:moment-matching-minimum-distance}

\begin{definition}[Method of moments and minimum distance]
\label{def:moment-matching}
\indexnotation{m(\theta),\ \widehat m_n}{Model moment vector and empirical moment vector used in moment matching.}
For iid observations $X_1,\ldots,X_n$ from a law $P_\theta$, let
$\psi$ be a measurable $\R^k$-valued function whose coordinates
are integrable under every $P_\theta$. Write
\[
m(\theta)=\int\psi\,dP_\theta,\qquad
\widehat m_n(x)=\frac1n\sum_{i=1}^n\psi(x_i).
\]
A \emph{method of moments} estimator is a measurable choice of
a solution to $m(\theta)=\widehat m_n(x)$ on a common full-probability
set. Given a metric $d$ on $\R^k$, a minimum-distance construction
instead minimizes $d(\widehat m_n(x),m(\theta))$ over $\theta$ and
requires a measurable choice of minimizer on such a set.
\end{definition}

\begin{proposition}[Moment fitting over a compact parameter set]
\label{prop:moment-fitting-compact}
Let $\Theta$ be a nonempty compact subset of $\R^d$ with its Borel sets,
and let $m:\Theta\to\R^k$ be continuous and injective. Exact moment
matching has a unique solution at each record for which
$\widehat m_n$ lies in $m(\Theta)$. That solution is measurable
on the set of such records. If $m(\Theta)$ is convex, Euclidean
minimum-distance fitting gives a unique measurable estimator
at every record. It agrees with exact matching when a solution exists.
\end{proposition}
\begin{proof}
The continuous bijection from compact $\Theta$ to $C=m(\Theta)$
has a continuous inverse, since images of closed subsets are compact
and hence closed. The empirical moment vector is measurable and
$C$ is closed. Composing with the inverse on the preimage of $C$
gives the measurable exact solution.

For any $y\in\R^k$, a nearest point in $C$ exists by compactness.
If two distinct points minimize squared Euclidean distance to $y$,
their midpoint belongs to $C$ and has squared distance equal to
their common value minus a quarter of their squared separation.
Hence the nearest point is unique. The nearest-point map is
continuous: if $y_j\to y$, any convergent subsequence of the nearest
points at $y_j$ has limit minimizing distance to $y$, by taking
limits in the minimizing inequalities. Compactness and uniqueness
give convergence of the full sequence. Composing this map with
$\widehat m_n$ and $m^{-1}$ gives the estimator. The nearest
point to a member of $C$ is itself.
\end{proof}

\begin{example}[Matching the mean of a uniform law]
\label{ex:mom-uniform-endpoint}
For the uniform endpoint model, the first model moment is
$m(\theta)=\theta/2$. With $\psi(x)=x$, exact matching gives
$\widehat\theta=2\overline X$ almost surely at every label.
This estimator is unbiased, and its variance is $\theta^2/(3n)$
since the variance of one draw is $\theta^2/12$. Under the
closed-support density choice the maximum likelihood estimator is
the sample maximum, whose bias and risk were computed in
Example~\ref{ex:mse-uniform-endpoint-comparison}.
\end{example}

\begin{example}[A moment constraint]
\label{ex:mom-poisson-compact}
Let $X_i$ be iid Poisson with rate $\lambda\in[1,3]$. The first
model moment is $\lambda$. If the sample mean lies outside $[1,3]$,
no exact first-moment match exists. Minimum-distance fitting for the
usual distance on $\R$ clips the sample mean to $[1,3]$.
The log likelihood has the form
$n\overline x\log\lambda-n\lambda+C(x)$, where $C(x)$ depends
only on the record. Its derivative has the sign of
$\overline x-\lambda$, so it is maximized at the same clipped value.
\end{example}

\begin{example}[One moment and two labels]
\label{ex:mom-normal-first-moment}
Observe iid draws from either $N(0,1)$ or $N(0,2)$, where the second
argument is the variance. Both first moments are zero, so matching
the empirical first moment gives either both labels or none.
The second model moments are one and two. Minimum-distance fitting
using $\psi(x)=x^2$ chooses the first label when the empirical second
moment is less than $3/2$ and the second when it is greater.
At equality either choice minimizes the distance, and a fixed
tie rule is measurable.
\end{example}

\subsection{Exercises}
\label{subsec:practice-5-2}

\begin{practiceexercise}[A finite tie]
\label{pex:ml-bernoulli-finite-tie}
For two iid Bernoulli draws with $p\in\{1/4,3/4\}$, give the
maximizing set at each count. Choose a tie rule and state
the resulting estimator as a function of the count.
\end{practiceexercise}

\begin{practiceexercise}[An impossible record]
\label{pex:ml-impossible-record}
On $\{0,1,2\}$ take laws $(1/4,3/4,0)$ and $(3/4,1/4,0)$.
Using counting measure, find the maximizing set at each
observation. What choices can an estimator make at outcome two?
\end{practiceexercise}

\begin{practiceexercise}[A countable parameter set]
\label{pex:ml-countable-first-maximizer}
Let $\Theta$ be countable with all subsets measurable. Suppose
a maximizer exists on a measurable common full-probability set.
Show that choosing the first maximizer in an enumeration there,
and a fixed label elsewhere, gives a measurable estimator.
\end{practiceexercise}

\begin{practiceexercise}[Closed Bernoulli constraint]
\label{pex:ml-bernoulli-clipped-constraint}
For $n$ iid Bernoulli observations, maximize the likelihood over
$[a,b]$ with $0<a<b<1$. Include the two constant records.
\end{practiceexercise}

\begin{practiceexercise}[Gaussian location on a half-line]
\label{pex:ml-normal-half-line}
For iid $N(\theta,1)$ observations with $\theta\geq0$, find
the unique maximum likelihood estimator using the residual
sum of squares identity.
\end{practiceexercise}

\begin{practiceexercise}[Weighted Gaussian location]
\label{pex:ml-weighted-gaussian-location}
Let the independent $X_i$ have laws $N(\theta,v_i)$ with known
$v_i>0$ and $\theta\in\R$. Complete the square in the log
likelihood to find its unique maximizer.
\end{practiceexercise}

\begin{practiceexercise}[A Poisson rate with an open boundary]
\label{pex:ml-poisson-zero-count}
For $n$ iid Poisson draws with rate $\lambda>0$, find the
maximizing set when the total is positive and when it is zero.
Compute the probability of the second case.
\end{practiceexercise}

\begin{practiceexercise}[An exponential rate]
\label{pex:ml-exponential-rate}
For $n$ iid rate-$\lambda$ exponential draws, with $\lambda>0$,
find the maximum likelihood estimator on records with positive
sum and show that these records have probability one.
\end{practiceexercise}

\begin{practiceexercise}[Uniform support with a restricted endpoint]
\label{pex:ml-uniform-restricted-endpoint}
In the uniform endpoint model use closed-support densities and
restrict $\theta$ to $[1,2]$. Find the maximizer for nonnegative
records with maximum at most two.
\end{practiceexercise}

\begin{practiceexercise}[The normal variance boundary]
\label{pex:ml-normal-variance-floor}
In Example~\ref{ex:ml-normal-location-variance}, restrict the
variance to $v\geq v_0$ for known $v_0>0$. Find both maximizing
coordinates, including on constant records.
\end{practiceexercise}

\begin{practiceexercise}[Likelihood after taking the count]
\label{pex:ml-count-ratio-reference}
For $n$ iid Bernoulli trials with $0<p<1$, use the law at
$p=1/2$ as the reference. Compute the full and count likelihood
ratios relative to this law and its induced law, respectively.
\end{practiceexercise}

\begin{practiceexercise}[A noninjective parameter map]
\label{pex:ml-bernoulli-squared-parameter}
Parametrize the Bernoulli probability as $p=\eta^2$ for
$\eta\in[-1,1]$. Given $k$ successes in $n$ iid trials, find
all maximizing labels. Include $k=0$ and $k=n$.
\end{practiceexercise}

\begin{practiceexercise}[A density modification at a moving point]
\label{pex:ml-pointwise-density-modification}
For the model in Example~\ref{ex:factorization-parameterwise-null-sets},
compare the maximizing sets for $f_\theta(x)=1$ and
$f_\theta(x)=1+\mathbf 1_{\{x=\theta\}}$ at each $x\in[0,1]$.
\end{practiceexercise}

\begin{practiceexercise}[Matching the mean of an exponential law]
\label{pex:mom-exponential-rate}
With $\psi(x)=x$ in the exponential rate model, solve the
moment equation when the sample mean is positive. Compare the
result with the maximum likelihood estimator.
\end{practiceexercise}

\begin{practiceexercise}[Two normal moments]
\label{pex:mom-normal-two-moments}
For $N(u,v)$ with $u\in\R$ and $v>0$, match the first and second
raw moments. Determine when a solution exists and compare it
with the likelihood fit.
\end{practiceexercise}

\begin{practiceexercise}[Projection in moment space]
\label{pex:mom-simplex-line-projection}
Let a three-outcome law have masses $(p,p,1-2p)$ with
$0\leq p\leq1/2$. Match the first two empirical category
frequencies by Euclidean minimum distance. Give the fitted $p$.
\end{practiceexercise}

\begin{practiceexercise}[A noninjective moment map]
\label{pex:mom-even-moment-normal-sign}
For $N(\theta,1)$ with $\theta\in\R$, use $\psi(x)=x^2$.
Find all solutions to the moment equation as a function of
the empirical second moment.
\end{practiceexercise}

\begin{practiceexercise}[A moment image that is not closed]
\label{pex:mom-open-bernoulli-distance}
For the Bernoulli model with $p\in(0,1)$, consider the absolute
distance between $p$ and the empirical proportion. Find the
infimum and determine when it is attained.
\end{practiceexercise}

\begin{exercise}[A rectangular support with two unknown endpoints]
\label{exer:ml-uniform-two-endpoints}
For $n\geq2$ iid uniform draws from $[a,b]$ with $a<b$, use the
closed-support density. Determine the maximizing pair on records
with distinct minimum and maximum, and analyze the likelihood
when they coincide. Show that the maximizing pair gives an
estimator on a common full-probability set.
\end{exercise}

\begin{exercise}[Projection in probability space and likelihood]
\label{exer:ml-three-outcomes-two-criteria}
On three outcomes take masses $(p,2p,1-3p)$ for $0\leq p\leq1/3$.
For an iid record, find the maximum likelihood estimator
and the Euclidean minimum-distance fit to the full vector
of empirical frequencies. Give a record for which they differ.
\end{exercise}

\begin{exercise}[Measurable approximate maximizers]
\label{exer:ml-measurable-approximate-maximizers}
Let $\Theta$ be a nonempty subset of $\R^d$ with its Borel sets.
Suppose $\theta\mapsto L_x(\theta)$ is continuous and has finite
supremum for every $x$. For each $\varepsilon>0$, construct a
measurable choice with likelihood strictly greater than the
supremum minus $\varepsilon$. Use a countable dense sequence
in $\Theta$. Give an example in which the supremum is unattained.
\end{exercise}

\begin{exercise}[Sufficient reduction and exact likelihood values]
\label{exer:ml-poisson-exposure-reduction}
For the Poisson exposure model of
Example~\ref{ex:factorization-poisson-exposures}, use the law at
a fixed $\lambda_0>0$ as reference. Show directly that the
likelihood ratio agrees with that of the total count using
its induced reference law. Derive the maximizing sets when
the rate ranges over $(0,\infty)$ and when it ranges over
$[\lambda_0,2\lambda_0]$.
\end{exercise}

\begin{exercise}[Unbiasedness and risk of the normal likelihood fit]
\label{exer:ml-gaussian-variance-risk}
For $n\geq2$ iid $N(u,v)$ observations, compute the bias and
squared risk of the maximum likelihood variance estimator.
Use the chi-square moments in
Exercise~\ref{exer:mse-normal-variance-denominator}. Compare
with the unbiased estimator, identifying when their squared
risks are equal.
\end{exercise}

\begin{exercise}[Matching two moments of a uniform law]
\label{exer:mom-one-moment-discrepancy}
For the uniform model $\operatorname{Unif}[a,b]$ with $a<b$,
match the first two raw moments to the empirical moments.
Determine when the fitted interval exists and whether it
must contain every observation. Give a record for which
matching both moments yields an interval that excludes
an observation.
\end{exercise}



\clearpage
\section{Rao--Blackwellization and complete sufficient statistics}
\label{sec:rao-blackwell-lehmann-scheffe}

Fix an experiment $\mathcal E=(\mathsf X,\mathcal A,
(P_\theta)_{\theta\in\Theta})$ and a real target $g$.

\subsection{Conditioning an estimator}
\label{subsec:rao-blackwellization}

\begin{definition}[Rao--Blackwellization]
\label{def:rao-blackwellization}
\indexnotation{\delta_T}{Rao--Blackwellized estimator based on the statistic \(T\).}
Let $\delta$ be integrable under each model law and let $T$ be
sufficient. A \emph{Rao--Blackwellization} of $\delta$ given $T$ is
an estimator $\delta_T=h\circ T$, with $h$ real and measurable,
such that $\delta_T$ is a version of
$\mathbb E_\theta[\delta\mid\sigma(T)]$ for every $\theta$.
\end{definition}

\begin{theorem}[Rao--Blackwell theorem]
\label{thm:rao-blackwell-squared-risk}
Let $T$ be sufficient and $\delta$ integrable at every label.
A Rao--Blackwellization $\delta_T$ exists and satisfies
$\mathbb E_\theta\delta_T=\mathbb E_\theta\delta$ for all $\theta$.
At every label where $\delta$ is square integrable, $\delta_T$ is
also square integrable, and
\[
\operatorname{MSE}_\theta(\delta)
=\operatorname{MSE}_\theta(\delta_T)
 +\mathbb E_\theta[(\delta-\delta_T)^2].
\]
Equality of the two risks at such a label holds exactly when
$\delta=\delta_T$ almost surely under that label.
\end{theorem}
\begin{proof}
For a nonnegative finite simple function of the observation,
a linear combination of the common conditional probabilities in
Definition~\ref{def:sufficient-statistic} gives a common conditional
expectation as a function of $T$. For a nonnegative measurable
function $f$, take finite simple functions $f_j\uparrow f$ and choose
such versions $h_j(T)$ for $f_j$. At each label, these versions
increase almost surely, since conditional expectation preserves
order and there are only countably many comparisons. The measurable
function $\liminf_j h_j$ therefore gives a common version for $f$
by monotone convergence on each event in $\sigma(T)$.

Apply this construction to the positive and negative parts of $\delta$.
Both resulting functions of $T$ are finite almost surely under each
label by integrability. Subtract them where both are finite and take
zero elsewhere. This gives $\delta_T$, and taking expectations
gives the mean identity.

At a label with finite second moment, conditional Jensen gives
$\mathbb E_\theta[\delta_T^2]\leq\mathbb E_\theta[\delta^2]$.
In the expansion
\[
\delta-g(\theta)=(\delta-\delta_T)+(\delta_T-g(\theta)),
\]
the two terms have zero expected product, since
$\mathbb E_\theta[\delta-\delta_T\mid\sigma(T)]=0$ and the second
term is $\sigma(T)$-measurable. This conditional-expectation
identity for a square-integrable multiplier follows from the bounded
case by truncation and Cauchy--Schwarz. Taking expected squares
gives the risk identity and its equality criterion.
\end{proof}

\begin{example}[Conditioning the first Bernoulli observation]
\label{ex:rao-blackwell-first-bernoulli}
For $n$ iid Bernoulli draws let $T=\sum_i X_i$. Given count $k$ of
positive mass, the records are uniform over the count fiber by
Example~\ref{ex:sufficiency-bernoulli-count}. The sum of the $n$
conditional coordinate means is $k$, and symmetry makes them
equal. Thus the Rao--Blackwellization of $X_1$ is $T/n$.
For estimating $p$, the squared risk decreases from $p(1-p)$
to $p(1-p)/n$. The decrease is strict when $n>1$ and $0<p<1$.
\end{example}

\begin{example}[Conditioning an endpoint estimate]
\label{ex:rao-blackwell-uniform-two-draws}
For two iid uniform draws on $[0,\theta]$, the estimator $2X_1$
is unbiased for $\theta$. Given $T=\max(X_1,X_2)=t>0$, the kernel
in Example~\ref{ex:sufficiency-uniform-maximum-pair} gives $X_1=t$
half the time and draws it uniformly on $[0,t]$ the other half.
Its conditional mean is $3t/4$, so the conditioned estimator is
$3T/2$. The original risk is $\theta^2/3$ and the conditioned
risk is $\theta^2/8$ by Example~\ref{ex:mse-uniform-endpoint-comparison}.
The exact decrease is $5\theta^2/24$.
\end{example}

\begin{example}[Conditioning on one draw of a pair]
\label{ex:rao-blackwell-insufficient-bernoulli}
Take two iid Bernoulli draws with $0<p<1$. The conditional
expectation of $X_2$ given $X_1$ is the constant $p$ by independence.
Both values of $X_1$ have positive probability at every label,
so no single function of $X_1$ represents these expectations
throughout the model. This is the insufficient reduction in
Example~\ref{ex:sufficiency-first-bernoulli-coordinate}.
\end{example}

\subsection{Optimality among unbiased estimators}
\label{subsec:lehmann-scheffe}

\begin{definition}[Uniform minimum-variance unbiased estimator]
\label{def:umvu-estimator}
\indexnotation{\mathrm{UMVU}}{Uniformly minimum-variance unbiased.}
An unbiased estimator $\delta^\star$ of $g$, square integrable
at every label, is \emph{uniformly minimum-variance unbiased}
if $\operatorname{Var}_\theta(\delta^\star)\leq
\operatorname{Var}_\theta(\delta)$ for every label and every
unbiased estimator $\delta$ that is square integrable at every label.
\end{definition}

\begin{theorem}[Lehmann--Scheff\'e theorem]
\label{thm:lehmann-scheffe}
Suppose $T$ is complete and sufficient and that $g$ admits an
unbiased estimator $\delta$ square integrable at every label.
Its Rao--Blackwellization $\delta_T$ is UMVU. Every UMVU estimator
of $g$ agrees with $\delta_T$ almost surely under each $P_\theta$.
\end{theorem}
\begin{proof}
For any unbiased square-integrable competitor $U$, the Rao--Blackwell
theorem gives an unbiased square-integrable estimator $U_T$.
The two functions of $T$, $U_T$ and $\delta_T$, have the same
expectation at every label. By completeness and
Proposition~\ref{prop:complete-unbiased-uniqueness}, they agree
almost surely at each label. The risk identity and unbiasedness give
\[
\operatorname{Var}_\theta(U)
=\operatorname{Var}_\theta(\delta_T)
 +\mathbb E_\theta[(U-\delta_T)^2].
\]
This proves optimality. If $U$ is also UMVU, the variances are equal
at every label and the expected squared difference is zero.
\end{proof}

\begin{proposition}[Covariance criterion for UMVU]
\label{prop:umvu-covariance-criterion}
An unbiased estimator $\delta$ square integrable at every label
is UMVU exactly when $\operatorname{Cov}_\theta(\delta,Z)=0$
at every label for every estimator $Z$ square integrable and
of mean zero at all labels.
\end{proposition}
\begin{proof}
If $\delta$ is UMVU, then $\delta+tZ$ is an unbiased competitor for
each $t\in\R$. At each label the quadratic
\[
\operatorname{Var}_\theta(\delta+tZ)
=\operatorname{Var}_\theta(\delta)
 +2t\operatorname{Cov}_\theta(\delta,Z)+t^2\operatorname{Var}_\theta(Z)
\]
has a minimum at zero, so its linear coefficient vanishes.
Conversely, for any unbiased square-integrable competitor $U$, put
$Z=U-\delta$. If every such covariance vanishes, the same formula
at $t=1$ gives
$\operatorname{Var}_\theta(U)=\operatorname{Var}_\theta(\delta)
+\operatorname{Var}_\theta(Z)$ at each label.
\end{proof}

\begin{example}[Estimating a squared Bernoulli probability]
\label{ex:umvu-bernoulli-squared}
For $n\geq2$ iid Bernoulli draws, $X_1X_2$ is unbiased for $p^2$.
Given $T=k$, all ordered pairs of distinct coordinates have the
same chance of both being one. There are $k(k-1)$ such pairs
of ones out of $n(n-1)$ pairs, so the conditioned estimate is
\[
\delta_T=\frac{T(T-1)}{n(n-1)}.
\]
The count is complete and sufficient by
Example~\ref{ex:complete-bernoulli-count}. This is therefore the
unique UMVU estimator up to model null sets.
\end{example}

\begin{example}[An unbiased endpoint estimate]
\label{ex:umvu-uniform-endpoint}
For $n$ iid draws from the uniform endpoint model,
$U=(n+1)T/n$ is unbiased and square integrable by
Example~\ref{ex:mse-uniform-endpoint-comparison}, with
$T$ the maximum. The maximum is complete and sufficient by
Example~\ref{ex:complete-uniform-maximum}, so $U$ is UMVU with
risk $\theta^2/(n(n+2))$. The biased estimate $(n+2)T/(n+1)$
has squared risk $\theta^2/(n+1)^2$ by the same moment formulas.
Since $(n+1)^2=n(n+2)+1$, that risk is strictly smaller
for every $\theta>0$.
\end{example}

\begin{example}[An incomplete model with no UMVU estimator for a target]
\label{ex:umvu-may-not-exist}
Use the two laws $P=(1/2,1/3,1/6)$ and $Q=(1/6,1/3,1/2)$ from
Example~\ref{ex:minimal-sufficient-incomplete-three}, and take target
values $g(P)=1/2$ and $g(Q)=1/6$. The indicator of the first outcome
is unbiased. The vectors of masses of $P,Q$ are linearly independent,
so all unbiased estimators have the form $(1,0,0)+a(1,-2,1)$ for
$a\in\R$. Put $Z=(1,-2,1)$. Its mean is zero and its variance
is two under both laws. The covariance of an unbiased estimator
with $Z$ is $1/2+2a$ under $P$ and $1/6+2a$ under $Q$.
These cannot vanish at the same $a$, so the covariance criterion
rules out a UMVU estimator for this target.
\end{example}

\subsection{Exercises}
\label{subsec:exercises-5-3}

\begin{exercise}[Preserving bias]
\label{exer:rao-blackwell-preserves-bias}
Show that conditioning on a sufficient statistic preserves bias at
each label, including when the original estimator is biased.
\end{exercise}

\begin{exercise}[A constant estimator]
\label{exer:rao-blackwell-constant}
Find the Rao--Blackwellization of a constant estimator given an
arbitrary sufficient statistic. What is the squared-risk decrease?
\end{exercise}

\begin{exercise}[An estimator already based on the statistic]
\label{exer:rao-blackwell-fixed-function}
Suppose $\delta=h(T)$ is integrable at every label. Show that it is
its own Rao--Blackwellization.
\end{exercise}

\begin{exercise}[Uniqueness of the common version]
\label{exer:rao-blackwell-versions}
Show that two Rao--Blackwellizations of the same estimator given
the same sufficient statistic agree off a common model null set.
\end{exercise}

\begin{exercise}[A constant sufficient statistic]
\label{exer:rao-blackwell-one-law}
Suppose all labels have the same observation law. Condition an
integrable estimator on a constant statistic. If its variance is
finite, give the squared-risk decrease.
\end{exercise}

\begin{exercise}[The first two Bernoulli draws]
\label{exer:rao-blackwell-two-bit-sum}
For $n\geq3$ iid Bernoulli draws, condition $(X_1+X_2)/2$ on
the total count and compute the squared-risk decrease for estimating $p$.
\end{exercise}

\begin{exercise}[A particular ordered pair of bits]
\label{exer:rao-blackwell-ordered-pair}
For $n\geq2$ iid Bernoulli draws, condition $X_1(1-X_2)$ on
the total count. Identify its target and explain why the result is UMVU.
\end{exercise}

\begin{exercise}[The all-zero event]
\label{exer:umvu-bernoulli-no-success}
For $1\leq m\leq n$, use the indicator that the first $m$ of $n$
Bernoulli draws are all zero. Condition it on the count to estimate
$(1-p)^m$.
\end{exercise}

\begin{exercise}[A count with zero probability]
\label{exer:rao-blackwell-null-count}
For Bernoulli parameters $p\in\{0,1\}$ and $n=2$, condition
$X_1$ on the count. Explain the freedom to specify the estimate
at count one.
\end{exercise}

\begin{exercise}[An infinite-variance estimator]
\label{exer:rao-blackwell-integrable-not-square}
In a one-law experiment, let $X$ be uniform on $(0,1)$ and
$\delta=X^{-1/2}$. Condition on a constant sufficient statistic.
Compare the two squared risks for the target $\mathbb E\delta$.
\end{exercise}

\begin{exercise}[A mean under a normal model]
\label{exer:rao-blackwell-normal-first}
For iid $N(\theta,1)$ observations, use the residual kernel from
Example~\ref{ex:sufficiency-normal-mean} to condition $X_1$ on the
sample mean. Compute the squared-risk decrease.
\end{exercise}

\begin{exercise}[A uniform endpoint with two observations]
\label{exer:rao-blackwell-uniform-conditional-variance}
In Example~\ref{ex:rao-blackwell-uniform-two-draws}, compute the
conditional variance of $2X_1$ given the maximum. Integrate it
to recover the risk decrease.
\end{exercise}

\begin{exercise}[Adding an ancillary term]
\label{exer:rao-blackwell-ancillary-noise}
Let $T$ be complete and sufficient, and $V$ ancillary. Suppose
$h(T)$ is square integrable and $a(V)$ is bounded with mean zero.
Condition $h(T)+a(V)$ on $T$ and compute the squared-risk decrease.
\end{exercise}

\begin{exercise}[The UMVU estimator of a constant]
\label{exer:umvu-constant}
For $g(\theta)=c$ at every label, give a UMVU estimator and prove
its uniqueness up to common model null sets.
\end{exercise}

\begin{exercise}[Linear combinations of optimal estimators]
\label{exer:umvu-linear-combination}
If $U$ and $V$ are UMVU for two targets, show that $aU+bV$
is UMVU for the corresponding linear combination of targets.
Use the covariance criterion.
\end{exercise}

\begin{exercise}[Optimality under a smaller model]
\label{exer:umvu-model-restriction}
Start with an unbiased estimator of $p$ based on one Bernoulli draw
for $p\in[0,1]$. Restrict the model to $\{1/2\}$. Determine whether
the original UMVU estimator remains UMVU for the restricted model.
\end{exercise}

\begin{exercise}[Label-specific variance minimizers]
\label{exer:umvu-no-common-minimizer}
In Example~\ref{ex:umvu-may-not-exist}, find the value of $a$
that minimizes the variance at each of the two labels separately.
\end{exercise}

\begin{exercise}[Uniqueness without sufficiency assumptions]
\label{exer:umvu-uniqueness-midpoint}
Suppose $U,V$ are both UMVU for the same target. Apply the
covariance criterion with $Z=U-V$ and deduce that the
estimators agree almost surely at every label.
\end{exercise}

\begin{exercise}[Polynomial targets of a Bernoulli model]
\label{exer:umvu-bernoulli-polynomial}
For a polynomial $g(p)=\sum_{j=0}^n a_j p^j$, derive its
UMVU estimator from the total count. Condition the products
of the first $j$ draws on the count, taking the empty product
to be one.
\end{exercise}

\begin{exercise}[Poisson estimation from the total]
\label{exer:rao-blackwell-poisson-zero}
For $n\geq2$ iid Poisson draws of mean $\lambda>0$, start with
$\mathbf 1_{\{X_1=0\}}$ to estimate $e^{-\lambda}$. Condition
on the total using the conditional distribution in
Exercise~\ref{exer:sufficiency-poisson-total}. Prove that the
total is complete by uniqueness of the coefficients of a power series.
Use integrability for all $\lambda>0$ to justify absolute convergence.
\end{exercise}

\begin{exercise}[The mean of a uniform sample given its maximum]
\label{exer:rao-blackwell-uniform-full-mean}
For $n\geq2$ iid uniform draws on $[0,\theta]$, condition
$2\overline X$ on the maximum using
Exercise~\ref{exer:sufficiency-uniform-maximum-general}.
Compute the squared-risk decrease and characterize equality.
\end{exercise}

\begin{exercise}[Conditioning under convex loss]
\label{exer:rao-blackwell-convex-loss}
Let $\delta$ be bounded and $T$ sufficient. For every $\theta$,
let $L(\theta,\cdot)$ be finite and convex on $\R$. Prove that
$\delta_T$ has no larger risk at any label by conditional Jensen.
Show that if $L(\theta,\cdot)$ is strictly convex, equality at that
label requires $\delta=\delta_T$ almost surely. One may use the
equality condition in Jensen's inequality.
\end{exercise}

\begin{exercise}[Bounded completeness and bounded estimators]
\label{exer:lehmann-scheffe-bounded-class}
Assume $T$ is boundedly complete and sufficient and that a
bounded unbiased estimator exists. Prove the Lehmann--Scheff\'e
conclusion within the class of bounded unbiased estimators.
\end{exercise}

\begin{exercise}[Incomplete sufficiency with a UMVU estimator]
\label{exer:umvu-incomplete-sufficient-example}
Observe $X=(Y,V)$ with $Y$ Bernoulli of parameter $p\in(0,1)$
and $V$ an independent fair bit. Verify that $X$ is sufficient
but incomplete. Use the covariance criterion to prove that
$Y$ is UMVU for $p$ in this experiment.
\end{exercise}



\clearpage
\section{Fisher information and information lower bounds}
\label{sec:fisher-information-cramer-rao}

\subsection{The score and Fisher information}
\label{subsec:fisher-score-information}

\begin{definition}[Score and Fisher information]
\label{def:fisher-score-information}
\indexnotation{s_\theta,\ I(\theta)}{Likelihood score and scalar Fisher information at parameter \(\theta\).}
Let $\Theta$ be an open interval in $\R$. Fix density versions
$f_\theta$ relative to a common $\sigma$-finite measure $\mu$.
Suppose that for $P_\theta$-almost every $x$ the density is positive
at $\theta$ and differentiable in the parameter there. The \emph{score} at $\theta$ and the \emph{Fisher
information} in the observation are
\[
s_\theta(x)=\partial_\theta\log f_\theta(x),
\qquad
I(\theta)=\mathbb E_\theta[s_\theta^2]\in[0,\infty].
\]
Use any measurable version of the score on the exceptional set.
On the differentiability set it is measurable as a limit of
difference quotients along a fixed sequence of steps tending to zero.
\end{definition}

\begin{proposition}[Differentiating expectations at a label]
\label{prop:score-differentiate-expectation}
Fix $\theta\in\Theta$. Suppose there is a measurable set $S$
on which, for all $t$ in a neighborhood of $\theta$, the densities
are positive and outside which they vanish, off a single
$\mu$-null set. Assume $f_t(x)$ is differentiable
at $t=\theta$ for $\mu$-almost every $x\in S$. Suppose there is a
nonnegative $\mu$-integrable $H$ satisfying
\begin{equation}
\left|\frac{f_{\theta+h}(x)-f_\theta(x)}{h}\right|
\leq H(x)
\label{eq:score-difference-domination}
\end{equation}
$\mu$-almost everywhere for each sufficiently small $h\ne0$.
Then $s_\theta$ is integrable and $\mathbb E_\theta s_\theta=0$.
If $\delta$ is integrable under $P_\theta$ and
$\int|\delta|H\,d\mu<\infty$, then $m(t)=\mathbb E_t\delta$
is finite near $\theta$ and
\[
m'(\theta)=\mathbb E_\theta[\delta s_\theta].
\]
\end{proposition}
\begin{proof}
Let $\dot f_\theta$ be the density derivative on $S$ and zero
elsewhere. Then $\dot f_\theta=s_\theta f_\theta$ almost
everywhere. Along any sequence of nonzero steps tending to zero,
the difference quotients converge almost everywhere to
$\dot f_\theta$. The integrable bound in
\eqref{eq:score-difference-domination} holds simultaneously
along the sequence, so $\lvert\dot f_\theta\rvert\leq H$
almost everywhere. In particular, the score is integrable.
Dominated convergence applied to the density integrals gives
$\int\dot f_\theta\,d\mu=0$, since every density integrates to one.

For small $h$, $f_{\theta+h}\leq f_\theta+|h|H$ almost everywhere,
which gives integrability of $\delta$ under $P_{\theta+h}$.
Apply dominated convergence to the difference quotients weighted
by $\delta$, using $|\delta|H$. Since this applies to every
sequence of steps tending to zero, the derivative exists and
equals $\int\delta\dot f_\theta\,d\mu$.
\end{proof}

\begin{proposition}[Additivity and change of parameter]
\label{prop:fisher-product-reparametrization}
For independent observations at a common scalar label $\theta$,
suppose their scores $s_{\theta,i}$ are centered and have finite
second moments. The product likelihood has score
$\sum_i s_{\theta,i}(X_i)$ and information $\sum_i I_i(\theta)$.
For iid observations the information is $nI(\theta)$.

If $\theta=\tau(\eta)$ for a differentiable map $\tau$ into
$\Theta$, the score of the reparametrized family is
$\tau'(\eta)s_{\tau(\eta)}$ almost surely at $\eta$.
At finite information this gives
$I_\eta(\eta)=\tau'(\eta)^2 I(\tau(\eta))$.
\end{proposition}
\begin{proof}
The product log likelihood is a sum by
Proposition~\ref{prop:product-densities}. Its derivative is the
sum of the scores almost surely under the product law. The
cross terms in the expectation of the squared sum vanish by
independence and centering. For the change of parameter apply
the chain rule and take the expected square.
\end{proof}

\begin{proposition}[Information from a second derivative]
\label{prop:fisher-expected-curvature}
At a fixed $\theta$, suppose the support is common and the densities
are positive and twice differentiable in the parameter on that
support, almost everywhere on a neighborhood of $\theta$.
Suppose $\lvert\partial_t f_t(x)\rvert$ and
$\lvert\partial_t^2 f_t(x)\rvert$ on that neighborhood are bounded
by $\mu$-integrable functions and $I(\theta)<\infty$. Then
\[
I(\theta)=-\mathbb E_\theta[\partial_\theta^2\log f_\theta(X)] .
\]
\end{proposition}
\begin{proof}
The mean-value theorem and the derivative bounds allow two
differentiations of $\int f_t\,d\mu=1$ by dominated convergence.
In particular $\int\partial_\theta^2 f_\theta\,d\mu=0$.
On the common support,
\[
\partial_\theta^2\log f_\theta
=\frac{\partial_\theta^2 f_\theta}{f_\theta}-s_\theta^2 .
\]
The first term is integrable under $P_\theta$ by the bound on
the second derivative, and the second by finite information.
Taking expectations gives the result.
\end{proof}

\subsection{The scalar Cram\'er--Rao bound}
\label{subsec:scalar-cramer-rao}

\begin{theorem}[Cram\'er--Rao bound]
\label{thm:scalar-cramer-rao}
Let $\delta$ have finite second moment at $\theta$ and suppose the
hypotheses of Proposition~\ref{prop:score-differentiate-expectation}
hold for $\delta$ there. If $0<I(\theta)<\infty$, then
\[
\operatorname{Var}_\theta(\delta)\geq\frac{m'(\theta)^2}{I(\theta)},
\qquad m(t)=\mathbb E_t\delta .
\]
Equality holds exactly when
$\delta-m(\theta)=m'(\theta)s_\theta/I(\theta)$ $P_\theta$-almost surely.
For a real target $g$, the corresponding bound on squared risk is
\[
\operatorname{MSE}_\theta(\delta)
\geq b_\theta(\delta)^2+\frac{m'(\theta)^2}{I(\theta)} .
\]
If $\delta$ is unbiased for $g$ near $\theta$, then $m'(\theta)=g'(\theta)$.
\end{theorem}
\begin{proof}
By Proposition~\ref{prop:score-differentiate-expectation},
the score is centered and
\[
m'(\theta)=\mathbb E_\theta[(\delta-m(\theta))s_\theta].
\]
Apply Cauchy--Schwarz to the centered estimator and score.
Because the score has positive squared norm, equality holds
exactly when the centered estimator is its scalar multiple.
Taking the inner product with the score gives the stated
coefficient. Add squared bias using
Proposition~\ref{prop:estimator-mse-decomposition}.
\end{proof}

\subsection{Examples and failure with moving support}
\label{subsec:fisher-examples}

\begin{example}[Bernoulli information]
\label{ex:fisher-bernoulli}
For one Bernoulli observation with $0<p<1$,
\[
s_p(X)=\frac{X-p}{p(1-p)},\qquad I(p)=\frac1{p(1-p)} .
\]
The first formula follows by differentiating the two log masses,
and the second uses $\operatorname{Var}_p(X)=p(1-p)$.
For $n$ iid draws, the sample mean is unbiased and its variance
$p(1-p)/n$ attains the bound. All density derivatives and
weighted bounds needed for the theorem exist locally because
the sample space is finite and the masses are polynomials.
\end{example}

\begin{example}[Normal location with known variance]
\label{ex:fisher-normal-location}
For $X_i$ iid $N(\theta,v)$ with known $v>0$, the joint score is
$n(\overline X-\theta)/v$ and the information is $n/v$.
The sample mean is unbiased and attains the bound $v/n$.
To justify differentiation, on a compact neighborhood of a
fixed label the first two density derivatives are bounded in
absolute value by a constant times
\[
(1+\textstyle\sum_i x_i^2)\exp(-\textstyle\sum_i x_i^2/(4v)).
\]
Indeed the derivatives multiply the density by polynomials of
degree at most two in the records, and
$(x_i-t)^2\geq x_i^2/2-t^2$. This bound stays integrable after
multiplication by $|\overline x|$. The second log likelihood
derivative is $-n/v$, in agreement with the curvature identity.
\end{example}

\begin{example}[A biased normal location estimator]
\label{ex:fisher-normal-affine-bias}
In the same normal model, let $\delta=a\overline X+c$ for fixed
$a,c\in\R$. Then $m'(\theta)=a$ and the bias for estimating
$\theta$ is $(a-1)\theta+c$. The squared risk is
\[
\frac{a^2v}{n}+\big((a-1)\theta+c\big)^2,
\]
which attains the biased bound for every label. The centered
estimator is $a$ times the centered sample mean, hence a
multiple of the score.
\end{example}

\begin{example}[Poisson observations with unequal exposure]
\label{ex:fisher-poisson-exposure}
In the exposure model of Example~\ref{ex:factorization-poisson-exposures},
put $A=\sum_i a_i$ and $T=\sum_i X_i$. The joint score is $T/\lambda-A$.
Since the independent counts have means and variances $a_i\lambda$,
the score is centered and the information is $A/\lambda$.
The estimator $T/A$ is unbiased for $\lambda$ with variance
$\lambda/A$, attaining the bound. On any interval $[l,u]$
with $0<l<\lambda<u$, the joint mass derivative is bounded by
$(T/l+A)e^{Au}$ times the joint mass at $u$. The bound is
summable also after multiplying by $T/A$, since $T$ has a
finite second moment at $u$.
\end{example}

\begin{example}[Exponential rate and mean]
\label{ex:fisher-exponential-rate-mean}
For one exponential observation of rate $\lambda>0$, the score
is $1/\lambda-X$ and the information is $1/\lambda^2$. In the
mean parameter $\mu_0=1/\lambda$, the information per observation
is $1/\mu_0^2$ by Proposition~\ref{prop:fisher-product-reparametrization}.
For $n$ iid draws, $\overline X$ estimates $\mu_0$ without bias
and attains the bound $\mu_0^2/n$. On any compact rate interval
$[l,u]\subset(0,\infty)$, the product density and its rate
derivative are bounded by a constant times
$(1+\sum_i x_i)\exp(-l\sum_i x_i)$ on the positive orthant.
Multiplying by the sample mean still gives an integrable bound.
Under the mean parametrization, the rate derivative is multiplied
by $-1/\mu_0^2$, which is bounded on a compact neighborhood of
the fixed mean.
\end{example}

\begin{example}[A moving uniform endpoint]
\label{ex:fisher-uniform-moving-support}
For $n$ iid observations uniform on $[0,\theta]$, $\theta>0$,
the pointwise score at $\theta$ is $-n/\theta$ on a set of full
$P_\theta$-probability. Its second moment is $n^2/\theta^2$,
and its mean is $-n/\theta$. The unbiased rescaled maximum in
Example~\ref{ex:mse-uniform-endpoint-comparison} has variance
\[
\frac{\theta^2}{n(n+2)}<\frac{\theta^2}{n^2}.
\]
The common-support hypothesis fails. Differentiating the
normalization $\int_0^\theta\theta^{-1}\,dx=1$ for one
observation includes a boundary contribution $1/\theta$,
which cancels the integral of the interior derivative.
\end{example}

\subsection{Exercises}
\label{subsec:exercises-5-4}

\begin{exercise}[Centering on a finite space]
\label{exer:fisher-finite-centering}
For positive masses $p_\theta(x)$ on a finite space, differentiate
their sum to verify $\mathbb E_\theta s_\theta=0$.
\end{exercise}

\begin{exercise}[Two-point information]
\label{exer:fisher-binary-linear-probability}
For $P_\theta(X=1)=(1+\theta)/2$ with $-1<\theta<1$,
compute the score and information. Find an unbiased estimator
of $\theta$ that attains the bound for one draw.
\end{exercise}

\begin{exercise}[The score of a constant model]
\label{exer:fisher-constant-model}
If every density version is the same function, compute the score
and information. Which targets can admit an unbiased estimator?
\end{exercise}

\begin{exercise}[A sample score at its maximizer]
\label{exer:fisher-score-bernoulli-mle}
For $n$ Bernoulli draws with $0<k<n$ successes, compute the
sample score and verify that it vanishes at the likelihood
maximizer. Compare with the all-zero record.
\end{exercise}

\begin{exercise}[Log odds]
\label{exer:fisher-bernoulli-log-odds}
Reparametrize Bernoulli by $\eta=\log(p/(1-p))$. Compute the
score and information in the new parameter.
\end{exercise}

\begin{exercise}[Normal location in different units]
\label{exer:fisher-normal-location-units}
Take $X\sim N(a\eta,v)$, with known $a\ne0$ and $v>0$.
Compute information for $\eta$ and check the bound for $X/a$.
\end{exercise}

\begin{exercise}[Information and independent repetition]
\label{exer:fisher-independent-repetition}
For iid draws with one-observation information $I(\theta)$,
compare the variance bounds for an unbiased estimate of a
differentiable target based on $n$ and on $2n$ draws.
\end{exercise}

\begin{exercise}[Gaussian measurements of unequal precision]
\label{exer:fisher-gaussian-unequal-precision}
Let independent $X_i$ have laws $N(\theta,v_i)$ with known
$v_i>0$. Compute the joint score and information.
Check equality for the inverse-variance weighted average.
\end{exercise}

\begin{exercise}[Second derivative for Bernoulli]
\label{exer:fisher-bernoulli-expected-curvature}
Compute the negative second derivative of the Bernoulli log
likelihood at both outcomes and average it. Compare with the
information in Example~\ref{ex:fisher-bernoulli}.
\end{exercise}

\begin{exercise}[Observed curvature for Poisson]
\label{exer:fisher-poisson-observed-curvature}
For one Poisson draw, compute the negative second log
likelihood derivative at a realized count and its expectation.
For which realized count is that derivative zero?
\end{exercise}

\begin{exercise}[Exponential log rate]
\label{exer:fisher-exponential-log-rate}
Set $\eta=\log\lambda$ in the exponential rate model.
Compute the score and information for one draw.
\end{exercise}

\begin{exercise}[Normal variance with known mean]
\label{exer:fisher-normal-variance-known-mean}
For iid $N(0,v)$ observations with $v>0$, compute the score
and information using $\mathbb E_v[X_i^4]=3v^2$.
Check the bound for $n^{-1}\sum_iX_i^2$.
\end{exercise}

\begin{exercise}[Normal location and a squared target]
\label{exer:fisher-normal-squared-target}
In the known-variance normal location model, compute the
unbiased variance bound for estimating $\theta^2$. At which
label does this bound vanish?
\end{exercise}

\begin{exercise}[Bias under shrinkage]
\label{exer:fisher-bernoulli-affine-bias}
For $\delta=a\overline X+c$ in the iid Bernoulli model,
use the biased bound and decide when equality holds.
\end{exercise}

\begin{exercise}[Stationary mean and equality]
\label{exer:fisher-zero-target-derivative}
If $m'(\theta)=0$ and $0<I(\theta)<\infty$, use the equality
condition to characterize the estimators attaining the
variance bound at that label.
\end{exercise}

\begin{exercise}[Information after squaring a parameter]
\label{exer:fisher-zero-information-reparameterization}
Let $X\sim N(\eta^2,1)$ for $\eta\in\R$. Compute $I_\eta(\eta)$,
including at zero. Explain which hypothesis of the bound
fails there.
\end{exercise}

\begin{exercise}[A moving uniform location]
\label{exer:fisher-uniform-location-score}
For $X\sim\operatorname{Unif}[\theta,\theta+1]$, compute the
score at interior records. Show that the family is identifiable
and identify the failed common-support hypothesis.
\end{exercise}

\begin{exercise}[Taking logarithms of the uniform endpoint]
\label{exer:fisher-uniform-endpoint-log-parameter}
Write $\theta=e^\eta$ in Example~\ref{ex:fisher-uniform-moving-support}.
Compute the score and its second moment for $n$ draws.
Does this change of parameter give common support?
\end{exercise}

\begin{exercise}[Variance of an unbiased squared normal mean]
\label{exer:cramer-rao-normal-square-gap}
For $n$ iid $N(\theta,v)$ draws with known $v$, take
$\delta=\overline X^{\,2}-v/n$. Prove unbiasedness for $\theta^2$,
compute its variance from the normal moments, and find the gap
above the Cram\'er--Rao bound. Explain the gap using the
equality condition.
\end{exercise}

\begin{exercise}[Which Bernoulli targets can attain the bound?]
\label{exer:cramer-rao-bernoulli-affine-targets}
Suppose an estimator is unbiased for a differentiable $g(p)$
in the $n$-draw Bernoulli model with $0<p<1$. Show that
attaining the bound at even one interior label forces the
estimator to be affine in the count on the whole sample space.
Deduce which targets admit such attainment.
\end{exercise}

\begin{exercise}[A uniform boundary term for an expectation]
\label{exer:fisher-uniform-expectation-boundary}
For $X\sim\operatorname{Unif}[0,\theta]$ and continuous
$\delta:[0,\infty)\to\R$, differentiate
$m(\theta)=\theta^{-1}\int_0^\theta\delta(x)\,dx$.
Find the term missing from $\mathbb E_\theta[\delta s_\theta]$,
and show explicitly how it centers the derivative for
$\delta\equiv1$.
\end{exercise}

\begin{exercise}[Information in a finite reduction]
\label{exer:fisher-finite-statistic-contraction}
For a finite observation space with strictly positive
differentiable masses and a statistic $T$, show that the score
of the induced law equals $\mathbb E_\theta[s_\theta\mid T]$.
Compare its information with that of the full observation
and characterize equality. Show that sufficiency implies equality.
\end{exercise}

\begin{exercise}[Differentiation for a Poisson functional]
\label{exer:fisher-poisson-polynomial-envelope}
For $X\sim\operatorname{Poisson}(\lambda)$ and an estimator
$\delta(X)$ bounded in absolute value by a polynomial in $X$,
construct an envelope in
\eqref{eq:score-difference-domination} that also works after
multiplication by $|\delta|$. Derive
$\partial_\lambda\mathbb E_\lambda\delta=
\operatorname{Cov}_\lambda(\delta,X)/\lambda$.
\end{exercise}

\begin{exercise}[A bound for a nonlinear Poisson target]
\label{exer:cramer-rao-poisson-zero-probability}
For $n$ iid Poisson draws of rate $\lambda>0$, compute the
bound for an unbiased estimate of $e^{-\lambda}$. Use the
conditional estimator from Exercise~\ref{exer:rao-blackwell-poisson-zero}
for $n\geq2$ and the equality condition to decide whether it
attains the bound at any $\lambda>0$.
\end{exercise}



\chapter{Tests and confidence procedures}
\label{chap:tests-confidence}
\notationgroup{Tests and confidence procedures}
\section{Hypotheses, randomized tests, size, and power}
\label{sec:hypotheses-size-power}

Fix an experiment $\mathcal E=(\mathsf X,\mathcal A,
(P_\theta)_{\theta\in\Theta})$.

\subsection{Hypotheses and test functions}
\label{subsec:test-functions}

\begin{definition}[Hypotheses and tests]
\label{def:hypotheses-test-function}
\indexnotation{\Theta_0,\ \Theta_1,\ \phi}{Null and alternative parameter sets, and conditional probability of rejection.}
Let $\Theta_0,\Theta_1$ be nonempty disjoint subsets of $\Theta$,
called the \emph{null} and \emph{alternative} parameter sets.
A hypothesis is \emph{simple} if it specifies a single observation
law and \emph{composite} if it specifies more than one.
A \emph{test function} is a measurable $\phi:\mathsf X\to[0,1]$.
At record $x$ the probability of rejecting the null is $\phi(x)$.
A test is \emph{nonrandomized} if $\phi$ takes only the values zero
and one. Its \emph{rejection region} is $\{x:\phi(x)=1\}$.
\end{definition}

\begin{definition}[Power, errors, size, and level]
\label{def:test-size-power}
\indexnotation{\beta_\phi(\theta),\ \alpha(\phi)}{Power function and size of a test.}
The \emph{power function} is $\beta_\phi(\theta)=\mathbb E_\theta\phi(X)$.
At $\theta\in\Theta_0$ this is the probability of a \emph{type I error}.
At $\theta\in\Theta_1$, the probability of a \emph{type II error} is
$1-\beta_\phi(\theta)$. The \emph{size} is
\[
\alpha(\phi)=\sup_{\theta\in\Theta_0}\beta_\phi(\theta).
\]
For $0\leq\alpha\leq1$, a test has \emph{level} $\alpha$ if its
size is at most $\alpha$.
\end{definition}

\begin{proposition}[Binary decisions and randomized tests]
\label{prop:binary-decision-test}
Test functions correspond bijectively to decision kernels with
actions $0$ (retain the null) and $1$ (reject it). The test
$\phi$ is implemented by rejecting when $U<\phi(X)$, where
$U$ is an independent uniform random variable on $[0,1]$.
A mixture $t\phi+(1-t)\psi$ of two test functions, for $t\in[0,1]$,
has power $t\beta_\phi+(1-t)\beta_\psi$. In particular,
the level-$\alpha$ tests form a convex set.
\end{proposition}
\begin{proof}
In Definition~\ref{def:randomized-rule}, the two masses of the
kernel at $x$ are $1-\phi(x)$ and $\phi(x)$. These determine the
kernel, and any kernel on these actions gives a measurable $\phi$.
The independent uniform variable implements this kernel since the
Lebesgue measure of $\{u\in[0,1]:u<\phi(x)\}$ is $\phi(x)$.
Linearity of integration gives the power of a mixture. Under each
null label, a mixture of two powers at most $\alpha$ is at most
$\alpha$.
\end{proof}

\begin{example}[Two labels with the same observation law]
\label{ex:test-bernoulli-aliased}
Let $P_\theta$ be the Bernoulli law with probability $\theta^2$,
where $\theta\in[-1,1]$. For null $\{-1/2\}$ and alternative
$\{1/2\}$, every test has the same rejection probability under
both hypotheses. Thus every level-$\alpha$ test has power at
most $\alpha$ at the alternative. The constant test
$\phi(x)=\alpha$ attains this power.
\end{example}

\subsection{Calibrating rejection probabilities}
\label{subsec:test-calibration}

\begin{proposition}[Randomization at a threshold]
\label{prop:test-threshold-randomization}
For a measurable real statistic $T$, threshold $c$, and
$\gamma\in[0,1]$, define
\[
\phi_{c,\gamma}(x)=\mathbf 1_{\{T(x)>c\}}
+\gamma\mathbf 1_{\{T(x)=c\}} .
\]
Its power is $P_\theta(T>c)+\gamma P_\theta(T=c)$.
For a simple null with law $P_0$, if
$P_0(T>c)\leq\alpha\leq P_0(T\geq c)$ and $P_0(T=c)>0$, then
$\gamma=(\alpha-P_0(T>c))/P_0(T=c)$ gives size $\alpha$.
If $P_0(T=c)=0$ and the same inequalities hold, every $\gamma$
gives size $\alpha$.
\end{proposition}
\begin{proof}
Integrate the two indicators. When the atom has positive mass,
the inequalities on $\alpha$ put the stated $\gamma$ in $[0,1]$.
Substitution gives size $\alpha$. When the atom has mass zero
the two tail probabilities are equal, so the inequalities already
force the size to be $\alpha$.
\end{proof}

\begin{proposition}[Calibration at the upper endpoint of a binomial null]
\label{prop:test-binomial-composite-calibration}
Let $S$ be the number of successes in $n\geq1$ iid Bernoulli
trials with parameter $p\in[0,1]$. Every test of the form $r(S)$
with $r$ nondecreasing has nondecreasing power in $p$.
For $\Theta_0=[0,p_0]$, its size is its power at $p_0$.
For every $\alpha\in[0,1]$, a test of the form $\phi_{c,\gamma}$
with $T=S$ attains size $\alpha$.
\end{proposition}
\begin{proof}
On a common probability space take independent uniforms
$U_1,\ldots,U_n$ on $[0,1]$. For each $p$ the sum
$\sum_i\mathbf 1_{\{U_i<p\}}$ has the required binomial law.
This sum is nondecreasing in $p$ for every realization, hence so is
its image under $r$. Taking expectations proves monotonicity.
The null includes its upper endpoint, so the supremum is attained
there. Under $p_0$ the tail probabilities $P_{p_0}(S\geq j)$ for
$j=0,\ldots,n+1$ decrease from one to zero. Some consecutive pair
brackets $\alpha$. Use that pair as $P_{p_0}(S\geq c)$ and
$P_{p_0}(S>c)$ in Proposition~\ref{prop:test-threshold-randomization}.
\end{proof}

\begin{example}[Exact size for two Bernoulli trials]
\label{ex:test-binomial-randomized-tie}
For two iid Bernoulli trials, take null $p\leq1/2$ and target
size $1/2$. With $S$ the count, reject at $S=2$, reject with
probability $1/2$ at $S=1$, and retain the null at $S=0$.
The power is $p^2+\tfrac12(2p(1-p))=p$, so the size is $1/2$.
The nonrandomized upper-tail test rejecting only at $S=2$
has size $1/4$. Rejecting at both $S=1$ and $S=2$ gives size
$3/4$.
\end{example}

\begin{example}[A continuous upper-tail test]
\label{ex:test-normal-upper-tail}
Let $X_1,\ldots,X_n$ be iid $N(\theta,v)$ for known $v>0$,
and take null $\theta\leq\theta_0$. Let $\Phi$ be the standard
normal distribution function and let $\Phi(z)=1-\alpha$ for
$0<\alpha<1$. Reject when
$\overline X>\theta_0+z\sqrt{v/n}$. The power is
\[
1-\Phi\bigl(z+(\theta_0-\theta)\sqrt{n/v}\bigr).
\]
It increases with $\theta$, equals $\alpha$ at $\theta_0$,
and exceeds $\alpha$ at every $\theta>\theta_0$. The normal
law gives zero mass to the threshold, so a tie rule has no
effect.
\end{example}

\begin{example}[A null whose largest rejection probability is a limit]
\label{ex:test-open-null-size}
Observe one Bernoulli bit and reject on success. Take null
$p<1/2$ within $[0,1]$ and alternative $p\geq1/2$. The power
at $p$ is $p$. The size is $1/2$, though the probability of
rejection is strictly below $1/2$ at every null label.
\end{example}

\subsection{Testing and separation of laws}
\label{subsec:test-separation-laws}

\begin{theorem}[Separation of two laws by tests]
\label{thm:test-total-variation-gap}
For probability laws $P,Q$ on $(\mathsf X,\mathcal A)$,
\[
\sup_{0\leq\phi\leq1}\left(\int\phi\,dQ-\int\phi\,dP\right)
=d_{\mathrm{TV}}(P,Q),
\]
where the supremum ranges over measurable tests and is attained by
a nonrandomized test. Consequently the smallest possible sum
of the two error probabilities for testing $P$ against $Q$ is
$1-d_{\mathrm{TV}}(P,Q)$. A test with type I error at most $\alpha$
has power at most $\alpha+d_{\mathrm{TV}}(P,Q)$ at $Q$.
\end{theorem}
\begin{proof}
Take densities $p,q$ relative to $P+Q$. The difference is
$\int\phi(q-p)\,d(P+Q)$. It is at most
$\int(q-p)_+\,d(P+Q)$, and equality holds for the indicator
of $\{q>p\}$. The positive and negative parts of $q-p$ have
equal integrals because both laws have mass one. Thus the positive
part integrates to $d_{\mathrm{TV}}(P,Q)$ by
Proposition~\ref{prop:total-variation-density}.
The sum of errors is one minus the displayed difference, which
proves the minimum. The power bound follows by adding the type I
error to the bound on the difference.
\end{proof}

\begin{example}[The test attaining total variation]
\label{ex:test-tv-four-outcomes}
Use the laws of Example~\ref{ex:statistic-increased-testing-risk},
\[
P=(1/2,1/8,1/4,1/8),\qquad Q=(1/8,1/4,1/8,1/2).
\]
Reject $P$ on $\{2,4\}$, where $Q$ has larger mass. The type I
error is $1/4$ and the power at $Q$ is $3/4$, so the sum of
errors is $1/2$. The total variation is $1/2$, which certifies
optimality for this sum. The statistic grouping the first two
and last two outcomes yields total variation $1/4$, so the
smallest error sum based on that statistic is $3/4$.
\end{example}

\begin{example}[A test with both errors zero]
\label{ex:test-disjoint-support}
For null $P=\operatorname{Unif}[0,1]$ and alternative
$Q=\operatorname{Unif}[1,2]$, reject when $X>1$. Under $P$ the
rejection probability is zero. Under $Q$ the only possible
retention in its support is at $1$, which has probability zero,
so the power is one. In this case $d_{\mathrm{TV}}(P,Q)=1$.
\end{example}

\subsection{Exercises}
\label{subsec:exercises-6-1}

\begin{exercise}[A test with one bit]
\label{exer:test-one-bit-power}
Let $X\sim\operatorname{Bernoulli}(p)$ and set $\phi(0)=a$ and
$\phi(1)=b$, with $a,b\in[0,1]$. Compute the power at every $p$.
When is the power nondecreasing?
\end{exercise}

\begin{exercise}[Two types of error]
\label{exer:test-two-error-probabilities}
In the preceding exercise take $a=0$, $b=3/4$, null $p=1/4$,
and alternative $p=3/4$. Compute both error probabilities.
\end{exercise}

\begin{exercise}[Size and level]
\label{exer:test-size-versus-level}
A test has rejection probabilities $1/8$ and $1/4$ at its two null
labels. Give its size and the set of levels at which it qualifies.
\end{exercise}

\begin{exercise}[A composite alternative]
\label{exer:test-composite-alternative-power}
For one Bernoulli draw, reject on success. Take null $p\leq1/3$
and alternative $p>1/3$. Find the size and the supremum of the
type II error over the alternative.
\end{exercise}

\begin{exercise}[An uninformative test]
\label{exer:test-constant-level}
For an arbitrary experiment give a test whose power at every label
is $\alpha$. What are its errors at the null and alternative labels?
\end{exercise}

\begin{exercise}[A deterministic test as a kernel]
\label{exer:test-indicator-kernel}
Given a measurable set $A$, write the two action probabilities of
the kernel that rejects precisely on $A$. Verify measurability.
\end{exercise}

\begin{exercise}[Mixing rejection regions]
\label{exer:test-mixture-overlap}
Let $A,B$ be measurable rejection regions. Express the test that
chooses $A$ with probability $t$ and $B$ otherwise at each of the
four possible membership patterns. Compute its power.
\end{exercise}

\begin{exercise}[The size of a mixture]
\label{exer:test-size-mixture-strict}
Under the null observe a point mass at either zero or one. Take
two tests that reject at zero and one respectively. Compare the
size of their equal mixture with the average of their sizes.
\end{exercise}

\begin{exercise}[Calibration in a finite law]
\label{exer:test-finite-threshold-gamma}
Let a statistic have masses $(1/2,1/3,1/6)$ at $0,1,2$ under
the simple null. Choose an upper threshold and tie probability
for exact size $1/3$.
\end{exercise}

\begin{exercise}[A threshold of zero mass]
\label{exer:test-zero-threshold-atom}
Let $T$ be uniform on $[0,1]$ under the null. Find an upper
threshold for exact size $\alpha\in(0,1)$, and show that the tie
probability has no effect.
\end{exercise}

\begin{exercise}[An exact lower-tail test]
\label{exer:test-binomial-lower-tail}
For the Bernoulli count with $n=2$, take null $p\geq1/2$.
Give a nonincreasing count test of size $1/2$ and compute its power.
\end{exercise}

\begin{exercise}[A binomial upper-tail calculation]
\label{exer:test-binomial-three-calibration}
Take $n=3$ Bernoulli observations and null $p\leq1/2$.
Construct an upper-tail count test of size $1/4$ and compute
its power at $p=3/4$.
\end{exercise}

\begin{exercise}[An open binomial null]
\label{exer:test-open-binomial-boundary}
For $n\geq1$ and null $0\leq p<p_0<1$, show that the size of
a nondecreasing count test equals its power evaluated at $p_0$.
Use continuity of the binomial power function.
\end{exercise}

\begin{exercise}[Normal sample size]
\label{exer:test-normal-power-sample-size}
In Example~\ref{ex:test-normal-upper-tail}, fix $\theta_1>\theta_0$,
$0<\alpha<1/2$, and $1/2<b<1$. Find the smallest integer $n\geq1$
giving power at least $b$ at $\theta_1$.
\end{exercise}

\begin{exercise}[The opposite orientation]
\label{exer:test-total-variation-symmetry}
Swap the roles of $P$ and $Q$ in Theorem~\ref{thm:test-total-variation-gap}.
Relate an optimal test for one orientation to an optimal test for
the other.
\end{exercise}

\begin{exercise}[Two overlapping uniforms]
\label{exer:test-overlapping-uniforms}
Take $P$ uniform on $[0,1]$ and $Q$ uniform on $[a,a+1]$ with
$0<a<1$. Find their total variation distance and a rejection
region attaining the minimum sum of errors.
\end{exercise}

\begin{exercise}[A coarsened observation]
\label{exer:test-tv-constant-statistic}
When a test is based on a constant statistic, characterize all
possible pairs of type I and type II errors for simple hypotheses.
Compare their sum with the optimum using the full record.
\end{exercise}

\begin{exercise}[Power and proximity of laws]
\label{exer:test-power-lipschitz}
For any test $\phi$, prove
$|\beta_\phi(\theta)-\beta_\phi(\vartheta)|\leq d_{\mathrm{TV}}(P_\theta,P_\vartheta)$.
Deduce continuity of power at a label where the laws vary
continuously in total variation.
\end{exercise}

\begin{exercise}[The value of randomization for one trial]
\label{exer:test-one-bernoulli-optimal-power}
Test $p=p_0$ against $p=p_1$ from one Bernoulli draw, with
$0<p_0<p_1<1$. Maximize power subject to size at most $\alpha$
by optimizing over $\phi(0),\phi(1)\in[0,1]$. Derive the
power as a function of $\alpha$ and identify the levels at
which a nonrandomized test attains it.
\end{exercise}

\begin{exercise}[Convexity of the achievable error pairs]
\label{exer:test-error-region-finite}
For simple hypotheses on a finite space, show that the set of
achievable error pairs is the convex hull of the pairs from
nonrandomized tests. Represent any randomized test as a
mixture over rejection regions by choosing each outcome for
the region with the specified rejection probability.
\end{exercise}

\begin{exercise}[All tests attaining the total variation gap]
\label{exer:test-tv-equality-set}
Use densities relative to $P+Q$ to characterize every test attaining
the supremum in Theorem~\ref{thm:test-total-variation-gap}.
Explain the freedom on the set where the two densities agree.
\end{exercise}

\begin{exercise}[Testing against a common mixture]
\label{exer:test-composite-mixture-tv}
Let $P_1,\ldots,P_m$ be null laws and $Q$ a simple alternative.
If $\int\phi\,dP_i\leq\alpha$ for each $i$, show that
\[
\int\phi\,dQ\leq\alpha+
\inf_{w}\,d_{\mathrm{TV}}\!\left(Q,\sum_i w_iP_i\right),
\]
where the infimum is over probability weights. Give an example
where $Q$ equals a mixture of distinct null laws while it is
different from every individual null law.
\end{exercise}

\begin{exercise}[A bound from relative entropy]
\label{exer:test-error-sum-pinsker}
Combine Theorems~\ref{thm:test-total-variation-gap}
and~\ref{thm:pinsker} to bound the minimum sum of errors from
below in terms of $D(P\|Q)$. Apply the bound to $n$ independent
$N(0,1)$ observations against $n$ independent $N(h,1)$
observations. Evaluate the result for $h=c/\sqrt n$.
\end{exercise}

\begin{exercise}[Sufficiency and the attainable power functions]
\label{exer:test-sufficiency-power-region}
Suppose a statistic $T$ has a common conditional kernel as in
Definition~\ref{def:common-conditional-kernel}. Given a test on
the full record, construct a randomized test based on $T$ with
the same power at every label. Prove that the two experiments
have identical attainable power functions.
\end{exercise}



\clearpage
\section{Neyman--Pearson theory and composite hypotheses}
\label{sec:neyman-pearson-composite}

\subsection{Optimal tests for two simple hypotheses}
\label{subsec:neyman-pearson-lemma}

\begin{definition}[Most powerful tests]
\label{def:most-powerful-test}
A level-$\alpha$ test is \emph{most powerful} against a specified
alternative law if its power at that law is at least that of
every level-$\alpha$ test for the same null. It is \emph{uniformly
most powerful} (UMP) against an alternative family if the
same test is most powerful at every law in that family.
\end{definition}

For a simple null law $P$ and alternative law $Q$, write $p,q$
for their densities relative to $\mu=P+Q$. Define the extended
likelihood ratio $\Lambda=q/p$ on $\{p>0\}$, set $\Lambda=+\infty$
on $\{p=0<q\}$, and set $\Lambda=0$ where both densities vanish.
\indexnotation{\Lambda=q/p}{Likelihood ratio of the alternative law to the null law.}

\begin{theorem}[Neyman--Pearson lemma]
\label{thm:neyman-pearson}
Let $0<\alpha<1$. There are $k\geq0$ finite and $\gamma\in[0,1]$
such that
\[
\phi^\star=\mathbf 1_{\{\Lambda>k\}}+
\gamma\mathbf 1_{\{\Lambda=k\}},
\qquad \int\phi^\star\,dP=\alpha .
\]
Every such test is most powerful for testing $P$ against $Q$
at level $\alpha$. A level-$\alpha$ test $\phi$ has the same
power as $\phi^\star$ exactly when it agrees with $\phi^\star$
$\mu$-almost everywhere on $\{q\ne kp\}$ and
$k(\alpha-\int\phi\,dP)=0$.
\end{theorem}
\begin{proof}
Under $P$, $\Lambda$ is finite almost surely. Let $F$ be its
distribution function and take
$k=\inf\{t\geq0:F(t)\geq1-\alpha\}$. Since $F$ tends to one,
$k$ is finite. Right continuity gives $F(k)\geq1-\alpha$ and
the left limit satisfies $F(k-)\leq1-\alpha$, including at zero.
Thus $P(\Lambda>k)\leq\alpha\leq P(\Lambda\geq k)$. If the tie
has positive mass, choose $\gamma$ to fill the difference as in
Proposition~\ref{prop:test-threshold-randomization}. If the
tie has mass zero, the upper tail already has mass $\alpha$.

For any level-$\alpha$ test $\phi$, subtract its power from
that of $\phi^\star$ and write the difference as
\[
\int(\phi^\star-\phi)(q-kp)\,d\mu
+k\left(\alpha-\int\phi\,dP\right).
\]
The first integrand is nonnegative. Where $q>kp$ the threshold
test is one, and where $q<kp$ it is zero. The second term is
nonnegative by the size constraint. The sum is zero exactly
when both terms vanish. The first term vanishes precisely
when the tests agree almost everywhere on $\{q\ne kp\}$.
\end{proof}

\begin{example}[Randomization at a binomial likelihood ratio]
\label{ex:np-two-bernoulli}
For two iid Bernoulli observations, test $p=1/4$ against $p=3/4$
at level $\alpha=1/4$. The likelihood ratios at counts zero,
one, and two are $1/9,1,9$. Under the null the upper count
has mass $1/16$ and the middle count has mass $6/16$, so a
most powerful test rejects at count two and rejects with
probability $1/2$ at count one. Its power is
$9/16+(1/2)(6/16)=3/4$.
\end{example}

\begin{example}[An alternative with mass outside the null support]
\label{ex:np-singular-alternative}
Let $P$ be uniform on $[0,1]$ and let $Q$ be the equally weighted
mixture of $P$ and the uniform law on $[2,3]$. The ratio is
infinite on $[2,3]$ and $1/2$ on $[0,1]$, up to $\mu$-null sets.
A most powerful level-$\alpha$ test rejects on $[2,3]$ and
rejects with probability $\alpha$ on $[0,1]$. Its power is
$(1+\alpha)/2$.
\end{example}

\begin{example}[A most powerful test with size below its level]
\label{ex:np-unused-size}
On three outcomes take masses $P=(1/2,1/2,0)$ and
$Q=(1/4,0,3/4)$. Rejecting on the first and third outcomes
has size $1/2$ and power one. At level $3/4$ one can raise the
size to $3/4$ by rejecting with probability $1/2$ on the
second outcome, without changing the power. The exact-size
Neyman--Pearson test has threshold $k=0$ here.
\end{example}

\subsection{Monotone likelihood ratios and one-sided tests}
\label{subsec:monotone-likelihood-ratio}

\begin{definition}[Monotone likelihood ratio]
\label{def:monotone-likelihood-ratio}
Let $\Theta\subseteq\R$ be ordered by the usual order. Suppose the
model has densities relative to a common $\sigma$-finite measure
and a common measurable set $S$ on which each density is positive
almost everywhere. Each density vanishes almost everywhere off $S$.
The model has a \emph{monotone likelihood ratio} in a real statistic
$T$ if, for each $\theta_1<\theta_2$, the ratio
$f_{\theta_2}/f_{\theta_1}$ equals $r_{\theta_1,\theta_2}(T)$
almost everywhere on $S$ for some nonnegative nondecreasing
Borel function $r_{\theta_1,\theta_2}:\R\to[0,\infty]$.
\end{definition}

\begin{theorem}[One-sided tests with a monotone likelihood ratio]
\label{thm:ump-monotone-likelihood-ratio}
Suppose the model has a monotone likelihood ratio in $T$, and
let $\theta_0\in\Theta$. For $0<\alpha<1$, choose $c\in\R$ and
$\gamma\in[0,1]$ so the threshold test of
Proposition~\ref{prop:test-threshold-randomization} has rejection
probability $\alpha$ under $P_{\theta_0}$. Such a choice exists.
The test is UMP of size $\alpha$ for the null $\theta\leq\theta_0$
against $\theta>\theta_0$.
\end{theorem}
\begin{proof}
The same quantile argument as in Theorem~\ref{thm:neyman-pearson},
now for the real statistic $T$ under $P_{\theta_0}$, gives a finite
$c$ whose upper tails bracket $\alpha$. Choose $\gamma$ using
Proposition~\ref{prop:test-threshold-randomization}. Write $\phi^\star$
for the resulting test.

For $\theta<\theta_0$, let $r(T)=f_{\theta_0}/f_\theta$ on $S$.
Under $P_\theta$, $r(T)$ is integrable with mean one. Write
$\phi^\star=a(T)$ with $a$ nondecreasing. If $T'$ is an independent
copy of $T$ under $P_\theta$, then
\[
2\operatorname{Cov}_\theta(a(T),r(T))
=\mathbb E_\theta[(a(T)-a(T'))(r(T)-r(T'))]\geq0 .
\]
The product is integrable since $a$ is bounded. Changing
measure, $\mathbb E_{\theta_0}\phi^\star=\mathbb E_\theta[a(T)r(T)]$
is therefore at least $\mathbb E_\theta\phi^\star$. Hence the
size is attained at $\theta_0$.

Fix $\theta>\theta_0$ and let $r(T)=f_\theta/f_{\theta_0}$.
Put $k=r(c)$. This is finite since $r$ is nondecreasing,
finite $P_{\theta_0}$-almost surely at $T$, and
$P_{\theta_0}(T\geq c)\geq\alpha>0$. On $S$, the test is
one wherever $r(T)>k$ and zero wherever $r(T)<k$, almost
everywhere. Thus for any composite-null level-$\alpha$
test $\phi$, the same difference-of-power identity as in the
Neyman--Pearson proof applies with $p=f_{\theta_0}$ and
$q=f_\theta$. Its first term is nonnegative by these comparisons.
Its second term is nonnegative because $\theta_0$ is in the null.
This proves optimality at each $\theta>\theta_0$.
\end{proof}

\begin{example}[A binomial upper-tail test is UMP]
\label{ex:ump-binomial-upper-tail}
For iid Bernoulli trials with $0<p<1$, the likelihood ratio for
$p_2>p_1$ is
\[
\left(\frac{1-p_2}{1-p_1}\right)^n
\left(\frac{p_2(1-p_1)}{p_1(1-p_2)}\right)^S
\]
where $S$ is the count. The base of the second factor exceeds
one, so the ratio increases in $S$. For null $p\leq p_0$, the
upper-tail count test calibrated at $p_0$ is UMP by
Theorem~\ref{thm:ump-monotone-likelihood-ratio}.
\end{example}

\begin{example}[A normal upper-tail test is UMP]
\label{ex:ump-normal-upper-tail}
For iid $N(\theta,v)$ observations with known $v>0$, the log
likelihood ratio for $\theta_2>\theta_1$ is
\[
\frac{n(\theta_2-\theta_1)}{v}\overline X
-\frac{n(\theta_2^2-\theta_1^2)}{2v}.
\]
The ratio increases in $\overline X$, so the calibrated test
in Example~\ref{ex:test-normal-upper-tail} is UMP against
$\theta>\theta_0$.
\end{example}

\subsection{A two-sided alternative}
\label{subsec:np-no-ump-two-sided}

\begin{example}[No UMP test for two opposite normal shifts]
\label{ex:no-ump-opposite-normal-shifts}
Observe $X\sim N(\theta,1)$. Test $\theta=0$ against
$\theta\in\{-a,a\}$ with $a>0$ and $0<\alpha<1$.
Against $a$, the likelihood ratio is $\exp(ax-a^2/2)$,
strictly increasing in $x$. The unique most powerful test
up to null sets rejects on the right tail of null mass $\alpha$,
by the equality condition in Theorem~\ref{thm:neyman-pearson}
and the absence of atoms. Against $-a$, the ratio is strictly
decreasing, so the unique most powerful test rejects on the
left tail of null mass $\alpha$. The two regions differ on
a set of positive null mass. Therefore no level-$\alpha$ test
is most powerful against both alternatives.
\end{example}

\subsection{Exercises}
\label{subsec:exercises-6-2}

\begin{exercise}[A likelihood ratio on four outcomes]
\label{exer:np-four-ratios}
Let $P=(1/4,1/4,1/4,1/4)$ and $Q=(1/8,1/8,1/4,1/2)$.
Find an exact-size test at level $1/4$ and compute its power.
\end{exercise}

\begin{exercise}[A tie at the critical ratio]
\label{exer:np-finite-tie}
For the laws of Exercise~\ref{exer:np-four-ratios}, find all
most powerful level-$5/8$ tests. Identify where they can differ.
\end{exercise}

\begin{exercise}[An alternative equal to the null]
\label{exer:np-identical-laws}
For $P=Q$, compute the maximum power at level $\alpha$ and
characterize all tests attaining it.
\end{exercise}

\begin{exercise}[Disjoint supports]
\label{exer:np-mutually-singular}
Suppose $P$ and $Q$ are mutually singular. Construct a
nonrandomized test of size zero and power one. Describe an
exact-size randomized test at an arbitrary level $0<\alpha<1$.
\end{exercise}

\begin{exercise}[A zero ratio at the threshold]
\label{exer:np-zero-threshold-condition}
In Example~\ref{ex:np-unused-size}, determine all most powerful
tests at level $3/4$ on the three outcomes.
\end{exercise}

\begin{exercise}[One Bernoulli trial revisited]
\label{exer:np-one-bit-ratio}
For $0<p_0<p_1<1$, order the two likelihood ratios of one
Bernoulli draw and use a randomized threshold to derive
the maximum power for $0<\alpha<1$.
\end{exercise}

\begin{exercise}[Changing which law is the null]
\label{exer:np-swap-null-alternative}
Reverse the hypotheses in Example~\ref{ex:np-two-bernoulli}.
At level $1/4$, find the critical count, tie probability, and power.
\end{exercise}

\begin{exercise}[A positive critical ratio]
\label{exer:np-positive-critical-ratio}
Suppose $P$ and $Q$ have positive masses on the same finite space.
Show that every most powerful test at level $0<\alpha<1$ has
size exactly $\alpha$.
\end{exercise}

\begin{exercise}[Testing a normal scale]
\label{exer:np-normal-scale}
For one draw, test $N(0,1)$ against $N(0,4)$. Show that the
likelihood ratio increases with $|x|$ and give the most powerful
level-$\alpha$ rule in terms of the standard normal distribution.
\end{exercise}

\begin{exercise}[Testing an exponential rate]
\label{exer:np-exponential-rate-simple}
For one exponential draw, test rate $\lambda_0$ against rate
$\lambda_1>\lambda_0>0$. Determine the direction of the
threshold, calibrate it, and compute the power.
\end{exercise}

\begin{exercise}[Poisson likelihood ratios]
\label{exer:np-poisson-mlr}
For $n$ independent Poisson counts of common rate $\lambda>0$,
show that the family has a monotone likelihood ratio in the
total count. Specify the UMP rule for $\lambda\leq\lambda_0$
against $\lambda>\lambda_0$ in terms of a null quantile.
\end{exercise}

\begin{exercise}[A lower-tail normal test]
\label{exer:np-normal-lower-tail}
In the known-variance normal location model, construct a UMP
test for $\theta\geq\theta_0$ against $\theta<\theta_0$ by
reversing the order of the parameter and the sample mean.
\end{exercise}

\begin{exercise}[Monotonicity by a change of measure]
\label{exer:np-monotone-power-functions}
Under the hypotheses of Theorem~\ref{thm:ump-monotone-likelihood-ratio},
show that every bounded nondecreasing function of $T$ has
nondecreasing expectation in $\theta$.
\end{exercise}

\begin{exercise}[A change of scale for the statistic]
\label{exer:np-mlr-increasing-transform}
Suppose the statistic takes finitely many real values and a model
has a monotone likelihood ratio in that statistic. Show that any
strictly increasing real transformation of its values preserves
the likelihood ratio property and the threshold tests.
\end{exercise}

\begin{exercise}[An endpoint not in the null]
\label{exer:np-open-normal-null}
For a normal mean with known variance, take null $\theta<0$
and alternative $\theta>0$. Show that every level-$\alpha$
test also has rejection probability at most $\alpha$ at $\theta=0$.
Use continuity in total variation to conclude that the usual
upper-tail test calibrated at zero is UMP.
\end{exercise}

\begin{exercise}[A shared rejection region]
\label{exer:np-common-region-different-ratios}
Give two distinct alternative laws on a finite space for which
the same nonrandomized level-$\alpha$ test is most powerful
against either one at a specified simple null.
\end{exercise}

\begin{exercise}[A necessary condition for two-sided optimality]
\label{exer:np-two-sided-uniqueness}
In Example~\ref{ex:no-ump-opposite-normal-shifts}, write the two
critical regions using the standard normal quantiles. Show
that they differ on a set of positive null probability even
when $\alpha>1/2$.
\end{exercise}

\begin{exercise}[Testing after observing a sufficient count]
\label{exer:np-full-vs-count}
For $n$ iid Bernoulli observations and two distinct interior
parameters, compare the likelihood ratio on the record with
the ratio on the count. Explain how this gives an optimal
test on the full record.
\end{exercise}

\begin{exercise}[Power as a function of level]
\label{exer:np-power-concavity}
For a fixed simple null and simple alternative, show that
the supremum of the power over level-$\alpha$ tests is a
nondecreasing concave function of $\alpha$ on $[0,1]$.
Use mixtures of tests to prove concavity.
\end{exercise}

\begin{exercise}[Optimal tests for a finite family of alternatives]
\label{exer:np-finite-common-order}
Let $P$ have positive masses on a finite space and let
$Q_1,\ldots,Q_m$ be alternatives. Suppose one ordering of the
outcomes makes all the ratios $q_j/p$ nonincreasing. Construct
a test that is most powerful against every $Q_j$ at level
$\alpha$, and prove its optimality even when some ratios tie.
\end{exercise}

\begin{exercise}[Finite composite null and a mixture certificate]
\label{exer:np-null-mixture-certificate}
Let the null consist of laws $P_1,\ldots,P_m$ and let $Q$
be a simple alternative. Suppose a test is most powerful at
level $\alpha$ for the simple null $\sum_i w_iP_i$, where
the $w_i$ are probability weights. If it also has level
$\alpha$ for the composite null, show that it is most powerful
for the composite problem. Give a finite example with at least
two positive weights where the certificate applies.
\end{exercise}

\begin{exercise}[A unique nonrandomized test and a randomized tie]
\label{exer:np-finite-uniqueness}
Assume a simple null and alternative have positive masses
on the same finite space and all likelihood ratios are distinct.
For $0<\alpha<1$, prove that the most powerful test is unique.
Give a necessary and sufficient condition on $\alpha$ for that
test to be nonrandomized.
\end{exercise}

\begin{exercise}[A maximin test for opposite normal shifts]
\label{exer:np-maximin-opposite-shifts}
In Example~\ref{ex:no-ump-opposite-normal-shifts}, maximize the
smaller of the powers at $-a$ and $a$ at level $\alpha$.
First apply Neyman--Pearson to the equal mixture of the two
alternatives. Use the symmetry of the maximizing test to show
that it solves the maximin problem.
\end{exercise}

\begin{exercise}[A moving-support one-sided test]
\label{exer:np-uniform-one-sided}
Let $X_1,\ldots,X_n$ be iid uniform on $[0,\theta]$, $\theta>0$.
Test $\theta\leq\theta_0$ against $\theta>\theta_0$ by rejecting
for large sample maximum. Calibrate to size $\alpha\in(0,1)$
and prove the test is UMP by applying Neyman--Pearson at the
boundary null for each alternative. Include the region outside
the support at $\theta_0$.
\end{exercise}



\clearpage
\section{Confidence sets and inversion of tests}
\label{sec:confidence-sets-test-inversion}

Fix an experiment $\mathcal E=(\mathsf X,\mathcal A,
(P_\theta)_{\theta\in\Theta})$ and equip $\Theta$ with a
sigma algebra $\mathcal T$.

\subsection{Coverage and test inversion}
\label{subsec:confidence-coverage-inversion}

\begin{definition}[Confidence set and coverage]
\label{def:confidence-set-coverage}
\indexnotation{C(X),\ \operatorname{cov}_C(\theta)}{Confidence set and its coverage at parameter \(\theta\).}
A confidence procedure assigns a subset $C(x)\subseteq\Theta$ to
each record $x$, with measurable graph
\[
G_C=\{(x,t):t\in C(x)\}\in\mathcal A\otimes\mathcal T .
\]
Its coverage at $\theta$ is
$\operatorname{cov}_C(\theta)=P_\theta\{\theta\in C(X)\}$.
It has confidence level $1-\alpha$ if
$\inf_{\theta\in\Theta}\operatorname{cov}_C(\theta)\geq1-\alpha$.
It has \emph{pointwise exact coverage} $1-\alpha$ if equality
holds at every $\theta$. A randomized procedure may use an
independent $U$ uniform on $[0,1]$ as part of the observation.
Coverage then uses $P_\theta\otimes\operatorname{Unif}[0,1]$.
\end{definition}

\begin{theorem}[Test--confidence duality]
\label{thm:test-confidence-duality}
Suppose $\phi(x,t)\in[0,1]$ is product measurable and
$\mathbb E_t[\phi(X,t)]\leq\alpha$ for every $t\in\Theta$.
Using a single independent uniform $U$, form
\[
C(x,u)=\{t\in\Theta:u\geq\phi(x,t)\}.
\]
This is a confidence procedure with
$\operatorname{cov}_C(\theta)=1-\mathbb E_\theta[\phi(X,\theta)]$.
If the tests are nonrandomized one can use
$C(x)=\{t:\phi(x,t)=0\}$. Conversely, a confidence procedure of
level $1-\alpha$ gives a level-$\alpha$ test of each point
null $t$ by rejecting when $t\notin C$. For a randomized
confidence procedure the corresponding test function averages
this rejection indicator over $U$.
\end{theorem}
\begin{proof}
The graph is measurable since it is defined by the comparison
$u\geq\phi(x,t)$ of two product-measurable functions.
At $\theta$, conditional on $X=x$, the probability that $C$
excludes $\theta$ is $\phi(x,\theta)$. Integration gives the
coverage identity. For nonrandomized tests, the graph given by
$\phi(x,t)=0$ gives the same identity.

Conversely, the test function at $t$ is the indicator of
$t\notin C(x)$ in the nonrandomized case. In the randomized case it is
$\int_0^1\mathbf 1_{\{t\notin C(x,u)\}}\,du$, product measurable by
integration of the graph indicator. Its rejection probability
at $t$ is $1-\operatorname{cov}_C(t)\leq\alpha$ by Tonelli.
\end{proof}

\begin{example}[Exact coverage with one Bernoulli observation]
\label{ex:confidence-bernoulli-shared-randomizer}
For $X$ Bernoulli of parameter $t$, fix $0<\alpha<1$.
Against larger success probabilities, the exact-size test
of $t$ rejects at $X=1$ with probability $\min\{1,\alpha/t\}$
for $t>0$ and with probability one for $t=0$. At $X=0$ it
rejects with probability $(\alpha-t)/(1-t)$ for $t<\alpha$
and with probability zero for $t\geq\alpha$.
For $t\geq\alpha$ its rejection probability under $t$ is
$t(\alpha/t)=\alpha$. For $t<\alpha$ the probability is
$t+(1-t)(\alpha-t)/(1-t)=\alpha$. Inversion gives pointwise
exact coverage using the same $U$ for every $t$.
\end{example}

\subsection{Pivots and continuous models}
\label{subsec:confidence-pivots}

\begin{definition}[Pivot]
\label{def:confidence-pivot}
\indexnotation{W(X,\theta)}{Pivot whose law at the true parameter is parameter independent.}
A real pivot is a product-measurable function $W(x,t)$ whose
law under $P_\theta$ at $t=\theta$ is the same for all $\theta$.
\end{definition}

\begin{proposition}[Inverting a pivot]
\label{prop:confidence-pivot-inversion}
Let $W$ be a pivot with common law $\nu$ at the true label.
If $A$ is Borel and $\nu(A)\geq1-\alpha$, then
\[
C(x)=\{t:W(x,t)\in A\}
\]
has coverage $\nu(A)$ at every label.
\end{proposition}
\begin{proof}
The graph is the inverse image of $A$ under $W$. The event
that $C(X)$ contains $\theta$ is $\{W(X,\theta)\in A\}$, of
probability $\nu(A)$.
\end{proof}

\begin{example}[A normal mean with known variance]
\label{ex:confidence-normal-known-variance}
For iid $N(\theta,v)$ observations with known $v>0$,
$W(X,\theta)=(\overline X-\theta)\sqrt{n/v}$ has law $N(0,1)$.
Write $z_u=\Phi^{-1}(u)$ for $0<u<1$. The region
$[-z_{1-\alpha/2},z_{1-\alpha/2}]$ has mass $1-\alpha$.
Inversion gives
\[
\left[\overline X-z_{1-\alpha/2}\sqrt{v/n},\
      \overline X+z_{1-\alpha/2}\sqrt{v/n}\right].
\]
The same interval results from inverting the two-tailed
tests that reject $t$ when $|(\overline X-t)\sqrt{n/v}|$
exceeds $z_{1-\alpha/2}$.
\end{example}

\begin{example}[A uniform endpoint]
\label{ex:confidence-uniform-endpoint}
For iid $\operatorname{Unif}[0,\theta]$ observations let $T$ be
the maximum. By Example~\ref{ex:statistic-uniform-maximum},
$(T/\theta)^n$ is uniform on $[0,1]$. Inverting the pivot
region $[\alpha,1]$ yields
\[
C(X)=[T,T\alpha^{-1/n}]\cap(0,\infty).
\]
Its coverage is exactly $1-\alpha$ for every $\theta>0$.
The parameter at the true endpoint enters through the support,
and the argument uses the distribution of the maximum directly.
\end{example}


\begin{example}[Changing the allocation of the normal tails]
\label{ex:confidence-normal-tail-allocation}
In the known-variance normal model, assign probability $a$
to the lower pivot tail and $b$ to the upper tail, where
$a,b>0$ and $a+b=\alpha$. The interval is
\[
[\overline X-z_{1-b}\sqrt{v/n},\
 \overline X-z_a\sqrt{v/n}].
\]
Its coverage is $1-\alpha$. For $b<a$ its lower endpoint
is farther below $\overline X$ than in the equal-tail
interval, while the upper endpoint also moves downward.
Both facts follow from strict monotonicity of the normal quantiles.
\end{example}

\subsection{Discrete tails and binomial intervals}
\label{subsec:confidence-discrete}

\begin{definition}[$p$-value]
\label{def:test-p-value}
\indexnotation{p_t(X)}{\(p\)-value for the null label \(t\).}
A $p$-value for the null label $t$ is a measurable
$p_t:\mathsf X\to[0,1]$ satisfying
\[
P_t\{p_t(X)\leq u\}\leq u\qquad(0\leq u\leq1).
\]
The same definition applies on the augmented observation space
when the $p$-value uses independent randomization.
\end{definition}

\begin{proposition}[Discrete upper tails]
\label{prop:discrete-tail-pvalues}
Let $S$ take finitely many ordered real values and set
\[
p_t^+(s)=P_t(S\geq s),\qquad
\widetilde p_t^+(s,u)=P_t(S>s)+uP_t(S=s).
\]
The first expression evaluated at $S$ is a $p$-value.
For independent $U$ uniform on $[0,1]$, the second expression
evaluated at $(S,U)$ is uniform on $[0,1]$ under $P_t$.
The corresponding lower-tail statements follow by replacing
$S$ with $-S$.
\end{proposition}
\begin{proof}
For the first statement, if there is a possible value $s$ with
$p_t^+(s)\leq u$, let $s_0$ be the smallest such value.
Then the event $\{p_t^+(S)\leq u\}$ is $\{S\geq s_0\}$ and has
probability $p_t^+(s_0)\leq u$. If there is no such value the
event is empty.

For the second statement, conditional on a value $s$ with
positive mass, $\widetilde p_t^+(s,U)$ is uniform on
$[P_t(S>s),P_t(S\geq s)]$. In descending order of $s$, these
intervals cover $[0,1]$ with disjoint interiors. Each interval
is selected with probability equal to its length. The
resulting distribution is uniform, as is seen by summing
the lengths of the intersections with any Borel set.
\end{proof}

\begin{proposition}[Equal-tail binomial interval]
\label{prop:clopper-pearson-interval}
Let $S$ be the count of $n\geq1$ iid Bernoulli draws with
parameter $p\in[0,1]$. For $0<\alpha<1$, put $a=\alpha/2$.
At observed count $k$, define $L_0=0$, $U_n=1$, and the
remaining endpoints by
\[
P_{L_k}(S\geq k)=a\quad (1\leq k\leq n),\qquad
P_{U_k}(S\leq k)=a\quad (0\leq k<n).
\]
The equations have unique solutions in $(0,1)$. The set
$[L_k,U_k]$, called the Clopper--Pearson interval, has
coverage at least $1-\alpha$ at every $p$.
\end{proposition}
\begin{proof}
The binomial upper tail at $k\geq1$ is continuous in $p$ and
strictly increases from zero to one. For strict increase
between two interior probabilities $p<q$, use the shared
uniform construction in Proposition~\ref{prop:test-binomial-composite-calibration}.
There is positive probability that exactly $k-1$ uniforms
are below $p$, one is in $(p,q)$, and all others exceed $q$.
The count crosses $k$ on this event. The lower tail at
$k<n$ is one minus the upper tail at $k+1$, so it strictly
decreases. This proves existence and uniqueness of the roots.
For $0<k<n$, at $p=L_k$ the lower tail is at least
$1-P_{L_k}(S\geq k)=1-a>a$. Hence $L_k<U_k$, and the endpoint
cases also satisfy $L_k\leq U_k$.

At the true $p$, exclusion by the lower endpoint requires
$P_p(S\geq k)<a$, by monotonicity. Exclusion by the upper
endpoint requires $P_p(S\leq k)<a$. Each event has probability
at most $a$ by Proposition~\ref{prop:discrete-tail-pvalues}.
The union has probability at most $2a=\alpha$.
\end{proof}

\begin{example}[Coverage on a finite space]
\label{ex:confidence-binomial-boundary}
At count zero the binomial interval is
$[0,1-(\alpha/2)^{1/n}]$ since $P_p(S=0)=(1-p)^n$.
At count $n$ it is $[(\alpha/2)^{1/n},1]$.
At $p=0$ and $p=1$ the true label is included with probability
one, so the interval's uniform coverage guarantee can be
strict.
\end{example}

\begin{example}[An exact randomized one-sided bound]
\label{ex:confidence-binomial-randomized-bound}
Let $S$ be the binomial count and let $U$ be independent and
uniform on $[0,1]$. Put
\[
C(S,U)=\{p\in[0,1]:\widetilde p_p^+(S,U)>\alpha\}.
\]
The graph is measurable since for each $s$ the randomized
tail is a polynomial in $p$ with coefficients depending
continuously on $U$. Under the true $p$ it is uniform by
Proposition~\ref{prop:discrete-tail-pvalues}, so the coverage
is $1-\alpha$ even at $p=0,1$. For fixed $(s,u)$ the tail is
$(1-u)P_p(S>s)+uP_p(S\geq s)$, a nondecreasing function of
$p$. Thus $C$ is an upper interval in the parameter space,
possibly empty or the whole space.
\end{example}

\subsection{Exercises}
\label{subsec:exercises-6-3}

\begin{exercise}[Reading coverage]
\label{exer:coverage-from-acceptance}
If a point-null test rejects the true label $t$ with probability
$0.04$ at every $t$, find the coverage of its inversion.
\end{exercise}

\begin{exercise}[Restricting the parameter space]
\label{exer:confidence-restricted-parameter}
For the normal mean interval in Example~\ref{ex:confidence-normal-known-variance},
intersect the interval with a known parameter space $[0,\infty)$.
Prove that coverage remains exact at every allowed label.
\end{exercise}

\begin{exercise}[Randomization on a singleton model]
\label{exer:confidence-singleton-randomization}
For a model with one label and a constant observation, describe
all possible coverages of nonrandomized confidence sets. Construct
a randomized confidence set with coverage $1-\alpha$.
\end{exercise}

\begin{exercise}[When the set is empty]
\label{exer:confidence-empty-set-bound}
Show that a level-$1-\alpha$ confidence procedure produces the
empty set with probability at most $\alpha$ at each label,
provided the empty-set event is measurable.
\end{exercise}

\begin{exercise}[Inversion on two labels]
\label{exer:confidence-two-label-inversion}
Let the labels be $0,1$ and let $P_t(X=t)=3/4$ for a binary
observation. Invert the tests that each reject their own label
when $X\ne t$. Find the resulting set and its coverage.
\end{exercise}

\begin{exercise}[An upper bound for a normal mean]
\label{exer:confidence-normal-upper-bound}
Invert the lower-tail point-null tests at level $\alpha$ for
a normal mean with known variance. Write the resulting
one-sided confidence set and verify coverage directly.
\end{exercise}

\begin{exercise}[Interval length and sample size]
\label{exer:confidence-normal-fixed-width}
In Example~\ref{ex:confidence-normal-known-variance}, find
the smallest $n\geq1$ giving interval length at most $w>0$.
\end{exercise}

\begin{exercise}[A pivot for a uniform minimum]
\label{exer:confidence-uniform-minimum}
For $n$ iid draws from $\operatorname{Unif}[0,\theta]$, $\theta>0$,
find the law of the sample minimum divided by $\theta$. Invert a
central pivot region to obtain an interval for $\theta$.
\end{exercise}

\begin{exercise}[Reallocating pivot tails]
\label{exer:confidence-uniform-pivot-tail-allocation}
In Example~\ref{ex:confidence-uniform-endpoint}, invert
$[a,1-b]$ for $a>0$, $b\geq0$, and $a+b=\alpha$.
Compare the lower and upper endpoints with those for $b=0$.
\end{exercise}

\begin{exercise}[A transformed normal interval]
\label{exer:confidence-normal-exp-transform}
For a normal mean with known variance, transform both
endpoints of the usual interval by the exponential map.
Find its coverage for the target $e^\theta$.
\end{exercise}

\begin{exercise}[A finite upper-tail p-value]
\label{exer:confidence-discrete-tail-table}
Let $S$ have masses $(1/2,1/3,1/6)$ at $0,1,2$ under a
point null. Tabulate $P(S\geq s)$ and compute the size of
the test that rejects when this tail is at most $1/4$.
\end{exercise}

\begin{exercise}[Randomization intervals]
\label{exer:confidence-discrete-randomization-table}
For the law in Exercise~\ref{exer:confidence-discrete-tail-table},
list the three intervals used by the randomized upper-tail
$p$-value. Compute its rejection probability at level $1/4$.
\end{exercise}

\begin{exercise}[Omitting the observed atom]
\label{exer:confidence-strict-tail-invalid}
For one fair Bernoulli draw, show that $P(S>s)$ evaluated
at the observed $s=S$ is not a valid $p$-value. Exhibit a
level at which it rejects too often.
\end{exercise}

\begin{exercise}[A two-tailed p-value]
\label{exer:confidence-two-tail-union-bound}
Suppose $p^-,p^+$ are valid $p$-values under the same law.
Show that $\min\{1,2\min(p^-,p^+)\}$ is also valid.
Which inequality in the proof can be strict?
\end{exercise}

\begin{exercise}[Endpoints with one draw]
\label{exer:confidence-cp-one-draw}
Write the Clopper--Pearson intervals at counts zero and one
when $n=1$. Determine the coverage function on $[0,1]$,
including at the two interior interval endpoints.
\end{exercise}

\begin{exercise}[Binomial symmetry]
\label{exer:confidence-cp-symmetry}
Show that the equal-tail binomial endpoints satisfy
$L_k=1-U_{n-k}$ for $0\leq k\leq n$.
\end{exercise}

\begin{exercise}[A randomized lower bound at the boundary]
\label{exer:confidence-randomized-bound-zero}
In Example~\ref{ex:confidence-binomial-randomized-bound}, take
the true parameter to be zero. Determine the event in $U$
on which zero belongs to the set. When does the set contain
the entire parameter space?
\end{exercise}

\begin{exercise}[Conservatism at an interior label]
\label{exer:confidence-binomial-interior-conservatism}
For $n=2$ and nominal confidence level $3/4$, use the
binomial tail probabilities at $p=1/2$ to compute the
coverage of the equal-tail interval there.
\end{exercise}

\begin{exercise}[Nested tests and confidence sets]
\label{exer:confidence-inversion-nesting}
Suppose $\phi_\alpha(x,t)$ is nondecreasing in $\alpha$ for
every $(x,t)$, with level $\alpha$ at each point null.
Use one common $U$ to prove nesting of the inverted
confidence sets as $\alpha$ varies.
\end{exercise}

\begin{exercise}[Expected length from inclusion probabilities]
\label{exer:confidence-expected-length-tonelli}
Let $\Theta\subseteq\R$ be Borel and let $C$ have a
measurable graph. Use Tonelli to write the expected Lebesgue
measure of $C(X)$ under $P_\theta$ as an integral of
inclusion probabilities. Interpret the integrand for a false
label in terms of the inverted test.
\end{exercise}

\begin{exercise}[Exclusion probabilities for one-sided intervals]
\label{exer:confidence-normal-ump-lower-bound}
In the known-variance normal location model, restrict to
confidence procedures whose sets are upper intervals in $\R$.
For fixed $t$, compare exclusion probabilities under $\theta>t$
using Theorem~\ref{thm:ump-monotone-likelihood-ratio}.
Show that inverting its tests maximizes those exclusion
probabilities simultaneously for all $t<\theta$.
\end{exercise}

\begin{exercise}[An exact two-sided randomized binomial procedure]
\label{exer:confidence-binomial-two-sided-randomized}
Use $R=\widetilde p_p^+(S,U)$ and reject the point null $p$
when $R\leq\alpha/2$ or $R\geq1-\alpha/2$. Prove that
the inverted set has exact coverage at every label and
is an interval when nonempty. Allow open endpoints.
\end{exercise}

\begin{exercise}[Testing power and confidence-set discrimination]
\label{exer:confidence-false-label-power}
For two normal means $t<\theta$, invert the family of
level-$\alpha$ upper-tail tests. Compute $P_\theta(t\notin C(X))$.
Show that no other level-$1-\alpha$ confidence procedure
excludes $t$ with greater probability under $P_\theta$.
\end{exercise}

\begin{exercise}[Inverting tests for a function of the parameter]
\label{exer:confidence-target-composite-null}
Let $g:\Theta\to\R$ be measurable. For each possible target
value $r$, suppose a product-measurable test family has
size at most $\alpha$ for the null $\{\theta:g(\theta)=r\}$.
Construct a confidence set for $g(\theta)$ and prove its
uniform coverage. Explain how to use a fixed independent
uniform variable for all target values.
\end{exercise}



\clearpage
\section{Multiple decisions and error criteria}
\label{sec:multiple-decisions-error-criteria}

\subsection{Choosing one of several labels}
\label{subsec:multiclass-decision-error}

\begin{definition}[Classification error]
\label{def:multiclass-error}
Let $\Theta=\{1,\ldots,m\}$ with observation laws $P_i$. A
classification rule is a decision kernel with the same set of
actions. Write $q_j(x)$ for the probability of choosing $j$ at $x$.
The zero-one risk at $i$ is $e_i(q)=1-\int q_i\,dP_i$.
For a prior $\pi$ on the labels the average error is
$e_\pi(q)=\sum_i\pi_i e_i(q)$ and the worst-case error is
$\max_i e_i(q)$.
\end{definition}

\begin{proposition}[Minimum average classification error]
\label{prop:multiclass-bayes-error}
Let $f_i=dP_i/d\mu$ for a common finite dominating measure.
The minimum average error is
\[
1-\int \max_i\{\pi_i f_i(x)\}\,\mu(dx).
\]
A rule attains it exactly when it puts all its action mass on
maximizers of $\pi_i f_i(x)$ for $\mu$-almost every $x$.
\end{proposition}
\begin{proof}
The prior probability of a correct decision is
$\int\sum_i q_i(x)\pi_i f_i(x)\,\mu(dx)$.
The integrand is a convex combination bounded by the maximum.
Choosing the first maximizing index is measurable by finite
comparisons and attains that bound. For equality, the nonnegative
gap must vanish almost everywhere, which holds exactly when
each index with positive action mass maximizes the weighted density.
\end{proof}

\begin{example}[Average and worst-case errors]
\label{ex:multiclass-no-observation}
Suppose all three labels give the same observation law. Every
rule has a common marginal action distribution $(r_1,r_2,r_3)$
and error vector $(1-r_1,1-r_2,1-r_3)$. For prior
$(1/2,1/3,1/6)$ the average error is minimized by always choosing
label one, giving average error $1/2$ and worst-case error one.
The minimum worst-case error is $2/3$, attained by choosing
each label with probability $1/3$. Indeed, at least one $r_i$
is at most $1/3$.
\end{example}

\begin{example}[Three labels with a shared observation]
\label{ex:multiclass-shared-observation}
For $i=1,2,3$, let $P_i$ assign mass $3/4$ to outcome $i$
and $1/4$ to outcome $0$. Under the uniform prior, choose $i$
at outcome $i$ and randomize uniformly at zero. At each label
the error is $(1/4)(2/3)=1/6$. This is Bayes by
Proposition~\ref{prop:multiclass-bayes-error}. The worst-case
error of any rule is at least its uniform prior average,
so the same rule is minimax.
\end{example}

\subsection{Rejecting several hypotheses}
\label{subsec:multiple-testing-errors}

\begin{definition}[Family-wise error and false discovery rate]
\label{def:fwer-fdr}
\indexnotation{\mathcal I_0(\theta),\ R,\ V}{Indices of true nulls, number of rejections, and number of false rejections.}
\indexnotation{\operatorname{FWER},\ \operatorname{FDR}}{Family-wise error rate and false discovery rate.}
Let $H_1,\ldots,H_m$ be null hypotheses, each a subset of the
parameter space. A multiple testing procedure returns a measurable
random subset of rejected indices. Under label $\theta$ put
$\mathcal I_0(\theta)=\{i:\theta\in H_i\}$. Let $R$ count all
rejections and $V$ count rejections in $\mathcal I_0(\theta)$. Define
\[
\operatorname{FWER}_\theta=P_\theta(V\geq1),\qquad
\operatorname{FDR}_\theta=\mathbb E_\theta\!\left[\frac{V}{\max\{R,1\}}\right].
\]
If the procedure uses auxiliary randomness these probabilities and
expectations include that randomness. Control at level $\alpha$
means that the specified rate is at most $\alpha$ at every label.
\end{definition}

\begin{proposition}[Comparing error criteria]
\label{prop:fdr-fwer-comparison}
At each $\theta$, $\operatorname{FDR}_\theta\leq
\operatorname{FWER}_\theta$. The two rates agree when every
null is true. If $\mathcal I_0(\theta)=\varnothing$, both
rates are zero.
\end{proposition}
\begin{proof}
The fraction $V/\max\{R,1\}$ is bounded by $\mathbf 1_{\{V\geq1\}}$.
If all nulls are true, $V=R$ and this bound is an equality.
If there are no true nulls then $V=0$.
\end{proof}

\begin{example}[One false rejection among several rejections]
\label{ex:fdr-fwer-separated}
For a fixed label, suppose the procedure always rejects four
false nulls and also rejects the only true null with probability
$1/4$. Then $\operatorname{FWER}_\theta=1/4$. On a false
rejection the fraction is $1/5$, and otherwise it is zero.
Thus $\operatorname{FDR}_\theta=1/20$.
\end{example}

\begin{example}[Separate level guarantees]
\label{ex:multiple-independent-null-unadjusted}
Under the global null, let the $m$ $p$-values be independent
uniform random variables on $[0,1]$. Rejecting each at threshold
$\alpha$ gives marginal rejection probability $\alpha$ and
family-wise error $1-(1-\alpha)^m$. For $m>1$ and $0<\alpha<1$
this exceeds $\alpha$.
\end{example}

\subsection{Family-wise control from marginal \texorpdfstring{$p$}{p}-values}
\label{subsec:bonferroni-holm}

For each $i$, let $p_i$ satisfy $P_\theta(p_i\leq u)\leq u$
for every $\theta\in H_i$ and $0\leq u\leq1$. This extends
Definition~\ref{def:test-p-value} to a composite null.

\begin{theorem}[Bonferroni control]
\label{thm:bonferroni-control}
Let $w_i\geq0$ with $\sum_i w_i=1$, and let $0\leq\alpha\leq1$.
Reject $i$ when $p_i\leq\alpha w_i$. The procedure controls
family-wise error at level $\alpha$.
\end{theorem}
\begin{proof}
At $\theta$ a false rejection means that
$p_i\leq\alpha w_i$ for at least one $i\in\mathcal I_0(\theta)$.
By the union bound and marginal validity its probability is at
most $\alpha\sum_{i\in\mathcal I_0(\theta)}w_i\leq\alpha$.
\end{proof}

\begin{theorem}[Holm's step-down procedure]
\label{thm:holm-control}
Order the $p$-values as $p_{(1)}\leq\cdots\leq p_{(m)}$,
breaking ties by index. Reject in order as long as
$p_{(j)}\leq\alpha/(m-j+1)$, stopping at the first failure.
This procedure controls family-wise error at level $\alpha$
and rejects every hypothesis rejected by equal-weight Bonferroni.
\end{theorem}
\begin{proof}
Fix $\theta$ and let $m_0$ be the number of true nulls.
When $m_0=0$ there is no false rejection. Otherwise let $j$
be the rank of the first true null. A false rejection requires
rejection at rank $j$. Since at least $m_0$ hypotheses remain
at that rank, the smallest true-null $p$-value must be at most
$\alpha/(m-j+1)\leq\alpha/m_0$. A union bound over the $m_0$
true nulls bounds this probability by $\alpha$.

All ranks at or before an equal-weight Bonferroni rejection
have $p$-value at most $\alpha/m$. Each step-down threshold
is at least $\alpha/m$, so none of those ranks stops the
procedure.
\end{proof}

\begin{example}[Holm can reject more hypotheses]
\label{ex:holm-extra-rejections}
Take level $0.06$ and ordered values
\[(0.01,0.025,0.04).\]
Equal-weight Bonferroni rejects only the first, at threshold
$0.02$. Holm compares the three values successively with
$0.02,0.03,0.06$ and rejects all three.
\end{example}

\begin{example}[Simultaneous confidence sets]
\label{ex:bonferroni-simultaneous-confidence}
Suppose $C_1,\ldots,C_m$ have marginal coverage at least
$1-\alpha/m$ for targets $g_1(\theta),\ldots,g_m(\theta)$,
respectively. The product set $\prod_i C_i$ covers the full
target vector with probability at least $1-\alpha$. Indeed,
its failure is the union of the marginal failures. For $m$
normal means with known variances, this gives simultaneous
intervals by using the quantile $z_{1-\alpha/(2m)}$ in each
interval from Example~\ref{ex:confidence-normal-known-variance}.
\end{example}

\subsection{False discovery rate under independence}
\label{subsec:benjamini-hochberg}

\begin{theorem}[Benjamini--Hochberg control]
\label{thm:benjamini-hochberg-independence}
At a label $\theta$, suppose each true-null $p$-value is independent
of the vector of the other $p$-values. For $0<\alpha<1$, set
\[
r=\max\bigl(\{j:1\leq j\leq m,\ p_{(j)}\leq\alpha j/m\}\cup\{0\}\bigr).
\]
Reject the first $r$ hypotheses in the ordering. If $m_0$ nulls
are true, then $\operatorname{FDR}_\theta\leq\alpha m_0/m$.
Equality holds if every true-null $p$-value is uniform.
\end{theorem}
\begin{proof}
Write $N(t)=\#\{i:p_i\leq t\}$. The condition at rank $j$
is $N(\alpha j/m)\geq j$. If $r>0$, then
$N(\alpha r/m)=r$, since a larger count would make rank $r+1$
pass its threshold. Thus the rejected indices are exactly those
with $p_i\leq\alpha r/m$ when $r>0$.

For a true null $i$, let $r_i$ be the rejection count after
setting $p_i=0$. Then $r_i\geq1$ and $r_i\geq r$. If $i$ is
rejected originally, replacing its value by zero leaves all
counts at thresholds of ranks greater than $r$ unchanged.
Hence $r_i=r$. Conversely, if $p_i\leq\alpha r_i/m$, the
count at threshold $\alpha r_i/m$ is unchanged by the replacement,
so $r\geq r_i$. Thus $r=r_i$ and $i$ is rejected. Its contribution
to the expected false discovery proportion is therefore
\[
\sum_{k=1}^m \frac{1}{k}
 P_\theta(r_i=k,\ p_i\leq\alpha k/m)
\leq \sum_{k=1}^m \frac{\alpha}{m}P_\theta(r_i=k)
=\frac{\alpha}{m}.
\]
The inequality uses independence since $r_i$ depends only on
the other $p$-values. It is an equality for a uniform $p_i$.
Summing over the true nulls proves both assertions.
\end{proof}

\subsection{Exercises}
\label{subsec:exercises-6-4}

\begin{exercise}[A binary rule as a classifier]
\label{exer:multiple-binary-classification}
For two labels, express the two classification errors in terms of the power of the test that rejects label one by choosing label two.
\end{exercise}

\begin{exercise}[Equal laws and a uniform prior]
\label{exer:multiple-equal-laws}
For $m$ labels with identical observation laws, compute the Bayes error under a uniform prior. Characterize all Bayes rules.
\end{exercise}

\begin{exercise}[A point mass and a shared atom]
\label{exer:multiple-shared-atom}
Extend Example~\ref{ex:multiclass-shared-observation} to $m$ labels with mass $a$ at zero and $1-a$ at its own label. Compute the minimax error.
\end{exercise}

\begin{exercise}[Weighted densities on two outcomes]
\label{exer:multiple-weighted-densities}
Use laws $(1/4,3/4)$, $(1/2,1/2)$, and $(3/4,1/4)$ on two outcomes, with equal priors. Find the Bayes action at each outcome and its average error.
\end{exercise}

\begin{exercise}[A label with zero prior mass]
\label{exer:multiple-zero-prior}
In Proposition~\ref{prop:multiclass-bayes-error}, explain when a label of prior mass zero can be a maximizing action.
\end{exercise}

\begin{exercise}[No rejections]
\label{exer:multiple-no-rejections}
For a procedure that rejects nothing, compute $R$, $V$, FDR, and FWER at every label.
\end{exercise}

\begin{exercise}[Two true nulls]
\label{exer:multiple-global-null}
When both nulls are true, tabulate the false discovery proportion for all possible rejection sets and recover the equality of the two error rates.
\end{exercise}

\begin{exercise}[One true and one false null]
\label{exer:multiple-two-null-rates}
Suppose the false null is always rejected and the true null is rejected with probability $b$. Compute both rates.
\end{exercise}

\begin{exercise}[A bound involving the expected number of false rejections]
\label{exer:multiple-expected-false-count}
Show that $P_\theta(V\geq1)\leq\mathbb E_\theta V$, and characterize equality.
\end{exercise}

\begin{exercise}[Independent uniform null values]
\label{exer:multiple-independent-bonferroni}
For Example~\ref{ex:multiple-independent-null-unadjusted}, compute FWER when the marginal threshold is $\alpha/m$.
\end{exercise}

\begin{exercise}[Identical null values]
\label{exer:multiple-identical-pvalues}
Take $p_i=U$ for all $i$, with a single uniform $U$ under the global null. Compute FWER for both Bonferroni and Holm.
\end{exercise}

\begin{exercise}[A choice of weights]
\label{exer:multiple-weights}
For $m=3$ and weights $(1/2,1/4,1/4)$, tabulate the rejection thresholds at level $0.04$.
\end{exercise}

\begin{exercise}[A zero weight]
\label{exer:multiple-zero-weight}
Show that a valid $p$-value equals zero with probability zero under its null. Explain what a zero weight can do under the alternative.
\end{exercise}

\begin{exercise}[First failure in Holm's rule]
\label{exer:multiple-holm-first-failure}
At level $0.05$, compare Holm, Bonferroni, and Benjamini--Hochberg
on $(0.012,0.028,0.029)$. Explain the role of stopping.
\end{exercise}

\begin{exercise}[An order-invariant tie]
\label{exer:multiple-holm-tie}
Show that the Holm rejection set does not depend on the ordering within a block of equal $p$-values.
\end{exercise}

\begin{exercise}[Only one true null]
\label{exer:multiple-single-true-null}
At a label with exactly one true null, relate FWER to the rejection probability of that null. Give an example with strict inequality between FDR and FWER.
\end{exercise}

\begin{exercise}[The probability of joint coverage]
\label{exer:multiple-joint-coverage}
For $m$ independent intervals each of exact coverage $1-\alpha/m$, compute the coverage of their product.
\end{exercise}

\begin{exercise}[Weighted simultaneous intervals]
\label{exer:multiple-weighted-coverage}
Generalize Example~\ref{ex:bonferroni-simultaneous-confidence} to marginal noncoverage bounds $\alpha w_i$.
\end{exercise}

\begin{exercise}[Dependent null values making Bonferroni exact]
\label{exer:multiple-exact-bonferroni}
For $0<\alpha<1$, construct $m$ jointly distributed uniform $p$-values for which the events $\{p_i\leq\alpha/m\}$ are pairwise disjoint up to null sets. Verify exact FWER.
\end{exercise}

\begin{exercise}[Subset coverage from a family of intervals]
\label{exer:multiple-subset-coverage}
Given simultaneous confidence sets $C_i$ for $g_i(\theta)$, reject the null $g_i(\theta)=t_i$ when $t_i\notin C_i$. Prove FWER control even if some nulls are false.
\end{exercise}

\begin{exercise}[Holm adjusted values]
\label{exer:multiple-holm-adjusted}
Show that Holm rejects an index at rank $j$ exactly when $\min\{1,\max_{k\leq j}(m-k+1)p_{(k)}\}\leq\alpha$ for $0\leq\alpha<1$.
\end{exercise}

\begin{exercise}[Counting false rejections]
\label{exer:multiple-false-count-tails}
Let $m_0$ nulls be true. If each is rejected with probability at most $\alpha/m$, bound $\mathbb E_\theta V$ and $P_\theta(V\geq k)$ for every integer $k\geq1$.
\end{exercise}

\begin{exercise}[Minimax randomization on a shared atom]
\label{exer:multiple-deterministic-minimax}
In Example~\ref{ex:multiclass-shared-observation}, compute the smallest worst-case error among deterministic rules. Compare it with the randomized minimax value.
\end{exercise}

\begin{exercise}[Worst-case error and a prior bound]
\label{exer:multiple-equal-risk-prior}
Show that any classifier whose errors are equal at all labels and whose prior average error is Bayes for some prior is minimax. Explain which inequalities become equalities.
\end{exercise}


\chapter{Bayesian inference}
\label{chap:bayesian-inference}
\notationgroup{Bayesian inference}

\section{Priors, posterior laws, and existence conditions}
\label{sec:bayesian-posteriors}

\subsection{Posterior kernels and their existence}
\label{subsec:bayesian-posterior-existence}

Let $P$ be a probability kernel from $(\Theta,\mathcal T)$ to
$(\mathsf X,\mathcal A)$ and let $\pi$ be a prior. On
$(\Theta\times\mathsf X,\mathcal T\otimes\mathcal A)$ with law
$J_\pi=\pi\ltimes P$, write $U$ for the parameter coordinate and $X$
for the observation coordinate. Their marginal laws are $\pi$ and
$M_\pi$ (Definition~\ref{def:prior-predictive}). A posterior $\Pi$ is
the conditional law of $U$ given $X$ in the sense of
Definition~\ref{def:posterior-kernel}.

\begin{definition}[Standard Borel space]
\label{def:standard-borel-space}
A measurable space is \emph{standard Borel} if it admits a bijection
onto a Borel subset of $\R$, with its trace Borel sigma algebra, such
that both the bijection and its inverse are measurable. A countable
space with all subsets measurable and a Borel subset of $\R$ with
its Borel sigma algebra are examples.
\end{definition}

We use the distribution-function theorem on $\R$. Every nondecreasing,
right-continuous function with limits zero at $-\infty$ and one at
$+\infty$ determines a unique Borel probability law having that
distribution function.

\begin{theorem}[Existence and uniqueness of a posterior]
\label{thm:bayesian-posterior-existence}
If $(\Theta,\mathcal T)$ is standard Borel, then $\pi$ and $P$
admit a posterior kernel. Any two posterior kernels agree as
probability measures outside a single $M_\pi$-null set of records.
The observation space may be arbitrary.
\end{theorem}
\begin{proof}
Choose a measurable isomorphism $h:\Theta\to S$ onto a Borel subset
$S$ of $\R$. For $q\in\mathbb Q$, the measure
\[
A\longmapsto J_\pi\{X\in A,\ h(U)\leq q\}
\]
is bounded by $M_\pi$. Let $f_q:\mathsf X\to[0,1]$ be a
Radon--Nikodym density. Off a single measurable $M_\pi$-null set,
\[
f_q\leq f_r\ (q<r),\qquad
f_{q+1/n}\downarrow f_q,\qquad
f_{-n}\downarrow 0,\qquad f_n\uparrow 1 .
\]
The inequalities follow from the ordering of the measures. For each
displayed limit, integrate over any $A\in\mathcal A$ and use bounded
convergence and the corresponding limit of events. Uniqueness of the
densities identifies the limits. There are countably many conditions.

At each record outside that set, put $F_x(t)=\inf_{q>t,\ q\in\mathbb Q}
f_q(x)$. It is nondecreasing and right-continuous. Indeed, if $t_k\downarrow
t$, each rational $q>t$ is greater than $t_k$ eventually, so
$\lim_k F_x(t_k)\leq f_q(x)$. Taking the infimum and using monotonicity
gives $\lim_k F_x(t_k)=F_x(t)$. The rational limits above give
$F_x(q)=f_q(x)$ and the limits zero and one at the two ends of $\R$.
Thus $F_x$ determines a probability law $Q(x,\cdot)$ on $\R$ by the
distribution-function theorem. On the exceptional set use a fixed point
mass at $h(\theta_0)$, where $\theta_0\in\Theta$.

The sets $C$ for which $x\mapsto Q(x,C)$ is measurable form a
Dynkin system containing the rational lower half-lines and $\R$.
They therefore include all Borel sets. For each $A\in\mathcal A$,
the two measures
\[
C\longmapsto \int_A Q(x,C)\,M_\pi(dx),\qquad
C\longmapsto J_\pi\{X\in A,\ h(U)\in C\}
\]
agree on the same generating system by the defining densities and
hence on all Borel sets. In particular, with $A=\mathsf X$ and
$C=S$, their common mass is one, so $Q(x,S)=1$ almost surely.
Where $Q(x,S)=1$, set $\Pi(x,B)=Q(x,h(B))$. Elsewhere set
$\Pi(x,\cdot)=\delta_{\theta_0}$. Each $h(B)$ is Borel, so $\Pi$
is a kernel and satisfies \eqref{eq:posterior-rectangles}.

The inverse images of the rational lower half-lines, together with
$\Theta$, form a countable generating system closed under finite
intersections. Uniqueness now follows from
Proposition~\ref{prop:conditional-law-versions}.
\end{proof}

For the regular conditional distribution theorem, see J.~Norris,
\href{https://www.statslab.cam.ac.uk/~jrn10/Lectures/ap.pdf}{\emph{Advanced
Probability}}, Section~1.6.

\begin{proposition}[Posterior integration]
\label{prop:bayesian-posterior-integration}
\indexnotation{\Pi g(x)}{Posterior integral of a parameter function \(g\).}
Suppose a posterior exists. For measurable $g:\Theta\to[0,\infty]$,
or real $g$ with $\int |g|\,d\pi<\infty$, put
$\Pi g(x)=\int g(\theta)\,\Pi(x,d\theta)$. In the integrable case use
a measurable real extension where the posterior absolute moment is
infinite.
Then $(\Pi g)(X)$ is a version of $\mathbb E_{J_\pi}[g(U)\mid\sigma(X)]$,
and
\[
\int_{\mathsf X}\Pi g(x)\,M_\pi(dx)=\int_\Theta g\,d\pi .
\]
\end{proposition}
\begin{proof}
Apply \eqref{eq:kernel-iterated-integral} to the reversed joint law
$M_\pi\ltimes\Pi$ with integrand $(x,\theta)\mapsto\mathbf 1_A(x)g(\theta)$
for each $A\in\mathcal A$. For integrable $g$, the same formula for
$|g|$ gives finiteness of its posterior integral almost surely, so one
may take the difference of the positive and negative parts. Taking
$A=\mathsf X$ gives the displayed identity.
\end{proof}

\begin{example}[Two Bernoulli rates]
\label{ex:bayesian-unequal-bernoulli-rates}
Give the rates $1/5$ and $3/5$ prior masses $3/4$ and $1/4$.
For one Bernoulli observation, $M_\pi(\{1\})=3/10$ and
$M_\pi(\{0\})=7/10$. By Proposition~\ref{prop:posterior-from-densities},
the posterior mass of $3/5$ is
\[
\frac{(1/4)(3/5)}{3/10}=\frac12\quad\text{after success},\qquad
\frac{(1/4)(2/5)}{7/10}=\frac17\quad\text{after failure}.
\]
Averaging these two masses with the predictive weights gives $1/4$.
\end{example}

\begin{example}[A continuous observation with moving support]
\label{ex:bayesian-uniform-endpoint-posterior}
Give $\theta$ the uniform prior on $(0,1)$ and take
$X\mid U=\theta$ uniform on $(0,\theta)$. For $0<x<1$ the
predictive density is $m(x)=\int_x^1 d\theta/\theta=-\log x$.
The posterior density at $x$ is
\[
\theta\longmapsto \frac{\mathbf 1_{\{x<\theta<1\}}}{-\theta\log x}.
\]
For $x\leq t\leq1$ its probability of $(0,t]$ is
$\log(t/x)/\log(1/x)$. In particular, its median is $\sqrt{x}$.
\end{example}

\begin{example}[An exact continuous observation]
\label{ex:bayesian-exact-continuous-record}
Let $\Theta=\mathsf X=[0,1]$, with Borel sets, take the uniform prior,
and put $P_\theta=\delta_\theta$. Then $M_\pi=\pi$ and
$\Pi(x,\cdot)=\delta_x$, since both sides of
\eqref{eq:posterior-rectangles} are $\pi(A\cap B)$. The posterior is
a point mass even though $M_\pi(\{x\})=0$ for every $x$. No
$\sigma$-finite measure dominates all the $P_\theta$: such a measure
would have to assign positive mass to uncountably many singletons,
whereas a finite-measure set contains at most countably many
positive-mass singletons.
\end{example}

\begin{example}[A magnitude without its sign]
\label{ex:bayesian-magnitude-posterior}
Take $U$ uniform on $[-1,1]$ and $X=|U|$. The predictive law is
uniform on $[0,1]$. For Borel $A\subseteq[0,1]$ and
$B\subseteq[-1,1]$,
\[
J_\pi\{X\in A,\ U\in B\}
=\tfrac12\int_A\bigl(\mathbf 1_B(x)+\mathbf 1_B(-x)\bigr)\,dx.
\]
Thus $\tfrac12\delta_x+\tfrac12\delta_{-x}$ is a posterior at $x$.
At $x=0$ its value is $\delta_0$. Replacing this value by $\delta_1$
also gives a posterior kernel, since the predictive mass of
$\{0\}$ is zero.
\end{example}

\begin{exercise}[A record of positive mass]
\label{exer:bayesian-posterior-atom}
Suppose $\{x\}\in\mathcal A$ and $M_\pi(\{x\})>0$. Derive the
posterior at $x$ by taking $A=\{x\}$ in
\eqref{eq:posterior-rectangles}. Show that its density with respect
to $\pi$ is $P_\theta(\{x\})/M_\pi(\{x\})$.
\end{exercise}

\begin{exercise}[One uninformative experiment]
\label{exer:bayesian-constant-kernel}
Suppose $P_\theta=Q$ at every label. Verify that the constant kernel
$\Pi(x,\cdot)=\pi$ is a posterior. Under the hypotheses of
Theorem~\ref{thm:bayesian-posterior-existence}, show that every
posterior agrees with it outside a single $Q$-null set.
\end{exercise}

\begin{exercise}[An erased binary observation]
\label{exer:bayesian-erased-label-posterior}
Put prior mass $s\in(0,1)$ on label one and mass $1-s$ on label zero.
The observation gives the label with probability $1-e$ and an erasure
with probability $e\in(0,1)$ at both labels. Compute the predictive law
and all posterior laws at the three records.
\end{exercise}

\begin{exercise}[Posterior moments in the endpoint model]
\label{exer:bayesian-endpoint-posterior-moments}
In Example~\ref{ex:bayesian-uniform-endpoint-posterior}, compute
$\Pi g(x)$ for $g(\theta)=\theta^a$ and $a>0$. Integrate the result
against $m(x)\,dx$ to recover the prior moment.
\end{exercise}

\begin{exercise}[A change on a null set]
\label{exer:bayesian-posterior-null-replacement}
Given a posterior $\Pi$, replace it on a measurable $M_\pi$-null
set by any other probability kernel. Show that the resulting kernel
is a posterior. Give a replacement on a set of positive predictive
probability that fails the rectangle identity.
\end{exercise}

\begin{exercise}[A biased sign]
\label{exer:bayesian-biased-sign-posterior}
Let $V$ be uniform on $[0,1]$. Given $V=v$, let $U=v$ with
probability $v$ and $U=-v$ with probability $1-v$. For $X=|U|$,
find a posterior for $U$ and verify it on rectangles. Compute
$\mathbb E[U\mid X]$.
\end{exercise}

\begin{exercise}[Exact observation mixed with noise]
\label{exer:bayesian-exact-or-noise}
Let $U$ and $Z$ be independent uniforms on $[0,1]$. Independently
draw a Bernoulli variable with success probability $r\in(0,1)$ and
record $X=U$ on success and $X=Z$ on failure. Find $M_\pi$ and
verify that the posterior is $r\delta_x+(1-r)\pi$.
\end{exercise}

\begin{exercise}[Observation events and posterior probabilities]
\label{exer:bayesian-average-event-posterior}
For $A\in\mathcal A$ with $M_\pi(A)>0$, express the conditional
probability of $U\in B$ given $X\in A$ as an average of
$\Pi(x,B)$. Evaluate it for $A=(1/2,1]$ and $B=(3/4,1]$ in
Example~\ref{ex:bayesian-magnitude-posterior}.
\end{exercise}

\subsection{Sequential Bayesian updating}
\label{subsec:bayesian-sequential-updating}

Let $X$ be an initial observation with law $P_\theta$ at parameter
$\theta$. Let $\Pi$ be its posterior under $\pi$ and $M_\pi$ its
predictive law. At a pair $(\theta,x)$ the next observation $Y$ has
law $K((\theta,x),\cdot)$ on $(\mathsf Y,\mathcal Y)$, where $K$ is
a probability kernel. The joint law of $(U,X,Y)$ is
$(\pi\ltimes P)\ltimes K$.

\begin{definition}[Conditional predictive law]
\label{def:bayesian-conditional-predictive}
\indexnotation{M(\cdot\mid x)}{Predictive law of the next observation after record \(x\).}
The \emph{conditional predictive law} of $Y$ after $x$ is
\[
M(C\mid x)=\int_\Theta K((\theta,x),C)\,\Pi(x,d\theta)
\qquad(C\in\mathcal Y).
\]
\end{definition}

\begin{theorem}[Updating the posterior with a conditional likelihood]
\label{thm:bayesian-conditional-update}
\indexnotation{\Pi^+(x,y;\cdot),\ m(y\mid x)}{Updated posterior and conditional predictive density.}
Suppose $K((\theta,x),\cdot)$ has a jointly measurable
nonnegative density $f(\theta,x,y)$ with respect to a fixed
$\sigma$-finite measure $\nu$ on $\mathsf Y$. Put
\[
m(y\mid x)=\int_\Theta f(\theta,x,y)\,\Pi(x,d\theta).
\]
Then $M(\cdot\mid x)$ is a conditional law of $Y$ given $X$ and has
density $m(\cdot\mid x)$ for each $x$. A posterior for $U$ given
$(X,Y)$ is
\begin{equation}
\Pi^+(x,y;B)=
\frac{\int_B f(\theta,x,y)\,\Pi(x,d\theta)}{m(y\mid x)}
\quad\text{when }0<m(y\mid x)<\infty .
\label{eq:bayesian-conditional-update}
\end{equation}
The other records form a null set for the law of $(X,Y)$.
On that set one may use $\Pi(x,\cdot)$.
\end{theorem}
\begin{proof}
The kernel integral formula in Proposition~\ref{prop:joint-action-law}
gives measurability of $M(C\mid x)$ by using the kernel
$x\mapsto\Pi(x,\cdot)$ with integrand $K((\theta,x),C)$.
For each $x$, monotone convergence gives countable additivity and
the total mass is one. For $A\in\mathcal A$ and $C\in\mathcal Y$,
reversing the first two coordinates of $\pi\ltimes P$ gives
\[
\mathbb P(X\in A,\ Y\in C)
=\int_A\int_\Theta K((\theta,x),C)\,\Pi(x,d\theta)\,M_\pi(dx).
\]
Thus the observation law is $M_\pi\ltimes M$, with $M$ the kernel
$x\mapsto M(\cdot\mid x)$. Tonelli gives the density $m(y\mid x)$ and
$\int m(y\mid x)\,\nu(dy)=1$ for each $x$. In particular, $m=0$
and $m=\infty$ both occur on sets of observation probability zero.

For fixed $B$, the numerator in \eqref{eq:bayesian-conditional-update}
is measurable in $(x,y)$ by the kernel integral formula, using
$(x,y)\mapsto\Pi(x,\cdot)$. On the stated set it defines, after
normalization, a probability measure in $B$. The fallback is a
kernel on the exceptional set. To verify the posterior identity,
integrate $\Pi^+(x,y;B)$ over an observation rectangle $A\times C$.
Cancellation of the density gives
\[
\int_A\int_C\int_B f(\theta,x,y)\,\Pi(x,d\theta)\,\nu(dy)\,M_\pi(dx)
=
\int_A\int_B K((\theta,x),C)\,\Pi(x,d\theta)\,M_\pi(dx).
\]
The inner integral on the left is zero where $m=0$, and the set
$m=\infty$ is $\nu$-null for each $x$. By reversal of the first
two coordinates, the right side is $\mathbb P(U\in B,X\in A,Y\in C)$.
For fixed $B$ the two finite measures on the observation space
agree on rectangles, so they agree everywhere.
\end{proof}

\begin{proposition}[Sequential and batch updates in a product experiment]
\label{prop:bayesian-batch-sequential}
\indexnotation{\Pi_t,\ m_{1:t}}{Posterior kernel and predictive density for the first \(t\) observations.}
For the product experiment of Definition~\ref{def:product-experiment},
suppose each $P_{\theta,i}$ has a jointly measurable finite density
$f_i(\theta,x_i)$ with respect to a fixed $\sigma$-finite measure
$\mu_i$. For $1\leq t\leq n$ put
\[
L_t(\theta;x_{1:t})=\prod_{i=1}^t f_i(\theta,x_i),\qquad
m_{1:t}(x_{1:t})=\int_\Theta L_t(\theta;x_{1:t})\,\pi(d\theta),
\qquad m_{1:0}=1.
\]
Here $x_{1:t}=(x_1,\ldots,x_t)$. Starting from $\Pi_0=\pi$,
there is a set of full predictive probability for the complete
record on which every prefix has $0<m_{1:t}<\infty$. On this set
the successive normalized updates give, for every $t$,
\[
\Pi_t(x_{1:t},d\theta)
=\frac{L_t(\theta;x_{1:t})}{m_{1:t}(x_{1:t})}\,\pi(d\theta),
\qquad
m(x_t\mid x_{1:t-1})
=\frac{m_{1:t}(x_{1:t})}{m_{1:t-1}(x_{1:t-1})}.
\]
In particular, the final update equals the batch update as a
probability measure on this set.
\end{proposition}
\begin{proof}
The product of the densities is the joint density at each label
because it gives the product law on rectangles by Tonelli.
Proposition~\ref{prop:posterior-from-densities} gives the batch formula
and implies that each prefix has a positive finite predictive density
almost surely. Intersect the finitely many prefix events under the full
predictive law.

At a record in this intersection, suppose the batch formula holds at
$t-1$. Integrating $f_t(\theta,x_t)$ against that posterior gives
$m_{1:t}/m_{1:t-1}$. Multiplying the same posterior by $f_t$ and
dividing by this ratio gives the batch formula at $t$ by
\eqref{eq:bayesian-conditional-update}. For $t=1$ the update is
Proposition~\ref{prop:posterior-from-densities}. Induction completes
the proof.
\end{proof}

\begin{example}[A success followed by a failure]
\label{ex:bayesian-sequential-two-rates}
In Example~\ref{ex:bayesian-unequal-bernoulli-rates}, observe two
conditionally independent Bernoulli trials. A success first leaves
posterior masses $1/2,1/2$ on the two rates. The next predictive
failure probability is $(4/5+2/5)/2=3/5$. A failure changes the mass
of $3/5$ to $(1/2)(2/5)/(3/5)=1/3$. In the batch update the two
unnormalized masses are
\[
(3/4)(1/5)(4/5)=3/25,\qquad
(1/4)(3/5)(2/5)=3/50.
\]
Their sum $9/50$ is also the product $(3/10)(3/5)$ of the
successive predictive probabilities.
\end{example}

\begin{example}[Two records with moving support]
\label{ex:bayesian-sequential-endpoint}
In Example~\ref{ex:bayesian-uniform-endpoint-posterior}, take two
conditionally independent draws $x,y$ and set $s=\max(x,y)$.
For $0<x,y<1$, the second likelihood multiplies the first posterior
density by $\theta^{-1}\mathbf 1_{\{y<\theta<1\}}$. The resulting
posterior and next predictive densities are
\[
\frac{\theta^{-2}\mathbf 1_{\{s<\theta<1\}}}{1/s-1},
\qquad m(y\mid x)=\frac{1/s-1}{-\log x}.
\]
Indeed, $\int_s^1\theta^{-2}\,d\theta=1/s-1$, which is the joint
predictive density by the product formula.
\end{example}

\begin{example}[The first observation is a magnitude]
\label{ex:bayesian-sequential-noisy-sign}
In Example~\ref{ex:bayesian-magnitude-posterior}, for $x>0$ the first
posterior gives equal mass to $x$ and $-x$. A second observation
reports the sign correctly with probability $3/4$ under either
sign. Its predictive probability of reporting positive is $1/2$.
After a positive report the posterior is
$\tfrac34\delta_x+\tfrac14\delta_{-x}$.
The update uses counting measure for the sign report. The first
observation was handled by its posterior kernel.
\end{example}

\begin{example}[Copying an earlier observation]
\label{ex:bayesian-sequential-copy}
Use the prior of Example~\ref{ex:bayesian-unequal-bernoulli-rates},
observe a Bernoulli bit $X$, and set $Y=X$. The second density is
$f(\theta,x,y)=\mathbf 1_{\{y=x\}}$ with respect to counting
measure. At $(1,1)$ it leaves the posterior mass of $3/5$ at $1/2$.
Treating the copy as a fresh Bernoulli draw would give
\[
\frac{(1/2)(3/5)}{(1/2)(1/5)+(1/2)(3/5)}=\frac34.
\]
At $(1,0)$ the conditional predictive mass is zero, so the update
in \eqref{eq:bayesian-conditional-update} uses the fallback there.
\end{example}

\begin{exercise}[Two successive successes]
\label{exer:bayesian-sequential-successes}
In Example~\ref{ex:bayesian-sequential-two-rates}, compute the
predictive probability of two successes and the corresponding
posterior mass of $3/5$. Check both by a batch calculation.
\end{exercise}

\begin{exercise}[Dependence under the predictive law]
\label{exer:bayesian-sequential-predictive-dependence}
Under any prior on the Bernoulli rate $p$, prove that the predictive
covariance of two conditionally independent trials is
$\operatorname{Var}_\pi(p)$. Determine when the two trials are
independent under their predictive law.
\end{exercise}

\begin{exercise}[Reversing the order of observations]
\label{exer:bayesian-sequential-order}
For a product experiment with finite jointly measurable densities,
fix a permutation of the observations. Show that the permuted and
original updates give the same final posterior on a common set
of full predictive probability. Compare the intermediate
posteriors in Example~\ref{ex:bayesian-sequential-two-rates}.
\end{exercise}

\begin{exercise}[A likelihood that is constant on the posterior support]
\label{exer:bayesian-sequential-constant-likelihood}
At a fixed $(x,y)$, suppose $f(\theta,x,y)$ is equal to a common
positive finite value for $\Pi(x,\cdot)$-almost every $\theta$.
Show that the posterior at this record stays at $\Pi(x,\cdot)$.
\end{exercise}

\begin{exercise}[A second experiment chosen from the first record]
\label{exer:bayesian-sequential-adaptive-experiment}
Use the prior from Example~\ref{ex:bayesian-unequal-bernoulli-rates}.
After observing the first Bernoulli bit $X$, draw a fresh
Bernoulli bit with rate $p$ if $X=1$. If $X=0$, draw a fair bit
independently of $p$. Write the conditional likelihood and find
the predictive mass and posterior at each of the four pairs.
\end{exercise}

\begin{exercise}[The posterior mean changes by a mean-zero increment]
\label{exer:bayesian-sequential-mean-increment}
Let $g$ be integrable under $\pi$. Under the hypotheses of
Theorem~\ref{thm:bayesian-conditional-update}, show that
$\mathbb E[\Pi^+g(X,Y)\mid \sigma(X)]=\Pi g(X)$ almost surely.
Here $\Pi^+g(x,y)=\int g(\theta)\,\Pi^+(x,y;d\theta)$ where
the integral is finite. Derive the identity by integrating over
sets $\{X\in A\}$.
\end{exercise}

\begin{exercise}[A decreasing upper tail after more draws]
\label{exer:bayesian-sequential-endpoint-tail}
In Example~\ref{ex:bayesian-sequential-endpoint}, for
$s\leq u\leq1$ compute the posterior probability of $\theta>u$.
If $0<y\leq x$, compare it with the probability after $x$ alone.
\end{exercise}

\begin{exercise}[The probability of the record]
\label{exer:bayesian-sequential-predictive-product}
Under Proposition~\ref{prop:bayesian-batch-sequential}, prove that the
product of the one-step predictive densities equals $m_{1:n}$
on the stated set of full predictive probability. With finite observation spaces and counting measures, show that
this product is the probability of the record.
\end{exercise}

% BEGIN NOTATION INDEX
\chapter*{Notation index}
\label{chap:notation-index}
\addcontentsline{toc}{chapter}{Notation index}

The linked page numbers lead to definitions or first uses. A symbol may
have a different meaning in a different part of the book.

\printnotationentries
% END NOTATION INDEX
\end{document}
